\input{preamble}

% OK, start here.
%
\begin{document}

\title{Compléments sur la platitude}

\maketitle

\phantomsection
\label{section-phantom}

\tableofcontents



\section{Introduction}
\label{section-introduction}

\noindent
Dans ce chapitre, nous exposons des résultats avancés sur les modules plats et les
morphismes plats de schémas, ainsi que des applications. La plupart des résultats de platitude
figurent dans l'article \cite{GruRay} de Raynaud et Gruson.

\medskip\noindent
Avant de lire ce chapitre, nous conseillons au lecteur de prendre connaissance
des résultats suivants (cette liste sert aussi à retrouver
les résultats précédents) :
\begin{enumerate}
\item Généralités sur les modules plats dans
Algèbre, Section \ref{algebra-section-flat}.
\item Relations entre les groupes $\text{Tor}$ et la platitude, voir
Algèbre, Section \ref{algebra-section-tor}.
\item Critères de platitude, voir
Algèbre, Section \ref{algebra-section-criteria-flatness}
(cas noethérien),
Algèbre, Section \ref{algebra-section-flatness-artinian}
(cas artinien),
Algèbre, Section \ref{algebra-section-more-flatness-criteria}
(cas non noethérien), et enfin
Compléments sur les morphismes, Section \ref{more-morphisms-section-criterion-flat-fibres}.
\item Platitude générique, voir
Algèbre, Section \ref{algebra-section-generic-flatness}
et
Morphismes, Section \ref{morphisms-section-generic-flatness}.
\item Ouverture du lieu de platitude, voir
Algèbre, Section \ref{algebra-section-open-flat}
et
Compléments sur les morphismes, Section \ref{more-morphisms-section-open-flat}.
\item Aplatissement, voir
Compléments d'algèbre, Sections\newline
\ref{more-algebra-section-flattening},
\ref{more-algebra-section-flattening-artinian},
\ref{more-algebra-section-flattening-local-base},
\ref{more-algebra-section-flattening-local-source-base}, et
\ref{more-algebra-section-flattening-Noetherian-complete-local}.
\item Résultats supplémentaires dans
Compléments d'algèbre, Sections \ref{more-algebra-section-descent-flatness-integral},
\ref{more-algebra-section-torsion-flat},
\ref{more-algebra-section-flat-finite}, et
\ref{more-algebra-section-blowup-flat}.
\end{enumerate}
Comme applications des résultats de platitude, nous traitons les sujets
suivants : une version non noethérienne du théorème d'existence de Grothendieck,
les éclatements et la platitude, le théorème de compactification de Nagata,
la topologie h, les carrés d'éclatement et la descente, la normalisation faible,
la descente des fibrés vectoriels en caractéristique positive et
la structure locale des complexes parfaits pour la topologie h.






\section{Lemmes sur la localisation \'etale}
\label{section-etale-localization}

\noindent
Nous réunissons dans cette section quelques lemmes de localisation \'etale qui seront
utiles dans la suite de ce chapitre. On peut omettre cette section en première lecture.

\begin{lemma}
\label{lemma-lift-etale}
Soit $i : Z \to X$ une immersion fermée de schémas affines.
Soit $Z' \to Z$ un morphisme \'etale, avec $Z'$ affine.
Il existe alors un morphisme \'etale $X' \to X$, avec $X'$
affine, tel que $Z' \cong Z \times_X X'$ comme schémas sur $Z$.
\end{lemma}

\begin{proof}
Voir
Algèbre, Lemme \ref{algebra-lemma-lift-etale}.
\end{proof}

\begin{lemma}
\label{lemma-etale-at-point}
Soit
$$
\xymatrix{
X \ar[d] & X' \ar[l] \ar[d] \\
S & S' \ar[l]
}
$$
un diagramme commutatif de schémas où $X' \to X$ et $S' \to S$ sont \'etales.
Soit $s' \in S'$ un point. Alors
$$
X' \times_{S'} \Spec(\mathcal{O}_{S', s'})
\longrightarrow
X \times_S \Spec(\mathcal{O}_{S', s'})
$$
est \'etale.
\end{lemma}

\begin{proof}
Cela résulte du fait que $X' \to X_{S'}$ est \'etale, en tant que morphisme de
schémas \'etales sur $X$, voir
Morphismes, Lemme \ref{morphisms-lemma-etale-permanence},
et du fait que le changement de base d'un morphisme \'etale est \'etale, voir
Morphismes, Lemme \ref{morphisms-lemma-base-change-etale}.
\end{proof}

\begin{lemma}
\label{lemma-etale-flat-up-down}
Soient $X \to T \to S$ des morphismes de schémas, avec $T \to S$ \'etale.
Soit $\mathcal{F}$ un $\mathcal{O}_X$-module quasi-cohérent.
Soit $x \in X$ un point. Alors
$$
\mathcal{F}\text{ flat over }S\text{ at }x
\Leftrightarrow
\mathcal{F}\text{ flat over }T\text{ at }x
$$
En particulier, $\mathcal{F}$ est plat sur $S$ si et seulement si $\mathcal{F}$
est plat sur $T$.
\end{lemma}

\begin{proof}
Puisqu'un morphisme \'etale est un morphisme plat (voir
Morphismes, Lemme \ref{morphisms-lemma-etale-flat}),
l'implication ``$\Leftarrow$'' résulte de
Algèbre, Lemme \ref{algebra-lemma-composition-flat}.
Réciproquement, supposons que $\mathcal{F}$ soit plat en $x$ sur $S$.
Notons $\tilde x \in X \times_S T$ le point au-dessus de $x$ dans $X$
et de l'image de $x$ dans $T$, prise comme point de $T$.
Alors $(X \times_S T \to X)^*\mathcal{F}$ est plat en $\tilde x$ sur $T$
pour $\text{pr}_2 : X \times_S T \to T$, voir
Morphismes, Lemme \ref{morphisms-lemma-base-change-module-flat}.
La diagonale $\Delta_{T/S} : T \to T \times_S T$ est une immersion ouverte ;
il suffit de combiner
Morphismes, Lemmes \ref{morphisms-lemma-diagonal-unramified-morphism} et
\ref{morphisms-lemma-etale-smooth-unramified}.
Ainsi $X$ s'identifie à un sous-schéma ouvert de $X \times_S T$,
la restriction de $\text{pr}_2$ à cet ouvert est le morphisme donné $X \to T$,
le point $\tilde x$ correspond au point $x$ de cet ouvert et
la restriction de $(X \times_S T \to X)^*\mathcal{F}$ à cet ouvert est $\mathcal{F}$.
On en déduit que $\mathcal{F}$ est plat en $x$ sur $T$.
\end{proof}

\begin{lemma}
\label{lemma-etale-flat-up-down-local-ring}
Soit $T \to S$ un morphisme \'etale. Soit $t \in T$ d'image $s \in S$.
Soit $M$ un $\mathcal{O}_{T, t}$-module. Alors
$$
M\text{ flat over }\mathcal{O}_{S, s}
\Leftrightarrow
M\text{ flat over }\mathcal{O}_{T, t}.
$$
\end{lemma}

\begin{proof}
On peut remplacer $S$ par un voisinage affine de $s$, puis
$T$ par un voisinage affine de $t$.
Posons $\mathcal{F} = (\Spec(\mathcal{O}_{T, t}) \to T)_*\widetilde M$.
C'est un faisceau quasi-cohérent (voir
Schémas, Lemme \ref{schemes-lemma-push-forward-quasi-coherent},
ou raisonner directement)
sur $T$, dont la fibre en $t$ est $M$ (détails omis).
On applique
le Lemme \ref{lemma-etale-flat-up-down}.
\end{proof}

\begin{lemma}
\label{lemma-flat-up-down-henselization}
Soient $S$ un schéma et $s \in S$ un point. Notons $\mathcal{O}_{S, s}^h$
(resp.\ $\mathcal{O}_{S, s}^{sh}$) le hensélisé (resp.\ le hensélisé
strict), voir
Algèbre, Définition \ref{algebra-definition-henselization}.
Soit $M^{sh}$ un $\mathcal{O}_{S, s}^{sh}$-module.
Les conditions suivantes sont équivalentes :
\begin{enumerate}
\item $M^{sh}$ est plat sur $\mathcal{O}_{S, s}$,
\item $M^{sh}$ est plat sur $\mathcal{O}_{S, s}^h$, et
\item $M^{sh}$ est plat sur $\mathcal{O}_{S, s}^{sh}$.
\end{enumerate}
Si $M^{sh} = M^h \otimes_{\mathcal{O}_{S, s}^h} \mathcal{O}_{S, s}^{sh}$,
elles sont aussi équivalentes à
\begin{enumerate}
\item[(4)] $M^h$ est plat sur $\mathcal{O}_{S, s}$, et
\item[(5)] $M^h$ est plat sur $\mathcal{O}_{S, s}^h$.
\end{enumerate}
Si $M^h = M \otimes_{\mathcal{O}_{S, s}} \mathcal{O}_{S, s}^h$,
elles sont aussi équivalentes à
\begin{enumerate}
\item[(6)] $M$ est plat sur $\mathcal{O}_{S, s}$.
\end{enumerate}
\end{lemma}

\begin{proof}
D'après Compléments d'algèbre, Lemme
\ref{more-algebra-lemma-dumb-properties-henselization},
les homomorphismes locaux
$\mathcal{O}_{S, s} \to \mathcal{O}_{S, s}^h \to \mathcal{O}_{S, s}^{sh}$
sont fidèlement plats.
Ainsi (3) $\Rightarrow$ (2) $\Rightarrow$ (1) et
(5) $\Rightarrow$ (4) résultent de
Algèbre, Lemme \ref{algebra-lemma-composition-flat}.
Par fidèle platitude, les équivalences (6) $\Leftrightarrow$ (5) et
(5) $\Leftrightarrow$ (3) résultent de
Algèbre, Lemme \ref{algebra-lemma-flatness-descends}.
Il suffit donc de montrer que
(1) $\Rightarrow$ (2) $\Rightarrow$ (3) et
(4) $\Rightarrow$ (5).
Pour cela, on peut supposer que $S$ est un schéma affine.

\medskip\noindent
Supposons (1). D'après
le Lemme \ref{lemma-etale-flat-up-down-local-ring},
on voit que $M^{sh}$ est plat sur $\mathcal{O}_{T, t}$ pour
tout voisinage \'etale $(T, t) \to (S, s)$. Puisque $\mathcal{O}_{S, s}^h$
et $\mathcal{O}_{S, s}^{sh}$ sont des limites inductives filtrantes d'anneaux locaux
de la forme $\mathcal{O}_{T, t}$ (voir
Algèbre, Lemmes \ref{algebra-lemma-henselization-different}
et \ref{algebra-lemma-strict-henselization-different}),
on en conclut que $M^{sh}$ est plat sur $\mathcal{O}_{S, s}^h$
et $\mathcal{O}_{S, s}^{sh}$ d'après
Algèbre, Lemme \ref{algebra-lemma-colimit-rings-flat}.
Ainsi (1) implique (2) et (3). Cela entraîne aussi
(2) $\Rightarrow$ (3) en remplaçant $\mathcal{O}_{S, s}$ par
$\mathcal{O}_{S, s}^h$. Le même raisonnement démontre
(4) $\Rightarrow$ (5).
\end{proof}

\begin{lemma}
\label{lemma-tor-amplitude-up-down-henselization}
Soient $S$ un schéma et $s \in S$ un point. Notons $\mathcal{O}_{S, s}^h$
(resp.\ $\mathcal{O}_{S, s}^{sh}$) le hensélisé (resp.\ le hensélisé
strict), voir
Algèbre, Définition \ref{algebra-definition-henselization}.
Soit $M^{sh}$ un objet de $D(\mathcal{O}_{S, s}^{sh})$.
Soient $a, b \in \mathbf{Z}$.
Les conditions suivantes sont équivalentes :
\begin{enumerate}
\item $M^{sh}$ est de Tor-amplitude contenue dans $[a, b]$ sur $\mathcal{O}_{S, s}$,
\item $M^{sh}$ est de Tor-amplitude contenue dans $[a, b]$ sur $\mathcal{O}_{S, s}^h$, et
\item $M^{sh}$ est de Tor-amplitude contenue dans $[a, b]$ sur $\mathcal{O}_{S, s}^{sh}$.
\end{enumerate}
Si $M^{sh} =
M^h \otimes_{\mathcal{O}_{S, s}^h}^\mathbf{L} \mathcal{O}_{S, s}^{sh}$
pour $M^h \in D(\mathcal{O}_{S, s}^h)$, elles sont aussi équivalentes à
\begin{enumerate}
\item[(4)] $M^h$ est de Tor-amplitude contenue dans $[a, b]$ sur $\mathcal{O}_{S, s}$, et
\item[(5)] $M^h$ est de Tor-amplitude contenue dans $[a, b]$ sur $\mathcal{O}_{S, s}^h$.
\end{enumerate}
Si $M^h = M \otimes_{\mathcal{O}_{S, s}}^\mathbf{L} \mathcal{O}_{S, s}^h$
pour $M \in D(\mathcal{O}_{S, s})$,
elles sont aussi équivalentes à
\begin{enumerate}
\item[(6)] $M$ est de Tor-amplitude contenue dans $[a, b]$ sur $\mathcal{O}_{S, s}$.
\end{enumerate}
\end{lemma}

\begin{proof}
D'après Compléments d'algèbre, Lemme
\ref{more-algebra-lemma-dumb-properties-henselization},
les homomorphismes locaux
$\mathcal{O}_{S, s} \to \mathcal{O}_{S, s}^h \to \mathcal{O}_{S, s}^{sh}$
sont fidèlement plats.
Ainsi (3) $\Rightarrow$ (2) $\Rightarrow$ (1) et
(5) $\Rightarrow$ (4) résultent de
Compléments d'algèbre, Lemme \ref{more-algebra-lemma-flat-push-tor-amplitude}.
Par fidèle platitude, les équivalences (6) $\Leftrightarrow$ (5) et
(5) $\Leftrightarrow$ (3) résultent de
Compléments d'algèbre, Lemme \ref{more-algebra-lemma-flat-descent-tor-amplitude}.
Il suffit donc de montrer que
(1) $\Rightarrow$ (3), (2) $\Rightarrow$ (3), et (4) $\Rightarrow$ (5).

\medskip\noindent
Supposons (1). En particulier, la cohomologie de $M^{sh}$ est nulle
en degrés $< a$ et $> b$. On peut donc représenter $M^{sh}$ par un complexe
$P^\bullet$ de $\mathcal{O}_{X, x}^{sh}$-modules libres tel que
$P^i = 0$ pour $i > b$
(voir, par exemple, le résultat très général
Catégories dérivées, Lemme \ref{derived-lemma-subcategory-left-resolution}).
Notons que $P^n$ est plat sur $\mathcal{O}_{S, s}$ pour tout $n$.
Considérons $\Coker(d_P^{a - 1})$. D'après Compléments d'algèbre, Lemme
\ref{more-algebra-lemma-last-one-flat},
c'est un $\mathcal{O}_{S, s}$-module plat.
Le Lemme \ref{lemma-flat-up-down-henselization} montre donc que
c'est un $\mathcal{O}_{S, s}^{sh}$-module plat.
Ainsi $\tau_{\geq a}P^\bullet$ est un complexe de
$\mathcal{O}_{S, s}^{sh}$-modules plats représentant $M^{sh}$
dans $D(\mathcal{O}_{S, s}^{sh}$ et l'on en déduit que
$M^{sh}$ est de Tor-amplitude contenue dans $[a, b]$, voir
Compléments d'algèbre, Lemme \ref{more-algebra-lemma-tor-amplitude}.
Ainsi (1) implique (3). Cela entraîne aussi
(2) $\Rightarrow$ (3) en remplaçant $\mathcal{O}_{S, s}$ par
$\mathcal{O}_{S, s}^h$. Le même raisonnement démontre
(4) $\Rightarrow$ (5).
\end{proof}

\begin{lemma}
\label{lemma-finite-flat-weak-assassin-up-down}
Soit $g : T \to S$ un morphisme fini et plat de schémas.
Soit $\mathcal{G}$ un $\mathcal{O}_S$-module quasi-cohérent.
Soit $t \in T$ un point d'image $s \in S$. Alors
$$
t \in \text{WeakAss}(g^*\mathcal{G})
\Leftrightarrow
s \in \text{WeakAss}(\mathcal{G})
$$
\end{lemma}

\begin{proof}
L'implication ``$\Leftarrow$'' résulte immédiatement de
Diviseurs, Lemme \ref{divisors-lemma-weakly-ass-pullback}.
Supposons $t \in \text{WeakAss}(g^*\mathcal{G})$.
Soit $\Spec(A) \subset S$ un voisinage ouvert affine de $s$.
Soit $\mathcal{G}$ le faisceau quasi-cohérent associé au $A$-module $M$.
Soit $\mathfrak p \subset A$ l'idéal premier correspondant à $s$.
Puisque $g$ est fini et plat, on a $g^{-1}(\Spec(A)) = \Spec(B)$
pour une certaine $A$-algèbre finie et plate $B$. Notons que
$g^*\mathcal{G}$ est le $\mathcal{O}_{\Spec(B)}$-module quasi-cohérent
associé au $B$-module $M \otimes_A B$ et que $g_*g^*\mathcal{G}$ est le
$\mathcal{O}_{\Spec(A)}$-module quasi-cohérent associé au
$A$-module $M \otimes_A B$. D'après
Algèbre, Lemme \ref{algebra-lemma-finite-flat-local},
on a $B_{\mathfrak p} \cong A_{\mathfrak p}^{\oplus n}$
pour un entier $n \geq 0$. Notons que $n \geq 1$, puisqu'on a supposé qu'il
existe au moins un point de $T$ au-dessus de $s$. En considérant
les fibres, on voit donc que
$$
s \in \text{WeakAss}(\mathcal{G})
\Leftrightarrow
s \in \text{WeakAss}(g_*g^*\mathcal{G})
$$
Or l'hypothèse $t \in \text{WeakAss}(g^*\mathcal{G})$
entraîne $s \in \text{WeakAss}(g_*g^*\mathcal{G})$ d'après
Diviseurs, Lemme \ref{divisors-lemma-weakly-associated-finite},
et donc, d'après ce qui précède, $s \in  \text{WeakAss}(\mathcal{G})$.
\end{proof}

\begin{lemma}
\label{lemma-etale-weak-assassin-up-down}
Soit $h : U \to S$ un morphisme \'etale de schémas.
Soit $\mathcal{G}$ un $\mathcal{O}_S$-module quasi-cohérent.
Soit $u \in U$ un point d'image $s \in S$. Alors
$$
u \in \text{WeakAss}(h^*\mathcal{G})
\Leftrightarrow
s \in \text{WeakAss}(\mathcal{G})
$$
\end{lemma}

\begin{proof}
En remplaçant $S$ et $U$ par des voisinages affines de $s$ et $u$,
on peut supposer que $g$ est un morphisme \'etale standard de schémas affines, voir
Morphismes, Lemme \ref{morphisms-lemma-etale-locally-standard-etale}.
On peut donc supposer $S = \Spec(A)$ et\newline
$X = \Spec(A[x, 1/g]/(f))$, où $f$ est unitaire et $f'$
est inversible dans $A[x, 1/g]$.
Notons que $A[x, 1/g]/(f) = (A[x]/(f))_g$ est aussi un localisé
de la $A$-algèbre finie et libre $A[x]/(f)$. On peut donc considérer
$U$ comme un sous-schéma ouvert du schéma $T = \Spec(A[x]/(f))$,
qui est fini localement libre sur $S$. On est ainsi ramené au
Lemme \ref{lemma-finite-flat-weak-assassin-up-down}
ci-dessus.
\end{proof}

\begin{lemma}
\label{lemma-weakly-associated-henselization}
Soient $S$ un schéma et $s \in S$ un point. Notons $\mathcal{O}_{S, s}^h$
(resp.\ $\mathcal{O}_{S, s}^{sh}$) le hensélisé (resp.\ le hensélisé
strict), voir
Algèbre, Définition \ref{algebra-definition-henselization}.
Soit $\mathcal{F}$ un $\mathcal{O}_S$-module quasi-cohérent.
Les conditions suivantes sont équivalentes :
\begin{enumerate}
\item $s$ est un point faiblement associé à $\mathcal{F}$,
\item $\mathfrak m_s$ est un idéal premier faiblement associé à $\mathcal{F}_s$,
\item $\mathfrak m_s^h$ est un idéal premier faiblement associé à
$\mathcal{F}_s \otimes_{\mathcal{O}_{S, s}} \mathcal{O}_{S, s}^h$, et
\item $\mathfrak m_s^{sh}$ est un idéal premier faiblement associé à
$\mathcal{F}_s \otimes_{\mathcal{O}_{S, s}} \mathcal{O}_{S, s}^{sh}$.
\end{enumerate}
\end{lemma}

\begin{proof}
L'équivalence de (1) et (2) est la définition, voir
Diviseurs, Définition \ref{divisors-definition-weakly-associated}.
Les implications (2) $\Rightarrow$ (3) $\Rightarrow$ (4)
résultent de Diviseurs, Lemme \ref{divisors-lemma-weakly-ass-pullback},
appliqué aux morphismes plats (Compléments d'algèbre, Lemme
\ref{more-algebra-lemma-dumb-properties-henselization})
suivants :
$$
\Spec(\mathcal{O}_{S, s}^{sh}) \to
\Spec(\mathcal{O}_{S, s}^h) \to
\Spec(\mathcal{O}_{S, s})
$$
et aux points fermés. Pour montrer (4) $\Rightarrow$ (2), on peut remplacer
$S$ par un voisinage affine. Supposons que
$x \in \mathcal{F}_s \otimes_{\mathcal{O}_{S, s}} \mathcal{O}_{S, s}^{sh}$
soit un élément dont l'annulateur a pour radical $\mathfrak m_s^{sh}$.
(Voir Algèbre, Lemme \ref{algebra-lemma-weakly-ass-local}.)
Puisque $\mathcal{O}_{S, s}^{sh}$ est la limite
des $\mathcal{O}_{U, u}$ sur les voisinages \'etales
$f : (U, u) \to (S, s)$ d'après Algèbre, Lemme
\ref{algebra-lemma-strict-henselization-different},
on peut supposer que $x$ est l'image d'un élément
$x' \in \mathcal{F}_s \otimes_{\mathcal{O}_{S, s}} \mathcal{O}_{U, u}$.
L'homomorphisme local $\mathcal{O}_{U, u} \to \mathcal{O}_{S, s}^{sh}$
est fidèlement plat (puisqu'il s'agit du hensélisé strict), donc
universellement injectif
(Algèbre, Lemme \ref{algebra-lemma-faithfully-flat-universally-injective}).
Il s'ensuit que l'annulateur de $x'$ est l'image réciproque de
l'annulateur de $x$. Le radical de cet annulateur
est donc égal à $\mathfrak m_u$.
Ainsi $u$ est un point faiblement associé à $f^*\mathcal{F}$.
D'après le Lemme \ref{lemma-etale-weak-assassin-up-down},
on voit que $s$ est un point faiblement associé à $\mathcal{F}$.
\end{proof}





\section{Structure locale d'un module de type fini}
\label{section-local-structure-module}

\noindent
Le lemme technique essentiel sur lequel reposent de nombreux raisonnements de ce
chapitre est le lemme géométrique
\ref{lemma-elementary-devissage}.

\begin{lemma}
\label{lemma-sheaf-lives-on-subscheme}
Soit $f : X \to S$ un morphisme de type fini de schémas affines.
Soit $\mathcal{F}$ un $\mathcal{O}_X$-module quasi-cohérent de type fini.
Soit $x \in X$ d'image $s = f(x)$ dans $S$.
Posons $\mathcal{F}_s = \mathcal{F}|_{X_s}$.
Il existe alors une immersion fermée $i : Z \to X$ de présentation finie
et un $\mathcal{O}_Z$-module quasi-cohérent de type fini $\mathcal{G}$
tels que $i_*\mathcal{G} = \mathcal{F}$ et
$Z_s = \text{Supp}(\mathcal{F}_s)$.
\end{lemma}

\begin{proof}
Supposons que le morphisme $f : X \to S$ soit donné par l'homomorphisme d'anneaux
$A \to B$ et que $\mathcal{F}$ soit le faisceau quasi-cohérent
associé au $B$-module $M$. D'après
Morphismes, Lemme \ref{morphisms-lemma-locally-finite-type-characterize},
on sait que $A \to B$ est un homomorphisme d'anneaux de type fini, et d'après
Propriétés, Lemme \ref{properties-lemma-finite-type-module},
on sait que $M$ est un $B$-module de type fini. En particulier, le
support de $\mathcal{F}$ est le sous-schéma fermé de $\Spec(B)$
défini par l'annulateur
$I = \{x \in B \mid xm = 0\ \forall m \in M\}$ de $M$, voir
Algèbre, Lemme \ref{algebra-lemma-support-closed}.
Soit $\mathfrak q \subset B$ l'idéal premier correspondant à $x$
et soit $\mathfrak p \subset A$ l'idéal premier correspondant à $s$.
Notons que $X_s = \Spec(B \otimes_A \kappa(\mathfrak p))$ et
que $\mathcal{F}_s$ est le faisceau quasi-cohérent associé au
$B \otimes_A \kappa(\mathfrak p)$-module $M \otimes_A \kappa(\mathfrak p)$. D'après
Morphismes, Lemme \ref{morphisms-lemma-support-finite-type},
le support de $\mathcal{F}_s$ est égal à
$V(I(B \otimes_A \kappa(\mathfrak p)))$. Puisque
$B \otimes_A \kappa(\mathfrak p)$ est de type fini sur $\kappa(\mathfrak p)$,
il existe un nombre fini d'éléments $f_1, \ldots, f_m \in I$
tels que
$$
I(B \otimes_A \kappa(\mathfrak p)) =
(f_1, \ldots, f_n)(B \otimes_A \kappa(\mathfrak p)).
$$
Notons $i : Z \to X$ le sous-schéma fermé défini par
$(f_1, \ldots, f_m)$, autrement dit $Z = \Spec(B/(f_1, \ldots, f_m))$.
Puisque $M$ est annulé par $I$, on peut considérer $M$ comme un
$B/(f_1, \ldots, f_m)$-module. En d'autres termes, $\mathcal{F}$ est
l'image directe d'un module de type fini sur $Z$.
Comme $Z_s = \text{Supp}(\mathcal{F}_s)$ par construction, cela
démontre le lemme.
\end{proof}

\begin{lemma}
\label{lemma-elementary-devissage}
Soit $f : X \to S$ un morphisme de schémas localement de type fini.
Soit $\mathcal{F}$ un $\mathcal{O}_X$-module quasi-cohérent de type fini.
Soit $x \in X$ d'image $s = f(x)$ dans $S$.
Posons $\mathcal{F}_s = \mathcal{F}|_{X_s}$ et
$n = \dim_x(\text{Supp}(\mathcal{F}_s))$.
On peut alors construire
\begin{enumerate}
\item des voisinages \'etales élémentaires $g : (X', x') \to (X, x)$,
$e : (S', s') \to (S, s)$,
\item un diagramme commutatif
$$
\xymatrix{
X \ar[dd]_f & X' \ar[dd] \ar[l]^g & Z' \ar[l]^i \ar[d]^\pi \\
& & Y' \ar[d]^h \\
S & S' \ar[l]_e & S' \ar@{=}[l]
}
$$
\item un point $z' \in Z'$ tel que $i(z') = x'$, $y' = \pi(z')$, $h(y') = s'$,
\item un $\mathcal{O}_{Z'}$-module quasi-cohérent de type fini $\mathcal{G}$,
\end{enumerate}
de sorte que les propriétés suivantes soient vérifiées :
\begin{enumerate}
\item $X'$, $Z'$, $Y'$, $S'$ sont des schémas affines,
\item $i$ est une immersion fermée de présentation finie,
\item $i_*(\mathcal{G}) \cong g^*\mathcal{F}$,
\item $\pi$ est fini et $\pi^{-1}(\{y'\}) = \{z'\}$,
\item l'extension $\kappa(y')/\kappa(s')$ est purement transcendante,
\item $h$ est lisse de dimension relative $n$
et à fibres géométriquement intègres.
\end{enumerate}
\end{lemma}

\begin{proof}
Soit $V \subset S$ un voisinage affine de $s$.
Soit $U \subset f^{-1}(V)$ un voisinage affine de $x$.
Il suffit alors de démontrer le lemme pour $f|_U : U \to V$ et
$\mathcal{F}|_U$. Dans la suite de la démonstration, on suppose donc que
$X$ et $S$ sont affines.

\medskip\noindent
Supposons d'abord que $X_s = \text{Supp}(\mathcal{F}_s)$ ; en particulier,
$n = \dim_x(X_s)$. Appliquons
Compléments sur les morphismes,
Lemmes \ref{more-morphisms-lemma-local-local-structure-finite-type} et
\ref{more-morphisms-lemma-local-local-structure-finite-type-affine}.
On obtient un diagramme commutatif
$$
\xymatrix{
X \ar[dd] & X' \ar[l]^g \ar[d]^\pi \\
& Y' \ar[d]^h  \\
S & S' \ar[l]_e
}
$$
et un point $x' \in X'$. On pose $Z' = X'$, $i = \text{id}$ et
$\mathcal{G} = g^*\mathcal{F}$, ce qui donne une solution dans ce cas.

\medskip\noindent
Dans le cas général, choisissons une immersion fermée $Z \to X$ et un faisceau
$\mathcal{G}$ sur $Z$ comme dans
le Lemme \ref{lemma-sheaf-lives-on-subscheme}.
En appliquant le résultat du paragraphe précédent à $Z \to S$ et
$\mathcal{G}$, on obtient un diagramme
$$
\xymatrix{
X \ar[dd]_f & Z \ar[l] \ar[dd]_{f|_Z} & Z' \ar[l]^g \ar[d]^\pi \\
& & Y' \ar[d]^h \\
S & S \ar@{=}[l] & S' \ar[l]_e
}
$$
et un point $z' \in Z'$ vérifiant toutes les propriétés requises.
Nous allons utiliser
le Lemme \ref{lemma-lift-etale}
pour plonger $Z'$ dans un schéma \'etale sur $X$. On ne peut appliquer directement le lemme,
car on veut que $X'$ soit un schéma sur $S'$. On
considère donc les morphismes
$$
\xymatrix{
Z' \ar[r] & Z \times_S S' \ar[r] & X \times_S S'
}
$$
Le premier morphisme est \'etale d'après
Morphismes, Lemme \ref{morphisms-lemma-etale-permanence}.
Le second est une immersion fermée, puisqu'il s'obtient par changement de base d'une immersion fermée.
Enfin, puisque $X$, $S$, $S'$, $Z$, $Z'$ sont tous affines, on peut appliquer
le Lemme \ref{lemma-lift-etale}
pour obtenir un morphisme \'etale de schémas affines $X' \to X \times_S S'$ tel que
$$
Z' = (Z \times_S S') \times_{(X \times_S S')} X' = Z \times_X X'.
$$
Puisque $Z \to X$ est une immersion fermée de présentation finie, il en est de même de $Z' \to X'$.
Soit $x' \in X'$ le point correspondant à $z' \in Z'$.
Le diagramme complété
$$
\xymatrix{
X \ar[dd] & X' \ar[dd] \ar[l] & Z' \ar[l]^i \ar[d]^\pi \\
& & Y' \ar[d]^h \\
S & S' \ar[l]_e & S' \ar@{=}[l]
}
$$
fournit alors une solution du problème initial.
\end{proof}

\begin{lemma}
\label{lemma-devissage-finite-presentation}
Sous les hypothèses et avec les notations du
Lemme \ref{lemma-elementary-devissage},
si $f$ est localement de présentation finie,
alors $\pi$ est de présentation finie.
Dans ce cas, les conditions suivantes sont équivalentes :
\begin{enumerate}
\item $\mathcal{F}$ est un $\mathcal{O}_X$-module de présentation finie
au voisinage de $x$,
\item $\mathcal{G}$ est un $\mathcal{O}_{Z'}$-module de présentation finie
au voisinage de $z'$, et
\item $\pi_*\mathcal{G}$ est un $\mathcal{O}_{Y'}$-module
de présentation finie au voisinage de $y'$.
\end{enumerate}
Toujours sous l'hypothèse que $f$ est localement de présentation finie, les conditions
suivantes sont équivalentes :
\begin{enumerate}
\item[(a)] $\mathcal{F}_x$ est un $\mathcal{O}_{X, x}$-module de présentation
finie,
\item[(b)] $\mathcal{G}_{z'}$ est un $\mathcal{O}_{Z', z'}$-module
de présentation finie, et
\item[(c)] $(\pi_*\mathcal{G})_{y'}$ est un $\mathcal{O}_{Y', y'}$-module
de présentation finie.
\end{enumerate}
\end{lemma}

\begin{proof}
Supposons $f$ localement de présentation finie. Alors $Z' \to S$ est localement
de présentation finie, comme composé de tels morphismes, voir
Morphismes, Lemme \ref{morphisms-lemma-composition-finite-presentation}.
Notons que $Y' \to S$ est aussi localement de présentation finie, comme composé
d'un morphisme lisse et d'un morphisme \'etale. Ainsi
Morphismes, Lemme \ref{morphisms-lemma-finite-presentation-permanence},
entraîne que $\pi$ est localement de présentation finie.
Puisque $\pi$ est fini, il est aussi séparé et
quasi-compact ; par suite, $\pi$ est bien de présentation finie.

\medskip\noindent
Pour démontrer l'équivalence de (1), (2) et (3),\newline considérons aussi :
(4) $g^*\mathcal{F}$ est un $\mathcal{O}_{X'}$-module de présentation finie
au voisinage de $x'$. L'image réciproque d'un module de présentation finie
est de présentation finie, voir
Modules, Lemme \ref{modules-lemma-pullback-finite-presentation}.
Ainsi (1) $\Rightarrow$ (4).
Le morphisme \'etale $g$ est ouvert, voir
Morphismes, Lemme \ref{morphisms-lemma-etale-open}.
Pour tout voisinage ouvert $U' \subset X'$ de $x'$, l'image
$g(U')$ est donc un voisinage ouvert de $x$ et
$\{U' \to g(U')\}$ est un recouvrement \'etale. Ainsi (4) $\Rightarrow$ (1) d'après
Descente, Lemme \ref{descent-lemma-finite-presentation-descends}.
En utilisant
Descente, Lemme \ref{descent-lemma-finite-finitely-presented-module},
et quelques arguments topologiques élémentaires (voir
Compléments sur les morphismes,
Lemme \ref{more-morphisms-lemma-finite-morphism-single-point-in-fibre}),
on voit que
(4) $\Leftrightarrow$ (2) $\Leftrightarrow$ (3).

\medskip\noindent
Pour démontrer l'équivalence de (a), (b), (c), considérons les homomorphismes d'anneaux
$$
\mathcal{O}_{X, x} \to
\mathcal{O}_{X', x'} \to
\mathcal{O}_{Z', z'} \leftarrow
\mathcal{O}_{Y', y'}
$$
Le premier homomorphisme est fidèlement plat. Par suite,
$\mathcal{F}_x$ est de présentation finie sur $\mathcal{O}_{X, x}$
si et seulement si $g^*\mathcal{F}_{x'}$ est de présentation finie sur
$\mathcal{O}_{X', x'}$, voir
Algèbre, Lemme \ref{algebra-lemma-descend-properties-modules}.
Le deuxième homomorphisme est surjectif (donc fini) et
de présentation finie par hypothèse ; par suite,
$g^*\mathcal{F}_{x'}$ est de présentation finie sur $\mathcal{O}_{X', x'}$
si et seulement si $\mathcal{G}_{z'}$ est de présentation finie sur
$\mathcal{O}_{Z', z'}$, voir
Algèbre, Lemme \ref{algebra-lemma-finite-finitely-presented-extension}.
Puisque $\pi$ est fini, de présentation finie et que
$\pi^{-1}(\{y'\}) = \{x'\}$, l'homomorphisme d'anneaux
$\mathcal{O}_{Y', y'} \leftarrow \mathcal{O}_{Z', z'}$ est fini
et de présentation finie, voir
Compléments sur les morphismes,
Lemme \ref{more-morphisms-lemma-finite-morphism-single-point-in-fibre}.
Par suite, $\mathcal{G}_{z'}$ est de présentation finie sur $\mathcal{O}_{Z', z'}$
si et seulement si $\pi_*\mathcal{G}_{y'}$ est de présentation finie sur
$\mathcal{O}_{Y', y'}$, voir
Algèbre, Lemme \ref{algebra-lemma-finite-finitely-presented-extension}.
\end{proof}

\begin{lemma}
\label{lemma-devissage-flat}
Sous les hypothèses et avec les notations du
Lemme \ref{lemma-elementary-devissage},
les conditions suivantes sont équivalentes :
\begin{enumerate}
\item $\mathcal{F}$ est plat sur $S$ au voisinage de $x$,
\item $\mathcal{G}$ est plat sur $S'$ au voisinage de $z'$, et
\item $\pi_*\mathcal{G}$ est plat sur $S'$ au voisinage de $y'$.
\end{enumerate}
Les conditions suivantes sont aussi équivalentes :
\begin{enumerate}
\item[(a)] $\mathcal{F}_x$ est plat sur $\mathcal{O}_{S, s}$,
\item[(b)] $\mathcal{G}_{z'}$ est plat sur $\mathcal{O}_{S', s'}$, et
\item[(c)] $(\pi_*\mathcal{G})_{y'}$ est plat sur $\mathcal{O}_{S', s'}$.
\end{enumerate}
\end{lemma}

\begin{proof}
Pour démontrer l'équivalence de (1), (2) et (3),\newline considérons aussi :
(4) $g^*\mathcal{F}$ est plat sur $S$ au voisinage de $x'$.
Nous utiliserons
le Lemme \ref{lemma-etale-flat-up-down}
pour identifier la platitude sur $S$ et sur $S'$ sans le mentionner de nouveau.
Le morphisme \'etale $g$ est plat et ouvert, voir
Morphismes, Lemme \ref{morphisms-lemma-etale-open}.
Pour tout voisinage ouvert $U' \subset X'$ de $x'$, l'image
$g(U')$ est donc un voisinage ouvert de $x$ et le morphisme
$U' \to g(U')$ est surjectif et plat.
Ainsi (4) $\Leftrightarrow$ (1) d'après
Morphismes, Lemme \ref{morphisms-lemma-flat-permanence}.
Notons que
$$
\Gamma(X', g^*\mathcal{F}) =
\Gamma(Z', \mathcal{G}) =
\Gamma(Y', \pi_*\mathcal{G})
$$
Les conditions de platitude de $g^*\mathcal{F}$, de $\mathcal{G}$ et de $\pi_*\mathcal{G}$
sur $S'$ sont donc équivalentes (on utilise ici le fait que $X'$, $Z'$, $Y'$ et
$S'$ sont tous affines). Des arguments topologiques laissés au lecteur (comparer
Compléments sur les morphismes,
Lemme \ref{more-morphisms-lemma-finite-morphism-single-point-in-fibre})
portant sur les voisinages affines montrent alors que
(4) $\Leftrightarrow$ (2) $\Leftrightarrow$ (3).

\medskip\noindent
Pour démontrer l'équivalence de (a), (b), (c), considérons le diagramme commutatif
d'homomorphismes locaux
$$
\xymatrix{
\mathcal{O}_{X', x'} \ar[r]_\iota &
\mathcal{O}_{Z', z'} &
\mathcal{O}_{Y', y'} \ar[l]^\alpha &
\mathcal{O}_{S', s'} \ar[l]^\beta \\
\mathcal{O}_{X, x} \ar[u]^\gamma & & &
\mathcal{O}_{S, s} \ar[lll]_\varphi \ar[u]_\epsilon
}
$$
Nous utiliserons
le Lemme \ref{lemma-etale-flat-up-down-local-ring}
pour identifier la platitude sur $\mathcal{O}_{S, s}$ et sur $\mathcal{O}_{S', s'}$
sans le mentionner de nouveau.
L'homomorphisme $\gamma$ est fidèlement plat. Par suite,
$\mathcal{F}_x$ est plat sur $\mathcal{O}_{S, s}$
si et seulement si $g^*\mathcal{F}_{x'}$ est plat sur
$\mathcal{O}_{S', s'}$, voir
Algèbre, Lemme \ref{algebra-lemma-flatness-descends-more-general}.
Comme $\mathcal{O}_{S', s'}$-modules, les modules
$g^*\mathcal{F}_{x'}$, $\mathcal{G}_{z'}$ et
$\pi_*\mathcal{G}_{y'}$ sont tous isomorphes, voir
Compléments sur les morphismes,
Lemme \ref{more-morphisms-lemma-finite-morphism-single-point-in-fibre}.
Cela achève la démonstration.
\end{proof}









\section{D\'evissage en une étape}
\label{section-one-step-devissage}

\noindent
Dans cette section, nous expliquons ce qu'est un d\'evissage en une étape d'un
module. Un d\'evissage en une étape existe localement pour la topologie \'etale sur la base et le but.
Nous étudions le changement de base, la restriction à des ouverts de Zariski et la localisation \'etale
d'un d\'evissage en une étape.

\begin{definition}
\label{definition-one-step-devissage}
Soit $S$ un schéma.
Soit $X$ localement de type fini sur $S$.
Soit $\mathcal{F}$ un $\mathcal{O}_X$-module quasi-cohérent de type fini.
Soit $s \in S$ un point.
Un {\it d\'evissage en une étape de $\mathcal{F}/X/S$ au-dessus de $s$}
est donné par des morphismes de schémas sur $S$
$$
\xymatrix{
X & Z \ar[l]_i \ar[r]^\pi & Y
}
$$
et un $\mathcal{O}_Z$-module quasi-cohérent $\mathcal{G}$ de type fini
tels que
\begin{enumerate}
\item $X$, $S$, $Z$ et $Y$ soient affines,
\item $i$ soit une immersion fermée de présentation finie,
\item $\mathcal{F} \cong i_*\mathcal{G}$,
\item $\pi$ soit fini, et
\item le morphisme structural $Y \to S$ soit lisse à
fibres géométriquement irréductibles de
dimension $\dim(\text{Supp}(\mathcal{F}_s))$.
\end{enumerate}
On dit alors que $(Z, Y, i, \pi, \mathcal{G})$ est un
d\'evissage en une étape de $\mathcal{F}/X/S$ au-dessus de $s$.
\end{definition}

\noindent
Un tel d\'evissage en une étape ne peut exister que si $X$ et $S$
sont affines. Dans la définition ci-dessus, on exige seulement que $X$ soit
(localement) de type fini sur $S$, et l'on conserve ce cadre
dans la suite. Dans \cite{GruRay}, les auteurs se placent systématiquement dans le cas
où $X \to S$ est localement de présentation finie et où $\mathcal{F}$
est un $\mathcal{O}_X$-module quasi-cohérent de type fini. Ce choix présente
l'avantage de donner un sens à la condition que $\mathcal{F}$ soit
de présentation finie comme $\mathcal{O}_X$-module, tandis que, dans notre
cadre, cette condition n'a pas de sens intrinsèque. Voir
Compléments sur les morphismes, Section
\ref{more-morphisms-section-finite-type-finite-presentation},
pour une discussion ; les observations qui y sont faites montrent que, dans notre cadre,
on peut considérer la condition que $\mathcal{F}$ soit « localement de présentation
finie relativement à $S$ », et l'on pourrait employer systématiquement cette
notion. Nous préférons cependant nous appuyer sur les résultats du
Lemme \ref{lemma-devissage-finite-presentation}
et les observations de la
Remarque \ref{remark-finite-presentation}
pour traiter cette question au cas par cas lorsqu'elle se présente.

\begin{definition}
\label{definition-one-step-devissage-at-x}
Soit $S$ un schéma.
Soit $X$ localement de type fini sur $S$.
Soit $\mathcal{F}$ un $\mathcal{O}_X$-module quasi-cohérent de type fini.
Soit $x \in X$ un point d'image $s$ dans $S$.
Un {\it d\'evissage en une étape de $\mathcal{F}/X/S$ en $x$}
est un système $(Z, Y, i, \pi, \mathcal{G}, z, y)$, où
$(Z, Y, i, \pi, \mathcal{G})$ est un d\'evissage en une étape de
$\mathcal{F}/X/S$ au-dessus de $s$ et où
\begin{enumerate}
\item $\dim_x(\text{Supp}(\mathcal{F}_s)) = \dim(\text{Supp}(\mathcal{F}_s))$,
\item $z \in Z$ est un point tel que $i(z) = x$ et $\pi(z) = y$,
\item on a $\pi^{-1}(\{y\}) = \{z\}$,
\item l'extension $\kappa(y)/\kappa(s)$ est purement
transcendante.
\end{enumerate}
\end{definition}

\noindent
Un d\'evissage en une étape de $\mathcal{F}/X/S$ en $x$ ne peut exister que si
$X$ et $S$ sont affines. La condition (1) assure que $Y \to S$ est de
dimension relative égale à $\dim_x(\text{Supp}(\mathcal{F}_s))$,
en vertu de la condition (5) de la
Définition \ref{definition-one-step-devissage}.

\begin{lemma}
\label{lemma-elementary-devissage-variant}
Soit $f : X \to S$ un morphisme de schémas localement de type fini.
Soit $\mathcal{F}$ un $\mathcal{O}_X$-module quasi-cohérent de type fini.
Soit $x \in X$ d'image $s = f(x)$ dans $S$.
Il existe alors un diagramme commutatif de schémas pointés
$$
\xymatrix{
(X, x) \ar[d]_f & (X', x') \ar[l]^g \ar[d] \\
(S, s) & (S', s') \ar[l] \\
}
$$
tel que $(S', s') \to (S, s)$ et $(X', x') \to (X, x)$
soient des voisinages \'etales élémentaires et tel que
$g^*\mathcal{F}/X'/S'$ admette un d\'evissage en une étape en $x'$.
\end{lemma}

\begin{proof}
Cela résulte immédiatement de la
Définition \ref{definition-one-step-devissage-at-x}
et du
Lemme \ref{lemma-elementary-devissage}.
\end{proof}

\begin{lemma}
\label{lemma-base-change-one-step}
Soient $S$, $X$, $\mathcal{F}$, $s$ comme dans la
Définition \ref{definition-one-step-devissage}.
Soit $(Z, Y, i, \pi, \mathcal{G})$ un d\'evissage en une étape
de $\mathcal{F}/X/S$ au-dessus de $s$.
Soit $(S', s') \to (S, s)$ un morphisme quelconque de schémas pointés.\newline
Ces données étant fixées, soient $X', Z', Y', i', \pi'$ les objets obtenus par changement
de base à partir de $X, Z, Y, i, \pi$ suivant $S' \to S$.
Soit $\mathcal{F}'$ l'image réciproque de $\mathcal{F}$ sur $X'$
et soit $\mathcal{G}'$ l'image réciproque de $\mathcal{G}$ sur $Z'$.
Si $S'$ est affine, alors $(Z', Y', i', \pi', \mathcal{G}')$
est un d\'evissage en une étape de $\mathcal{F}'/X'/S'$ au-dessus de $s'$.
\end{lemma}

\begin{proof}
Les produits fibrés de schémas affines sont affines, voir
Schémas, Lemme \ref{schemes-lemma-fibre-product-affines}.
Le changement de base préserve
les immersions fermées,
les morphismes de présentation finie,
les morphismes finis,
les morphismes lisses,
les morphismes à fibres géométriquement irréductibles et
les morphismes de dimension relative $n$, voir
Morphismes, Lemmes \ref{morphisms-lemma-base-change-closed-immersion},
\ref{morphisms-lemma-base-change-finite-presentation},
\ref{morphisms-lemma-base-change-finite},
\ref{morphisms-lemma-base-change-smooth},
\ref{morphisms-lemma-base-change-relative-dimension-d}, et
Compléments sur les morphismes, Lemme
\ref{more-morphisms-lemma-base-change-fibres-geometrically-irreducible}.
On a $i'_*\mathcal{G}' \cong \mathcal{F}'$, car l'image directe
par le morphisme fini $i$ commute au changement de base, voir
Cohomologie des schémas, Lemme \ref{coherent-lemma-affine-base-change}.
On a
$\dim(\text{Supp}(\mathcal{F}_s)) = \dim(\text{Supp}(\mathcal{F}'_{s'}))$
d'après
Morphismes, Lemme \ref{morphisms-lemma-dimension-fibre-after-base-change},
puisque
$$
\text{Supp}(\mathcal{F}_s) \times_s s' = \text{Supp}(\mathcal{F}'_{s'}).
$$
Cela démontre le lemme.
\end{proof}

\begin{lemma}
\label{lemma-base-change-one-step-at-x}
Soient $S$, $X$, $\mathcal{F}$, $x$, $s$ comme dans la
Définition \ref{definition-one-step-devissage-at-x}.\newline
Soit $(Z, Y, i, \pi, \mathcal{G}, z, y)$ un d\'evissage en une étape
de $\mathcal{F}/X/S$ en $x$.
Soit $(S', s') \to (S, s)$ un morphisme de schémas pointés
induisant un isomorphisme $\kappa(s) = \kappa(s')$.
Soit $(Z', Y', i', \pi', \mathcal{G}')$ le système construit dans le
Lemme \ref{lemma-base-change-one-step},
et soit $x' \in X'$ (resp.\ $z' \in Z'$, $y' \in Y'$) l'unique
point s'envoyant à la fois sur $x \in X$ (resp.\ $z \in Z$, $y \in Y$)
et sur $s' \in S'$.
Si $S'$ est affine, alors $(Z', Y', i', \pi', \mathcal{G}', z', y')$
est un d\'evissage en une étape de $\mathcal{F}'/X'/S'$ en $x'$.
\end{lemma}

\begin{proof}
D'après le
Lemme \ref{lemma-base-change-one-step},
$(Z', Y', i', \pi', \mathcal{G}')$ est un d\'evissage en une étape de
$\mathcal{F}'/X'/S'$ au-dessus de $s'$. Les propriétés (1) -- (4) de la
Définition \ref{definition-one-step-devissage-at-x}
sont vérifiées pour $(Z', Y', i', \pi', \mathcal{G}', z', y')$,
car l'hypothèse $\kappa(s) = \kappa(s')$ assure que les fibres
$X'_{s'}$, $Z'_{s'}$ et $Y'_{s'}$ sont respectivement isomorphes à
$X_s$, $Z_s$ et $Y_s$.
\end{proof}

\begin{definition}
\label{definition-shrink}
Soient $S$, $X$, $\mathcal{F}$, $x$, $s$ comme dans la
Définition \ref{definition-one-step-devissage-at-x}.\newline
Soit $(Z, Y, i, \pi, \mathcal{G}, z, y)$ un d\'evissage en une étape
de $\mathcal{F}/X/S$ en $x$. Appelons
{\it restriction standard} de cette situation la donnée
d'ouverts standards $S' \subset S$, $X' \subset X$, $Z' \subset Z$
et $Y' \subset Y$ tels que $s \in S'$, $x \in X'$, $z \in Z'$ et
$y \in Y'$, et tels que
$$
(Z', Y', i|_{Z'}, \pi|_{Z'}, \mathcal{G}|_{Z'}, z, y)
$$
soit un d\'evissage en une étape de $\mathcal{F}|_{X'}/X'/S'$ en $x$.
\end{definition}

\begin{lemma}
\label{lemma-shrink}
Sous les hypothèses et avec les notations de la
Définition \ref{definition-shrink},
on a :
\begin{enumerate}
\item
\label{item-shrink-base}
Si $S' \subset S$ est un voisinage ouvert standard de $s$, alors,
en posant $X' = X_{S'}$, $Z' = Z_{S'}$ et $Y' = Y_{S'}$, on obtient une
restriction standard.
\item
\label{item-shrink-on-Y}
Soit $W \subset Y$ un voisinage ouvert standard de $y$.
Il existe alors une restriction standard telle que $Y' = W \times_S S'$.
\item
\label{item-shrink-on-X}
Soit $U \subset X$ un voisinage ouvert de $x$.
Il existe alors une restriction standard telle que $X' \subset U$.
\end{enumerate}
\end{lemma}

\begin{proof}
L'assertion (1) résulte immédiatement du
Lemme \ref{lemma-base-change-one-step-at-x}
et du fait que l'image réciproque d'un ouvert standard par un morphisme
de schémas affines est un ouvert standard, voir
Algèbre, Lemme \ref{algebra-lemma-spec-functorial}.

\medskip\noindent
Soit $W \subset Y$ comme dans (2). Puisque $Y \to S$ est lisse, il est ouvert, voir
Morphismes, Lemme \ref{morphisms-lemma-smooth-open}.
On peut donc trouver un voisinage ouvert standard $S'$ de $s$
contenu dans l'image de $W$. Les fibres de $W_{S'} \to S'$
sont alors des sous-schémas ouverts non vides des fibres de $Y \to S$ au-dessus de $S'$
et sont donc elles aussi géométriquement irréductibles. En posant $Y' = W_{S'}$
et $Z' = \pi^{-1}(Y')$, on voit que $Z' \subset Z$ est un voisinage ouvert standard
de $z$. Soit $\overline{h} \in \Gamma(Z, \mathcal{O}_Z)$
une fonction telle que $Z' = D(\overline{h})$. Puisque $i : Z \to X$
est une immersion fermée, on peut trouver une fonction $h \in \Gamma(X, \mathcal{O}_X)$
telle que $i^\sharp(h) = \overline{h}$. Prenons $X' = D(h) \subset X$.
On obtient ainsi une restriction standard comme dans (2).

\medskip\noindent
Soit $U \subset X$ comme dans (3). Quitte à restreindre $U$, on peut supposer que
$U$ est un ouvert standard. D'après
Compléments sur les morphismes,
Lemme \ref{more-morphisms-lemma-finite-morphism-single-point-in-fibre},
il existe un ouvert standard $W \subset Y$, voisinage de $y$, tel
que $\pi^{-1}(W) \subset i^{-1}(U)$. Appliquons (2) pour obtenir une restriction
standard $X', S', Z', Y'$ avec $Y' = W_{S'}$. Puisque
$Z' \subset \pi^{-1}(W) \subset i^{-1}(U)$, on peut remplacer $X'$ par
$X' \cap U$ (qui est encore un ouvert standard, puisque $U$ l'est aussi)
sans modifier aucune des conditions définissant une restriction standard.
On obtient ainsi le résultat voulu.
\end{proof}

\begin{lemma}
\label{lemma-elementary-etale-neighbourhood}
Soient $S$, $X$, $\mathcal{F}$, $x$, $s$ comme dans la
Définition \ref{definition-one-step-devissage-at-x}.\newline
Soit $(Z, Y, i, \pi, \mathcal{G}, z, y)$ un d\'evissage en une étape
de $\mathcal{F}/X/S$ en $x$. Soit
$$
\xymatrix{
(Y, y) \ar[d] & (Y', y') \ar[l] \ar[d] \\
(S, s) & (S', s') \ar[l]
}
$$
un diagramme commutatif de schémas pointés dont les flèches horizontales
sont des voisinages \'etales élémentaires. Il existe alors
un diagramme commutatif
$$
\xymatrix{
& & (X'', x'') \ar[lld] \ar[d] & (Z'', z'') \ar[l] \ar[lld] \ar[d] \\
(X, x) \ar[d] & (Z, z) \ar[l] \ar[d] &
(S'', s'') \ar[lld] & (Y'', y'') \ar[lld] \ar[l] \\
(S, s) & (Y, y) \ar[l]
}
$$
de schémas pointés vérifiant les propriétés suivantes :
\begin{enumerate}
\item $(S'', s'') \to (S', s')$ est un voisinage \'etale élémentaire et
le morphisme $S'' \to S$ est le composé $S'' \to S' \to S$,
\item $Y''$ est un sous-schéma ouvert de $Y' \times_{S'} S''$,
\item $Z'' = Z \times_Y Y''$,
\item $(X'', x'') \to (X, x)$ est un voisinage \'etale élémentaire, et
\item $(Z'', Y'', i'', \pi'', \mathcal{G}'', z'', y'')$ est un
d\'evissage en une étape en $x''$ du faisceau $\mathcal{F}''$.
\end{enumerate}
Ici $\mathcal{F}''$ (resp.\ $\mathcal{G}''$) est l'image réciproque de
$\mathcal{F}$ (resp.\ $\mathcal{G}$) par le morphisme $X'' \to X$
(resp.\ $Z'' \to Z$), et $i'' : Z'' \to X''$ et $\pi'' : Z'' \to Y''$
sont les morphismes du diagramme.
\end{lemma}

\begin{proof}
Soit $(S'', s'') \to (S', s')$ un voisinage \'etale élémentaire quelconque
avec $S''$ affine. Soit $Y'' \subset Y' \times_{S'} S''$ un voisinage
ouvert affine quelconque du point $y'' = (y', s'')$. On
obtient alors un schéma affine $(Z'', z'')$ par (3). De plus, $Z_{S''} \to X_{S''}$
est une immersion fermée et $Z'' \to Z_{S''}$ est un morphisme \'etale.
Le
Lemme \ref{lemma-lift-etale}
s'applique donc, et l'on peut trouver un morphisme \'etale $X'' \to X_{S'}$ de schémas affines
tel que $Z'' \cong X'' \times_{X_{S'}} Z_{S'}$. Notons $i'' : Z'' \to X''$
l'immersion fermée correspondante. En posant $x'' = i''(z'')$, on obtient un
diagramme commutatif comme dans le lemme.
Les propriétés (1), (2), (3) et (4) sont vérifiées par construction.
Il suffit donc de montrer que (5) est vérifiée pour un choix convenable de
$(S'', s'') \to (S', s')$ et de $Y''$.

\medskip\noindent
Énumérons d'abord les propriétés vérifiées pour tout choix de
$(S'', s'') \to (S', s')$ et de $Y''$ comme dans le premier paragraphe.
Puisqu'on a $Z'' = X'' \times_X Z$ par construction, on voit que
$i''_*\mathcal{G}'' = \mathcal{F}''$ (avec les notations de
l'énoncé du lemme), voir
Cohomologie des schémas, Lemme \ref{coherent-lemma-affine-base-change}.
Posons $n = \dim(\text{Supp}(\mathcal{F}_s)) = \dim_x(\text{Supp}(\mathcal{F}_s))$.
Le morphisme $Y'' \to S''$ est lisse de dimension relative $n$
(car $Y' \to S'$ est lisse de dimension relative $n$,
comme composé $Y' \to Y_{S'} \to S'$ d'un morphisme \'etale et
d'un morphisme lisse de dimension relative $n$, et parce que le changement de base
préserve les morphismes lisses de dimension relative $n$).
On a $\kappa(y'') = \kappa(y)$ et $\kappa(s) = \kappa(s'')$ ;
ainsi $\kappa(y'')$ est une extension purement transcendante de $\kappa(s'')$.
Le morphisme de fibres $X''_{s''} \to X_s$ est un morphisme \'etale de schémas
affines sur $\kappa(s) = \kappa(s'')$, envoyant le point $x''$ sur le
point $x$ et dont l'image réciproque de $\mathcal{F}_s$ est $\mathcal{F}''_{s''}$.
Par suite,
$$
\dim(\text{Supp}(\mathcal{F}''_{s''})) =
\dim(\text{Supp}(\mathcal{F}_s)) = n =
\dim_x(\text{Supp}(\mathcal{F}_s)) =
\dim_{x''}(\text{Supp}(\mathcal{F}''_{s''}))
$$
car la dimension est invariante par localisation \'etale, voir
Descente, Lemme \ref{descent-lemma-dimension-at-point-local}.
Puisque $\pi'' : Z'' \to Y''$ est obtenu par changement de base à partir de $\pi$, on voit que
$\pi''$ est fini et, puisque $\kappa(y) = \kappa(y'')$, on voit que
$\pi^{-1}(\{y''\}) = \{z''\}$.

\medskip\noindent
À ce stade, on a vérifié toutes les conditions de la
Définition \ref{definition-one-step-devissage},
sauf que l'on n'a pas vérifié que $Y'' \to S''$ est à fibres géométriquement
irréductibles. Cela n'est évidemment pas vrai en général,
et c'est ici que l'on utilise le fait que
$\kappa(s) \subset \kappa(y)$ est purement transcendante. En effet,
soit $T \subset Y'_{s'}$ la composante irréductible de
$Y'_{s'}$ contenant $y' = (y, s')$. Notons que $T$ est un sous-schéma ouvert
de $Y'_{s'}$, car ce dernier est un schéma lisse sur $\kappa(s')$. D'après
Variétés,
Lemme \ref{varieties-lemma-geometrically-connected-if-connected-and-point},
on voit que $T$ est géométriquement connexe, puisque $\kappa(s') = \kappa(s)$
est algébriquement fermé dans $\kappa(y') = \kappa(y)$.
Puisque $T$ est lisse, on voit que $T$ est géométriquement irréductible. Le
Lemme de Compléments sur les morphismes,
\ref{more-morphisms-lemma-normal-morphism-irreducible},
s'applique donc, et l'on peut trouver un morphisme \'etale élémentaire
$(S'', s'') \to (S', s')$ et un ouvert affine $Y'' \subset Y'_{S''}$
tels que toutes les fibres de $Y'' \to S''$ soient géométriquement irréductibles
et que $T = Y''_{s''}$. Quitte à restreindre d'abord $Y''$, puis $S''$,
on peut supposer que $Y''$ et $S''$ sont tous deux affines.
Cela achève la démonstration du lemme.
\end{proof}

\begin{lemma}
\label{lemma-existence-alpha}
Soient $S$, $X$, $\mathcal{F}$, $s$ comme dans la
Définition \ref{definition-one-step-devissage}.
Soit $(Z, Y, i, \pi, \mathcal{G})$ un d\'evissage en une étape
de $\mathcal{F}/X/S$ au-dessus de $s$.
Soit $\xi \in Y_s$ l'unique point générique.
Il existe alors un entier $r > 0$ et un homomorphisme de $\mathcal{O}_Y$-modules
$$
\alpha : \mathcal{O}_Y^{\oplus r} \longrightarrow \pi_*\mathcal{G}
$$
tels que
$$
\alpha :
\kappa(\xi)^{\oplus r}
\longrightarrow
(\pi_*\mathcal{G})_\xi \otimes_{\mathcal{O}_{Y, \xi}} \kappa(\xi)
$$
soit un isomorphisme. De plus, on a alors
$$
\dim(\text{Supp}(\Coker(\alpha)_s)) < \dim(\text{Supp}(\mathcal{F}_s)).
$$
\end{lemma}

\begin{proof}
Par hypothèse, les schémas $S$ et $Y$ sont affines.
Écrivons $S = \Spec(A)$ et $Y = \Spec(B)$.
Puisque $\pi$ est fini, le $\mathcal{O}_Y$-module $\pi_*\mathcal{G}$
est un $\mathcal{O}_Y$-module quasi-cohérent de type fini.
On a donc $\pi_*\mathcal{G} = \widetilde{N}$ pour un certain $B$-module de type fini $N$.
Soit $\mathfrak p \subset B$ l'idéal premier correspondant à $\xi$.
Pour obtenir $\alpha$, posons
$r = \dim_{\kappa(\mathfrak p)} N \otimes_B \kappa(\mathfrak p)$
et choisissons $x_1, \ldots, x_r \in N$ dont les images forment une base de
$N \otimes_B \kappa(\mathfrak p)$. Prenons pour $\alpha : B^{\oplus r} \to N$
l'homomorphisme donné par la formule $\alpha(b_1, \ldots, b_r) = \sum b_ix_i$.
Il est clair que
$\alpha : \kappa(\mathfrak p)^{\oplus r} \to N \otimes_B \kappa(\mathfrak p)$
est un isomorphisme, comme souhaité. Enfin, soit $\alpha$ un homomorphisme quelconque ayant cette
propriété. Alors $N' = \Coker(\alpha)$ est un $B$-module de type fini
tel que $N' \otimes \kappa(\mathfrak p) = 0$. Par le lemme de Nakayama
(Algèbre, Lemme \ref{algebra-lemma-NAK}),
on voit que $N'_{\mathfrak p} = 0$. Puisque la fibre $Y_s$ est
géométriquement irréductible de dimension $n$, de point générique $\xi$,
et que l'on vient de voir que $\xi$ n'appartient pas au support de
$\Coker(\alpha)$, la dernière assertion du lemme est vérifiée.
\end{proof}


\section{D\'evissage complet}
\label{section-complete-devissage}

\noindent
Dans cette section, nous expliquons ce qu'est un d\'evissage complet d'un
module et démontrons l'existence de tels dévissages. Le contenu de cette
section consiste surtout à organiser les constructions précédentes.

\begin{definition}
\label{definition-complete-devissage}
Soit $S$ un schéma.
Soit $X$ localement de type fini sur $S$.
Soit $\mathcal{F}$ un $\mathcal{O}_X$-module quasi-cohérent de type fini.
Soit $s \in S$ un point.
Un {\it d\'evissage complet de $\mathcal{F}/X/S$ au-dessus de $s$} est donné par un
diagramme
$$
\xymatrix{
X & Z_1 \ar[l]^{i_1} \ar[d]^{\pi_1} \\
& Y_1 & Z_2 \ar[l]^{i_2} \ar[d]^{\pi_2} \\
& & Y_2 & Z_3 \ar[l] \ar[d] \\
& & & ... & ... \ar[l] \ar[d] \\
& & & & Y_n
}
$$
de schémas sur $S$, des $\mathcal{O}_{Z_k}$-modules quasi-cohérents de type fini
$\mathcal{G}_k$ et des homomorphismes de $\mathcal{O}_{Y_k}$-modules
$$
\alpha_k :
\mathcal{O}_{Y_k}^{\oplus r_k}
\longrightarrow
\pi_{k, *}\mathcal{G}_k,
\quad
k = 1, \ldots, n
$$
vérifiant les propriétés suivantes :
\begin{enumerate}
\item $(Z_1, Y_1, i_1, \pi_1, \mathcal{G}_1)$ est un
d\'evissage en une étape de $\mathcal{F}/X/S$ au-dessus de $s$,
\item l'homomorphisme $\alpha_k$ induit un isomorphisme
$$
\kappa(\xi_k)^{\oplus r_k} \longrightarrow
(\pi_{k, *}\mathcal{G}_k)_{\xi_k}
\otimes_{\mathcal{O}_{Y_k, \xi_k}} \kappa(\xi_k)
$$
où $\xi_k \in (Y_k)_s$ est l'unique point générique,
\item pour $k = 2, \ldots, n$, le système
$(Z_k, Y_k, i_k, \pi_k, \mathcal{G}_k)$
est un d\'evissage en une étape de $\Coker(\alpha_{k - 1})/Y_{k - 1}/S$
au-dessus de $s$,
\item $\Coker(\alpha_n) = 0$.
\end{enumerate}
On dit alors que
$(Z_k, Y_k, i_k, \pi_k, \mathcal{G}_k, \alpha_k)_{k = 1, \ldots, n}$
est un d\'evissage complet de $\mathcal{F}/X/S$ au-dessus de $s$.
\end{definition}

\begin{definition}
\label{definition-complete-devissage-at-x}
Soit $S$ un schéma.
Soit $X$ localement de type fini sur $S$.
Soit $\mathcal{F}$ un $\mathcal{O}_X$-module quasi-cohérent de type fini.
Soit $x \in X$ un point d'image $s \in S$.
Un {\it d\'evissage complet de $\mathcal{F}/X/S$ en $x$} est donné par un
système
$$
(Z_k, Y_k, i_k, \pi_k, \mathcal{G}_k, \alpha_k, z_k, y_k)_{k = 1, \ldots, n}
$$
tel que $(Z_k, Y_k, i_k, \pi_k, \mathcal{G}_k, \alpha_k)$ soit un
d\'evissage complet de $\mathcal{F}/X/S$ au-dessus de $s$, et tel que
\begin{enumerate}
\item $(Z_1, Y_1, i_1, \pi_1, \mathcal{G}_1, z_1, y_1)$ soit un
d\'evissage en une étape de $\mathcal{F}/X/S$ en $x$,
\item pour $k = 2, \ldots, n$, le système
$(Z_k, Y_k, i_k, \pi_k, \mathcal{G}_k, z_k, y_k)$
soit un d\'evissage en une étape de $\Coker(\alpha_{k - 1})/Y_{k - 1}/S$
en $y_{k - 1}$.
\end{enumerate}
\end{definition}

\noindent
Remarquons encore qu'un d\'evissage complet ne peut exister que si $X$ et
$S$ sont affines.

\begin{lemma}
\label{lemma-base-change-complete}
Soient $S$, $X$, $\mathcal{F}$, $s$ comme dans la
Définition \ref{definition-complete-devissage}.
Soit $(S', s') \to (S, s)$ un morphisme quelconque de schémas pointés.
Soit $(Z_k, Y_k, i_k, \pi_k, \mathcal{G}_k, \alpha_k)_{k = 1, \ldots, n}$
un d\'evissage complet de $\mathcal{F}/X/S$ au-dessus de $s$.
Ces données étant fixées,\newline soient $X', Z'_k, Y'_k, i'_k, \pi'_k$ les objets obtenus par changement
de base à partir de\newline $X, Z_k, Y_k, i_k, \pi_k$ suivant $S' \to S$.
Soit $\mathcal{F}'$ l'image réciproque de $\mathcal{F}$ sur $X'$
et soit $\mathcal{G}'_k$ l'image réciproque de $\mathcal{G}_k$ sur $Z'_k$.
Soit $\alpha'_k$ l'image réciproque de $\alpha_k$ sur $Y'_k$.
Si $S'$ est affine, alors
$(Z'_k, Y'_k, i'_k, \pi'_k, \mathcal{G}'_k, \alpha'_k)_{k = 1, \ldots, n}$
est un d\'evissage complet de $\mathcal{F}'/X'/S'$ au-dessus de $s'$.
\end{lemma}

\begin{proof}
D'après le
Lemme \ref{lemma-base-change-one-step},
on sait que le changement de base d'un d\'evissage en une étape est un
d\'evissage en une étape. Il suffit donc de montrer que la formation de
$\Coker(\alpha_k)$ commute au changement de base et que la
condition (2) de la
Définition \ref{definition-complete-devissage}
est préservée par changement de base. La première assertion résulte du fait que
$\pi'_{k, *}\mathcal{G}'_k$ est l'image réciproque de
$\pi_{k, *}\mathcal{G}_k$ (d'après
Cohomologie des schémas, Lemme \ref{coherent-lemma-affine-base-change})
et de l'exactitude à droite de $\otimes$. Pour la seconde,
le même raisonnement donne
$$
(\pi_{k, *}\mathcal{G}_k)_{\xi_k}
\otimes_{\mathcal{O}_{Y_k, \xi_k}} \kappa(\xi_k)
\otimes_{\kappa(\xi_k)} \kappa(\xi'_k)
\cong
(\pi'_{k, *}\mathcal{G}'_k)_{\xi'_k}
\otimes_{\mathcal{O}_{Y'_k, \xi'_k}} \kappa(\xi'_k)
$$
avec des notations évidentes.
\end{proof}

\begin{lemma}
\label{lemma-base-change-complete-at-x}
Soient $S$, $X$, $\mathcal{F}$, $x$, $s$ comme dans la
Définition \ref{definition-complete-devissage-at-x}.
Soit $(S', s') \to (S, s)$ un morphisme de schémas pointés
induisant un isomorphisme $\kappa(s) = \kappa(s')$. Soit
$(Z_k, Y_k, i_k, \pi_k, \mathcal{G}_k, \alpha_k, z_k, y_k)_{k = 1, \ldots, n}$
un d\'evissage complet de $\mathcal{F}/X/S$ en $x$.
Soit
$(Z'_k, Y'_k, i'_k, \pi'_k, \mathcal{G}'_k, \alpha'_k)_{k = 1, \ldots, n}$
le système construit dans le
Lemme \ref{lemma-base-change-complete},
et soit $x' \in X'$ (resp.\ $z'_k \in Z'$, $y'_k \in Y'$) l'unique
point s'envoyant à la fois sur $x \in X$ (resp.\ $z_k \in Z_k$, $y_k \in Y_k$)
et sur $s' \in S'$.\newline
Si $S'$ est affine, alors
$(Z'_k, Y'_k, i'_k, \pi'_k, \mathcal{G}'_k, \alpha'_k,
z'_k, y'_k)_{k = 1, \ldots, n}$
est un d\'evissage complet de $\mathcal{F}'/X'/S'$ en $x'$.
\end{lemma}

\begin{proof}
On combine
le Lemme \ref{lemma-base-change-complete}
et
le Lemme \ref{lemma-base-change-one-step-at-x}.
\end{proof}

\begin{definition}
\label{definition-shrink-complete}
Soient $S$, $X$, $\mathcal{F}$, $x$, $s$ comme dans la
Définition \ref{definition-complete-devissage-at-x}.
Considérons un d\'evissage complet
$(Z_k, Y_k, i_k, \pi_k, \mathcal{G}_k, \alpha_k, z_k, y_k)_{k = 1, \ldots, n}$
de $\mathcal{F}/X/S$ en $x$. Appelons
{\it restriction standard} de cette situation la donnée
d'ouverts standards $S' \subset S$, $X' \subset X$,
$Z'_k \subset Z_k$ et $Y'_k \subset Y_k$ tels que $s_k \in S'$,
$x_k \in X'$, $z_k \in Z'$ et $y_k \in Y'$, et tels que
$$
(Z'_k, Y'_k, i'_k, \pi'_k,
\mathcal{G}'_k, \alpha'_k, z_k, y_k)_{k = 1, \ldots, n}
$$
soit un d\'evissage en une étape de $\mathcal{F}'/X'/S'$ en $x$, où
$\mathcal{G}'_k = \mathcal{G}_k|_{Z'_k}$ et
$\mathcal{F}' = \mathcal{F}|_{X'}$.
\end{definition}

\begin{lemma}
\label{lemma-shrink-complete}
Sous les hypothèses et avec les notations de la
Définition \ref{definition-shrink-complete},
on a :
\begin{enumerate}
\item
\label{item-shrink-base-complete}
Si $S' \subset S$ est un voisinage ouvert standard de $s$, alors,
en posant $X' = X_{S'}$, $Z'_k = Z_{S'}$ et $Y'_k = Y_{S'}$, on obtient une
restriction standard.
\item
\label{item-shrink-on-Y-complete}
Soit $W \subset Y_n$ un voisinage ouvert standard de $y$.
Il existe alors une restriction standard telle que $Y'_n = W \times_S S'$.
\item
\label{item-shrink-on-X-complete}
Soit $U \subset X$ un voisinage ouvert de $x$.
Il existe alors une restriction standard telle que $X' \subset U$.
\end{enumerate}
\end{lemma}

\begin{proof}
L'assertion (1) résulte immédiatement des
Lemmes \ref{lemma-base-change-complete-at-x} et
\ref{lemma-shrink}.

\medskip\noindent
Démonstration de (2). Pour simplifier les notations, posons $X = Y_0$. On applique
le Lemme \ref{lemma-shrink} (\ref{item-shrink-on-Y})
pour trouver une restriction standard
$S', Y'_{n - 1}, Z'_n, Y'_n$
du d\'evissage en une étape de $\Coker(\alpha_{n - 1})/Y_{n - 1}/S$
en $y_{n - 1}$ avec $Y'_n = W \times_S S'$. On peut répéter cette opération
et trouver une restriction standard
$S'', Y''_{n - 2}, Z''_{n - 1}, Y''_{n - 1}$
du d\'evissage en une étape de $\Coker(\alpha_{n - 2})/Y_{n - 2}/S$
en $y_{n - 2}$ avec $Y''_{n - 1} = Y'_{n - 1} \times_S S''$.
On continue ainsi jusqu'à obtenir
$S^{(n)}, Y^{(n)}_0, Z^{(n)}_1, Y^{(n)}_1$.
Il est alors clair que l'on obtient la restriction standard cherchée
en prenant $S^{(n)}$, $X^{(n)}$, $Z_k^{(n - k)} \times_S S^{(n)}$ et
$Y_k^{(n - k)} \times_S S^{(n)}$, avec la propriété voulue.

\medskip\noindent
Démonstration de (3). On raisonne par récurrence sur la longueur du
d\'evissage complet. On applique d'abord
le Lemme \ref{lemma-shrink} (\ref{item-shrink-on-X})
pour trouver une restriction standard
$S', X', Z'_1, Y'_1$
du d\'evissage en une étape de $\mathcal{F}/X/S$ en $x$
avec $X' \subset U$. Si $n = 1$, la démonstration est terminée.
Si $n > 1$, l'hypothèse de récurrence fournit une restriction standard
$S''$, $Y''_1$, $Z''_k$ et $Y''_k$ du d\'evissage complet\newline
$(Z_k, Y_k, i_k, \pi_k, \mathcal{G}_k, \alpha_k, z_k, y_k)_{k = 2, \ldots, n}$
de $\Coker(\alpha_1)/Y_1/S$ en $x$, telle que
$Y''_1 \subset Y'_1$. En utilisant
le Lemme \ref{lemma-shrink} (\ref{item-shrink-on-Y}),
on peut trouver $S''' \subset S'$, $X''' \subset X'$, $Z'''_1$ et
$Y'''_1 = Y''_1 \times_S S'''$ qui constituent une restriction standard.
La solution du problème consiste à prendre
$$
S''', X''', Z'''_1, Y'''_1, Z''_2 \times_S S''',
Y''_2 \times_S S''', \ldots, Z''_n \times_S S''', Y''_n \times_S S'''
$$
Cela achève la démonstration du lemme.
\end{proof}

\begin{proposition}
\label{proposition-existence-complete-at-x}
Soit $S$ un schéma.
Soit $X$ localement de type fini sur $S$.
Soit $x \in X$ un point d'image $s \in S$.
Il existe un diagramme commutatif
$$
\xymatrix{
(X, x) \ar[d] & (X', x') \ar[l]^g \ar[d] \\
(S, s) & (S', s') \ar[l]
}
$$
de schémas pointés dont les flèches
horizontales sont des voisinages \'etales élémentaires
et tel que $g^*\mathcal{F}/X'/S'$ admette un
d\'evissage complet en $x$.
\end{proposition}

\begin{proof}
On démontre cette assertion par récurrence sur l'entier
$d = \dim_x(\text{Supp}(\mathcal{F}_s))$.
D'après le
Lemme \ref{lemma-elementary-devissage-variant},
il existe un diagramme
$$
\xymatrix{
(X, x) \ar[d] & (X', x') \ar[l]^g \ar[d] \\
(S, s) & (S', s') \ar[l]
}
$$
de schémas pointés dont les flèches
horizontales sont des voisinages \'etales élémentaires
et tel que $g^*\mathcal{F}/X'/S'$ admette un d\'evissage en une étape en $x'$.
Le caractère local du problème permet de remplacer
$(X, x) \to (S, s)$ par $(X', x') \to (S', s')$. Après ce remplacement,
on peut donc supposer qu'il existe un d\'evissage en une étape
$(Z_1, Y_1, i_1, \pi_1, \mathcal{G}_1)$ de $\mathcal{F}/X/S$ en $x$.

\medskip\noindent
On applique
le Lemme \ref{lemma-existence-alpha}
pour trouver un homomorphisme
$$
\alpha_1 :
\mathcal{O}_{Y_1}^{\oplus r_1}
\longrightarrow
\pi_{1, *}\mathcal{G}_1
$$
qui induit un isomorphisme d'espaces vectoriels sur $\kappa(\xi_1)$,
où $\xi_1 \in Y_1$ est l'unique point générique de la fibre de
$Y_1$ au-dessus de $s$. De plus,\newline
$\dim_{y_1}(\text{Supp}(\Coker(\alpha_1)_s)) < d$.
Il peut arriver que la fibre du faisceau $\Coker(\alpha_1)_s$
en $y_1$ soit nulle. On peut alors restreindre $Y_1$ en utilisant
le Lemme \ref{lemma-shrink} (\ref{item-shrink-on-Y})
et supposer que $\Coker(\alpha_1) = 0$, obtenant ainsi un
d\'evissage complet de longueur nulle.

\medskip\noindent
Supposons maintenant que la fibre du faisceau $\Coker(\alpha_1)_s$
en $y_1$ ne soit pas nulle. L'hypothèse de récurrence fournit alors un
diagramme commutatif
\begin{equation}
\label{equation-overcome-this}
\vcenter{
\xymatrix{
(Y_1, y_1) \ar[d] & (Y'_1, y'_1) \ar[l]^h \ar[d] \\
(S, s) & (S', s') \ar[l]
}
}
\end{equation}
de schémas pointés dont les flèches
horizontales sont des voisinages \'etales élémentaires
et tel que $h^*\Coker(\alpha_1)/Y'_1/S'$ admette un
d\'evissage complet
$$
(Z_k, Y_k, i_k, \pi_k, \mathcal{G}_k, \alpha_k, z_k, y_k)_{k = 2, \ldots, n}
$$
en $y'_1$. (En particulier, $i_2 : Z_2 \to Y'_1$ est une immersion fermée dans
$Y'_2$.) On applique alors
le Lemme \ref{lemma-elementary-etale-neighbourhood}
à $S, X, \mathcal{F}, x, s$, au système
$(Z_1, Y_1, i_1, \pi_1, \mathcal{G}_1)$ et au
diagramme (\ref{equation-overcome-this}). On obtient un diagramme
$$
\xymatrix{
& & (X'', x'') \ar[lld] \ar[d] & (Z''_1, z''_1) \ar[l] \ar[lld] \ar[d] \\
(X, x) \ar[d] & (Z_1, z_1) \ar[l] \ar[d] &
(S'', s'') \ar[lld] & (Y''_1, y''_1) \ar[lld] \ar[l] \\
(S, s) & (Y_1, y_1) \ar[l]
}
$$
possédant toutes les propriétés énumérées dans le lemme cité.
En particulier, $Y''_1 \subset Y'_1 \times_{S'} S''$. Posons
$X_1 = Y'_1 \times_{S'} S''$ et notons $\mathcal{F}_1$
l'image réciproque de $\Coker(\alpha_1)$. D'après le
Lemme \ref{lemma-base-change-complete-at-x},
le système
\begin{equation}
\label{equation-shrink-this}
(Z_k \times_{S'} S'',
Y_k \times_{S'} S'', i''_k, \pi''_k, \mathcal{G}''_k,
\alpha''_k, z''_k, y''_k)_{k = 2, \ldots, n}
\end{equation}
est un d\'evissage complet de $\mathcal{F}_1$
sur $X_1$. De nouveau, le caractère local du problème
permet de remplacer $(X, x) \to (S, s)$ par $(X'', x'') \to (S'', s'')$.
On voit ainsi que l'on peut supposer :
\begin{enumerate}
\item[(a)] Il existe un d\'evissage en une étape
$(Z_1, Y_1, i_1, \pi_1, \mathcal{G}_1)$ de $\mathcal{F}/X/S$ en $x$,
\item[(b)] il existe un $\alpha_1 : \mathcal{O}_{Y_1}^{\oplus r_1}
\to \pi_{1, *}\mathcal{G}_1$ tel que $\alpha \otimes \kappa(\xi_1)$
soit un isomorphisme,
\item[(c)] $Y_1 \subset X_1$ est ouvert, $y_1 = x_1$, et
$\mathcal{F}_1|_{Y_1} \cong \Coker(\alpha_1)$, et
\item[(d)] il existe un d\'evissage complet\newline
$(Z_k, Y_k, i_k, \pi_k, \mathcal{G}_k, \alpha_k, z_k, y_k)_{k = 2, \ldots, n}$
de $\mathcal{F}_1/X_1/S$ en $x_1$.
\end{enumerate}
Pour achever la démonstration, il suffit de restreindre le d\'evissage en une étape
et le d\'evissage complet de façon à les assembler en un
d\'evissage complet. (Nous suggérons au lecteur de le faire lui-même en utilisant
les Lemmes \ref{lemma-shrink} et
\ref{lemma-shrink-complete},
au lieu de lire la démonstration qui suit.) Puisque $Y_1 \subset X_1$
est un voisinage ouvert de $x_1$, on peut appliquer
le Lemme \ref{lemma-shrink-complete} (\ref{item-shrink-on-X-complete})
pour trouver une restriction standard $S', X'_1, Z'_2, Y'_2, \ldots, Y'_n$
de la donnée (d), telle que $X'_1 \subset Y_1$. Notons que $X'_1$ est aussi
un ouvert standard du schéma affine $Y_1$. On restreint ensuite la donnée
(a) comme suit : on restreint d'abord la base $S$ à $S'$, voir
le Lemme \ref{lemma-shrink} (\ref{item-shrink-base}), puis
on restreint le résultat à $S''$, $X''$, $Z''_1$, $Y''_1$ en utilisant
le Lemme \ref{lemma-shrink} (\ref{item-shrink-on-Y}),
de sorte que finalement $Y''_1 = X'_1 \times_S S''$ et $S'' \subset S'$.
On voit alors que
$$
Z''_1, Y''_1, Z'_2 \times_{S'} S'', Y'_2 \times_{S'} S'', \ldots,
Y'_n \times_{S'} S''
$$
fournit le d\'evissage complet cherché.
\end{proof}

\noindent
En organisant encore les constructions, on obtient la conséquence suivante.

\begin{lemma}
\label{lemma-existence-complete}
Soit $X \to S$ un morphisme de type fini de schémas.
Soit $\mathcal{F}$ un $\mathcal{O}_X$-module quasi-cohérent de type fini.
Soit $s \in S$ un point.
Il existe un voisinage \'etale élémentaire
$(S', s') \to (S, s)$ et des morphismes \'etales
$h_i : Y_i \to X_{S'}$, $i = 1, \ldots, n$, tels que, pour chaque
$i$, il existe un d\'evissage complet de $\mathcal{F}_i/Y_i/S'$ au-dessus de $s'$,
où $\mathcal{F}_i$ est l'image réciproque de $\mathcal{F}$ sur $Y_i$,
et tels que $X_s = (X_{S'})_{s'} \subset \bigcup h_i(Y_i)$.
\end{lemma}

\begin{proof}
Pour tout point $x \in X_s$, on peut trouver un diagramme
$$
\xymatrix{
(X, x) \ar[d] & (X', x') \ar[l]^g \ar[d] \\
(S, s) & (S', s') \ar[l]
}
$$
de schémas pointés dont les flèches horizontales sont des voisinages
\'etales élémentaires et tel que $g^*\mathcal{F}/X'/S'$ admette un
d\'evissage complet en $x'$. Puisque $X \to S$ est de type fini, la
fibre $X_s$ est quasi-compacte et, puisque chaque $g : X' \to X$ comme ci-dessus
est ouvert, on peut recouvrir $X_s$ par une réunion finie de $g(X'_{s'})$.
On peut donc trouver une famille finie de tels diagrammes
$$
\vcenter{
\xymatrix{
(X, x) \ar[d] & (X'_i, x'_i) \ar[l]^{g_i} \ar[d] \\
(S, s) & (S'_i, s'_i) \ar[l]
}
}
\quad i = 1, \ldots, n
$$
telle que $X_s = \bigcup g_i(X'_i)$. Posons
$S' = S'_1 \times_S \ldots \times_S S'_n$
et soit $Y_i = X_i \times_{S'_i} S'$ le changement de base de $X'_i$ à $S'$. D'après
le Lemme \ref{lemma-base-change-complete},
on voit que l'image réciproque de $\mathcal{F}$ sur $Y_i$ admet un d\'evissage complet
au-dessus de $s$, ce qui donne le résultat voulu.
\end{proof}



\section{Traduction en termes d'algèbre}
\label{section-translation}

\noindent
Il peut être utile d'expliciter en termes algébriques ce que signifie l'existence d'un
d\'evissage complet. Nous introduisons la notion suivante (qui n'est pas
très utile, d'où le nom démesurément long que nous lui donnons).

\begin{definition}
\label{definition-elementary-etale-neighbourhood}
Soit $R \to S$ un homomorphisme d'anneaux. Soit $\mathfrak q$ un idéal premier de $S$ au-dessus
de l'idéal premier $\mathfrak p$ de $R$. Une {\it localisation \'etale élémentaire de
l'homomorphisme d'anneaux $R \to S$ en $\mathfrak q$} est donnée par un diagramme commutatif
d'anneaux et par les idéaux premiers correspondants
$$
\xymatrix{
S \ar[r] & S' \\
R \ar[u] \ar[r] & R' \ar[u]
}
\quad\quad
\xymatrix{
\mathfrak q \ar@{-}[r] & \mathfrak q' \\
\mathfrak p \ar@{-}[u] \ar@{-}[r] & \mathfrak p' \ar@{-}[u]
}
$$
tels que $R \to R'$ et $S \to S'$ soient des homomorphismes d'anneaux \'etales et que
$\kappa(\mathfrak p) = \kappa(\mathfrak p')$ et
$\kappa(\mathfrak q) = \kappa(\mathfrak q')$.
\end{definition}

\begin{definition}
\label{definition-complete-devissage-algebra}
Soit $R \to S$ un homomorphisme d'anneaux de type fini.
Soit $\mathfrak r$ un idéal premier de $R$.
Soit $N$ un $S$-module de type fini.
Un {\it d\'evissage complet de $N/S/R$ au-dessus de $\mathfrak r$}
est donné par des homomorphismes de $R$-algèbres
$$
\xymatrix{
& A_1 & & A_2 & & ... & & A_n \\
S \ar[ru] & & B_1 \ar[lu] \ar[ru] & & ... \ar[lu] \ar[ru] & &
... \ar[lu] \ar[ru] & & B_n \ar[lu]
}
$$
des $A_i$-modules de type fini $M_i$ et des homomorphismes de $B_i$-modules
$\alpha_i : B_i^{\oplus r_i} \to M_i$ tels que
\begin{enumerate}
\item $S \to A_1$ soit surjectif et de présentation finie,
\item $B_i \to A_{i + 1}$ soit surjectif et de présentation finie,
\item $B_i \to A_i$ soit fini,
\item $R \to B_i$ soit lisse à fibres géométriquement irréductibles,
\item $N \cong M_1$ comme $S$-modules,
\item $\Coker(\alpha_i) \cong M_{i + 1}$ comme $B_i$-modules,
\item $\alpha_i : \kappa(\mathfrak p_i)^{\oplus r_i}
\to M_i \otimes_{B_i} \kappa(\mathfrak p_i)$ soit un isomorphisme,
où $\mathfrak p_i = \mathfrak rB_i$, et
\item $\Coker(\alpha_n) = 0$.
\end{enumerate}
On dit alors que
$(A_i, B_i, M_i, \alpha_i)_{i = 1, \ldots, n}$
est un d\'evissage complet de $N/S/R$ au-dessus de $\mathfrak r$.
\end{definition}

\begin{remark}
\label{remark-finite-presentation}
Notons que les $R$-algèbres $B_i$, pour tout $i$, et $A_i$, pour $i \geq 2$,
sont de présentation finie sur $R$. Si $S$ est de présentation finie sur
$R$, alors $A_1$ est également de présentation finie sur
$R$. Dans ce cas, tous les homomorphismes d'anneaux du d\'evissage complet sont de
présentation finie. Voir
Algèbre, Lemme \ref{algebra-lemma-compose-finite-type}.
Toujours sous l'hypothèse que $S$ est de présentation finie sur $R$,
les conditions suivantes sont équivalentes :
\begin{enumerate}
\item $M$ est de présentation finie sur $S$,
\item $M_1$ est de présentation finie sur $A_1$,
\item $M_1$ est de présentation finie sur $B_1$,
\item chaque $M_i$ est de présentation finie à la fois comme $A_i$-module
et comme $B_i$-module.
\end{enumerate}
Les équivalences (1) $\Leftrightarrow$ (2) et (2) $\Leftrightarrow$ (3)
résultent de
Algèbre, Lemme \ref{algebra-lemma-finite-finitely-presented-extension}.
Si $M_1$ est de présentation finie, il en est de même de $\Coker(\alpha_1)$ (voir
Algèbre, Lemme \ref{algebra-lemma-extension}),
donc de $M_2$, et ainsi de suite.
\end{remark}

\begin{definition}
\label{definition-complete-devissage-at-x-algebra}
Soit $R \to S$ un homomorphisme d'anneaux de type fini.
Soit $\mathfrak q$ un idéal premier de $S$ au-dessus de l'idéal premier $\mathfrak r$ de $R$.
Soit $N$ un $S$-module de type fini.
Un {\it d\'evissage complet de $N/S/R$ en $\mathfrak q$} est donné par un
d\'evissage complet $(A_i, B_i, M_i, \alpha_i)_{i = 1, \ldots, n}$
de $N/S/R$ au-dessus de $\mathfrak r$ et par des idéaux premiers $\mathfrak q_i \subset B_i$
au-dessus de $\mathfrak r$ tels que
\begin{enumerate}
\item $\kappa(\mathfrak r) \subset \kappa(\mathfrak q_i)$ soit purement
transcendante,
\item il existe un unique idéal premier $\mathfrak q'_i \subset A_i$
au-dessus de $\mathfrak q_i \subset B_i$,
\item $\mathfrak q = \mathfrak q'_1 \cap S$ et
$\mathfrak q_i = \mathfrak q'_{i + 1} \cap A_i$,
\item $R \to B_i$ soit de dimension relative
$\dim_{\mathfrak q_i}(\text{Supp}(M_i \otimes_R \kappa(\mathfrak r)))$.
\end{enumerate}
\end{definition}

\begin{remark}
\label{remark-same-notion}
Soit $A \to B$ un homomorphisme d'anneaux de type fini et soit $N$ un
$B$-module de type fini. Soit $\mathfrak q$ un idéal premier de $B$ au-dessus de l'idéal premier
$\mathfrak r$ de $A$. Posons $X = \Spec(B)$, $S = \Spec(A)$ et
$\mathcal{F} = \widetilde{N}$ sur $X$. Soit $x$ le point correspondant
à $\mathfrak q$ et soit $s \in S$ le point correspondant à
$\mathfrak p$. Alors
\begin{enumerate}
\item s'il existe un d\'evissage complet de $\mathcal{F}/X/S$
au-dessus de $s$, il existe un d\'evissage complet de
$N/B/A$ au-dessus de $\mathfrak p$, et
\item il existe un d\'evissage complet de $\mathcal{F}/X/S$
en $x$ si et seulement s'il existe un d\'evissage complet de
$N/B/A$ en $\mathfrak q$.
\end{enumerate}
La seule nuance est que nous avons omis la condition sur
la dimension relative dans la formulation « un d\'evissage complet de
$N/B/A$ au-dessus de $\mathfrak p$ », ce qui explique que l'implication (1)
ne soit donnée que dans un sens.
La notion de d\'evissage complet en
$\mathfrak q$ inclut bien cette condition. Quoi qu'il en soit, nous utiliserons
seulement le fait que l'existence pour $\mathcal{F}/X/S$
entraîne l'existence pour $N/B/A$.
\end{remark}

\begin{lemma}
\label{lemma-existence-algebra}
Soit $R \to S$ un homomorphisme d'anneaux de type fini.
Soit $M$ un $S$-module de type fini.
Soit $\mathfrak q$ un idéal premier de $S$.
Il existe une localisation \'etale élémentaire
$R' \to S', \mathfrak q', \mathfrak p'$ de
l'homomorphisme d'anneaux $R \to S$ en $\mathfrak q$ telle qu'il
existe un d\'evissage complet de
$(M \otimes_S S')/S'/R'$ en $\mathfrak q'$.
\end{lemma}

\begin{proof}
C'est une reformulation de la
Proposition \ref{proposition-existence-complete-at-x}
au moyen de la
Remarque \ref{remark-same-notion}
\end{proof}




\section{Localisation et homomorphismes universellement injectifs}
\label{section-localize-universally-injective}


\begin{lemma}
\label{lemma-homothety-spectrum}
Soit $R \to S$ un homomorphisme d'anneaux.
Soit $N$ un $S$-module.
Supposons que
\begin{enumerate}
\item $R$ soit un anneau local d'idéal maximal $\mathfrak m$,
\item $\overline{S} = S/\mathfrak m S$ soit noethérien, et
\item $\overline{N} = N/\mathfrak m_R N$ soit un $\overline{S}$-module de type fini.
\end{enumerate}
Soit $\Sigma \subset S$ la partie multiplicative des éléments qui ne sont pas
diviseurs de zéro sur $\overline{N}$. Alors $\Sigma^{-1}S$ est un anneau semi-local
dont le spectre est formé des idéaux premiers $\mathfrak q \subset S$ contenus dans un
élément de $\text{Ass}_S(\overline{N})$. De plus, tout idéal
maximal de $\Sigma^{-1}S$ correspond à un idéal premier associé à
$\overline{N}$ sur $\overline{S}$.
\end{lemma}

\begin{proof}
Notons que
$\text{Ass}_S(\overline{N}) = \text{Ass}_{\overline{S}}(\overline{N})$, voir
Algèbre, Lemme \ref{algebra-lemma-ass-quotient-ring}.
Cet ensemble est fini d'après
Algèbre, Lemme \ref{algebra-lemma-finite-ass}.
Écrivons\newline $\{\mathfrak q_1, \ldots, \mathfrak q_r\} = \text{Ass}_S(\overline{N})$.
On a $\Sigma = S \setminus (\bigcup \mathfrak q_i)$ d'après
Algèbre, Lemme \ref{algebra-lemma-ass-zero-divisors}.
D'après la description de $\Spec(\Sigma^{-1}S)$ donnée dans
Algèbre, Lemme \ref{algebra-lemma-spec-localization},
et d'après
Algèbre, Lemme \ref{algebra-lemma-silly},
on voit que les idéaux premiers de $\Sigma^{-1}S$ correspondent aux idéaux premiers de
$S$ contenus dans l'un des $\mathfrak q_i$.
Les idéaux maximaux de $\Sigma^{-1}S$ correspondent donc bijectivement aux
éléments maximaux (pour\ l'inclusion) de l'ensemble
$\{\mathfrak q_1, \ldots, \mathfrak q_r\}$. Cela démontre le lemme.
\end{proof}

\begin{lemma}
\label{lemma-homothety-universally-injective}
Sous les hypothèses et avec les notations du
Lemme \ref{lemma-homothety-spectrum},
supposons de plus que
\begin{enumerate}
\item $S$ soit local et que $R \to S$ soit un homomorphisme local,
\item $S$ soit essentiellement de présentation finie sur $R$,
\item $N$ soit de présentation finie sur $S$, et
\item $N$ soit plat sur $R$.
\end{enumerate}
Alors tout $s \in \Sigma$ définit un
homomorphisme universellement injectif de $R$-modules $s : N \to N$, et
l'homomorphisme $N \to \Sigma^{-1}N$ est $R$-universellement injectif.
\end{lemma}

\begin{proof}
D'après
Algèbre, Lemme \ref{algebra-lemma-mod-injective-general},
la suite\newline $0 \to N \to N \to N/sN \to 0$ est exacte et
$N/sN$ est plat sur $R$. Il en résulte que $s : N \to N$
est universellement injectif, voir
Algèbre, Lemme \ref{algebra-lemma-flat-tor-zero}.
L'homomorphisme $N \to \Sigma^{-1}N$ est universellement injectif comme limite
inductive filtrante des homomorphismes $s : N \to N$.
\end{proof}

\begin{lemma}
\label{lemma-base-change-universally-flat-local}
Soit $R \to S$ un homomorphisme d'anneaux.
Soit $N$ un $S$-module.
Soit $S \to S'$ un homomorphisme d'anneaux.
Supposons que
\begin{enumerate}
\item $R \to S$ soit un homomorphisme local d'anneaux locaux,
\item $S$ soit essentiellement de présentation finie sur $R$,
\item $N$ soit de présentation finie sur $S$,
\item $N$ soit plat sur $R$,
\item $S \to S'$ soit plat, et
\item l'image de $\Spec(S') \to \Spec(S)$ contienne
tous les idéaux premiers $\mathfrak q$ de $S$ au-dessus de $\mathfrak m_R$
tels que $\mathfrak q$ soit un idéal premier associé à $N/\mathfrak m_R N$.
\end{enumerate}
Alors $N \to N \otimes_S S'$ est $R$-universellement injectif.
\end{lemma}

\begin{proof}
Posons $N' = N \otimes_R S'$. Considérons le diagramme commutatif
$$
\xymatrix{
N \ar[d] \ar[r] & N' \ar[d] \\
\Sigma^{-1}N \ar[r] & \Sigma^{-1}N'
}
$$
où $\Sigma \subset S$ est l'ensemble des éléments qui ne sont pas
diviseurs de zéro sur $N/\mathfrak m_R N$. S'il est établi que l'homomorphisme
$N \to \Sigma^{-1}N'$ est universellement injectif, alors $N \to N'$
l'est aussi (voir
Algèbre, Lemme \ref{algebra-lemma-universally-injective-permanence}).

\medskip\noindent
D'après le
Lemme \ref{lemma-homothety-spectrum},
l'anneau $\Sigma^{-1}S$ est un anneau semi-local dont les idéaux maximaux
correspondent aux idéaux premiers associés à $N/\mathfrak m_R N$.
L'image de
$\Spec(\Sigma^{-1}S') \to \Spec(\Sigma^{-1}S)$
contient donc tous ces idéaux maximaux par hypothèse. D'après
Algèbre, Lemme \ref{algebra-lemma-ff-rings},
l'homomorphisme d'anneaux $\Sigma^{-1}S \to \Sigma^{-1}S'$ est fidèlement plat.
Ainsi $\Sigma^{-1}N \to \Sigma^{-1}N'$, qui est l'homomorphisme
$$
N \otimes_S \Sigma^{-1}S \longrightarrow N \otimes_S \Sigma^{-1}S'
$$
est universellement injectif, voir
Algèbre, Lemmes \ref{algebra-lemma-faithfully-flat-universally-injective} et
\ref{algebra-lemma-universally-injective-tensor}.
Enfin, on applique
le Lemme \ref{lemma-homothety-universally-injective}
pour voir que $N \to \Sigma^{-1}N$ est universellement injectif.
Puisque le composé d'homomorphismes de modules universellement injectifs est universellement
injectif (voir
Algèbre, Lemme \ref{algebra-lemma-composition-universally-injective}),
on en conclut que $N \to \Sigma^{-1}N'$ est universellement injectif, ce qui donne le résultat voulu.
\end{proof}

\begin{lemma}
\label{lemma-base-change-universally-flat}
Soit $R \to S$ un homomorphisme d'anneaux.
Soit $N$ un $S$-module.
Soit $S \to S'$ un homomorphisme d'anneaux.
Supposons que
\begin{enumerate}
\item $R \to S$ soit de présentation finie et que $N$ soit de présentation finie
sur $S$,
\item $N$ soit plat sur $R$,
\item $S \to S'$ soit plat, et
\item l'image de $\Spec(S') \to \Spec(S)$ contienne
tous les idéaux premiers $\mathfrak q$ tels que $\mathfrak q$ soit un idéal premier associé
à $N \otimes_R \kappa(\mathfrak p)$, où $\mathfrak p$ est l'image
réciproque de $\mathfrak q$ dans $R$.
\end{enumerate}
Alors $N \to N \otimes_S S'$ est $R$-universellement injectif.
\end{lemma}

\begin{proof}
D'après
Algèbre, Lemme \ref{algebra-lemma-universally-injective-check-stalks},
il suffit de montrer que $N_{\mathfrak q} \to (N \otimes_R S')_{\mathfrak q}$
est $R_{\mathfrak p}$-universellement injectif pour tout idéal premier $\mathfrak q$
de $S$ au-dessus de $\mathfrak p$ dans $R$. On peut donc appliquer
le Lemme \ref{lemma-base-change-universally-flat-local}
aux homomorphismes d'anneaux
$R_{\mathfrak p} \to S_{\mathfrak q} \to S'_{\mathfrak q}$
et au module $N_{\mathfrak q}$.
\end{proof}

\noindent
Le lecteur pourra comparer le lemme suivant à
Algèbre, Lemmes \ref{algebra-lemma-mod-injective} et
\ref{algebra-lemma-mod-injective-general}, ainsi qu'aux résultats de la
Section \ref{section-variants-mod-injective}.
Dans chaque cas, la conclusion est que l'homomorphisme $u : M \to N$ est
universellement injectif, de conoyau plat.

\begin{lemma}
\label{lemma-universally-injective-local}
Soit $(R, \mathfrak m)$ un anneau local. Soit $u : M \to N$ un homomorphisme de $R$-modules.
Si $M$ est un $R$-module projectif, si $N$ est un $R$-module plat et si
$\overline{u} : M/\mathfrak mM \to N/\mathfrak mN$ est injectif,
alors $u$ est universellement injectif.
\end{lemma}

\begin{proof}
D'après
Algèbre, Théorème \ref{algebra-theorem-projective-free-over-local-ring},
le module $M$ est libre. Si l'on établit le résultat pour tout facteur direct de type
fini de $M$, le lemme en résultera. On peut donc
supposer que $M$ est libre de type fini. Écrivons $N = \colim_i N_i$ comme
limite inductive filtrante de modules libres de type fini, voir
Algèbre, Théorème \ref{algebra-theorem-lazard}.
Notons que $u : M \to N$ se factorise par $N_i$ pour un certain $i$ (puisque $M$ est libre
de type fini). Notons $u_i : M \to N_i$ l'homomorphisme de $R$-modules correspondant.
Puisque $\overline{u}$ est injectif, on voit que
$\overline{u_i} : M/\mathfrak mM \to N_i/\mathfrak mN_i$ est
injectif et le reste après composition avec les homomorphismes
$N_i/\mathfrak mN_i \to N_{i'}/\mathfrak mN_{i'}$ pour tout $i' \geq i$.
Puisque $M$ et $N_{i'}$ sont libres de type fini sur l'anneau local $R$, cela entraîne
que $M \to N_{i'}$ est une injection scindée pour tout $i' \geq i$. Par suite,
pour tout $R$-module $Q$, on voit que $M \otimes_R Q \to N_{i'} \otimes_R Q$
est injectif pour tout $i' \geq i$. Puisque $- \otimes_R Q$ commute aux
limites inductives, on en conclut que $M \otimes_R Q \to N_{i'} \otimes_R Q$
est injectif, comme souhaité.
\end{proof}

\begin{lemma}
\label{lemma-invert-universally-injective}
Sous les hypothèses et avec les notations du
Lemme \ref{lemma-homothety-spectrum},
supposons de plus que $N$ soit projectif comme $R$-module.
Alors tout $s \in \Sigma$ définit un
homomorphisme universellement injectif de $R$-modules $s : N \to N$, et
l'homomorphisme $N \to \Sigma^{-1}N$ est $R$-universellement injectif.
\end{lemma}

\begin{proof}
Choisissons $s \in \Sigma$. D'après le
Lemme \ref{lemma-universally-injective-local},
l'homomorphisme $s : N \to N$ est universellement injectif.
L'homomorphisme $N \to \Sigma^{-1}N$ est universellement injectif comme limite
inductive filtrante des homomorphismes $s : N \to N$.
\end{proof}






\section{Complétion et modules de Mittag-Leffler}
\label{section-completion-ML}


\begin{lemma}
\label{lemma-completed-direct-sum-ML}
Soit $R$ un anneau. Soit $I \subset R$ un idéal. Soit $A$ un ensemble.
Supposons que $R$ soit noethérien et complet pour la topologie définie par $I$. Le complété
$(\bigoplus\nolimits_{\alpha \in A} R)^\wedge$
est plat et de Mittag-Leffler.
\end{lemma}

\begin{proof}
D'après
Compléments d'algèbre, Lemme
\ref{more-algebra-lemma-ui-completion-direct-sum-into-product},\newline
l'homomorphisme\newline $(\bigoplus\nolimits_{\alpha \in A} R)^\wedge
\to \prod_{\alpha \in A} R$ est universellement injectif.
Ainsi, d'après
Algèbre, Lemmes \ref{algebra-lemma-ui-flat-domain} et
\ref{algebra-lemma-pure-submodule-ML},
il suffit de montrer que $\prod_{\alpha \in A} R$ est plat et de Mittag-Leffler.
D'après
Algèbre, Proposition \ref{algebra-proposition-characterize-coherent}
(et
Algèbre, Lemme \ref{algebra-lemma-Noetherian-coherent}),
on voit que $\prod_{\alpha \in A} R$ est plat.
On conclut donc, car un produit de copies de $R$ est de Mittag-Leffler, voir
Algèbre, Lemme \ref{algebra-lemma-product-over-Noetherian-ring}.
\end{proof}

\begin{lemma}
\label{lemma-lift-ML}
Soit $R$ un anneau. Soit $I \subset R$ un idéal.
Soit $M$ un $R$-module.
Supposons que
\begin{enumerate}
\item $R$ soit noethérien et complet pour la topologie $I$-adique,
\item $M$ soit plat sur $R$, et
\item $M/IM$ soit un $R/I$-module projectif.
\end{enumerate}
Alors le complété $I$-adique $M^\wedge$ est un
$R$-module plat de Mittag-Leffler.
\end{lemma}

\begin{proof}
Choisissons une surjection $F \to M$, où $F$ est un $R$-module libre. D'après
Algèbre, Lemme \ref{algebra-lemma-split-completed-sequence},
le module $M^\wedge$ est un facteur direct du module $F^\wedge$.
Il suffit donc de démontrer le lemme pour $F$.
Dans ce cas, le lemme résulte du
Lemme \ref{lemma-completed-direct-sum-ML}.
\end{proof}

\noindent
Dans les
Lemmes \ref{lemma-universally-injective-to-completion} et
\ref{lemma-universally-injective-to-completion-flat},
l'hypothèse que $S$ est noethérien est satisfaite si $R \to S$ est de type fini, voir
Algèbre, Lemme \ref{algebra-lemma-Noetherian-permanence}.

\begin{lemma}
\label{lemma-universally-injective-to-completion}
Soit $R$ un anneau.
Soit $I \subset R$ un idéal.
Soient $R \to S$ un homomorphisme d'anneaux et $N$ un $S$-module.
Supposons que
\begin{enumerate}
\item $R$ soit un anneau noethérien,
\item $S$ soit un anneau noethérien,
\item $N$ soit un $S$-module de type fini, et
\item pour tout $R$-module de type fini $Q$, tout
$\mathfrak q \in \text{Ass}_S(Q \otimes_R N)$
vérifie $IS + \mathfrak q \not = S$.
\end{enumerate}
Alors l'homomorphisme $N \to N^\wedge$ de $N$ dans le complété $I$-adique de $N$
est universellement injectif comme homomorphisme de $R$-modules.
\end{lemma}

\begin{proof}
Il faut montrer que, pour tout $R$-module de type fini $Q$, l'homomorphisme
$Q \otimes_R N \to Q \otimes_R N^\wedge$ est injectif, voir
Algèbre, Théorème \ref{algebra-theorem-universally-exact-criteria}.
Puisqu'il existe un homomorphisme canonique $Q \otimes_R N^\wedge \to (Q \otimes_R N)^\wedge$,
il suffit de démontrer que l'homomorphisme canonique
$Q \otimes_R N \to (Q \otimes_R N)^\wedge$ est injectif.
On peut donc remplacer $N$ par $Q \otimes_R N$, et il suffit de démontrer
l'injectivité de l'homomorphisme $N \to N^\wedge$.

\medskip\noindent
Soit $K = \Ker(N \to N^\wedge)$. Il suffit de montrer que
$K_{\mathfrak q} = 0$ pour $\mathfrak q \in \text{Ass}(N)$, puisque $N$ est un
sous-module de $\prod_{\mathfrak q \in \text{Ass}(N)} N_{\mathfrak q}$, voir
Algèbre, Lemme \ref{algebra-lemma-zero-at-ass-zero}.
Choisissons $\mathfrak q \in \text{Ass}(N)$. La dernière hypothèse montre
qu'il existe un idéal premier $\mathfrak q' \supset IS + \mathfrak q$.
Puisque $K_{\mathfrak q}$ est un localisé de $K_{\mathfrak q'}$,
il suffit de démontrer la nullité de $K_{\mathfrak q'}$.
Notons que $K = \bigcap I^nN$ ; par suite,
$K_{\mathfrak q'} \subset \bigcap I^nN_{\mathfrak q'}$.
Ainsi $K_{\mathfrak q'} = 0$ d'après
Algèbre, Lemme \ref{algebra-lemma-intersect-powers-ideal-module-zero}.
\end{proof}

\begin{lemma}
\label{lemma-universally-injective-to-completion-flat}
Soit $R$ un anneau.
Soit $I \subset R$ un idéal.
Soient $R \to S$ un homomorphisme d'anneaux et $N$ un $S$-module.
Supposons que
\begin{enumerate}
\item $R$ soit un anneau noethérien,
\item $S$ soit un anneau noethérien,
\item $N$ soit un $S$-module de type fini,
\item $N$ soit plat sur $R$, et
\item pour tout idéal premier $\mathfrak q \subset S$ associé à
$N \otimes_R \kappa(\mathfrak p)$, où $\mathfrak p = R \cap \mathfrak q$,
on ait $IS + \mathfrak q \not = S$.
\end{enumerate}
Alors l'homomorphisme $N \to N^\wedge$ de $N$ dans le complété $I$-adique de $N$
est universellement injectif comme homomorphisme de $R$-modules.
\end{lemma}

\begin{proof}
Cela résulte du
Lemme \ref{lemma-universally-injective-to-completion},
car
Algèbre, Lemme
\ref{algebra-lemma-bourbaki-fibres} et
Remarque \ref{algebra-remark-bourbaki}
assurent que les ensembles d'idéaux premiers associés aux produits tensoriels
$N \otimes_R Q$ sont contenus dans l'ensemble des idéaux premiers associés aux
modules $N \otimes_R \kappa(\mathfrak p)$.
\end{proof}




\section{Modules projectifs}
\label{section-projective}

\noindent
Le lemme suivant permet de démontrer la projectivité par
récurrence noethérienne sur la base, voir
le Lemme \ref{lemma-fibres-irreducible-flat-projective}.

\begin{lemma}
\label{lemma-flat-pure-over-complete-projective}
Soit $R$ un anneau.
Soit $I \subset R$ un idéal.
Soient $R \to S$ un homomorphisme d'anneaux et $N$ un $S$-module.
Supposons que
\begin{enumerate}
\item $R$ soit noethérien et complet pour la topologie $I$-adique,
\item $R \to S$ soit de type fini,
\item $N$ soit un $S$-module de type fini,
\item $N$ soit plat sur $R$,
\item $N/IN$ soit projectif comme $R/I$-module, et
\item pour tout idéal premier $\mathfrak q \subset S$ associé à
$N \otimes_R \kappa(\mathfrak p)$, où $\mathfrak p = R \cap \mathfrak q$,
on ait $IS + \mathfrak q \not = S$.
\end{enumerate}
Alors $N$ est projectif comme $R$-module.
\end{lemma}

\begin{proof}
D'après le
Lemme \ref{lemma-universally-injective-to-completion-flat},
l'homomorphisme $N \to N^\wedge$ est universellement injectif.
D'après le
Lemme \ref{lemma-lift-ML},
le module $N^\wedge$ est de Mittag-Leffler.
D'après
Algèbre, Lemme \ref{algebra-lemma-pure-submodule-ML},
on en conclut que $N$ est de Mittag-Leffler.
Ainsi $N$ est engendré par une famille dénombrable, plat et de Mittag-Leffler comme $R$-module,
donc projectif d'après
Algèbre, Lemme \ref{algebra-lemma-countgen-projective}.
\end{proof}

\begin{lemma}
\label{lemma-fibres-irreducible-flat-projective}
Soit $R$ un anneau.
Soit $R \to S$ un homomorphisme d'anneaux.
Supposons que
\begin{enumerate}
\item $R$ soit noethérien,
\item $R \to S$ soit de type fini et plat, et
\item tout anneau de fibre $S \otimes_R \kappa(\mathfrak p)$ soit
géométriquement intègre sur $\kappa(\mathfrak p)$.
\end{enumerate}
Alors $S$ est projectif comme $R$-module.
\end{lemma}

\begin{proof}
Considérons l'ensemble
$$
\{I \subset R \mid S/IS\text{ not projective as }R/I\text{-module}\}
$$
Il faut montrer que cet ensemble est vide. Raisonnons par l'absurde en le supposant
non vide. Il contient alors un élément maximal $I$.
Soit $J = \sqrt{I}$ son radical. Si $I \not = J$, alors
$S/JS$ est projectif comme $R/J$-module, $S/IS$ est plat sur $R/I$
et $J/I$ est un idéal nilpotent de $R/I$. En appliquant
Algèbre, Lemme \ref{algebra-lemma-lift-projective},
on voit que $S/IS$ est un $R/I$-module projectif, ce qui est contradictoire.
On peut donc supposer que $I$ est un idéal radical. Autrement dit, on
est ramené à démontrer le lemme lorsque $R$ est un anneau réduit et que
$S/IS$ est un $R/I$-module projectif pour tout idéal non nul $I$
de $R$.

\medskip\noindent
Supposons que $R$ soit un anneau réduit et que $S/IS$ soit un $R/I$-module projectif
pour tout idéal non nul $I$ de $R$. Par platitude générique,
Algèbre, Lemme \ref{algebra-lemma-generic-flatness-Noetherian}
(appliqué à un localisé $R_g$ qui est intègre), ou par le résultat plus général
Algèbre, Lemme \ref{algebra-lemma-generic-flatness-reduced},
il existe un élément non nul $f \in R$ tel que $S_f$ soit libre comme
$R_f$-module. Notons $R^\wedge = \lim R/(f^n)$ le complété $(f)$-adique
de $R$. Notons que l'homomorphisme d'anneaux
$$
R \longrightarrow R_f \times R^\wedge
$$
est fidèlement plat, voir
Algèbre, Lemme \ref{algebra-lemma-completion-flat}.
Par descente fidèlement plate de la projectivité, voir
Algèbre, Théorème \ref{algebra-theorem-ffdescent-projectivity},
il suffit donc de démontrer que $S \otimes_R R^\wedge$ est un
$R^\wedge$-module projectif. Pour cela, on utilisera le critère du
Lemme \ref{lemma-flat-pure-over-complete-projective}.
Tout d'abord, notons que $S/fS = (S \otimes_R R^\wedge)/f(S \otimes_R R^\wedge)$
est un $R/(f)$-module projectif et que $S \otimes_R R^\wedge$ est plat
et de type fini sur $R^\wedge$, par changement de base.
Supposons ensuite que $\mathfrak p^\wedge$ soit un idéal premier
de $R^\wedge$. Soit $\mathfrak p \subset R$ l'idéal premier correspondant
de $R$. Puisque $R \to S$ est à anneaux de fibres géométriquement intègres,
il en est de même après tout changement de base. Par suite,
$\mathfrak q^\wedge = \mathfrak p^\wedge(S \otimes_R R^\wedge)$,
est un idéal premier au-dessus de $\mathfrak p^\wedge$
et c'est l'unique idéal premier associé à
$S \otimes_R \kappa(\mathfrak p^\wedge)$. Il suffit donc de montrer que
$f(S \otimes_R R^\wedge) + \mathfrak q^\wedge \not = S \otimes_R R^\wedge$.
Cela résulte de $\mathfrak p^\wedge + fR^\wedge \not = R^\wedge$, puisque
$f$ appartient au radical de Jacobson de l'anneau $f$-adiquement complet $R^\wedge$,
et du fait que $R^\wedge \to S \otimes_R R^\wedge$ est surjectif sur les spectres,
ses fibres étant non vides (les espaces irréductibles sont non vides).
\end{proof}

\begin{lemma}
\label{lemma-fibres-irreducible-flat-projective-nonnoetherian}
Soit $R$ un anneau. Soit $R \to S$ un homomorphisme d'anneaux.
Supposons que
\begin{enumerate}
\item $R \to S$ soit de présentation finie et plat, et
\item tout anneau de fibre $S \otimes_R \kappa(\mathfrak p)$ soit
géométriquement intègre sur $\kappa(\mathfrak p)$.
\end{enumerate}
Alors $S$ est projectif comme $R$-module.
\end{lemma}

\begin{proof}
On peut trouver un diagramme cocartésien d'anneaux
$$
\xymatrix{
S_0 \ar[r] & S \\
R_0 \ar[u] \ar[r] & R \ar[u]
}
$$
tel que $R_0$ soit de type fini sur $\mathbf{Z}$ et que l'homomorphisme
$R_0 \to S_0$ soit de type fini et plat, à fibres géométriquement
intègres, voir
Compléments sur les morphismes,
Lemmes \ref{more-morphisms-lemma-Noetherian-approximation-flat},
\ref{more-morphisms-lemma-Noetherian-approximation-geometrically-reduced},
\ref{more-morphisms-lemma-Noetherian-approximation-geometrically-irreducible},
et \ref{more-morphisms-lemma-Noetherian-approximation-combine}.
D'après le
Lemme \ref{lemma-fibres-irreducible-flat-projective},
on voit que $S_0$ est un $R_0$-module projectif. Ainsi $S = S_0 \otimes_{R_0} R$
est un $R$-module projectif, voir
Algèbre, Lemme \ref{algebra-lemma-ascend-properties-modules}.
\end{proof}

\begin{remark}
\label{remark-how-in-RG}
Le Lemme \ref{lemma-fibres-irreducible-flat-projective-nonnoetherian}
est une étape essentielle dans l'établissement des résultats de ce chapitre. Son analogue
dans \cite{GruRay} est \cite[I Proposition 3.3.1]{GruRay} :
si $R \to S$ est lisse à fibres géométriquement intègres, alors $S$
est projectif comme $R$-module. C'est un cas particulier du
Lemme \ref{lemma-fibres-irreducible-flat-projective-nonnoetherian},
mais, puisque nous améliorerons encore ce lemme dans la suite, il est peu utile
de disposer dès à présent d'un résultat plus fort.
Esquissons brièvement la démonstration donnée dans \cite{GruRay}.
\begin{enumerate}
\item On se ramène d'abord au cas où $R$ est noethérien, comme ci-dessus.
\item Puisque la projectivité descend par les homomorphismes d'anneaux fidèlement plats, voir
Algèbre, Théorème \ref{algebra-theorem-ffdescent-projectivity},
on peut travailler localement pour la topologie fppf sur $R$, donc supposer
que $R \to S$ admet une section $\sigma : S \to R$. (Il suffit, comme d'habitude,
d'effectuer le changement de base à $S$.) Posons $I = \Ker(S \to R)$.
\item En localisant encore sur $R$, on peut supposer que $I/I^2$ est un
$R$-module libre et que le complété $S^\wedge$ de $S$ pour la topologie définie par $I$
est isomorphe à $R[[t_1, \ldots, t_n]]$, voir
Morphismes, Lemme \ref{morphisms-lemma-section-smooth-morphism}.
On utilise ici le fait que $R \to S$ est lisse.
\item Pour démontrer que $S$ est projectif comme $R$-module, il suffit de
montrer que $S$ est plat, engendré par une famille dénombrable et de Mittag-Leffler comme
$R$-module, voir
Algèbre, Lemme \ref{algebra-lemma-countgen-projective}.
Les deux premières propriétés sont évidentes. Il suffit donc de montrer que $S$
est de Mittag-Leffler comme $R$-module. D'après
Algèbre, Lemme \ref{algebra-lemma-power-series-ML},
le module $R[[t_1, \ldots, t_n]]$ est de Mittag-Leffler sur $R$. Ainsi
Algèbre, Lemme \ref{algebra-lemma-pure-submodule-ML},
montre qu'il suffit d'établir que
$S \to S^\wedge$ est universellement injectif comme homomorphisme de $R$-modules.
\item On applique
le Lemme \ref{lemma-base-change-universally-flat}
pour voir que $S \to S^\wedge$ est $R$-universellement injectif.
En effet, puisque $R \to S$ est à fibres géométriquement intègres, tout point associé
de tout anneau de fibre est précisément le point générique de cet anneau de fibre,
qui appartient à l'image de $\Spec(S^\wedge) \to \Spec(S)$.
\end{enumerate}
Il existe une analogie entre la démonstration ainsi esquissée et
les arguments qui conduisent à la démonstration du
Lemme \ref{lemma-fibres-irreducible-flat-projective-nonnoetherian}.
Dans les deux cas, une complétion joue un rôle essentiel, et l'hypothèse
que les fibres sont géométriquement intègres assure chaque fois que
l'homomorphisme de $S$ dans le complété de $S$ est $R$-universellement injectif.
\end{remark}













\section{Modules plats de type fini, première partie}
\label{section-finite-type-flat-I}

\noindent
Dans certains cas, étant donnés un homomorphisme d'anneaux $R \to S$ de présentation finie et
un $S$-module de type fini $N$, la platitude de $N$ sur $R$ entraîne que $N$
est de présentation finie. Nous démontrons dans cette section que cette assertion est vraie
« ponctuellement ». Remarquons que la première démonstration de la
Proposition \ref{proposition-finite-type-flat-at-point}
utilise les résultats géométriques de la
Section \ref{section-local-structure-module},
mais non l'existence d'un d\'evissage complet.

\begin{lemma}
\label{lemma-induction-step}
Soit $(R, \mathfrak m)$ un anneau local. Soit $R \to S$ un homomorphisme d'anneaux de présentation finie,
plat et à fibres géométriquement intègres. Posons
$\mathfrak p = \mathfrak mS$. Soit $\mathfrak q \subset S$ un idéal premier
au-dessus de $\mathfrak m$. Soit $N$ un $S$-module de type fini.
Il existe $r \geq 0$ et un homomorphisme de $S$-modules
$$
\alpha : S^{\oplus r} \longrightarrow N
$$
tels que
$\alpha : \kappa(\mathfrak p)^{\oplus r} \to N \otimes_S \kappa(\mathfrak p)$
soit un isomorphisme. Pour tout tel $\alpha$, les conditions suivantes sont équivalentes :
\begin{enumerate}
\item $N_{\mathfrak q}$ est $R$-plat,
\item $\alpha$ est $R$-universellement injectif et
$\Coker(\alpha)_{\mathfrak q}$ est $R$-plat,
\item $\alpha$ est injectif et
$\Coker(\alpha)_{\mathfrak q}$ est $R$-plat,
\item $\alpha_{\mathfrak p}$ est un isomorphisme et
$\Coker(\alpha)_{\mathfrak q}$ est $R$-plat, et
\item $\alpha_{\mathfrak q}$ est injectif et
$\Coker(\alpha)_{\mathfrak q}$ est $R$-plat.
\end{enumerate}
\end{lemma}

\begin{proof}
Pour obtenir $\alpha$, posons
$r = \dim_{\kappa(\mathfrak p)} N \otimes_S \kappa(\mathfrak p)$ et choisissons
$x_1, \ldots, x_r \in N$ dont les images forment une base de\newline
$N \otimes_S \kappa(\mathfrak p)$. Définissons
$\alpha(s_1, \ldots, s_r) = \sum s_i x_i$. Cela démontre l'existence.

\medskip\noindent
Fixons un tel $\alpha$. L'implication la plus intéressante est
(1) $\Rightarrow$ (2), que nous démontrons d'abord. Supposons (1).
Puisque $S/\mathfrak mS$ est un anneau intègre de corps des fractions $\kappa(\mathfrak p)$,
on voit que
$(S/\mathfrak mS)^{\oplus r} \to
N_{\mathfrak p}/\mathfrak mN_{\mathfrak p} = N \otimes_S \kappa(\mathfrak p)$
est injectif. Par suite, d'après les
Lemmes \ref{lemma-universally-injective-local} et
\ref{lemma-fibres-irreducible-flat-projective-nonnoetherian},
l'homomorphisme $S^{\oplus r} \to N_{\mathfrak p}$ est $R$-universellement injectif.
Il en résulte que $S^{\oplus r} \to N$ est $R$-universellement injectif, voir
Algèbre, Lemme \ref{algebra-lemma-universally-injective-permanence}.
Le localisé $\alpha_{\mathfrak q}$ est alors lui aussi $R$-universellement
injectif, voir
Algèbre, Lemme \ref{algebra-lemma-universally-injective-localize}.
On en conclut que $\Coker(\alpha)_{\mathfrak q}$ est $R$-plat d'après
Algèbre, Lemme \ref{algebra-lemma-ui-flat-domain}.

\medskip\noindent
L'implication (2) $\Rightarrow$ (3) est immédiate. Si (3) est vérifiée, alors
$\alpha_{\mathfrak p}$ est injectif, comme localisé d'un homomorphisme injectif
de modules. Par le lemme de Nakayama
(Algèbre, Lemme \ref{algebra-lemma-NAK}),
$\alpha_{\mathfrak p}$ est aussi surjectif. Ainsi (3) $\Rightarrow$ (4).
Si (4) est vérifiée, alors $\alpha_{\mathfrak p}$ est un isomorphisme ; par suite,
$\alpha$ est injectif, puisque $S_{\mathfrak q} \to S_{\mathfrak p}$ est injectif.
En effet, les éléments de $S \setminus \mathfrak p$ ne sont pas diviseurs de zéro sur $S$,
comme on le voit en combinant les
Lemmes \ref{lemma-invert-universally-injective} et
\ref{lemma-fibres-irreducible-flat-projective-nonnoetherian}.
Ainsi (4) $\Rightarrow$ (5). Enfin, si (5) est vérifiée, alors
$N_{\mathfrak q}$ est $R$-plat, comme extension de modules plats, voir
Algèbre, Lemme \ref{algebra-lemma-flat-ses}.
Ainsi (5) $\Rightarrow$ (1), ce qui achève la démonstration.
\end{proof}

\begin{lemma}
\label{lemma-complete-devissage-flat-finite-type-module}
Soit $(R, \mathfrak m)$ un anneau local.
Soit $R \to S$ un homomorphisme d'anneaux de présentation finie.
Soit $N$ un $S$-module de type fini.
Soit $\mathfrak q$ un idéal premier de $S$ au-dessus de $\mathfrak m$.
Supposons que $N_{\mathfrak q}$ soit plat sur $R$ et
qu'il existe un d\'evissage complet de $N/S/R$ en $\mathfrak q$.
Alors $N$ est un $S$-module de présentation finie, libre comme $R$-module,
et il existe un isomorphisme
$$
N \cong B_1^{\oplus r_1} \oplus \ldots \oplus B_n^{\oplus r_n}
$$
de $R$-modules, où chaque $B_i$ est une $R$-algèbre lisse à fibres géométriquement
irréductibles.
\end{lemma}

\begin{proof}
Soit $(A_i, B_i, M_i, \alpha_i, \mathfrak q_i)_{i = 1, \ldots, n}$
le d\'evissage complet donné. On démontre le lemme par récurrence sur $n$.
Notons que $N$ est de présentation finie comme $S$-module si et seulement si
$M_1$ est de présentation finie comme $B_1$-module, voir
la Remarque \ref{remark-finite-presentation}.
Notons que $N_{\mathfrak q} \cong (M_1)_{\mathfrak q_1}$ comme $R$-modules,
car (a) $N_{\mathfrak q} \cong (M_1)_{\mathfrak q'_1}$, où
$\mathfrak q'_1$ est l'unique idéal premier de $A_1$ au-dessus de $\mathfrak q_1$,
et (b) $(A_1)_{\mathfrak q'_1} = (A_1)_{\mathfrak q_1}$ d'après
Algèbre, Lemme \ref{algebra-lemma-unique-prime-over-localize-below},
d'où (c) $(M_1)_{\mathfrak q'_1} \cong (M_1)_{\mathfrak q_1}$.
Ainsi $(M_1)_{\mathfrak q_1}$ est un $R$-module plat.\newline On peut donc remplacer
$(S, N)$ par $(B_1, M_1)$ pour démontrer le lemme. D'après le
Lemme \ref{lemma-induction-step},
l'homomorphisme $\alpha_1 : B_1^{\oplus r_1} \to M_1$ est $R$-universellement injectif
et $\Coker(\alpha_1)_{\mathfrak q}$ est $R$-plat.
Notons que $(A_i, B_i, M_i, \alpha_i, \mathfrak q_i)_{i = 2, \ldots, n}$
est un d\'evissage complet de $\Coker(\alpha_1)/B_1/R$ en
$\mathfrak q_1$. L'hypothèse de récurrence
entraîne donc que $\Coker(\alpha_1)$ est de présentation finie comme
$B_1$-module, libre comme $R$-module et possède une décomposition comme dans le lemme.
Il en résulte que $M_1$ est de présentation finie comme $B_1$-module, voir
Algèbre, Lemme \ref{algebra-lemma-extension}.
Il en résulte aussi que
$M_1 \cong B_1^{\oplus r_1} \oplus \Coker(\alpha_1)$
comme $R$-modules, d'où une décomposition comme dans le lemme.
Enfin, $B_1$ est projectif comme $R$-module d'après le
Lemme \ref{lemma-fibres-irreducible-flat-projective-nonnoetherian},
donc libre comme $R$-module d'après
Algèbre, Théorème \ref{algebra-theorem-projective-free-over-local-ring}.
Cela achève la démonstration.
\end{proof}

\begin{proposition}
\label{proposition-finite-type-flat-at-point}
Soit $f : X \to S$ un morphisme de schémas.
Soit $\mathcal{F}$ un faisceau quasi-cohérent sur $X$.
Soit $x \in X$ d'image $s \in S$.
Supposons que
\begin{enumerate}
\item $f$ soit localement de présentation finie,
\item $\mathcal{F}$ soit de type fini, et
\item $\mathcal{F}$ soit plat en $x$ sur $S$.
\end{enumerate}
Il existe alors un voisinage \'etale élémentaire $(S', s') \to (S, s)$
et un sous-schéma ouvert
$$
V \subset X \times_S \Spec(\mathcal{O}_{S', s'})
$$
qui contient l'unique point de
$X \times_S \Spec(\mathcal{O}_{S', s'})$ s'envoyant sur $x$,
tels que l'image réciproque de $\mathcal{F}$ sur $V$ soit un $\mathcal{O}_V$-module
de présentation finie et plat sur $\mathcal{O}_{S', s'}$.
\end{proposition}

\begin{proof}[Première démonstration]
Cette démonstration est plus longue, mais n'utilise pas l'existence d'un d\'evissage complet.
Le problème est local au voisinage de $x$ et de $s$ ; on peut donc supposer que $X$
et $S$ sont affines. Au cours de la démonstration, on remplacera un nombre fini de fois
$S$ par un voisinage \'etale élémentaire de $(S, s)$. Le but sera alors de trouver,
après un tel remplacement, un ouvert
$V \subset X \times_S \Spec(\mathcal{O}_{S, s})$ contenant $x$
tel que $\mathcal{F}|_V$ soit plat sur $S$ et de présentation finie.
On peut évidemment aussi remplacer $S$ par $\Spec(\mathcal{O}_{S, s})$
à tout moment de la démonstration, c'est-à-dire supposer que $S$ est un schéma local.
On démontrera la proposition par récurrence sur l'entier
$n = \dim_x(\text{Supp}(\mathcal{F}_s))$.

\medskip\noindent
On peut choisir
\begin{enumerate}
\item des voisinages \'etales élémentaires $g : (X', x') \to (X, x)$,
$e : (S', s') \to (S, s)$,
\item un diagramme commutatif
$$
\xymatrix{
X \ar[dd]_f & X' \ar[dd] \ar[l]^g & Z' \ar[l]^i \ar[d]^\pi \\
& & Y' \ar[d]^h \\
S & S' \ar[l]_e & S' \ar@{=}[l]
}
$$
\item un point $z' \in Z'$ tel que $i(z') = x'$, $y' = \pi(z')$, $h(y') = s'$,
\item un $\mathcal{O}_{Z'}$-module quasi-cohérent de type fini $\mathcal{G}$,
\end{enumerate}
comme dans
le Lemme \ref{lemma-elementary-devissage}.
Nous allons remplacer $S$ par $\Spec(\mathcal{O}_{S', s'})$, voir
les remarques du premier paragraphe de la démonstration. Considérons le diagramme
$$
\xymatrix{
X_{\mathcal{O}_{S', s'}} \ar[ddr]_f &
X'_{\mathcal{O}_{S', s'}} \ar[dd] \ar[l]^g &
Z'_{\mathcal{O}_{S', s'}} \ar[l]^i \ar[d]^\pi \\
& & Y'_{\mathcal{O}_{S', s'}} \ar[dl]^h \\
& \Spec(\mathcal{O}_{S', s'})
}
$$
On a ici effectué le changement de base des schémas $X', Z', Y'$ sur $S'$ suivant
$\Spec(\mathcal{O}_{S', s'}) \to S'$ et celui du schéma $X$ sur $S$ suivant
$\Spec(\mathcal{O}_{S', s'}) \to S$. Le morphisme
$g$ est encore \'etale, voir
le Lemme \ref{lemma-etale-at-point}.
Après avoir remplacé $X$ par $X_{\mathcal{O}_{S', s'}}$,
$X'$ par $X'_{\mathcal{O}_{S', s'}}$,
$Z'$ par $Z'_{\mathcal{O}_{S', s'}}$ et
$Y'$ par $Y'_{\mathcal{O}_{S', s'}}$,
on peut supposer que l'on dispose d'un diagramme comme dans le
Lemme \ref{lemma-elementary-devissage},
où, de plus, $S = S'$ est un schéma local de point fermé $s$. D'après les
Lemmes \ref{lemma-devissage-finite-presentation} et
\ref{lemma-devissage-flat},
le résultat pour $Y' \to S$, le faisceau $\pi_*\mathcal{G}$ et le
point $y'$ entraîne le résultat pour $X \to S$, $\mathcal{F}$ et $x$.
On peut donc supposer que $S$ est local et que $X \to S$ est un morphisme lisse
de schémas affines, à fibres géométriquement irréductibles de dimension $n$.

\medskip\noindent
Cas initial de la récurrence : $n = 0$.
Puisque $X \to S$ est lisse à fibres géométriquement
irréductibles de dimension $0$, on voit que $X \to S$ est une immersion
ouverte, voir
Descente, Lemme \ref{descent-lemma-universally-injective-etale-open-immersion}.
Puisque $S$ est local et que son point fermé appartient à l'image de $X \to S$,
on en conclut que $X = S$. Ainsi $\mathcal{F}$ correspond
à un $\mathcal{O}_{S, s}$-module plat de type fini. Dans ce cas, le résultat
découle de
Algèbre, Lemme \ref{algebra-lemma-finite-flat-local},
qui affirme que $\mathcal{F}$ est en fait libre de type fini.

\medskip\noindent
Étape de récurrence. Supposons le résultat établi lorsque la dimension
du support dans la fibre fermée est $< n$. Écrivons $S = \Spec(A)$,
$X = \Spec(B)$ et $\mathcal{F} = \widetilde{N}$ pour un certain $B$-module
$N$. Notons que $A$ est un anneau local ; désignons par $\mathfrak m$ son idéal maximal.
Alors $\mathfrak p = \mathfrak mB$ est l'unique idéal premier minimal au-dessus de
$\mathfrak m$, puisque $X \to S$ est à fibres géométriquement irréductibles. Enfin,
soit $\mathfrak q \subset B$ l'idéal premier correspondant à $x$. D'après le
Lemme \ref{lemma-induction-step},
on peut choisir un homomorphisme
$$
\alpha : B^{\oplus r} \to N
$$
tel que $\kappa(\mathfrak p)^{\oplus r} \to N \otimes_B \kappa(\mathfrak p)$
soit un isomorphisme. De plus, puisque $N_{\mathfrak q}$ est $A$-plat, le lemme
montre aussi que $\alpha$ est injectif et que
$\Coker(\alpha)_{\mathfrak q}$ est $A$-plat.
Posons $Q = \Coker(\alpha)$. Notons que le support de $Q/\mathfrak mQ$
ne contient pas $\mathfrak p$. On a donc certainement
$\dim_{\mathfrak q}(\text{Supp}(Q/\mathfrak mQ)) < n$.
En réunissant les propriétés connues de $Q$, on voit
que l'hypothèse de récurrence s'applique à $Q$. Il existe donc
un morphisme \'etale élémentaire $(S', s) \to (S, s)$ tel que la conclusion
soit vérifiée pour $Q \otimes_A A'$ sur $B \otimes_A A'$, où
$A' = \mathcal{O}_{S', s'}$. Après avoir remplacé $A$ par $A'$, on dispose d'une
suite exacte
$$
0 \to B^{\oplus r} \to N \to Q \to 0
$$
(on utilise ici l'injectivité de $\alpha$ indiquée plus haut)
de $B$-modules de type fini, et l'on obtient aussi un élément
$g \in B$, $g \not \in \mathfrak q$, tel que
$Q_g$ soit de présentation finie sur $B_g$ et plat sur $A$. La localisation
étant exacte, on voit que
$$
0 \to B_g^{\oplus r} \to N_g \to Q_g \to 0
$$
est encore exacte. Puisque $B_g$ et $Q_g$ sont plats sur $A$, on en conclut que
$N_g$ est plat sur $A$, voir
Algèbre, Lemme \ref{algebra-lemma-flat-ses},
et puisque $B_g$ et $Q_g$ sont de présentation finie sur $B_g$, il en est de même
de $N_g$, voir
Algèbre, Lemme \ref{algebra-lemma-extension}.
\end{proof}

\begin{proof}[Deuxième démonstration]
On applique la
Proposition \ref{proposition-existence-complete-at-x}
pour trouver un diagramme commutatif
$$
\xymatrix{
(X, x) \ar[d] & (X', x') \ar[l]^g \ar[d] \\
(S, s) & (S', s') \ar[l]
}
$$
de schémas pointés dont les flèches
horizontales sont des voisinages \'etales élémentaires
et tel que $g^*\mathcal{F}/X'/S'$ admette un
d\'evissage complet en $x$.
(En particulier, $S'$ et $X'$ sont affines.) D'après
Morphismes, Lemme \ref{morphisms-lemma-flat-permanence},
on voit que $g^*\mathcal{F}$ est plat en $x'$ sur $S$ et, d'après le
Lemme \ref{lemma-etale-flat-up-down},
on voit qu'il est plat en $x'$ sur $S'$. Au moyen de la
Remarque \ref{remark-same-notion},
on en déduit que
$$
\Gamma(X', g^*\mathcal{F})/
\Gamma(X', \mathcal{O}_{X'})/
\Gamma(S', \mathcal{O}_{S'})
$$
admet un d\'evissage complet en l'idéal premier de $\Gamma(X', \mathcal{O}_{X'})$
correspondant à $x'$. On peut effectuer le changement de base de ce
d\'evissage complet à l'anneau local $\mathcal{O}_{S', s'}$
de $\Gamma(S', \mathcal{O}_{S'})$ en l'idéal premier correspondant à
$s'$. Le
Lemme \ref{lemma-complete-devissage-flat-finite-type-module}
entraîne donc que
$$
\Gamma(X', \mathcal{F}')
\otimes_{\Gamma(S', \mathcal{O}_{S'})}
\mathcal{O}_{S', s'}
$$
est plat sur $\mathcal{O}_{S', s'}$ et de présentation finie sur
$\Gamma(X', \mathcal{O}_{X'})
\otimes_{\Gamma(S', \mathcal{O}_{S'})} \mathcal{O}_{S', s'}$.
Autrement dit, la restriction de $\mathcal{F}$ à
$X' \times_{S'} \Spec(\mathcal{O}_{S', s'})$
est de présentation finie et plate sur $\mathcal{O}_{S', s'}$.
Puisque le morphisme
$X' \times_{S'} \Spec(\mathcal{O}_{S', s'})
\to X \times_S \Spec(\mathcal{O}_{S', s'})$
est \'etale
(Lemme \ref{lemma-etale-at-point}),
son image $V \subset X \times_S \Spec(\mathcal{O}_{S', s'})$
est un sous-schéma ouvert et, par descente \'etale, la restriction
de $\mathcal{F}$ à $V$ est de présentation finie et plate sur
$\mathcal{O}_{S', s'}$. (Résultats utilisés :
Morphismes, Lemme \ref{morphisms-lemma-etale-open},
Descente, Lemme \ref{descent-lemma-finite-presentation-descends}, et
Morphismes, Lemme \ref{morphisms-lemma-flat-permanence}.)
\end{proof}

\begin{lemma}
\label{lemma-open-in-fibre-where-flat}
Soit $f : X \to S$ un morphisme de schémas localement de type fini.
Soit $\mathcal{F}$ un $\mathcal{O}_X$-module quasi-cohérent de type fini.
Soit $s \in S$. Alors l'ensemble
$$
\{x \in X_s \mid \mathcal{F} \text{ flat over }S\text{ at }x\}
$$
est ouvert dans la fibre $X_s$.
\end{lemma}

\begin{proof}
Supposons $x \in U$. Choisissons un voisinage \'etale élémentaire
$(S', s') \to (S, s)$ et un ouvert
$V \subset X \times_S \Spec(\mathcal{O}_{S', s'})$ comme dans la
Proposition \ref{proposition-finite-type-flat-at-point}.
Notons que $X_{s'} = X_s$, puisque $\kappa(s) = \kappa(s')$.
Si $x' \in V \cap X_{s'}$, alors l'image réciproque de $\mathcal{F}$ sur
$X \times_S S'$ est plate sur $S'$ en $x'$. Ainsi $\mathcal{F}$ est
plat en $x'$ sur $S$, voir
Morphismes, Lemme \ref{morphisms-lemma-flat-permanence}.
Autrement dit, $X_s \cap V \subset U$ est un voisinage ouvert
de $x$ dans $U$.
\end{proof}

\begin{lemma}
\label{lemma-finite-type-flat-at-point}
Soit $f : X \to S$ un morphisme de schémas.
Soit $\mathcal{F}$ un faisceau quasi-cohérent sur $X$.
Soit $x \in X$ d'image $s \in S$.
Supposons que
\begin{enumerate}
\item $f$ soit localement de type fini,
\item $\mathcal{F}$ soit de type fini, et
\item $\mathcal{F}$ soit plat en $x$ sur $S$.
\end{enumerate}
Il existe alors un voisinage \'etale élémentaire $(S', s') \to (S, s)$
et un sous-schéma ouvert
$$
V \subset X \times_S \Spec(\mathcal{O}_{S', s'})
$$
qui contient l'unique point de
$X \times_S \Spec(\mathcal{O}_{S', s'})$ s'envoyant sur $x$,
tels que l'image réciproque de $\mathcal{F}$ sur $V$ soit plate sur
$\mathcal{O}_{S', s'}$.
\end{lemma}

\begin{proof}
(La seule différence avec la
Proposition \ref{proposition-finite-type-flat-at-point}
est que l'on ne suppose pas $f$ de présentation finie.)
La question est locale sur $X$ et sur $S$ ; on peut donc supposer que $X$ et $S$
sont affines. Écrivons $X = \Spec(B)$, $S = \Spec(A)$ et
$B = A[x_1, \ldots, x_n]/I$. Autrement dit, on obtient une immersion fermée
$i : X \to \mathbf{A}^n_S$. Notons $t = i(x) \in \mathbf{A}^n_S$.
On peut appliquer la
Proposition \ref{proposition-finite-type-flat-at-point}
à $\mathbf{A}^n_S \to S$, au faisceau $i_*\mathcal{F}$
et au point $t$. On obtient un voisinage
\'etale élémentaire $(S', s') \to (S, s)$ et un sous-schéma ouvert
$$
W \subset \mathbf{A}^n_{\mathcal{O}_{S', s'}}
$$
tels que l'image réciproque de $i_*\mathcal{F}$ sur $W$ soit plate sur
$\mathcal{O}_{S', s'}$. Cela signifie que
$V := W \cap \big(X \times_S \Spec(\mathcal{O}_{S', s'})\big)$
est le sous-schéma ouvert cherché.
\end{proof}

\begin{lemma}
\label{lemma-finite-type-flat-along-fibre}
Soit $f : X \to S$ un morphisme de schémas.
Soit $\mathcal{F}$ un faisceau quasi-cohérent sur $X$.
Soit $s \in S$.
Supposons que
\begin{enumerate}
\item $f$ soit de présentation finie,
\item $\mathcal{F}$ soit de type fini, et
\item $\mathcal{F}$ soit plat sur $S$ en tout point de la fibre $X_s$.
\end{enumerate}
Il existe alors un voisinage \'etale élémentaire $(S', s') \to (S, s)$
et un sous-schéma ouvert
$$
V \subset X \times_S \Spec(\mathcal{O}_{S', s'})
$$
qui contient la fibre $X_s = X \times_S s'$, tels que l'image réciproque
de $\mathcal{F}$ sur $V$ soit un $\mathcal{O}_V$-module
de présentation finie et plat sur $\mathcal{O}_{S', s'}$.
\end{lemma}

\begin{proof}
Pour tout point $x \in X_s$, on peut utiliser la
Proposition \ref{proposition-finite-type-flat-at-point}
pour trouver un voisinage \'etale élémentaire $(S_x, s_x) \to (S, s)$
et un ouvert $V_x \subset X \times_S \Spec(\mathcal{O}_{S_x, s_x})$
tels que $x \in X_s = X \times_S s_x$ appartienne à $V_x$ et que
l'image réciproque de $\mathcal{F}$ sur $V_x$ soit un
$\mathcal{O}_{V_x}$-module de présentation finie et plat sur
$\mathcal{O}_{S_x, s_x}$. En particulier, on peut considérer la fibre
$(V_x)_{s_x}$ comme un voisinage ouvert de $x$ dans $X_s$.
Puisque $X_s$ est quasi-compact, on peut trouver un nombre fini de points
$x_1, \ldots, x_n \in X_s$ tels que $X_s$ soit la réunion
des $(V_{x_i})_{s_{x_i}}$. Choisissons un voisinage \'etale élémentaire
$(S' , s') \to (S, s)$ qui domine chacun des voisinages
$(S_{x_i}, s_{x_i})$, voir
Compléments sur les morphismes,
Lemme \ref{more-morphisms-lemma-elementary-etale-neighbourhoods}.
Posons $V = \bigcup V_i$, où $V_i$ est l'image réciproque de l'ouvert
$V_{x_i}$ par le morphisme
$$
X \times_S \Spec(\mathcal{O}_{S', s'})
\longrightarrow
X \times_S \Spec(\mathcal{O}_{S_{x_i}, s_{x_i}})
$$
Par construction, $V$ contient $X_s$ et l'image réciproque
de $\mathcal{F}$ sur $V$ est un $\mathcal{O}_V$-module
de présentation finie et plat sur $\mathcal{O}_{S', s'}$.
\end{proof}

\begin{lemma}
\label{lemma-finite-type-flat-along-fibre-variant}
Soit $f : X \to S$ un morphisme de schémas.
Soit $\mathcal{F}$ un faisceau quasi-cohérent sur $X$.
Soit $s \in S$.
Supposons que
\begin{enumerate}
\item $f$ soit de type fini,
\item $\mathcal{F}$ soit de type fini, et
\item $\mathcal{F}$ soit plat sur $S$ en tout point de la fibre $X_s$.
\end{enumerate}
Il existe alors un voisinage \'etale élémentaire $(S', s') \to (S, s)$
et un sous-schéma ouvert
$$
V \subset X \times_S \Spec(\mathcal{O}_{S', s'})
$$
qui contient la fibre $X_s = X \times_S s'$, tels que l'image réciproque
de $\mathcal{F}$ sur $V$ soit plate sur $\mathcal{O}_{S', s'}$.
\end{lemma}

\begin{proof}
(La seule différence avec le
Lemme \ref{lemma-finite-type-flat-along-fibre}
est que l'on ne suppose pas $f$ de présentation finie.)
Pour tout point $x \in X_s$, on peut utiliser le
Lemme \ref{lemma-finite-type-flat-at-point}
pour trouver un voisinage \'etale élémentaire $(S_x, s_x) \to (S, s)$
et un ouvert $V_x \subset X \times_S \Spec(\mathcal{O}_{S_x, s_x})$
tels que $x \in X_s = X \times_S s_x$ appartienne à $V_x$ et que
l'image réciproque de $\mathcal{F}$ sur $V_x$ soit plate sur
$\mathcal{O}_{S_x, s_x}$. En particulier, on peut considérer la fibre
$(V_x)_{s_x}$ comme un voisinage ouvert de $x$ dans $X_s$.
Puisque $X_s$ est quasi-compact, on peut trouver un nombre fini de points
$x_1, \ldots, x_n \in X_s$ tels que $X_s$ soit la réunion
des $(V_{x_i})_{s_{x_i}}$. Choisissons un voisinage \'etale élémentaire
$(S' , s') \to (S, s)$ qui domine chacun des voisinages
$(S_{x_i}, s_{x_i})$, voir
Compléments sur les morphismes,
Lemme \ref{more-morphisms-lemma-elementary-etale-neighbourhoods}.
Posons $V = \bigcup V_i$, où $V_i$ est l'image réciproque de l'ouvert
$V_{x_i}$ par le morphisme
$$
X \times_S \Spec(\mathcal{O}_{S', s'})
\longrightarrow
X \times_S \Spec(\mathcal{O}_{S_{x_i}, s_{x_i}})
$$
Par construction, $V$ contient $X_s$ et l'image réciproque
de $\mathcal{F}$ sur $V$ est plate sur $\mathcal{O}_{S', s'}$.
\end{proof}

\begin{lemma}
\label{lemma-finite-type-flat-at-point-X}
Soit $S$ un schéma. Soit $X$ localement de type fini sur $S$.
Soit $x \in X$ d'image $s \in S$.
Si $X$ est plat en $x$ sur $S$, il existe alors un voisinage
\'etale élémentaire $(S', s') \to (S, s)$ et un sous-schéma ouvert
$$
V \subset X \times_S \Spec(\mathcal{O}_{S', s'})
$$
qui contient l'unique point de
$X \times_S \Spec(\mathcal{O}_{S', s'})$ s'envoyant sur $x$,
tels que $V \to \Spec(\mathcal{O}_{S', s'})$
soit plat et de présentation finie.
\end{lemma}

\begin{proof}
La question est locale sur $X$ et sur $S$ ; on peut donc supposer que $X$ et $S$
sont affines. Écrivons $X = \Spec(B)$, $S = \Spec(A)$ et
$B = A[x_1, \ldots, x_n]/I$. Autrement dit, on obtient une immersion fermée
$i : X \to \mathbf{A}^n_S$. Notons $t = i(x) \in \mathbf{A}^n_S$.
On peut appliquer la
Proposition \ref{proposition-finite-type-flat-at-point}
à $\mathbf{A}^n_S \to S$, au faisceau $\mathcal{F} = i_*\mathcal{O}_X$
et au point $t$. On obtient un voisinage
\'etale élémentaire $(S', s') \to (S, s)$ et un sous-schéma ouvert
$$
W \subset \mathbf{A}^n_{\mathcal{O}_{S', s'}}
$$
tels que l'image réciproque de $i_*\mathcal{O}_X$ soit plate et de présentation
finie. Cela signifie que
$V := W \cap \big(X \times_S \Spec(\mathcal{O}_{S', s'})\big)$
est le sous-schéma ouvert cherché.
\end{proof}

\begin{lemma}
\label{lemma-finite-type-flat-at-point-local}
Soit $f : X \to S$ un morphisme localement de présentation finie.
Soit $\mathcal{F}$ un $\mathcal{O}_X$-module quasi-cohérent de type fini.
Si $x \in X$ et si $\mathcal{F}$ est plat en $x$ sur $S$, alors
$\mathcal{F}_x$ est un $\mathcal{O}_{X, x}$-module de présentation finie.
\end{lemma}

\begin{proof}
Soit $s = f(x)$. D'après la
Proposition \ref{proposition-finite-type-flat-at-point},
il existe un voisinage \'etale élémentaire $(S', s') \to (S, s)$
tel que l'image réciproque de $\mathcal{F}$ sur
$X \times_S \Spec(\mathcal{O}_{S', s'})$ soit de
présentation finie au voisinage du point $x' \in X_{s'} = X_s$
correspondant à $x$. L'homomorphisme d'anneaux
$$
\mathcal{O}_{X, x} \longrightarrow
\mathcal{O}_{X \times_S \Spec(\mathcal{O}_{S', s'}), x'}
=
\mathcal{O}_{X \times_S S', x'}
$$
est plat et local, comme localisé d'un homomorphisme d'anneaux \'etale. Par suite,
$\mathcal{F}_x$ est de présentation finie sur $\mathcal{O}_{X, x}$
par descente, voir
Algèbre, Lemme \ref{algebra-lemma-descend-properties-modules}
(et le fait qu'un homomorphisme local plat est fidèlement plat, voir
Algèbre, Lemme \ref{algebra-lemma-local-flat-ff}).
\end{proof}

\begin{lemma}
\label{lemma-finite-type-flat-at-point-local-X}
Soit $f : X \to S$ un morphisme localement de type fini.
Soit $x \in X$ d'image $s \in S$. Si $f$ est plat en $x$ sur $S$, alors
$\mathcal{O}_{X, x}$ est essentiellement de présentation finie sur
$\mathcal{O}_{S, s}$.
\end{lemma}

\begin{proof}
On peut supposer $X$ et $S$ affines. Écrivons $X = \Spec(B)$,
$S = \Spec(A)$ et $B = A[x_1, \ldots, x_n]/I$.
Autrement dit, on obtient une immersion fermée $i : X \to \mathbf{A}^n_S$.
Notons $t = i(x) \in \mathbf{A}^n_S$. On peut appliquer
le Lemme \ref{lemma-finite-type-flat-at-point-local}
à $\mathbf{A}^n_S \to S$, au faisceau $\mathcal{F} = i_*\mathcal{O}_X$
et au point $t$. On en conclut que $\mathcal{O}_{X, x}$ est
de présentation finie sur $\mathcal{O}_{\mathbf{A}^n_S, t}$,
ce qui donne l'assertion cherchée.
\end{proof}







\section{Extension de propriétés à partir d'un ouvert}
\label{section-extending-properties}

\noindent
Nous réunissons dans cette section plusieurs résultats de la forme suivante : si $f : X \to S$
est un morphisme plat de schémas et si $f$ possède une certaine propriété au-dessus d'un
ouvert dense de $S$, alors $f$ possède cette même propriété au-dessus de tout $S$.

\begin{lemma}
\label{lemma-flat-finite-type-finitely-presented-over-dense-open}
\begin{slogan}
Les prolongements $S$-plats et de type fini de modules de présentation finie
sur un (bon) ouvert sont aussi de présentation finie sur $X$.
\end{slogan}
Soit $f : X \to S$ un morphisme de schémas. Soit $\mathcal{F}$ un
$\mathcal{O}_X$-module quasi-cohérent. Soit $U \subset S$ un ouvert.
Supposons que
\begin{enumerate}
\item $f$ soit localement de présentation finie,
\item $\mathcal{F}$ soit de type fini et plat sur $S$,
\item $U \subset S$ soit rétrocompact et schématiquement dense,
\item $\mathcal{F}|_{f^{-1}U}$ soit de présentation finie.
\end{enumerate}
Alors $\mathcal{F}$ est de présentation finie.
\end{lemma}

\begin{proof}
Le problème est local sur $X$ et sur $S$ ; on peut donc supposer $X$ et $S$ affines.
Écrivons $S = \Spec(A)$ et $X = \Spec(B)$. Soit $N$ un $B$-module de type fini tel
que $\mathcal{F}$ soit le faisceau quasi-cohérent associé à $N$.
On a $U = D(f_1) \cup \ldots \cup D(f_n)$ pour certains $f_i \in A$, voir
Algèbre, Lemme \ref{algebra-lemma-qc-open}.
Puisque $U$ est schématiquement dense, l'homomorphisme
$A \to A_{f_1} \times \ldots \times A_{f_n}$ est injectif.
Choisissons un idéal premier $\mathfrak q \subset B$ au-dessus de $\mathfrak p \subset A$,
correspondant à $x \in X$ s'envoyant sur $s \in S$.
D'après le Lemme \ref{lemma-finite-type-flat-at-point-local},
le module $N_\mathfrak q$ est de présentation finie sur $B_\mathfrak q$.
Choisissons une surjection $\varphi : B^{\oplus m} \to N$ de $B$-modules.
Choisissons $k_1, \ldots, k_t \in \Ker(\varphi)$ et
posons $N' = B^{\oplus m}/\sum Bk_j$. Il existe une surjection canonique
$N' \to N$, et $N$ est la limite inductive filtrante des $B$-modules $N'$
ainsi construits. On voit donc que l'on peut choisir
$k_1, \ldots, k_t$ de sorte que (a) $N'_{f_i} \cong N_{f_i}$, $i = 1, \ldots, n$,
et (b) $N'_\mathfrak q \cong N_\mathfrak q$.
Cela entraîne en particulier que $N'_\mathfrak q$ est plat sur $A$.
Par ouverture du lieu de platitude, voir
Algèbre, Théorème \ref{algebra-theorem-openness-flatness},
on en conclut qu'il existe $g \in B$, $g \not \in \mathfrak q$,
tel que $N'_g$ soit plat sur $A$. Considérons le diagramme commutatif
$$
\xymatrix{
N'_g \ar[r] \ar[d] & N_g \ar[d] \\
\prod N'_{gf_i} \ar[r] & \prod N_{gf_i}
}
$$
La flèche inférieure est un isomorphisme par le choix de $k_1, \ldots, k_t$.
La flèche verticale de gauche est injective, puisque
$A \to \prod A_{f_i}$ est injectif et que $N'_g$ est plat sur $A$.
La flèche horizontale supérieure est donc injective, donc un isomorphisme.
Cela démontre que $N_g$ est de présentation finie sur $B_g$.
On conclut en appliquant
Algèbre, Lemme \ref{algebra-lemma-cover}.
\end{proof}

\begin{lemma}
\label{lemma-flat-finite-type-finitely-presented-over-dense-open-X}
Soit $f : X \to S$ un morphisme de schémas. Soit $U \subset S$ un ouvert.
Supposons que
\begin{enumerate}
\item $f$ soit localement de type fini et plat,
\item $U \subset S$ soit rétrocompact et schématiquement dense,
\item $f|_{f^{-1}U} : f^{-1}U \to U$ soit localement de présentation finie.
\end{enumerate}
Alors $f$ est localement de présentation finie.
\end{lemma}

\begin{proof}
La question est locale sur $X$ et sur $S$ ; on peut donc supposer $X$ et
$S$ affines. Choisissons une immersion fermée $i : X \to \mathbf{A}^n_S$
et appliquons
le Lemme \ref{lemma-flat-finite-type-finitely-presented-over-dense-open}
à $i_*\mathcal{O}_X$. Nous omettons certains détails.
\end{proof}

\begin{lemma}
\label{lemma-flat-finite-presentation-dimension-over-dense-open}
Soit $f : X \to S$ un morphisme de schémas plat et localement
de type fini. Soit $U \subset S$ un ouvert dense tel que
$X_U \to U$ soit de dimension relative $\leq e$, voir
Morphismes, Définition \ref{morphisms-definition-relative-dimension-d}.
Si, de plus, l'une des conditions suivantes est vérifiée :
\begin{enumerate}
\item $f$ est localement de présentation finie, ou
\item $U \subset S$ est rétrocompact,
\end{enumerate}
alors $f$ est de dimension relative $\leq e$.
\end{lemma}

\begin{proof}
Démonstration dans le cas (1). Soit $W \subset X$ le sous-schéma ouvert construit
et étudié dans Compléments sur les morphismes, Lemmes
\ref{more-morphisms-lemma-flat-finite-presentation-CM-open} et
\ref{more-morphisms-lemma-flat-finite-presentation-CM-pieces}.
Notons que tout point générique de toute fibre appartient à $W$ ;
il suffit donc de démontrer le résultat pour $W$. Puisque
$W = \bigcup_{d \geq 0} U_d$, il suffit de montrer que $U_d = \emptyset$
pour $d > e$. Puisque $f$ est plat et localement de présentation finie, il est ouvert ;
ainsi $f(U_d)$ est ouvert (Morphismes, Lemme \ref{morphisms-lemma-fppf-open}).
Si $U_d$ est non vide, on a donc $f(U_d) \cap U \not = \emptyset$,
comme souhaité.

\medskip\noindent
Démonstration dans le cas (2). On peut remplacer $S$ par son réduit. Alors $U$ est
schématiquement dense. Par suite, $f$ est localement de présentation finie
d'après le Lemme \ref{lemma-flat-finite-type-finitely-presented-over-dense-open-X}.
On est ainsi ramené au cas (1).
\end{proof}

\begin{lemma}
\label{lemma-proper-flat-finite-over-dense-open}
Soit $f : X \to S$ un morphisme de schémas plat et propre.
Soit $U \subset S$ un ouvert dense tel que $X_U \to U$ soit fini.
Si, de plus, $f$ est localement de présentation finie ou si
$U \subset S$ est rétrocompact, alors $f$ est fini.
\end{lemma}

\begin{proof}
D'après le Lemme \ref{lemma-flat-finite-presentation-dimension-over-dense-open},
les fibres de $f$ sont de dimension nulle.
Ainsi $f$ est quasi-fini
(Morphismes, Lemme \ref{morphisms-lemma-locally-quasi-finite-rel-dimension-0}),
donc à fibres finies
(Morphismes, Lemme \ref{morphisms-lemma-quasi-finite}).
Par suite, $f$ est fini d'après
Compléments sur les morphismes, Lemme \ref{more-morphisms-lemma-characterize-finite}.
\end{proof}

\begin{lemma}
\label{lemma-zariski}
Soient $f : X \to S$ un morphisme de schémas et $U \subset S$ un ouvert. Si
\begin{enumerate}
\item $f$ est séparé, localement de type fini et plat,
\item $f^{-1}(U) \to U$ est un isomorphisme, et
\item $U \subset S$ est rétrocompact et schématiquement dense,
\end{enumerate}
alors $f$ est une immersion ouverte.
\end{lemma}

\begin{proof}
D'après le Lemme \ref{lemma-flat-finite-type-finitely-presented-over-dense-open-X},
le morphisme $f$ est localement de présentation finie.
L'image $f(X) \subset S$ est ouverte
(Morphismes, Lemme \ref{morphisms-lemma-fppf-open}) ;
on peut donc remplacer $S$ par $f(X)$. Il faut ainsi démontrer que
$f$ est un isomorphisme. On peut supposer $S$ affine. On peut se ramener
au cas où $X$ est quasi-compact, car il suffit de montrer
que tout ouvert quasi-compact $X' \subset X$ dont l'image est $S$
s'envoie isomorphiquement sur $S$. On peut donc supposer $f$ quasi-compact.
Toutes les fibres de $f$ sont de dimension $0$, voir
le Lemme \ref{lemma-flat-finite-presentation-dimension-over-dense-open}.
Ainsi $f$ est quasi-fini, voir
Morphismes, Lemme \ref{morphisms-lemma-locally-quasi-finite-rel-dimension-0}.
Soit $s \in S$. Choisissons un voisinage \'etale
élémentaire $g : (T, t) \to (S, s)$ tel que $X \times_S T = V \amalg W$
avec $V \to T$ fini et $W_t = \emptyset$, voir
Compléments sur les morphismes, Lemme
\ref{more-morphisms-lemma-etale-splits-off-quasi-finite-part}.
Notons $\pi : V \amalg W \to T$ le morphisme donné. Puisque $\pi$ est
plat et localement de présentation finie, on voit que $\pi(V)$ est
ouvert dans $T$ (Morphismes, Lemme \ref{morphisms-lemma-fppf-open}).
Quitte à restreindre $T$, on peut supposer que $T = \pi(V)$.
Puisque $f$ est un isomorphisme au-dessus de $U$, on voit que $\pi$ est un
isomorphisme au-dessus de $g^{-1}U$. Puisque $\pi(V) = T$, cela entraîne
que $\pi^{-1}g^{-1}U$ est contenu dans $V$.
D'après Morphismes, Lemme
\ref{morphisms-lemma-flat-morphism-scheme-theoretically-dense-open},
on voit que $\pi^{-1}g^{-1}U \subset V \amalg W$ est schématiquement
dense. On en déduit donc que $W = \emptyset$.
Ainsi $X \times_S T = V$ est fini sur $T$.
Il en résulte que $f$ est fini (après avoir remplacé $S$ par un
voisinage ouvert de $s$), par exemple d'après
Descente, Lemme \ref{descent-lemma-descending-property-finite}.
Alors $f$ est fini localement libre
(Morphismes, Lemme \ref{morphisms-lemma-finite-flat})
et, après avoir restreint $S$ à un voisinage ouvert plus petit de $s$,
on voit que $f$ est fini localement libre d'un certain degré $d$
(Morphismes, Lemme \ref{morphisms-lemma-finite-locally-free}).
Mais $d = 1$, puisque le degré vaut $1$ au-dessus
de l'ouvert dense $U$. Ainsi $f$ est un isomorphisme.
\end{proof}





\section{Modules plats de présentation finie}
\label{section-finitely-presented-flat}

\noindent
Dans certains cas, étant donnés un homomorphisme d'anneaux $R \to S$ de présentation finie et
un $S$-module de présentation finie $N$, la platitude de $N$ sur $R$ entraîne
que $N$ est projectif comme $R$-module, au moins après avoir remplacé $S$
par une extension \'etale. Nous réunissons dans cette section quelques résultats
de cette nature.

\begin{lemma}
\label{lemma-induction-step-fp}
Soit $R$ un anneau. Soit $R \to S$ un homomorphisme d'anneaux de présentation finie,
plat et à fibres géométriquement intègres. Soit
$\mathfrak q \subset S$ un idéal premier au-dessus de l'idéal premier
$\mathfrak r \subset R$. Posons $\mathfrak p = \mathfrak r S$.
Soit $N$ un $S$-module de présentation finie.
Il existe $r \geq 0$ et un homomorphisme de $S$-modules
$$
\alpha : S^{\oplus r} \longrightarrow N
$$
tels que
$\alpha : \kappa(\mathfrak p)^{\oplus r} \to N \otimes_S \kappa(\mathfrak p)$
soit un isomorphisme. Pour tout tel $\alpha$, les conditions suivantes sont équivalentes :
\begin{enumerate}
\item $N_{\mathfrak q}$ est $R$-plat,
\item il existe $f \in R$, $f \not \in \mathfrak r$, tel que
$\alpha_f : S_f^{\oplus r} \to N_f$ soit $R_f$-universellement injectif et
un $g \in S$, $g \not \in \mathfrak q$, tel que $\Coker(\alpha)_g$
soit $R$-plat,
\item $\alpha_{\mathfrak r}$ est $R_{\mathfrak r}$-universellement injectif et
$\Coker(\alpha)_{\mathfrak q}$ est $R$-plat,
\item $\alpha_{\mathfrak r}$ est injectif et
$\Coker(\alpha)_{\mathfrak q}$ est $R$-plat,
\item $\alpha_{\mathfrak p}$ est un isomorphisme et
$\Coker(\alpha)_{\mathfrak q}$ est $R$-plat, et
\item $\alpha_{\mathfrak q}$ est injectif et
$\Coker(\alpha)_{\mathfrak q}$ est $R$-plat.
\end{enumerate}
\end{lemma}

\begin{proof}
Pour obtenir $\alpha$, posons
$r = \dim_{\kappa(\mathfrak p)} N \otimes_S \kappa(\mathfrak p)$ et choisissons
$x_1, \ldots, x_r \in N$ dont les images forment une base de\newline
$N \otimes_S \kappa(\mathfrak p)$. Définissons
$\alpha(s_1, \ldots, s_r) = \sum s_i x_i$. Cela démontre l'existence.

\medskip\noindent
Fixons un choix de $\alpha$.
On peut appliquer
le Lemme \ref{lemma-induction-step}
à l'homomorphisme
$\alpha_{\mathfrak r} : S_{\mathfrak r}^{\oplus r} \to N_{\mathfrak r}$.
On voit donc que (1), (3), (4), (5) et (6) sont équivalentes.
Puisqu'il est aussi clair que (2) implique (3), il reste seulement
à montrer que (1) implique (2).

\medskip\noindent
Supposons (1). Par ouverture du lieu de platitude, voir
Algèbre, Théorème \ref{algebra-theorem-openness-flatness},
l'ensemble
$$
U_1 = \{\mathfrak q' \subset S \mid N_{\mathfrak q'}\text{ is flat over }R\}
$$
est ouvert dans $\Spec(S)$. Il contient $\mathfrak q$ par hypothèse,
donc $\mathfrak p$. Puisque $S^{\oplus r}$ et $N$ sont des
$S$-modules de présentation finie, l'ensemble
$$
U_2 = \{\mathfrak q' \subset S \mid
\alpha_{\mathfrak q'}\text{ is an isomorphism}\}
$$
est ouvert dans $\Spec(S)$, voir
Algèbre, Lemme \ref{algebra-lemma-map-between-finitely-presented}.
Il contient $\mathfrak p$ d'après (5). Puisque $R \to S$
est de présentation finie et plat, l'application
$\Phi : \Spec(S) \to \Spec(R)$ est ouverte, voir
Algèbre, Proposition \ref{algebra-proposition-fppf-open}.
Pour tout idéal premier $\mathfrak r' \in \Phi(U_1 \cap U_2)$, on voit
qu'il existe un idéal premier $\mathfrak q'$ au-dessus de $\mathfrak r'$ tel que
$N_{\mathfrak q'}$ soit plat et que $\alpha_{\mathfrak q'}$ soit
un isomorphisme, ce qui entraîne que $\alpha \otimes \kappa(\mathfrak p')$
est un isomorphisme, où $\mathfrak p' = \mathfrak r' S$. Ainsi
$\alpha_{\mathfrak r'}$ est $R_{\mathfrak r'}$-universellement injectif
par l'implication (1) $\Rightarrow$ (3).
Par suite, si l'on choisit $f \in R$, $f \not \in \mathfrak r$, tel que
$D(f) \subset \Phi(U_1 \cap U_2)$, on en conclut que
$\alpha_f$ est $R_f$-universellement injectif, voir
Algèbre, Lemme \ref{algebra-lemma-universally-injective-check-stalks}.
Le même raisonnement montre aussi que, pour tout
$\mathfrak q' \in U_1 \cap \Phi^{-1}(\Phi(U_1 \cap U_2))$,
le module $\Coker(\alpha)_{\mathfrak q'}$ est $R$-plat.
Notons que $\mathfrak q \in U_1 \cap \Phi^{-1}(\Phi(U_1 \cap U_2))$.
On peut donc trouver $g \in S$, $g \not \in \mathfrak q$, tel
que $D(g) \subset U_1 \cap \Phi^{-1}(\Phi(U_1 \cap U_2))$,
ce qui donne le résultat voulu.
\end{proof}

\begin{lemma}
\label{lemma-complete-devissage-flat-finitely-presented-module}
Soit $R \to S$ un homomorphisme d'anneaux de présentation finie.
Soit $N$ un $S$-module de présentation finie, plat sur $R$.
Soit $\mathfrak r \subset R$ un idéal premier.
Supposons qu'il existe un d\'evissage complet de $N/S/R$ au-dessus de $\mathfrak r$.
Il existe alors $f \in R$, $f \not \in \mathfrak r$,
tel que
$$
N_f \cong B_1^{\oplus r_1} \oplus \ldots \oplus B_n^{\oplus r_n}
$$
comme $R$-modules, où chaque $B_i$ est une $R_f$-algèbre lisse à fibres géométriquement
irréductibles. De plus, $N_f$ est projectif comme $R_f$-module.
\end{lemma}

\begin{proof}
Soit $(A_i, B_i, M_i, \alpha_i)_{i = 1, \ldots, n}$ le
d\'evissage complet donné. On démontre le lemme par récurrence sur $n$.
Notons que les assertions du lemme portent uniquement sur la structure
de $N$ comme $R$-module. On peut donc remplacer $N$ par $M_1$ et
considérer $M_1$ comme un $B_1$-module. Voir
la Remarque \ref{remark-finite-presentation}
pour vérifier que $M_1$ est de présentation finie comme $B_1$-module. D'après le
Lemme \ref{lemma-induction-step-fp},
on peut, après avoir remplacé $R$ par $R_f$ pour un certain
$f \in R$, $f \not \in \mathfrak r$, supposer que
l'homomorphisme $\alpha_1 : B_1^{\oplus r_1} \to M_1$ est $R$-universellement injectif.
Puisque $M_1$ et $B_1^{\oplus r_1}$ sont $R$-plats et de présentation finie comme
$B_1$-modules, on voit que $\Coker(\alpha_1)$ est $R$-plat
(Algèbre, Lemme \ref{algebra-lemma-ui-flat-domain})
et de présentation finie comme $B_1$-module. Notons que
$(A_i, B_i, M_i, \alpha_i)_{i = 2, \ldots, n}$ est un
d\'evissage complet de $\Coker(\alpha_1)$. L'hypothèse de récurrence
entraîne donc qu'après avoir remplacé
$R$ par $R_f$ pour un certain $f \in R$, $f \not \in \mathfrak r$,
on peut supposer que $\Coker(\alpha_1)$ possède une décomposition
comme dans le lemme et est projectif. En particulier,
$M_1 = B_1^{\oplus r_1} \oplus \Coker(\alpha_1)$.
Cela démontre l'assertion relative à la décomposition.
L'assertion de projectivité en résulte, puisque $B_1$ est projectif comme
$R$-module d'après le
Lemme \ref{lemma-fibres-irreducible-flat-projective-nonnoetherian}.
\end{proof}

\begin{remark}
\label{remark-complete-devissage-flat-finitely-presented-module}
Il existe une variante du
Lemme \ref{lemma-complete-devissage-flat-finitely-presented-module}
où l'on affaiblit la condition de platitude en supposant seulement que $N$
est plat en un idéal premier donné $\mathfrak q$ au-dessus de $\mathfrak r$,
mais où l'on renforce la condition de d\'evissage en supposant
l'existence d'un d\'evissage complet {\it en $\mathfrak q$}. Comparer avec le
Lemme \ref{lemma-complete-devissage-flat-finite-type-module}.
\end{remark}

\noindent
Voici le résultat principal de cette section.

\begin{proposition}
\label{proposition-finite-presentation-flat-at-point}
Soit $f : X \to S$ un morphisme de schémas.
Soit $\mathcal{F}$ un faisceau quasi-cohérent sur $X$.
Soit $x \in X$ d'image $s \in S$.
Supposons que
\begin{enumerate}
\item $f$ soit localement de présentation finie,
\item $\mathcal{F}$ soit de présentation finie, et
\item $\mathcal{F}$ soit plat en $x$ sur $S$.
\end{enumerate}
Il existe alors un diagramme commutatif de schémas pointés
$$
\xymatrix{
(X, x) \ar[d] & (X', x') \ar[l]^g \ar[d] \\
(S, s) & (S', s') \ar[l]
}
$$
dont les flèches horizontales sont des voisinages \'etales élémentaires,
tel que $X'$, $S'$ soient affines et que
$\Gamma(X', g^*\mathcal{F})$ soit un
$\Gamma(S', \mathcal{O}_{S'})$-module projectif.
\end{proposition}

\begin{proof}
Par ouverture du lieu de platitude, voir
Compléments sur les morphismes, Théorème \ref{more-morphisms-theorem-openness-flatness},
on peut remplacer $X$ par un voisinage ouvert de $x$ et supposer que
$\mathcal{F}$ est plat sur $S$. On applique ensuite la
Proposition \ref{proposition-existence-complete-at-x}
pour trouver un diagramme comme dans l'énoncé de la proposition,
tel que $g^*\mathcal{F}/X'/S'$ admette un d\'evissage complet au-dessus de $s'$.
(En particulier, $S'$ et $X'$ sont affines.) D'après
Morphismes, Lemme \ref{morphisms-lemma-flat-permanence},
on voit que $g^*\mathcal{F}$ est plat sur $S$ et, d'après le
Lemme \ref{lemma-etale-flat-up-down},
on voit qu'il est plat sur $S'$. Au moyen de la
Remarque \ref{remark-same-notion},
on en déduit que
$$
\Gamma(X', g^*\mathcal{F})/
\Gamma(X', \mathcal{O}_{X'})/
\Gamma(S', \mathcal{O}_{S'})
$$
admet un d\'evissage complet au-dessus de l'idéal premier de $\Gamma(S', \mathcal{O}_{S'})$
correspondant à $s'$. Le
Lemme \ref{lemma-complete-devissage-flat-finitely-presented-module}
entraîne donc que le résultat de la proposition est vérifié après avoir remplacé
$S'$ par un voisinage ouvert standard de $s'$.
\end{proof}

\noindent
Dans la suite de cette section, nous démontrons plusieurs variantes
de ce résultat. La première est une version « globale ».

\begin{lemma}
\label{lemma-finite-presentation-flat-along-fibre}
Soit $f : X \to S$ un morphisme de schémas.
Soit $\mathcal{F}$ un faisceau quasi-cohérent sur $X$.
Soit $s \in S$.
Supposons que
\begin{enumerate}
\item $f$ soit de présentation finie,
\item $\mathcal{F}$ soit de présentation finie, et
\item $\mathcal{F}$ soit plat sur $S$ en tout point de la fibre $X_s$.
\end{enumerate}
Il existe alors un voisinage \'etale élémentaire
$(S', s') \to (S, s)$ et un diagramme commutatif de schémas
$$
\xymatrix{
X \ar[d] & X' \ar[l]^g \ar[d] \\
S & S' \ar[l]
}
$$
tels que $g$ soit \'etale, $X_s \subset g(X')$, que les schémas
$X'$, $S'$ soient affines et que
$\Gamma(X', g^*\mathcal{F})$ soit un
$\Gamma(S', \mathcal{O}_{S'})$-module projectif.
\end{lemma}

\begin{proof}
Pour tout point $x \in X_s$, on peut utiliser la
Proposition \ref{proposition-finite-presentation-flat-at-point}
pour trouver un diagramme commutatif
$$
\xymatrix{
(X, x) \ar[d] & (Y_x, y_x) \ar[l]^{g_x} \ar[d] \\
(S, s) & (S_x, s_x) \ar[l]
}
$$
dont les flèches horizontales sont des voisinages \'etales élémentaires,
tel que $Y_x$, $S_x$ soient affines et que
$\Gamma(Y_x, g_x^*\mathcal{F})$ soit un
$\Gamma(S_x, \mathcal{O}_{S_x})$-module projectif. En particulier,
$g_x(Y_x) \cap X_s$ est un voisinage ouvert de $x$ dans $X_s$.
Puisque $X_s$ est quasi-compact, on peut trouver un nombre fini de points
$x_1, \ldots, x_n \in X_s$ tels que $X_s$ soit la réunion
des $g_{x_i}(Y_{x_i}) \cap X_s$. Choisissons un voisinage \'etale élémentaire
$(S' , s') \to (S, s)$ qui domine chacun des voisinages
$(S_{x_i}, s_{x_i})$, voir
Compléments sur les morphismes,
Lemme \ref{more-morphisms-lemma-elementary-etale-neighbourhoods}.
On peut aussi supposer $S'$ affine.
Posons $X' = \coprod Y_{x_i} \times_{S_{x_i}} S'$ et munissons-le du
morphisme évident $g : X' \to X$.
Par construction, $g(X')$ contient $X_s$ et
$$
\Gamma(X', g^*\mathcal{F})
=
\bigoplus \Gamma(Y_{x_i}, g_{x_i}^*\mathcal{F})
\otimes_{\Gamma(S_{x_i}, \mathcal{O}_{S_{x_i}})}
\Gamma(S', \mathcal{O}_{S'}).
$$
C'est un $\Gamma(S', \mathcal{O}_{S'})$-module projectif, voir
Algèbre, Lemme \ref{algebra-lemma-ascend-properties-modules}.
\end{proof}

\noindent
Les deux lemmes suivants reformulent les résultats
précédents lorsque $\mathcal{F} = \mathcal{O}_X$.

\begin{lemma}
\label{lemma-finite-presentation-flat-at-point-X}
Soit $f : X \to S$ localement de présentation finie.
Soit $x \in X$ d'image $s \in S$.
Si $f$ est plat en $x$ sur $S$, il existe alors un diagramme commutatif
de schémas pointés
$$
\xymatrix{
(X, x) \ar[d] & (X', x') \ar[l]^g \ar[d] \\
(S, s) & (S', s') \ar[l]
}
$$
dont les flèches horizontales sont des voisinages \'etales élémentaires,
tel que $X'$, $S'$ soient affines et que
$\Gamma(X', \mathcal{O}_{X'})$ soit un
$\Gamma(S', \mathcal{O}_{S'})$-module projectif.
\end{lemma}

\begin{proof}
C'est un cas particulier de la
Proposition \ref{proposition-finite-presentation-flat-at-point}.
\end{proof}

\begin{lemma}
\label{lemma-finite-presentation-flat-along-fibre-X}
Soit $f : X \to S$ de présentation finie.
Soit $s \in S$.
Si $X$ est plat sur $S$ en tout point de $X_s$,
il existe alors un voisinage \'etale élémentaire
$(S', s') \to (S, s)$ et un diagramme commutatif de schémas
$$
\xymatrix{
X \ar[d] & X' \ar[l]^g \ar[d] \\
S & S' \ar[l]
}
$$
avec $g$ \'etale, $X_s \subset g(X')$, tels que $X'$, $S'$
soient affines et que
$\Gamma(X', \mathcal{O}_{X'})$ soit un
$\Gamma(S', \mathcal{O}_{S'})$-module projectif.
\end{lemma}

\begin{proof}
C'est un cas particulier du
Lemme \ref{lemma-finite-presentation-flat-along-fibre}.
\end{proof}

\noindent
Les lemmes suivants donnent des conséquences de la
Proposition \ref{proposition-finite-presentation-flat-at-point}
lorsque l'on suppose seulement que le morphisme et le faisceau sont de type fini
(et pas nécessairement de présentation finie).

\begin{lemma}
\label{lemma-finite-type-flat-at-point-free}
Soit $f : X \to S$ un morphisme de schémas.
Soit $\mathcal{F}$ un faisceau quasi-cohérent sur $X$.
Soit $x \in X$ d'image $s \in S$.
Supposons que
\begin{enumerate}
\item $f$ soit localement de présentation finie,
\item $\mathcal{F}$ soit de type fini, et
\item $\mathcal{F}$ soit plat en $x$ sur $S$.
\end{enumerate}
Il existe alors un voisinage \'etale élémentaire $(S', s') \to (S, s)$
et un diagramme commutatif de schémas pointés
$$
\xymatrix{
(X, x) \ar[d] & (X', x') \ar[l]^g \ar[d] \\
(S, s) & (\Spec(\mathcal{O}_{S', s'}), s') \ar[l]
}
$$
tels que $X' \to X \times_S \Spec(\mathcal{O}_{S', s'})$
soit \'etale, $\kappa(x) = \kappa(x')$, que le schéma $X'$ soit
affine de présentation finie sur $\mathcal{O}_{S', s'}$,
que le faisceau $g^*\mathcal{F}$ soit de présentation finie sur $\mathcal{O}_{X'}$
et que $\Gamma(X', g^*\mathcal{F})$ soit un
$\mathcal{O}_{S', s'}$-module libre.
\end{lemma}

\begin{proof}
Pour démontrer le lemme, on peut remplacer $(S, s)$ par un voisinage \'etale
élémentaire quelconque, et l'on peut aussi remplacer $S$ par
$\Spec(\mathcal{O}_{S, s})$. Ainsi, d'après la
Proposition \ref{proposition-finite-type-flat-at-point},
on peut supposer que $\mathcal{F}$ est de présentation finie et plat sur
$S$ au voisinage de $x$. Dans ce cas, le résultat découle de la
Proposition \ref{proposition-finite-presentation-flat-at-point},
car
Algèbre, Théorème \ref{algebra-theorem-projective-free-over-local-ring},
assure que projectif $=$ libre sur un anneau local.
\end{proof}

\begin{lemma}
\label{lemma-finite-type-flat-at-point-free-variant}
Soit $f : X \to S$ un morphisme de schémas.
Soit $\mathcal{F}$ un faisceau quasi-cohérent sur $X$.
Soit $x \in X$ d'image $s \in S$.
Supposons que
\begin{enumerate}
\item $f$ soit localement de type fini,
\item $\mathcal{F}$ soit de type fini, et
\item $\mathcal{F}$ soit plat en $x$ sur $S$.
\end{enumerate}
Il existe alors un voisinage \'etale élémentaire $(S', s') \to (S, s)$
et un diagramme commutatif de schémas pointés
$$
\xymatrix{
(X, x) \ar[d] & (X', x') \ar[l]^g \ar[d] \\
(S, s) & (\Spec(\mathcal{O}_{S', s'}), s') \ar[l]
}
$$
tels que $X' \to X \times_S \Spec(\mathcal{O}_{S', s'})$
soit \'etale, $\kappa(x) = \kappa(x')$, que le schéma $X'$ soit
affine et que $\Gamma(X', g^*\mathcal{F})$ soit un
$\mathcal{O}_{S', s'}$-module libre.
\end{lemma}

\begin{proof}
(La seule différence avec le
Lemme \ref{lemma-finite-type-flat-at-point-free}
est que l'on ne suppose pas $f$ de présentation finie.)
Le problème est local sur $X$ et sur $S$. On peut donc supposer $X$ et
$S$ affines, et écrire $X = \Spec(B)$ et $S = \Spec(A)$.
Puisque $B$ est une $A$-algèbre de type fini, on peut trouver une surjection
$A[x_1, \ldots, x_n] \to B$. Autrement dit, on peut choisir une immersion
fermée $i : X \to \mathbf{A}^n_S$. Posons $t = i(x)$ et
$\mathcal{G} = i_*\mathcal{F}$. Notons que $\mathcal{G}_t \cong \mathcal{F}_x$
comme $\mathcal{O}_{S, s}$-modules. Ainsi $\mathcal{G}$ est plat sur $S$ en $t$.
On applique
le Lemme \ref{lemma-finite-type-flat-at-point-free}
au morphisme $\mathbf{A}^n_S \to S$, au point $t$ et au
faisceau $\mathcal{G}$. On peut donc trouver un
voisinage \'etale élémentaire $(S', s') \to (S, s)$
et un diagramme commutatif de schémas pointés
$$
\xymatrix{
(\mathbf{A}^n_S, t) \ar[d] & (Y, y) \ar[l]^h \ar[d] \\
(S, s) & (\Spec(\mathcal{O}_{S', s'}), s') \ar[l]
}
$$
tels que $Y \to \mathbf{A}^n_{\mathcal{O}_{S', s'}}$
soit \'etale, $\kappa(t) = \kappa(y)$, que le schéma $Y$ soit
affine et que $\Gamma(Y, h^*\mathcal{G})$ soit un
$\mathcal{O}_{S', s'}$-module projectif. Une solution du problème
initial est alors donnée par le sous-schéma fermé
$X' = Y \times_{\mathbf{A}^n_S} X$ de $Y$.
\end{proof}

\begin{lemma}
\label{lemma-finite-type-flat-along-fibre-free}
Soit $f : X \to S$ un morphisme de schémas.
Soit $\mathcal{F}$ un faisceau quasi-cohérent sur $X$.
Soit $s \in S$.
Supposons que
\begin{enumerate}
\item $f$ soit de présentation finie,
\item $\mathcal{F}$ soit de type fini, et
\item $\mathcal{F}$ soit plat sur $S$ en tout point de $X_s$.
\end{enumerate}
Il existe alors un voisinage \'etale élémentaire $(S', s') \to (S, s)$
et un diagramme commutatif de schémas
$$
\xymatrix{
X \ar[d] & X' \ar[l]^g \ar[d] \\
S & \Spec(\mathcal{O}_{S', s'}) \ar[l]
}
$$
tels que $X' \to X \times_S \Spec(\mathcal{O}_{S', s'})$
soit \'etale, $X_s = g((X')_{s'})$, que le schéma $X'$ soit
affine de présentation finie sur $\mathcal{O}_{S', s'}$,
que le faisceau $g^*\mathcal{F}$ soit de présentation finie sur $\mathcal{O}_{X'}$
et que $\Gamma(X', g^*\mathcal{F})$ soit un
$\mathcal{O}_{S', s'}$-module libre.
\end{lemma}

\begin{proof}
Pour tout point $x \in X_s$, on peut utiliser le
Lemme \ref{lemma-finite-type-flat-at-point-free}
pour trouver un voisinage \'etale élémentaire $(S_x , s_x) \to (S, s)$
et un diagramme commutatif
$$
\xymatrix{
(X, x) \ar[d] & (Y_x, y_x) \ar[l]^{g_x} \ar[d] \\
(S, s) & (\Spec(\mathcal{O}_{S_x, s_x}), s_x) \ar[l]
}
$$
tels que $Y_x \to X \times_S \Spec(\mathcal{O}_{S_x, s_x})$
soit \'etale, $\kappa(x) = \kappa(y_x)$, que le schéma $Y_x$ soit affine
de présentation finie sur $\mathcal{O}_{S_x, s_x}$, que le faisceau
$g_x^*\mathcal{F}$ soit de présentation finie sur $\mathcal{O}_{Y_x}$ et
que $\Gamma(Y_x, g_x^*\mathcal{F})$ soit un
$\mathcal{O}_{S_x, s_x}$-module libre. En particulier,
$g_x((Y_x)_{s_x})$ est un voisinage ouvert de $x$ dans $X_s$.
Puisque $X_s$ est quasi-compact, on peut trouver un nombre fini de points
$x_1, \ldots, x_n \in X_s$ tels que $X_s$ soit la réunion
des $g_{x_i}((Y_{x_i})_{s_{x_i}})$. Choisissons un voisinage \'etale élémentaire
$(S' , s') \to (S, s)$ qui domine chacun des voisinages
$(S_{x_i}, s_{x_i})$, voir
Compléments sur les morphismes,
Lemme \ref{more-morphisms-lemma-elementary-etale-neighbourhoods}.
Posons
$$
X' = \coprod Y_{x_i} \times_{\Spec(\mathcal{O}_{S_{x_i}, s_{x_i}})}
\Spec(\mathcal{O}_{S', s'})
$$
et munissons-le du morphisme évident $g : X' \to X$.
Par construction, $X_s = g(X'_{s'})$ et
$$
\Gamma(X', g^*\mathcal{F})
=
\bigoplus \Gamma(Y_{x_i}, g_{x_i}^*\mathcal{F})
\otimes_{\mathcal{O}_{S_{x_i}, s_{x_i}}}
\mathcal{O}_{S', s'}.
$$
C'est un $\mathcal{O}_{S', s'}$-module libre, comme somme directe
de changements de base de modules libres. Nous omettons quelques détails mineurs.
\end{proof}

\begin{lemma}
\label{lemma-finite-type-flat-along-fibre-free-variant}
Soit $f : X \to S$ un morphisme de schémas.
Soit $\mathcal{F}$ un faisceau quasi-cohérent sur $X$.
Soit $s \in S$.
Supposons que
\begin{enumerate}
\item $f$ soit de type fini,
\item $\mathcal{F}$ soit de type fini, et
\item $\mathcal{F}$ soit plat sur $S$ en tout point de $X_s$.
\end{enumerate}
Il existe alors un voisinage \'etale élémentaire $(S', s') \to (S, s)$
et un diagramme commutatif de schémas
$$
\xymatrix{
X \ar[d] & X' \ar[l]^g \ar[d] \\
S & \Spec(\mathcal{O}_{S', s'}) \ar[l]
}
$$
tels que $X' \to X \times_S \Spec(\mathcal{O}_{S', s'})$
soit \'etale, $X_s = g((X')_{s'})$, que le schéma $X'$ soit affine
et que $\Gamma(X', g^*\mathcal{F})$ soit un
$\mathcal{O}_{S', s'}$-module libre.
\end{lemma}

\begin{proof}
(La seule différence avec le
Lemme \ref{lemma-finite-type-flat-along-fibre-free}
est que l'on ne suppose pas $f$ de présentation finie.)
Pour tout point $x \in X_s$, on peut utiliser le
Lemme \ref{lemma-finite-type-flat-at-point-free-variant}
pour trouver un voisinage \'etale élémentaire $(S_x , s_x) \to (S, s)$
et un diagramme commutatif
$$
\xymatrix{
(X, x) \ar[d] & (Y_x, y_x) \ar[l]^{g_x} \ar[d] \\
(S, s) & (\Spec(\mathcal{O}_{S_x, s_x}), s_x) \ar[l]
}
$$
tels que $Y_x \to X \times_S \Spec(\mathcal{O}_{S_x, s_x})$
soit \'etale, $\kappa(x) = \kappa(y_x)$, que le schéma $Y_x$ soit affine et
que $\Gamma(Y_x, g_x^*\mathcal{F})$ soit un
$\mathcal{O}_{S_x, s_x}$-module libre. En particulier,
$g_x((Y_x)_{s_x})$ est un voisinage ouvert de $x$ dans $X_s$.
Puisque $X_s$ est quasi-compact, on peut trouver un nombre fini de points
$x_1, \ldots, x_n \in X_s$ tels que $X_s$ soit la réunion
des $g_{x_i}((Y_{x_i})_{s_{x_i}})$. Choisissons un voisinage \'etale élémentaire
$(S' , s') \to (S, s)$ qui domine chacun des voisinages
$(S_{x_i}, s_{x_i})$, voir
Compléments sur les morphismes,
Lemme \ref{more-morphisms-lemma-elementary-etale-neighbourhoods}.
Posons
$$
X' = \coprod Y_{x_i} \times_{\Spec(\mathcal{O}_{S_{x_i}, s_{x_i}})}
\Spec(\mathcal{O}_{S', s'})
$$
et munissons-le du morphisme évident $g : X' \to X$.
Par construction, $X_s = g(X'_{s'})$ et
$$
\Gamma(X', g^*\mathcal{F})
=
\bigoplus \Gamma(Y_{x_i}, g_{x_i}^*\mathcal{F})
\otimes_{\mathcal{O}_{S_{x_i}, s_{x_i}}}
\mathcal{O}_{S', s'}.
$$
C'est un $\mathcal{O}_{S', s'}$-module libre, comme somme directe
de changements de base de modules libres.
\end{proof}




\section{Modules plats de type fini, II}
\label{section-finite-type-flat-II}

\noindent
Nous aurons besoin du lemme suivant.

\begin{lemma}
\label{lemma-weak-bourbaki-pre-pre}
Soit $R \to S$ un homomorphisme d'anneaux de présentation finie. Soit $N$ un
$S$-module de présentation finie. Soit $\mathfrak q \subset S$ un idéal premier
au-dessus de $\mathfrak p \subset R$. Posons
$\overline{S} = S \otimes_R \kappa(\mathfrak p)$,
$\overline{\mathfrak q} = \mathfrak q \overline{S}$ et
$\overline{N} = N \otimes_R \kappa(\mathfrak p)$. Alors
il existe un $g \in S$ tel que
$g \not \in \mathfrak q$ et que
$\overline{g} \in \mathfrak r$ pour tout
$\mathfrak r \in \text{Ass}_{\overline{S}}(\overline{N})$
tel que $\mathfrak r \not \subset \overline{\mathfrak q}$.
\end{lemma}

\begin{proof}
En effet, si $\text{Ass}_{\overline{S}}(\overline{N}) =
\{\mathfrak r_1, \ldots, \mathfrak r_n\}$
(la finitude résulte d'Algèbre, lemme \ref{algebra-lemma-finite-ass}),
alors, quitte à renuméroter, nous pouvons supposer que
$$
\mathfrak r_1 \subset \overline{\mathfrak q},
\ldots,
\mathfrak r_r \subset \overline{\mathfrak q}, \quad
\mathfrak r_{r + 1} \not \subset \overline{\mathfrak q},
\ldots,
\mathfrak r_n \not \subset \overline{\mathfrak q}
$$
Comme $\overline{\mathfrak q}$ est un idéal premier, le produit
$\mathfrak r_{r + 1} \ldots \mathfrak r_n$ n'est pas contenu dans
$\overline{\mathfrak q}$ ; nous pouvons donc choisir un élément
$a$ de $\overline{S}$ appartenant à
$\mathfrak r_{r + 1}, \ldots, \mathfrak r_n$ mais non à
$\overline{\mathfrak q}$.
S'il existe $g \in S$ d'image $a$, alors $g$
convient. En général, nous pouvons trouver un élément non nul
$\lambda \in \kappa(\mathfrak p)$
tel que $\lambda a$ soit l'image d'un $g \in S$.
\end{proof}

\noindent
Le lemme suivant admet une variante légèrement plus forte,
le lemme \ref{lemma-weak-bourbaki}
ci-dessous.

\begin{lemma}
\label{lemma-weak-bourbaki-pre}
Soit $R \to S$ un homomorphisme d'anneaux de présentation finie.
Soit $N$ un $S$-module de présentation finie
qui est plat comme $R$-module. Soit $M$ un $R$-module.
Soit $\mathfrak q$ un idéal premier de $S$ au-dessus de $\mathfrak p \subset R$.
Alors
$$
\mathfrak q \in \text{WeakAss}_S(M \otimes_R N)
\Leftrightarrow
\Big(
\mathfrak p \in \text{WeakAss}_R(M)
\text{ and }
\overline{\mathfrak q} \in \text{Ass}_{\overline{S}}(\overline{N})
\Big)
$$
Ici, $\overline{S} = S \otimes_R \kappa(\mathfrak p)$,
$\overline{\mathfrak q} = \mathfrak q \overline{S}$ et
$\overline{N} = N \otimes_R \kappa(\mathfrak p)$.
\end{lemma}

\begin{proof}
Choisissons $g \in S$ comme au lemme \ref{lemma-weak-bourbaki-pre-pre}.
Appliquons la proposition \ref{proposition-finite-presentation-flat-at-point}
au morphisme de schémas $\Spec(S_g) \to \Spec(R)$, au module quasi-cohérent
associé à $N_g$ et aux points
correspondant aux idéaux premiers $\mathfrak qS_g$ et $\mathfrak p$. La traduction
algébrique donne un diagramme commutatif d'anneaux
$$
\xymatrix{
S \ar[r] & S_g \ar[r] & S' \\
& R \ar[lu] \ar[u] \ar[r] & R' \ar[u]
}
\quad\quad
\xymatrix{
\mathfrak q \ar@{-}[r] \ar@{-}[rd] &
\mathfrak qS_g \ar@{-}[d] \ar@{-}[r] & \mathfrak q' \ar@{-}[d] \\
& \mathfrak p \ar@{-}[r] & \mathfrak p'
}
$$
muni des idéaux premiers indiqués ; les flèches horizontales sont \'etales,
et $N \otimes_S S'$ est projectif comme $R'$-module. Posons
$N' = N \otimes_S S'$, $M' = M \otimes_R R'$,
$\overline{S}' = S' \otimes_{R'} \kappa(\mathfrak q')$,
$\overline{\mathfrak q}' = \mathfrak q' \overline{S}'$,
et
$$
\overline{N}' = N' \otimes_{R'} \kappa(\mathfrak p') =
\overline{N} \otimes_{\overline{S}} \overline{S}'
$$
D'après le lemme \ref{lemma-etale-weak-assassin-up-down}, nous avons
\begin{align*}
\text{WeakAss}_{S'}(M' \otimes_{R'} N') & =
(\Spec(S') \to \Spec(S))^{-1}\text{WeakAss}_S(M \otimes_R N) \\
\text{WeakAss}_{R'}(M') & =
(\Spec(R') \to \Spec(R))^{-1}\text{WeakAss}_R(M) \\
\text{Ass}_{\overline{S}'}(\overline{N}') & =
(\Spec(\overline{S}') \to \Spec(\overline{S}))^{-1}
\text{Ass}_{\overline{S}}(\overline{N})
\end{align*}
Utilisons Algèbre, lemme \ref{algebra-lemma-ass-weakly-ass},
pour $\overline{N}$ et $\overline{N}'$. En particulier, nous avons
\begin{align*}
\mathfrak q \in \text{WeakAss}_S(M \otimes_R N)
& \Leftrightarrow
\mathfrak q' \in \text{WeakAss}_{S'}(M' \otimes_{R'} N') \\
\mathfrak p \in \text{WeakAss}_R(M)
& \Leftrightarrow
\mathfrak p' \in \text{WeakAss}_{R'}(M') \\
\overline{\mathfrak q} \in \text{Ass}_{\overline{S}}(\overline{N})
& \Leftrightarrow
\overline{\mathfrak q}' \in \text{WeakAss}_{\overline{S}'}(\overline{N}')
\end{align*}
Notre choix précis de $g$ et la formule donnant
$\text{Ass}_{\overline{S}'}(\overline{N}')$ ci-dessus montrent que
\begin{equation}
\label{equation-key-observation}
\text{if }\mathfrak r' \in \text{Ass}_{\overline{S}'}(\overline{N}')
\text{ lies over }\mathfrak r \subset \overline{S}\text{ then }
\mathfrak r \subset \overline{\mathfrak q}
\end{equation}
Cette observation sera essentielle dans la suite de la démonstration. Nous utiliserons
la caractérisation des idéaux premiers faiblement associés donnée dans
Algèbre, lemme \ref{algebra-lemma-weakly-ass-local}, sans autre mention.

\medskip\noindent
Supposons que
$\overline{\mathfrak q} \not \in \text{Ass}_{\overline{S}}(\overline{N})$.
Alors
$\overline{\mathfrak q}' \not \in \text{Ass}_{\overline{S}'}(\overline{N}')$.
D'après
Algèbre, lemmes \ref{algebra-lemma-ass-zero-divisors},
\ref{algebra-lemma-finite-ass} et
\ref{algebra-lemma-silly},
il existe un élément $\overline{a}' \in \overline{\mathfrak q}'$
qui n'est pas diviseur de zéro sur $\overline{N}'$.
Quitte à remplacer $\overline{a}'$ par $\lambda \overline{a}'$ pour un élément non nul
$\lambda \in \kappa(\mathfrak p)$, nous pouvons trouver
$a' \in \mathfrak q'$ d'image $\overline{a}'$. D'après le
lemme \ref{lemma-invert-universally-injective},
l'application $a' : N'_{\mathfrak p'} \to N'_{\mathfrak p'}$ est
$R'_{\mathfrak p'}$-universellement injective. En particulier,
$a' : M' \otimes_{R'} N' \to M' \otimes_{R'} N'$ est
injective après localisation en $\mathfrak p'$, donc après
localisation en $\mathfrak q'$. Il en résulte clairement que
$\mathfrak q' \not \in \text{WeakAss}_{S'}(M' \otimes_{R'} N')$.
Nous concluons que $\mathfrak q \in \text{WeakAss}_S(M \otimes_R N)$ entraîne
$\overline{\mathfrak q} \in \text{Ass}_{\overline{S}}(\overline{N})$.

\medskip\noindent
Supposons $\mathfrak q \in \text{WeakAss}_S(M \otimes_R N)$. Nous voulons
montrer que $\mathfrak p \in \text{WeakAss}_S(M)$.
Soit $z \in M \otimes_R N$ un élément tel que $\mathfrak q$
soit minimal au-dessus de $J = \text{Ann}_S(z)$.
Soient $f_i \in \mathfrak p$, $i \in I$, des générateurs de
l'idéal $\mathfrak p$. Comme $\mathfrak q$ est au-dessus de $\mathfrak p$, pour tout $i$,
nous pouvons choisir $n_i \geq 1$ et $g_i \in S$, $g_i \not \in \mathfrak q$,
tels que $g_i f_i^{n_i} \in J$, c'est-à-dire $g_i f_i^{n_i} z = 0$.
Soit $z' \in (M' \otimes_{R'} N')_{\mathfrak p'}$ l'image de $z$.
Observons que $z'$ est non nul car $z$ a une image non nulle dans
$(M \otimes_R N)_\mathfrak q$ et car $S_\mathfrak q \to S'_{\mathfrak q'}$
est fidèlement plat. Nous affirmons que $f_i^{n_i} z' = 0$.

\medskip\noindent
Démontrons cette assertion : soit $g'_i \in S'$ l'image de $g_i$.
D'après l'observation essentielle (\ref{equation-key-observation}),
l'image $\overline{g}'_i \in \overline{S}'$
n'appartient à $\mathfrak r'$ pour aucun
$\mathfrak r' \in \text{Ass}_{\overline{S}'}(\overline{N})$.
Ainsi, d'après le lemme \ref{lemma-invert-universally-injective},
l'application $g'_i : N'_{\mathfrak p'} \to N'_{\mathfrak p'}$ est
$R'_{\mathfrak p'}$-universellement injective. En particulier,
$g'_i : M' \otimes_{R'} N' \to M' \otimes_{R'} N'$ est
injective après localisation en $\mathfrak p'$. L'assertion
en résulte puisque $g_i f_i^{n_i} z' = 0$.

\medskip\noindent
Cette assertion montre que l'annulateur de $z'$ dans $R'_{\mathfrak p'}$
contient les éléments $f_i^{n_i}$. Comme $R \to R'$ est \'etale, nous avons
$\mathfrak p'R'_{\mathfrak p'} = \mathfrak pR'_{\mathfrak p'}$
d'après Algèbre, lemme \ref{algebra-lemma-etale-at-prime}.
L'annulateur de $z'$ dans $R'_{\mathfrak p'}$ a donc pour radical
$\mathfrak p' R_{\mathfrak p'}$ (nous utilisons ici que $z'$ est non nul).
D'autre part,
$$
z' \in (M' \otimes_{R'} N')_{\mathfrak p'} =
M'_{\mathfrak p'} \otimes_{R'_{\mathfrak p'}} N'_{\mathfrak p'}
$$
Le module $N'_{\mathfrak p'}$ est projectif sur l'anneau local
$R'_{\mathfrak p'}$, donc libre
(Algèbre, théorème \ref{algebra-theorem-projective-free-over-local-ring}).
Nous pouvons donc trouver un facteur direct libre de type fini $F' \subset N'_{\mathfrak p'}$
tel que $z' \in M'_{\mathfrak p'} \otimes_{R'_{\mathfrak p'}} F'$.
Si $F'$ est de rang $n$, nous en déduisons que
$\mathfrak p' R'_{\mathfrak p'} \in
\text{WeakAss}_{R'_{\mathfrak p'}}({M'_{\mathfrak p'}}^{\oplus n})$.
Cela entraîne
$\mathfrak p'R'_{\mathfrak p'} \in \text{WeakAss}(M'_{\mathfrak p'})$,
par exemple d'après Algèbre, lemme \ref{algebra-lemma-weakly-ass}.
Alors $\mathfrak p' \in \text{WeakAss}_{R'}(M')$,
ce qui donne à son tour $\mathfrak p \in \text{WeakAss}_R(M)$.
Cela achève la démonstration de l'implication
``$\Rightarrow$'' de l'équivalence du lemme.

\medskip\noindent
Supposons que $\mathfrak p \in \text{WeakAss}_R(M)$ et
$\overline{\mathfrak q} \in \text{Ass}_{\overline{S}}(\overline{N})$.
Nous voulons montrer que $\mathfrak q$ est faiblement associé à $M \otimes_R N$.
Notons que $\overline{\mathfrak q}'$ est un élément maximal de
$\text{Ass}_{\overline{S}'}(\overline{N}')$.
Cela résulte de (\ref{equation-key-observation})
et du fait qu'il n'y a pas d'inclusions entre les idéaux premiers
de $\overline{S}'$ au-dessus de $\overline{\mathfrak q}$
(les fibres des morphismes \'etales étant discrètes,
Morphismes, lemme \ref{morphisms-lemma-etale-over-field}).
Ainsi, en remplaçant $R, S, \mathfrak p, \mathfrak q, M, N$ par
$R', S', \mathfrak p', \mathfrak q', M', N'$,
nous pouvons supposer, outre les hypothèses du lemme, que
\begin{enumerate}
\item $\mathfrak p \in \text{WeakAss}_R(M)$,
\item $\overline{\mathfrak q} \in \text{Ass}_{\overline{S}}(\overline{N})$,
\item $N$ est projectif comme $R$-module, et
\item $\overline{\mathfrak q}$ est maximal dans
$\text{Ass}_{\overline{S}}(\overline{N})$.
\end{enumerate}
Une réduction supplémentaire consiste à remplacer
$R, S, M, N$ par leurs localisés en $\mathfrak p$.
Cela fournit une condition supplémentaire, à savoir :
\begin{enumerate}
\item[(5)] $R$ est un anneau local d'idéal maximal $\mathfrak p$.
\end{enumerate}
Nous terminerons en montrant que (1) -- (5) entraînent
$\mathfrak q \in \text{WeakAss}(M \otimes_R N)$.

\medskip\noindent
Comme $R$ est local et $\mathfrak p \in \text{WeakAss}_R(M)$,
nous pouvons choisir un $y \in M$ dont l'annulateur $I$ a pour radical
$\mathfrak p$.
Écrivons $\overline{\mathfrak q} = (\overline{g}_1, \ldots, \overline{g}_n)$
pour des $\overline{g}_i \in \overline{S}$. Choisissons $g_i \in S$
d'image $\overline{g}_i$.
Alors $\mathfrak q = \mathfrak pS + g_1S + \ldots + g_nS$.
Considérons l'application
$$
\Psi : N/IN \longrightarrow (N/IN)^{\oplus n}, \quad
z \longmapsto (g_1z, \ldots, g_nz).
$$
C'est un homomorphisme de $R/I$-modules projectifs.
L'anneau local $R/I$ est auto-associé
(Compléments d'algèbre, définition \ref{more-algebra-definition-auto-ass})
car $\mathfrak p/I$ est localement nilpotent.
L'application $\Psi \otimes \kappa(\mathfrak p)$ n'est pas injective, car
$\overline{\mathfrak q} \in \text{Ass}_{\overline{S}}(\overline{N})$.
Ainsi, Compléments d'algèbre, lemme \ref{more-algebra-lemma-P-fPD-zero},
entraîne que $\Psi$ n'est pas injective. Choisissons $z \in N/IN$
non nul dans le noyau de $\Psi$. L'annulateur $J = \text{Ann}_S(z)$
contient $IS$ et $g_i$ par construction. Ainsi,
$\sqrt{J} \subset S$ contient $\mathfrak q$.
Soit $\mathfrak s \subset S$ un idéal premier minimal au-dessus de $J$.
Alors $\mathfrak q \subset \mathfrak s$,
$\mathfrak s$ est au-dessus de $\mathfrak p$, et
$\mathfrak s \in \text{WeakAss}_S(N/IN)$.
Ce dernier fait résulte de la définition des idéaux premiers faiblement associés.
Appliquons le sens ``$\Rightarrow$'' du lemme (déjà démontré)
à l'homomorphisme d'anneaux $R \to S$ et aux modules $R/I$ et $N$
pour en déduire que
$\overline{\mathfrak s} \in \text{Ass}_{\overline{S}}(\overline{N})$.
Comme $\overline{\mathfrak q} \subset \overline{\mathfrak s}$,
la maximalité de $\overline{\mathfrak q}$, voir la condition (4) ci-dessus,
entraîne que $\overline{\mathfrak q} = \overline{\mathfrak s}$.
Cela montre que $\mathfrak q = \mathfrak s$ et donne
la conclusion voulue.
\end{proof}

\begin{lemma}
\label{lemma-bourbaki-finite-type-general-base-at-point}
Soit $S$ un schéma.
Soit $f : X \to S$ localement de type fini.
Soit $x \in X$ d'image $s \in S$.
Soit $\mathcal{F}$ un faisceau quasi-cohérent de type fini sur $X$.
Soit $\mathcal{G}$ un faisceau quasi-cohérent sur $S$.
Si $\mathcal{F}$ est plat en $x$ sur $S$, alors
$$
x \in \text{WeakAss}_X(\mathcal{F} \otimes_{\mathcal{O}_X} f^*\mathcal{G})
\Leftrightarrow
s \in \text{WeakAss}_S(\mathcal{G})
\text{ and }
x \in \text{Ass}_{X_s}(\mathcal{F}_s).
$$
\end{lemma}

\begin{proof}
Dans ce paragraphe, nous nous ramenons au cas où $f$ est de présentation finie.
La question est locale sur $X$ et $S$ ; nous pouvons donc supposer $X$ et $S$
affines. Écrivons $X = \Spec(B)$, $S = \Spec(A)$ et
$B = A[x_1, \ldots, x_n]/I$. Autrement dit, nous obtenons une immersion fermée
$i : X \to \mathbf{A}^n_S$ sur $S$. Notons $t = i(x) \in \mathbf{A}^n_S$.
Notons que $i_*\mathcal{F}$ est un faisceau quasi-cohérent de type fini sur
$\mathbf{A}^n_S$, plat en $t$ sur $S$, et que
$$
i_*(\mathcal{F} \otimes_{\mathcal{O}_X} f^*\mathcal{G}) =
i_*\mathcal{F} \otimes_{\mathcal{O}_{\mathbf{A}^n_S}} p^*\mathcal{G}
$$
où $p : \mathbf{A}^n_S \to S$ est la projection. Notons que
$t$ est un point faiblement associé de
$i_*(\mathcal{F} \otimes_{\mathcal{O}_X} f^*\mathcal{G})$
si et seulement si $x$ est un point faiblement associé de
$\mathcal{F} \otimes_{\mathcal{O}_X} f^*\mathcal{G}$, voir
Diviseurs, lemme \ref{divisors-lemma-weakly-associated-finite}.
De même, $x \in \text{Ass}_{X_s}(\mathcal{F}_s)$ si et seulement
si $t \in \text{Ass}_{\mathbf{A}^n_s}((i_*\mathcal{F})_s)$ (voir
Algèbre, lemme \ref{algebra-lemma-ass-quotient-ring}).
Il suffit donc de démontrer le lemme lorsque $X = \mathbf{A}^n_S$.
Nous pouvons ainsi supposer que $X \to S$ est de présentation finie.

\medskip\noindent
Dans ce paragraphe, nous nous ramenons au cas où $\mathcal{F}$ est de présentation finie
et plat sur $S$.
Choisissons un voisinage \'etale élémentaire $e : (S', s') \to (S, s)$
et un ouvert $V \subset X \times_S \Spec(\mathcal{O}_{S', s'})$
comme dans la proposition \ref{proposition-finite-type-flat-at-point}.
Soit $x' \in X' = X \times_S S'$ l'unique point d'images $x$
et $s'$. Il suffit alors de démontrer l'énoncé pour
$X' \to S'$, $x'$, $s'$, $(X' \to X)^*\mathcal{F}$ et $e^*\mathcal{G}$, voir
le lemme \ref{lemma-etale-weak-assassin-up-down}.
Soit $v \in V$ l'unique point d'image $x'$
et soit $s' \in \Spec(\mathcal{O}_{S', s'})$ le point fermé.
Alors $\mathcal{O}_{V, v} = \mathcal{O}_{X', x'}$
et $\mathcal{O}_{\Spec(\mathcal{O}_{S', s'}), s'} =
\mathcal{O}_{S', s'}$, et de même pour les fibres des images réciproques de
$\mathcal{F}$ et $\mathcal{G}$.
De plus, $V_{s'} \subset X'_{s'}$ est un sous-schéma ouvert.
Comme la propriété d'être un point faiblement associé
ne dépend que de la fibre du faisceau, nous pouvons
remplacer
$X' \to S'$, $x'$, $s'$, $(X' \to X)^*\mathcal{F}$ et $e^*\mathcal{G}$
par
$V \to \Spec(\mathcal{O}_{S', s'})$, $v$, $s'$, $(V \to X)^*\mathcal{F}$
et $(\Spec(\mathcal{O}_{S', s'}) \to S)^*\mathcal{G}$.
Nous pouvons donc supposer que $f$ est de présentation finie et que
$\mathcal{F}$ est de présentation finie et plat sur $S$.

\medskip\noindent
Supposons $f$ de présentation finie et
$\mathcal{F}$ de présentation finie et plat sur $S$.
Après restriction de $X$ et $S$ à des voisinages affines
de $x$ et $s$, ce cas est traité par le
lemme \ref{lemma-weak-bourbaki-pre}.
\end{proof}

\begin{lemma}
\label{lemma-weak-bourbaki}
Soit $R \to S$ un homomorphisme d'anneaux essentiellement de type fini.
Soit $N$ un localisé d'un $S$-module de type fini, plat sur $R$.
Soit $M$ un $R$-module. Alors
$$
\text{WeakAss}_S(M \otimes_R N)
=
\bigcup\nolimits_{\mathfrak p \in \text{WeakAss}_R(M)}
\text{Ass}_{S \otimes_R \kappa(\mathfrak p)}(N \otimes_R \kappa(\mathfrak p))
$$
\end{lemma}

\begin{proof}
Ce lemme est la traduction
du lemme \ref{lemma-bourbaki-finite-type-general-base-at-point}
en termes algébriques. Les détails de cette traduction sont omis.
\end{proof}

\begin{lemma}
\label{lemma-bourbaki-finite-type-general-base}
Soit $f : X \to S$ un morphisme localement de type fini.
Soit $\mathcal{F}$ un faisceau quasi-cohérent de type fini sur $X$
qui est plat sur $S$. Soit $\mathcal{G}$ un faisceau quasi-cohérent sur $S$.
Alors nous avons
$$
\text{WeakAss}_X(\mathcal{F} \otimes_{\mathcal{O}_X} f^*\mathcal{G}) =
\bigcup\nolimits_{s \in \text{WeakAss}_S(\mathcal{G})}
\text{Ass}_{X_s}(\mathcal{F}_s)
$$
\end{lemma}

\begin{proof}
Conséquence immédiate du
lemme \ref{lemma-bourbaki-finite-type-general-base-at-point}.
\end{proof}

\begin{theorem}
\label{theorem-finite-type-flat}
\begin{slogan}
Le lieu de platitude est ouvert (version non noethérienne).
\end{slogan}
Soit $f : X \to S$ un morphisme de schémas.
Soit $\mathcal{F}$ un $\mathcal{O}_X$-module quasi-cohérent.
Supposons que
\begin{enumerate}
\item $X \to S$ soit localement de présentation finie,
\item $\mathcal{F}$ soit un $\mathcal{O}_X$-module de type fini, et
\item l'ensemble des points faiblement associés de $S$ soit localement fini dans $S$.
\end{enumerate}
Alors $U = \{x \in X \mid \mathcal{F}\text{ flat at }x\text{ over }S\}$
est ouvert dans $X$ et $\mathcal{F}|_U$ est un $\mathcal{O}_U$-module
de présentation finie et plat sur $S$.
\end{theorem}

\begin{proof}
Soit $x \in X$ tel que $\mathcal{F}$ soit plat en $x$ sur $S$.
Il faut trouver un voisinage ouvert de $x$ sur lequel $\mathcal{F}$ se restreigne
en un module de présentation finie plat sur $S$.
Le problème est local sur $X$ et $S$ ; nous pouvons donc supposer $X$ et $S$
affines. Comme $\mathcal{F}_x$ est un
$\mathcal{O}_{X, x}$-module de présentation finie d'après le
lemme \ref{lemma-finite-type-flat-at-point-local},
nous déduisons d'
Algèbre, lemme \ref{algebra-lemma-construct-fp-module-from-stalk},
qu'il existe un $\mathcal{O}_X$-module de présentation finie $\mathcal{F}'$
et une application $\varphi : \mathcal{F}' \to \mathcal{F}$ induisant
un isomorphisme $\varphi_x : \mathcal{F}'_x \to \mathcal{F}_x$. En particulier,
$\mathcal{F}'$ est plat sur $S$ en $x$ ; par ouverture
du lieu de platitude,
Compléments sur les morphismes, théorème \ref{more-morphisms-theorem-openness-flatness},
nous pouvons donc, après restriction de $X$, supposer que
$\mathcal{F}'$ est plat sur $S$. Comme $\mathcal{F}$ est de type fini,
nous pouvons, après restriction de $X$, supposer $\varphi$ surjective, voir
Modules, lemme \ref{modules-lemma-finite-type-surjective-on-stalk},
ou utiliser le lemme de Nakayama
(Algèbre, lemme \ref{algebra-lemma-NAK}).
D'après le
lemme \ref{lemma-bourbaki-finite-type-general-base},
nous avons
$$
\text{WeakAss}_X(\mathcal{F}') \subset
\bigcup\nolimits_{s \in \text{WeakAss}(S)} \text{Ass}_{X_s}(\mathcal{F}'_s)
$$
Comme $\text{WeakAss}(S)$ est fini par hypothèse et puisque
$\text{Ass}_{X_s}(\mathcal{F}'_s)$ est fini d'après
Diviseurs, lemme \ref{divisors-lemma-finite-ass},
nous concluons que $\text{WeakAss}_X(\mathcal{F}')$ est fini. À l'aide d'
Algèbre, lemme \ref{algebra-lemma-silly},
nous pouvons, après une nouvelle restriction de $X$, supposer que
$\text{WeakAss}_X(\mathcal{F}')$ est contenu dans l'ensemble des généralisations
de $x$. Considérons maintenant $\mathcal{K} = \Ker(\varphi)$. Nous avons
$\text{WeakAss}_X(\mathcal{K}) \subset \text{WeakAss}_X(\mathcal{F}')$
(d'après
Diviseurs, lemme \ref{divisors-lemma-ses-weakly-ass}),
mais, d'autre part, $\varphi_x$ est un isomorphisme, donc $\varphi_{x'}$
est aussi un isomorphisme pour tout $x' \leadsto x$. Nous concluons que
$\text{WeakAss}_X(\mathcal{K}) = \emptyset$, d'où
$\mathcal{K} = 0$ d'après
Diviseurs, lemme \ref{divisors-lemma-weakly-ass-zero}.
\end{proof}

\begin{lemma}
\label{lemma-finite-type-flat-algebra}
Soit $R \to S$ un homomorphisme d'anneaux de présentation finie.
Soit $M$ un $S$-module de type fini. Supposons $\text{WeakAss}_R(R)$ fini.
Alors
$$
U = \{\mathfrak q \subset S \mid M_{\mathfrak q}\text{ flat over }R\}
$$
est ouvert dans $\Spec(S)$ et, pour tout $g \in S$ tel que
$D(g) \subset U$, le localisé $M_g$ est un
$S_g$-module de présentation finie, plat sur $R$.
\end{lemma}

\begin{proof}
Cela résulte immédiatement du
théorème \ref{theorem-finite-type-flat}.
\end{proof}

\begin{lemma}
\label{lemma-finite-type-flat-X}
Soit $f : X \to S$ un morphisme de schémas localement de type
fini. Supposons l'ensemble des points faiblement associés de $S$ localement fini
dans $S$. Alors l'ensemble des points $x \in X$ où $f$ est plat est un sous-schéma
ouvert $U \subset X$ et $U \to S$ est plat et localement de présentation
finie.
\end{lemma}

\begin{proof}
Le problème est local sur $X$ et $S$ ; nous pouvons donc supposer
$X$ et $S$ affines. Alors $X \to S$ correspond à un homomorphisme d'anneaux
de type fini $A \to B$. Choisissons une surjection $A[x_1, \ldots, x_n] \to B$
et considérons $B$ comme un $A[x_1, \ldots, x_n]$-module. Une application du
lemme \ref{lemma-finite-type-flat-algebra}
achève la démonstration.
\end{proof}

\begin{lemma}
\label{lemma-finite-type-flat-over-integral}
Soit $f : X \to S$ un morphisme de schémas
localement de type fini et plat. Si $S$ est intègre, alors $f$
est localement de présentation finie.
\end{lemma}

\begin{proof}
Cas particulier du
lemme \ref{lemma-finite-type-flat-X}.
\end{proof}

\begin{proposition}
\label{proposition-flat-finite-type-finite-presentation-domain}
Soit $R$ un anneau intègre. Soit $R \to S$ un homomorphisme d'anneaux de type fini.
Soit $M$ un $S$-module de type fini.
\begin{enumerate}
\item Si $S$ est plat sur $R$, alors $S$ est une $R$-algèbre de présentation finie.
\item Si $M$ est plat comme $R$-module, alors $M$ est de présentation finie
comme $S$-module.
\end{enumerate}
\end{proposition}

\begin{proof}
La partie (1) est un cas particulier du
lemme \ref{lemma-finite-type-flat-over-integral}.
Pour la partie (2), choisissons une surjection $R[x_1, \ldots, x_n] \to S$.
D'après le lemme \ref{lemma-finite-type-flat-algebra}, $M$
est de présentation finie comme $R[x_1, \ldots, x_n]$-module.
Nous concluons par Algèbre, lemme
\ref{algebra-lemma-finitely-presented-over-subring}.
\end{proof}

\begin{lemma}[Version de type fini du théorème \ref{theorem-finite-type-flat}]
\label{lemma-finite-type-flat}
Soit $f : X \to S$ un morphisme de schémas.
Soit $\mathcal{F}$ un $\mathcal{O}_X$-module quasi-cohérent.
Supposons que
\begin{enumerate}
\item $X \to S$ soit localement de type fini,
\item $\mathcal{F}$ soit un $\mathcal{O}_X$-module de type fini, et
\item l'ensemble des points faiblement associés de $S$ soit localement fini dans $S$.
\end{enumerate}
Alors $U = \{x \in X \mid \mathcal{F}\text{ flat at }x\text{ over }S\}$
est ouvert dans $X$ et $\mathcal{F}|_U$ est plat sur $S$ et localement
de présentation finie relativement à $S$ (voir
Compléments sur les morphismes, définition
\ref{more-morphisms-definition-relatively-finitely-presented-sheaf}).
\end{lemma}

\begin{proof}
La question est locale sur $X$ et $S$. Nous pouvons donc supposer $X$ et $S$ affines.
Nous pouvons alors choisir une immersion fermée $i : X \to \mathbf{A}^n_S$.
Appliquons le théorème \ref{theorem-finite-type-flat} à $X' = \mathbf{A}^n_S \to S$
et au module quasi-cohérent $\mathcal{F}' = i_*\mathcal{F}$ de type fini ;
nous trouvons que
$$
U' = \{x' \in X' \mid \mathcal{F}'\text{ flat at }x'\text{ over }S\}
$$
est ouvert dans $X'$ et que $\mathcal{F}'|_{U'}$ est de présentation finie.
Comme $\mathcal{F}'$ se restreint à zéro sur $X' \setminus i(X)$ et
comme $\mathcal{F}'_{i(x)} \cong \mathcal{F}_x$ pour tout $x \in X$, nous voyons que
$$
U' = i(U) \amalg (X' \setminus i(X))
$$
Ainsi, $U = i^{-1}(U')$ est ouvert. De plus, il est clair que
$\mathcal{F}'|_{U'} = (i|_U)_*(\mathcal{F}|_U)$.
Nous en déduisons que $\mathcal{F}|_U$ est de présentation finie
relativement à $S$ d'après Compléments sur les morphismes, lemmes
\ref{more-morphisms-lemma-relative-finite-presentation} et
\ref{more-morphisms-lemma-finite-morphism-relative-finite-presentation}.
\end{proof}

\begin{lemma}
\label{lemma-finite-type-flat-algebra-bis}
Soit $R \to S$ un homomorphisme d'anneaux de type fini.
Soit $M$ un $S$-module de type fini.
Supposons $\text{WeakAss}_R(R)$ fini.
Alors
$$
U = \{\mathfrak q \subset S \mid M_{\mathfrak q}\text{ flat over }R\}
$$
est ouvert dans $\Spec(S)$ et, pour tout $g \in S$ tel que
$D(g) \subset U$, le localisé $M_g$ est plat sur $R$ et
est un $S_g$-module de présentation finie relativement à $R$ (voir
Compléments d'algèbre, définition
\ref{more-algebra-definition-relatively-finitely-presented}).
\end{lemma}

\begin{proof}
C'est le lemme \ref{lemma-finite-type-flat} traduit en termes algébriques.
\end{proof}








\section{Exemples de modules relativement purs}
\label{section-examples-pure-modules}

\noindent
Dans cette courte section, nous discutons quelques exemples de résultats qui motiveront
la notion de {\it module relativement pur} et le
concept d'{\it impureté} que nous introduirons plus loin. Chaque
exemple est énoncé sous forme de lemme. Notons l'analogie entre la condition sur
les idéaux premiers associés et les conditions figurant dans les
lemmes \ref{lemma-base-change-universally-flat},
\ref{lemma-universally-injective-to-completion},
\ref{lemma-universally-injective-to-completion-flat} et
\ref{lemma-flat-pure-over-complete-projective}.
Voir aussi
Algèbre, lemme \ref{algebra-lemma-compare-relative-assassins},
pour une discussion.

\begin{lemma}
\label{lemma-explain-why-pure}
Soit $R$ un anneau local d'idéal maximal $\mathfrak m$.
Soit $R \to S$ un homomorphisme d'anneaux. Soit $N$ un $S$-module.
Supposons que
\begin{enumerate}
\item $N$ soit projectif comme $R$-module, et
\item $S/\mathfrak mS$ soit noethérien et $N/\mathfrak mN$ soit un
$S/\mathfrak mS$-module de type fini.
\end{enumerate}
Alors, pour tout idéal premier $\mathfrak q \subset S$ associé à
$N \otimes_R \kappa(\mathfrak p)$, où $\mathfrak p = R \cap \mathfrak q$,
nous avons $\mathfrak q + \mathfrak m S \not = S$.
\end{lemma}

\begin{proof}
Notons que les hypothèses des
lemmes \ref{lemma-homothety-spectrum} et
\ref{lemma-invert-universally-injective}
sont satisfaites. Nous utiliserons leurs conclusions sans autre
mention. Soit $\Sigma \subset S$ l'ensemble multiplicatif des éléments
qui ne sont pas diviseurs de zéro sur $N/\mathfrak mN$. L'application
$N \to \Sigma^{-1}N$ est $R$-universellement injective. Ainsi,
tout $\mathfrak q \subset S$ qui est un idéal premier associé à
$N \otimes_R \kappa(\mathfrak p)$ est aussi un idéal premier associé à
$\Sigma^{-1}N \otimes_R \kappa(\mathfrak p)$. Il en résulte clairement que
$\mathfrak q$ correspond à un idéal premier de $\Sigma^{-1}S$.
Ainsi, $\mathfrak q \subset \mathfrak q'$, où $\mathfrak q'$
correspond à un idéal premier associé à $N/\mathfrak mN$, ce qui conclut.
\end{proof}

\noindent
Le lemme suivant donne un autre exemple, un peu trivial, de ce phénomène.

\begin{lemma}
\label{lemma-explain-why-pure-complete}
Soit $R$ un anneau. Soit $I \subset R$ un idéal.
Soit $R \to S$ un homomorphisme d'anneaux. Soit $N$ un $S$-module.
Si $N$ est complet pour la topologie $I$-adique, alors, pour tout $R$-module $M$ et
pour tout idéal premier $\mathfrak q \subset S$ associé à
$N \otimes_R M$, nous avons $\mathfrak q + I S \not = S$.
\end{lemma}

\begin{proof}
Notons $S^\wedge$ le complété $I$-adique de $S$.
Notons que $N$ est un $S^\wedge$-module, donc
$N \otimes_R M$ est aussi un $S^\wedge$-module.
Soit $z \in N \otimes_R M$ un élément tel que
$\mathfrak q = \text{Ann}_S(z)$. Comme $z \not = 0$, nous voyons
que $\text{Ann}_{S^\wedge}(z) \not = S^\wedge$. Ainsi,
$\mathfrak q S^\wedge \not = S^\wedge$. Il existe donc un
idéal maximal $\mathfrak m \subset S^\wedge$ tel que
$\mathfrak q S^\wedge \subset \mathfrak m$. Puisque
$IS^\wedge \subset \mathfrak m$ d'après
Algèbre, lemme \ref{algebra-lemma-radical-completion},
nous obtenons la conclusion.
\end{proof}

\noindent
Notons que le lemme suivant donne une autre démonstration du
lemme \ref{lemma-explain-why-pure},
car un module projectif sur un anneau local est libre, voir
Algèbre, théorème \ref{algebra-theorem-projective-free-over-local-ring}.

\begin{lemma}
\label{lemma-explain-why-pure-direct-sum-finite-modules}
Soit $R$ un anneau local d'idéal maximal $\mathfrak m$.
Soit $R \to S$ un homomorphisme d'anneaux. Soit $N$ un $S$-module.
Supposons $N$ isomorphe, comme $R$-module, à une somme
directe de $R$-modules de type fini. Alors, pour tout $R$-module $M$ et
pour tout idéal premier $\mathfrak q \subset S$ associé à
$N \otimes_R M$, nous avons $\mathfrak q + \mathfrak m S \not = S$.
\end{lemma}

\begin{proof}
Écrivons $N = \bigoplus_{i \in I} M_i$, chaque $M_i$ étant un $R$-module de type fini.
Soient $M$ un $R$-module et $\mathfrak q \subset S$ un idéal premier
associé à $N \otimes_R M$ tel que $\mathfrak q + \mathfrak m S = S$. Soit
$z \in N \otimes_R M$ un élément tel que $\mathfrak q = \text{Ann}_S(z)$.
En modifiant légèrement la décomposition en somme directe, nous pouvons supposer
$z \in M_1 \otimes_R M$ pour un élément $1 \in I$. Écrivons
$1 = f + \sum x_j g_j$ pour des $f \in \mathfrak q$, $x_j \in \mathfrak m$
et $g_j \in S$. Pour tout $g \in S$, notons $g'$ l'application $R$-linéaire
$$
M_1 \to N \xrightarrow{g} N \to M_1
$$
dont la première flèche est l'inclusion, la deuxième la multiplication
par $g$ et la troisième la projection. Comme chaque $x_j \in R$,
nous obtenons l'égalité
$$
f' + \sum x_j g'_j = \text{id}_{M_1} \in \text{End}_R(M_1)
$$
D'après le lemme de Nakayama
(Algèbre, lemme \ref{algebra-lemma-NAK}),
$f'$ est surjective ; ainsi, d'après
Algèbre, lemme \ref{algebra-lemma-fun},
$f'$ est un isomorphisme. En particulier, l'application
$$
M_1 \otimes_R M \to N \otimes_R M \xrightarrow{f} N \otimes_R M
\to M_1 \otimes_R M
$$
est un isomorphisme. Cela contredit l'hypothèse $fz = 0$.
\end{proof}

\begin{lemma}
\label{lemma-explain-why-pure-ML}
Soit $R$ un anneau local hensélien d'idéal maximal $\mathfrak m$.
Soit $R \to S$ un homomorphisme d'anneaux. Soit $N$ un $S$-module.
Supposons $N$ engendré par un ensemble dénombrable et de Mittag-Leffler comme $R$-module.
Alors, pour tout $R$-module $M$ et pour tout idéal premier $\mathfrak q \subset S$
associé à $N \otimes_R M$, nous avons
$\mathfrak q + \mathfrak m S \not = S$.
\end{lemma}

\begin{proof}
Ce lemme se ramène au
lemme \ref{lemma-explain-why-pure-direct-sum-finite-modules}
d'après
Algèbre, lemme \ref{algebra-lemma-split-ML-henselian}.
\end{proof}

\noindent
Supposons que $f : X \to S$ soit un morphisme de schémas et que
$\mathcal{F}$ soit un module quasi-cohérent sur $X$.
Soient $\xi \in \text{Ass}_{X/S}(\mathcal{F})$ et $Z = \overline{\{\xi\}}$.
Considérons le diagramme
$$
\xymatrix{
\xi \ar@{|->}[d] & Z \ar[r] \ar[rd] & X \ar[d]^f \\
f(\xi) & & S
}
$$
Notons que $f(Z) \subset \overline{\{f(\xi)\}}$ et que $f(Z)$ est fermé
si et seulement s'il y a égalité, c'est-à-dire $f(Z) = \overline{\{f(\xi)\}}$.
Il résulte du
lemme \ref{lemma-explain-why-pure}
que, si $S$, $X$ sont affines, si les fibres $X_s$ sont noethériennes,
si $\mathcal{F}$ est de type fini et si $\Gamma(X, \mathcal{F})$
est un $\Gamma(S, \mathcal{O}_S)$-module projectif, alors
$f(Z) = \overline{\{f(\xi)\}}$ est une partie fermée.
Des énoncés analogues légèrement différents valent dans les cas décrits aux
lemmes \ref{lemma-explain-why-pure-complete},
\ref{lemma-explain-why-pure-direct-sum-finite-modules} et
\ref{lemma-explain-why-pure-ML}.




\section{Impuretés}
\label{section-impure}

\noindent
Nous voulons formaliser le phénomène dont nous avons donné des exemples dans la
section \ref{section-examples-pure-modules}
en termes de spécialisations des points de $\text{Ass}_{X/S}(\mathcal{F})$.
Nous voulons aussi travailler localement au voisinage d'un point $s \in S$. À cette fin,
nous introduisons les définitions suivantes.

\begin{situation}
\label{situation-pre-pure}
Ici, $S$, $X$ sont des schémas et $f : X \to S$ est un morphisme de type fini.
De plus, $\mathcal{F}$ est un $\mathcal{O}_X$-module quasi-cohérent de type fini.
Enfin, $s$ est un point de $S$.
\end{situation}

\noindent
Dans cette situation, considérons un morphisme $g : T \to S$, un point $t \in T$
tel que $g(t) = s$, une spécialisation $t' \leadsto t$ et un point
$\xi \in X_T$ du changement de base de $X$, au-dessus de $t'$. Diagramme :
\begin{equation}
\label{equation-impurity}
\vcenter{
\xymatrix{
\xi \ar@{|->}[d] & \\
t' \ar@{~>}[r] & t \ar@{|->}[r] & s
}
}
\quad\quad
\vcenter{
\xymatrix{
X_T \ar[d] \ar[r] & X \ar[d] \\
T \ar[r]^g & S
}
}
\end{equation}
Notons en outre $\mathcal{F}_T$ l'image réciproque de $\mathcal{F}$ sur $X_T$.

\begin{definition}
\label{definition-impurity}
Dans la
situation \ref{situation-pre-pure},
nous disons qu'un diagramme (\ref{equation-impurity}) définit une
{\it impureté de $\mathcal{F}$ au-dessus de $s$}
si $\xi \in \text{Ass}_{X_T/T}(\mathcal{F}_T)$ et
$\overline{\{\xi\}} \cap X_t = \emptyset$. Nous exprimerons
cela en disant : « soit $(g : T \to S, t' \leadsto t, \xi)$ une
impureté de $\mathcal{F}$ au-dessus de $s$ ».
\end{definition}

\begin{lemma}
\label{lemma-impure-finite-presentation}
Plaçons-nous dans la situation \ref{situation-pre-pure}.
S'il existe une impureté de $\mathcal{F}$ au-dessus de $s$, alors
il existe une impureté $(g : T \to S, t' \leadsto t, \xi)$
de $\mathcal{F}$ au-dessus de $s$ telle que $g$ soit localement de présentation
finie et $t$ soit un point fermé de la fibre de $g$ au-dessus de $s$.
\end{lemma}

\begin{proof}
Soit $(g : T \to S, t' \leadsto t, \xi)$ une impureté quelconque de
$\mathcal{F}$ au-dessus de $s$. Appliquons
Limites, lemme \ref{limits-lemma-separate},
à $t \in T$ et $Z = \overline{\{\xi\}}$ pour obtenir un voisinage ouvert
$V \subset T$ de $t$, un diagramme commutatif
$$
\xymatrix{
V \ar[d] \ar[r]_a & T' \ar[d]^b \\
T \ar[r]^g & S,
}
$$
et un sous-schéma fermé $Z' \subset X_{T'}$ tels que
\begin{enumerate}
\item le morphisme $b : T' \to S$ soit localement de présentation finie,
\item on ait $Z' \cap X_{a(t)} = \emptyset$, et
\item $Z \cap X_V$ soit envoyé dans $Z'$ par le morphisme $X_V \to X_{T'}$.
\end{enumerate}
Comme $t'$ se spécialise en $t$, nous pouvons remplacer $T$ par le voisinage ouvert
$V$ de $t$. Nous avons donc un diagramme commutatif
$$
\xymatrix{
X_T \ar[d] \ar[r] &
X_{T'} \ar[d] \ar[r] &
X \ar[d] \\
T \ar[r]^a & T' \ar[r]^b & S
}
$$
où $b \circ a = g$. Notons $\xi' \in X_{T'}$
l'image de $\xi$. D'après
Diviseurs, lemme \ref{divisors-lemma-base-change-relative-assassin},
nous avons $\xi' \in \text{Ass}_{X_{T'}/T'}(\mathcal{F}_{T'})$.
De plus, par construction, l'adhérence de $\overline{\{\xi'\}}$
est contenue dans la partie fermée $Z'$, qui ne rencontre pas la fibre
$X_{a(t)}$. Ainsi, $(T' \to S, a(t') \leadsto a(t), \xi')$
est une impureté de $\mathcal{F}$ au-dessus de $s$.

\medskip\noindent
Nous pouvons donc supposer $g : T \to S$ localement de présentation finie.
Soit $Z = \overline{\{\xi\}}$. Par hypothèse, $Z_t = \emptyset$. D'après
Compléments sur les morphismes, lemme \ref{more-morphisms-lemma-empty-generic-fibre},
cela signifie que $Z_{t''} = \emptyset$ pour $t''$ dans une partie ouverte
de $\overline{\{t\}}$. Comme la fibre de
$T \to S$ au-dessus de $s$ est un schéma de Jacobson, voir
Morphismes, lemme \ref{morphisms-lemma-ubiquity-Jacobson-schemes},
il existe un point fermé $t'' \in \overline{\{t\}}$ tel que
$Z_{t''} = \emptyset$. Alors $(g : T \to S, t' \leadsto t'', \xi)$ est
l'impureté cherchée.
\end{proof}

\begin{lemma}
\label{lemma-impure-limit}
Plaçons-nous dans la situation \ref{situation-pre-pure}.
Soit $(g : T \to S, t' \leadsto t, \xi)$ une impureté de
$\mathcal{F}$ au-dessus de $s$. Supposons que $T = \lim_{i \in I} T_i$
soit une limite dirigée de schémas affines sur $S$. Alors, pour
un certain $i$, le triplet $(T_i \to S, t'_i \leadsto t_i, \xi_i)$
est une impureté de $\mathcal{F}$ au-dessus de $s$.
\end{lemma}

\begin{proof}
Les notations de l'énoncé ont le sens suivant : soient $p_i : T \to T_i$
les projections, $t_i = p_i(t)$ et $t'_i = p_i(t')$.
Enfin, $\xi_i \in X_{T_i}$ est l'image de $\xi$. D'après
Diviseurs, lemme \ref{divisors-lemma-base-change-relative-assassin},
$\xi_i$ est bien un point de l'assassin
relatif de $\mathcal{F}_{T_i}$ sur $T_i$. Il reste donc seulement à
montrer que $\overline{\{\xi_i\}} \cap X_{t_i} = \emptyset$ pour un certain $i$.

\medskip\noindent
Première démonstration. Soient $Z_i = \overline{\{\xi_i\}} \subset X_{T_i}$
et $Z = \overline{\{\xi\}} \subset X_T$
munis des structures de sous-schémas réduits induites.
Alors $Z = \lim Z_i$ d'après
Limites, lemme \ref{limits-lemma-inverse-limit-irreducibles}.
Choisissons un corps $k$ et un morphisme $\Spec(k) \to T$ d'image $t$.
Alors
$$
\emptyset =
Z \times_T \Spec(k) = (\lim Z_i) \times_{(\lim T_i)} \Spec(k)
= \lim Z_i \times_{T_i} \Spec(k)
$$
car les limites commutent aux produits fibrés (les limites commutent aux limites).
Chaque $Z_i \times_{T_i} \Spec(k)$ est quasi-compact, car $X_{T_i} \to T_i$
est de type fini, donc $Z_i \to T_i$ est de type fini.
Ainsi, $Z_i \times_{T_i} \Spec(k)$ est vide pour un certain $i$ d'après
Limites, lemme \ref{limits-lemma-limit-nonempty}.
Comme l'image du composé $\Spec(k) \to T \to T_i$ est $t_i$,
nous obtenons la conclusion voulue.

\medskip\noindent
Deuxième démonstration. Posons $Z = \overline{\{\xi\}}$. Appliquons
Limites, lemme \ref{limits-lemma-separate},
à cette situation pour obtenir un voisinage ouvert
$V \subset T$ de $t$, un diagramme commutatif
$$
\xymatrix{
V \ar[d] \ar[r]_a & T' \ar[d]^b \\
T \ar[r]^g & S,
}
$$
et un sous-schéma fermé $Z' \subset X_{T'}$ tels que
\begin{enumerate}
\item le morphisme $b : T' \to S$ soit localement de présentation finie,
\item on ait $Z' \cap X_{a(t)} = \emptyset$, et
\item $Z \cap X_V$ soit envoyé dans $Z'$ par le morphisme $X_V \to X_{T'}$.
\end{enumerate}
Nous pouvons supposer $V$ ouvert affine de $T$ ; d'après
Limites, lemmes \ref{limits-lemma-descend-opens} et
\ref{limits-lemma-limit-affine},
nous pouvons donc trouver un $i$ et un ouvert affine $V_i \subset T_i$ tels que
$V = f_i^{-1}(V_i)$. D'après
Limites,
proposition \ref{limits-proposition-characterize-locally-finite-presentation},
nous pouvons, quitte à augmenter $i$, trouver un morphisme
$a_i : V_i \to T'$ tel que $a = a_i \circ f_i|_V$.
Le morphisme induit $X_{V_i} \to X_{T'}$ envoie $\xi_i$ dans
$Z'$. Comme $Z' \cap X_{a(t)} = \emptyset$, nous concluons que
$(T_i \to S, t'_i \leadsto t_i, \xi_i)$ est une impureté de
$\mathcal{F}$ au-dessus de $s$.
\end{proof}

\begin{lemma}
\label{lemma-quasi-finite-impurity-elementary}
Plaçons-nous dans la situation \ref{situation-pre-pure}.
S'il existe une impureté $(g : T \to S, t' \leadsto t, \xi)$
de $\mathcal{F}$ au-dessus de $s$ avec $g$ quasi-fini en $t$, alors il
existe une impureté $(g : T \to S, t' \leadsto t, \xi)$ telle que
$(T, t) \to (S, s)$ soit un voisinage \'etale élémentaire.
\end{lemma}

\begin{proof}
Soit $(g : T \to S, t' \leadsto t, \xi)$ une impureté de
$\mathcal{F}$ au-dessus de $s$ telle que $g$ soit quasi-fini en $t$.
Après restriction de $T$, nous pouvons supposer $g$ localement de type fini.
Appliquons
Compléments sur les morphismes,
lemme \ref{more-morphisms-lemma-etale-makes-quasi-finite-finite-at-point},
à $T \to S$ et $t \mapsto s$. Nous obtenons un diagramme
$$
\xymatrix{
T \ar[d] & T \times_S U \ar[l] \ar[d] & V \ar[l] \ar[ld] \\
S & U \ar[l]
}
$$
où $(U, u) \to (S, s)$ est un voisinage \'etale élémentaire
et $V \subset T \times_S U$ est un voisinage ouvert de $v = (t, u)$
tel que $V \to U$ soit fini et que $v$ soit l'unique point de $V$
au-dessus de $u$. Comme le morphisme $V \to T$ est \'etale,
donc plat, il existe une spécialisation $v' \leadsto v$ telle
que $v' \mapsto t'$. Notons que $\kappa(t') \subset \kappa(v')$
est une extension finie séparable. Choisissons un point quelconque $\zeta \in X_{v'}$ d'image
$\xi \in X_{t'}$. D'après
Diviseurs, lemme \ref{divisors-lemma-base-change-relative-assassin},
nous avons $\zeta \in \text{Ass}_{X_V/V}(\mathcal{F}_V)$.
De plus, l'adhérence $\overline{\{\zeta\}}$ ne rencontre pas
la fibre $X_v$, puisque, par hypothèse, l'adhérence $\overline{\{\xi\}}$
ne rencontre pas $X_t$. Autrement dit, $(V \to S, v' \leadsto v, \zeta)$
est une impureté de $\mathcal{F}$ au-dessus de $S$.

\medskip\noindent
Ensuite, soient $u' \in U'$ l'image de $v'$ et
$\theta \in X_U$ l'image de $\zeta$.
Alors $\theta \mapsto u'$ et $u' \leadsto u$.
D'après
Diviseurs, lemme \ref{divisors-lemma-base-change-relative-assassin},
nous avons $\theta \in \text{Ass}_{X_U/U}(\mathcal{F})$.
De plus, comme $\pi : X_V \to X_U$ est fini, nous avons
$\pi\big(\overline{\{\zeta\}}\big) = \overline{\{\pi(\zeta)\}}$. Puisque
$v$ est l'unique point de $V$ au-dessus de $u$, nous avons
$X_u \cap \overline{\{\pi(\zeta)\}} = \emptyset$, car
$X_v \cap \overline{\{\zeta\}} = \emptyset$. Nous concluons ainsi que
$(U \to S, u' \leadsto u, \theta)$ est une impureté de
$\mathcal{F}$ au-dessus de $s$, ce qui achève la démonstration.
\end{proof}

\begin{lemma}
\label{lemma-Noetherian-impurity-quasi-finite}
Plaçons-nous dans la situation \ref{situation-pre-pure}.
Supposons $S$ localement noethérien.
S'il existe une impureté de $\mathcal{F}$ au-dessus de $s$, alors
il existe une impureté $(g : T \to S, t' \leadsto t, \xi)$
de $\mathcal{F}$ au-dessus de $s$ telle que $g$ soit quasi-fini en $t$.
\end{lemma}

\begin{proof}
Nous pouvons remplacer $S$ par un voisinage affine de $s$. D'après le
lemme \ref{lemma-impure-finite-presentation},
nous pouvons supposer donnée une impureté $(g : T \to S, t' \leadsto t, \xi)$
telle que $g$ soit localement de type fini et $t$ un point fermé de la
fibre de $g$ au-dessus de $s$. Nous pouvons remplacer $T$ par le sous-schéma réduit
induit sur $\overline{\{t'\}}$. Soit
$Z = \overline{\{\xi\}} \subset X_T$. Par hypothèse, $Z_t = \emptyset$
et l'image de $Z \to T$ contient $t'$. D'après
Compléments sur les morphismes,
lemme \ref{more-morphisms-lemma-relative-assassin-in-neighbourhood},
il existe un ouvert non vide $V \subset Z$ tel que, pour tout
$w \in f(V)$, tout point générique $\xi'$ de $V_w$ appartienne à
$\text{Ass}_{X_T/T}(\mathcal{F}_T)$. D'après
Compléments sur les morphismes, lemme \ref{more-morphisms-lemma-nonempty-generic-fibre},
il existe un ouvert non vide $W \subset T$ tel que $W \subset f(V)$. D'après
Compléments sur les morphismes, lemme
\ref{more-morphisms-lemma-quasi-finite-quasi-section-meeting-nearby-open-X},
il existe un sous-schéma fermé $T' \subset T$ tel que
$t \in T'$, que $T' \to S$ soit quasi-fini en $t$ et qu'il existe un point
$z \in T' \cap W$, $z \leadsto t$, qui ne s'envoie pas sur $s$.
Choisissons un point générique quelconque $\xi'$ du schéma non vide $V_z$.
Alors $(T' \to S, z \leadsto t, \xi')$ est l'impureté cherchée.
\end{proof}

\noindent
Dans la suite, nous utiliserons l'hensélisé
$S^h = \Spec(\mathcal{O}_{S, s}^h)$
de $S$ en $s$, voir
Cohomologie \'etale,
définition \ref{etale-cohomology-definition-etale-local-rings}.
Comme $S^h \to S$ envoie le point fermé de $S^h$ sur $s$ et
induit un isomorphisme des corps résiduels, nous noterons
aussi $s \in S^h$ ce point fermé. Ainsi, $(S^h, s) \to (S, s)$ est
un morphisme de schémas pointés.

\begin{lemma}
\label{lemma-impurity-on-henselization}
Plaçons-nous dans la situation \ref{situation-pre-pure}.
S'il existe une impureté $(S^h \to S, s' \leadsto s, \xi)$
de $\mathcal{F}$ au-dessus de $s$, alors il existe une impureté
$(T \to S, t' \leadsto t, \xi)$ de $\mathcal{F}$ au-dessus de $s$
où $(T, t) \to (S, s)$ est un voisinage \'etale élémentaire.
\end{lemma}

\begin{proof}
Nous pouvons remplacer $S$ par un voisinage affine de $s$.
Disons que $S = \Spec(A)$ et que $s$ correspond à l'idéal premier
$\mathfrak p \subset A$. Alors
$\mathcal{O}_{S, s}^h = \colim_{(T, t)} \Gamma(T, \mathcal{O}_T)$,
la limite étant prise sur la catégorie opposée à la
catégorie cofiltrante des voisinages \'etales élémentaires affines
$(T, t)$ de $(S, s)$, voir
Compléments sur les morphismes,
lemme \ref{more-morphisms-lemma-describe-henselization}
et sa démonstration. Ainsi, $S^h = \lim_i T_i$ et nous concluons par le
lemme \ref{lemma-impure-limit}.
\end{proof}

\begin{lemma}
\label{lemma-pure-along-X-s}
Dans la situation \ref{situation-pre-pure}, les conditions suivantes
sont équivalentes :
\begin{enumerate}
\item il existe une impureté $(S^h \to S, s' \leadsto s, \xi)$
de $\mathcal{F}$ au-dessus de $s$, où $S^h$ est l'hensélisé de $S$ en $s$,
\item il existe une impureté $(T \to S, t' \leadsto t, \xi)$
de $\mathcal{F}$ au-dessus de $s$ telle que $(T, t) \to (S, s)$ soit un
voisinage \'etale élémentaire, et
\item il existe une impureté $(T \to S, t' \leadsto t, \xi)$
de $\mathcal{F}$ au-dessus de $s$ telle que $T \to S$ soit quasi-fini en $t$.
\end{enumerate}
\end{lemma}

\begin{proof}
Comme un morphisme \'etale est localement quasi-fini, il est clair que
(2) entraîne (3). Nous avons vu que (3) entraîne (2) dans le
lemme \ref{lemma-quasi-finite-impurity-elementary}.
Nous avons vu que (1) entraîne (2) dans le
lemme \ref{lemma-impurity-on-henselization}.
Enfin, si $(T \to S, t' \leadsto t, \xi)$ est une impureté
de $\mathcal{F}$ au-dessus de $s$ telle que $(T, t) \to (S, s)$ soit un
voisinage \'etale élémentaire, nous pouvons choisir une factorisation
$S^h \to T \to S$ du morphisme structural $S^h \to S$.
Choisissons un point quelconque $s' \in S^h$ d'image $t'$ et un point quelconque
$\xi' \in X_{s'}$ d'image $\xi \in X_{t'}$. Alors
$(S^h \to S, s' \leadsto s, \xi')$ est une impureté de
$\mathcal{F}$ au-dessus de $s$. Nous omettons les détails.
\end{proof}





\section{Modules relativement purs}
\label{section-pure}

\noindent
La notion de module pur relativement à une base a été introduite dans \cite{GruRay}.

\begin{definition}
\label{definition-pure}
Soit $f : X \to S$ un morphisme de schémas de type fini.
Soit $\mathcal{F}$ un $\mathcal{O}_X$-module quasi-cohérent de type fini.
\begin{enumerate}
\item Soit $s \in S$. Nous disons que $\mathcal{F}$ est {\it pur le long de $X_s$}
s'il n'existe aucune impureté $(g : T \to S, t' \leadsto t, \xi)$
de $\mathcal{F}$ au-dessus de $s$ avec $(T, t) \to (S, s)$ un
voisinage \'etale élémentaire.
\item Nous disons que $\mathcal{F}$ est {\it universellement pur le long de $X_s$}
s'il n'existe aucune impureté de $\mathcal{F}$ au-dessus de $s$.
\item Nous disons que $X$ est {\it pur le long de $X_s$} si $\mathcal{O}_X$
est pur le long de $X_s$.
\item Nous disons que $\mathcal{F}$ est {\it universellement $S$-pur}, ou
{\it universellement pur relativement à $S$}, si $\mathcal{F}$ est universellement
pur le long de $X_s$ pour tout $s \in S$.
\item Nous disons que $\mathcal{F}$ est {\it $S$-pur}, ou
{\it pur relativement à $S$}, si $\mathcal{F}$ est pur le long de $X_s$
pour tout $s \in S$.
\item Nous disons que $X$ est {\it $S$-pur} ou {\it pur relativement à $S$}
si $\mathcal{O}_X$ est pur relativement à $S$.
\end{enumerate}
\end{definition}

\noindent
Nous nous restreignons ici délibérément aux morphismes
de type fini, et non aux morphismes seulement localement de
type fini ; voir la
remarque \ref{remark-discuss-finite-type}
pour une discussion. Dans la situation de la définition, le
lemme \ref{lemma-pure-along-X-s}
nous dit que les conditions suivantes sont équivalentes :
\begin{enumerate}
\item $\mathcal{F}$ est pur le long de $X_s$,
\item il n'existe aucune impureté $(g : T \to S, t' \leadsto t, \xi)$ avec $g$
quasi-fini en $t$,
\item il n'existe aucune impureté de la forme
$(S^h \to S, s' \leadsto s, \xi)$, où $S^h$ est l'hensélisé
de $S$ en $s$.
\end{enumerate}
Si nous notons $X^h = X \times_S S^h$ et $\mathcal{F}^h$ l'image réciproque
de $\mathcal{F}$ sur $X^h$, nous pouvons formuler cette dernière condition
sous la forme positive suivante :
\begin{enumerate}
\item[(4)] Tous les points de $\text{Ass}_{X^h/S^h}(\mathcal{F}^h)$ se spécialisent
en des points de $X_s$.
\end{enumerate}
En particulier, il est clair que $\mathcal{F}$ est pur le long de $X_s$
si et seulement si l'image réciproque de $\mathcal{F}$ sur
$X \times_S \Spec(\mathcal{O}_{S, s})$ est pure le long de $X_s$.

\begin{remark}
\label{remark-discuss-finite-type}
Soient $f : X \to S$ un morphisme localement de type fini
et $\mathcal{F}$ un $\mathcal{O}_X$-module quasi-cohérent de type fini.
Dans ce cas, il est encore vrai que (1) et (2) ci-dessus sont équivalentes,
car la démonstration du
lemme \ref{lemma-quasi-finite-impurity-elementary}
n'utilise pas la quasi-compacité de $f$. Il est également clair que
(3) et (4) sont équivalentes. Cependant, nous ignorons si (1) et (3) sont
équivalentes. Dans ce cas, il peut parfois être plus commode de définir
la pureté par les conditions équivalentes (3) et (4), comme dans \cite{GruRay}.
D'autre part, pour de nombreuses applications, il semble que la notion appropriée
soit bien celle de pureté universelle.
\end{remark}

\noindent
Il est naturel de se demander si la propriété de pureté relativement à
la base est conservée par changement de base, autrement dit si être pur revient
à être universellement pur. Il se trouve que c'est vrai
sur une base noethérienne (voir le
lemme \ref{lemma-Noetherian-base-change}),
ou si le faisceau est plat (voir les
lemmes \ref{lemma-finite-type-flat-pure-along-fibre-is-universal} et
\ref{lemma-finite-type-flat-pure-is-universal}).
Ce n'est pas vrai en général, même si le morphisme et le faisceau sont de
présentation finie ; voir
Exemples, section \ref{examples-section-pure-not-universally},
pour un contre-exemple. Vérifions d'abord que notre emploi de « universellement »
correspond à la notion usuelle.

\begin{lemma}
\label{lemma-base-change-universally}
Soit $f : X \to S$ un morphisme de schémas de type fini.
Soit $\mathcal{F}$ un $\mathcal{O}_X$-module quasi-cohérent de type fini.
Soit $s \in S$. Les conditions suivantes sont équivalentes :
\begin{enumerate}
\item $\mathcal{F}$ est universellement pur le long de $X_s$, et
\item pour tout morphisme de schémas pointés $(S', s') \to (S, s)$,
l'image réciproque $\mathcal{F}_{S'}$ est pure le long de $X_{s'}$.
\end{enumerate}
En particulier, $\mathcal{F}$ est universellement pur relativement à $S$ si et
seulement si tout changement de base $\mathcal{F}_{S'}$ de $\mathcal{F}$ est
pur relativement à $S'$.
\end{lemma}

\begin{proof}
C'est formel.
\end{proof}

\begin{lemma}
\label{lemma-quasi-finite-base-change}
Soit $f : X \to S$ un morphisme de schémas de type fini.
Soit $\mathcal{F}$ un $\mathcal{O}_X$-module quasi-cohérent de type fini.
Soit $s \in S$. Soit $(S', s') \to (S, s)$ un morphisme de schémas pointés.
Si $S' \to S$ est quasi-fini en $s'$ et $\mathcal{F}$ est pur le long de $X_s$,
alors $\mathcal{F}_{S'}$ est pur le long de $X_{s'}$.
\end{lemma}

\begin{proof}
Si $(T \to S', t' \leadsto t, \xi)$ est une impureté de
$\mathcal{F}_{S'}$ au-dessus de $s'$ avec $T \to S'$ quasi-fini en $t$,
alors $(T \to S, t' \to t, \xi)$ est une impureté de $\mathcal{F}$
au-dessus de $s$ avec $T \to S$ quasi-fini en $t$, voir
Morphismes, lemme \ref{morphisms-lemma-composition-quasi-finite}.
Le lemme résulte donc immédiatement de la caractérisation (2)
de la pureté donnée après la
définition \ref{definition-pure}.
\end{proof}

\begin{lemma}
\label{lemma-Noetherian-base-change}
Soit $f : X \to S$ un morphisme de schémas de type fini.
Soit $\mathcal{F}$ un $\mathcal{O}_X$-module quasi-cohérent de type fini.
Soit $s \in S$. Si $\mathcal{O}_{S, s}$ est noethérien, alors
$\mathcal{F}$ est pur le long de $X_s$ si et seulement si $\mathcal{F}$
est universellement pur le long de $X_s$.
\end{lemma}

\begin{proof}
Nous pouvons d'abord remplacer $S$ par $\Spec(\mathcal{O}_{S, s})$, c'est-à-dire
supposer $S$ noethérien. Utilisons ensuite le
lemme \ref{lemma-Noetherian-impurity-quasi-finite}
et la caractérisation (2) de la pureté donnée dans la discussion suivant la
définition \ref{definition-pure}
pour conclure.
\end{proof}

\noindent
La pureté satisfait à la descente plate.

\begin{lemma}
\label{lemma-flat-descend-pure}
Soit $f : X \to S$ un morphisme de schémas de type fini.
Soit $\mathcal{F}$ un $\mathcal{O}_X$-module quasi-cohérent de type fini.
Soit $s \in S$. Soit $(S', s') \to (S, s)$ un morphisme de schémas pointés.
Supposons $S' \to S$ plat en $s'$.
\begin{enumerate}
\item Si $\mathcal{F}_{S'}$ est pur le long de $X_{s'}$,
alors $\mathcal{F}$ est pur le long de $X_s$.
\item Si $\mathcal{F}_{S'}$ est universellement pur le long de $X_{s'}$,
alors $\mathcal{F}$ est universellement pur le long de $X_s$.
\end{enumerate}
\end{lemma}

\begin{proof}
Soit $(T \to S, t' \leadsto t, \xi)$ une impureté de
$\mathcal{F}$ au-dessus de $s$. Posons $T_1 = T \times_S S'$ et soit $t_1$
l'unique point de $T_1$ d'images $t$ et $s'$. Comme
$T_1 \to T$ est plat en $t_1$, voir
Morphismes, lemme \ref{morphisms-lemma-base-change-flat},
il existe une spécialisation $t'_1 \leadsto t_1$ au-dessus de
$t' \leadsto t$, voir
Algèbre, section \ref{algebra-section-going-up}.
Choisissons un point $\xi_1 \in X_{t'_1}$ correspondant à un point
générique de $\Spec(\kappa(t'_1) \otimes_{\kappa(t')} \kappa(\xi))$, voir
Schémas, lemme \ref{schemes-lemma-points-fibre-product}.
D'après
Diviseurs, lemme \ref{divisors-lemma-base-change-relative-assassin},
nous avons $\xi_1 \in \text{Ass}_{X_{T_1}/T_1}(\mathcal{F}_{T_1})$.
Comme l'adhérence de Zariski de $\{\xi_1\}$ dans $X_{T_1}$ s'envoie dans
l'adhérence de Zariski de $\{\xi\}$ dans $X_T$, nous concluons que
cette adhérence est disjointe de $X_{t_1}$. Ainsi,
$(T_1 \to S', t'_1 \leadsto t_1, \xi_1)$
est une impureté de $\mathcal{F}_{S'}$ au-dessus de $s'$.
Autrement dit, nous avons démontré la contraposée de la partie (2) du
lemme. Enfin, si $(T, t) \to (S, s)$ est un voisinage
\'etale élémentaire, alors $(T_1, t_1) \to (S', s')$ est aussi un
voisinage \'etale élémentaire ; nous voyons ainsi que (1) est vraie.
\end{proof}

\begin{lemma}
\label{lemma-supported-on-closed}
Soit $i : Z \to X$ une immersion fermée de schémas de type fini sur
un schéma $S$. Soit $s \in S$. Soit $\mathcal{F}$ un
faisceau quasi-cohérent de type fini sur $Z$. Alors $\mathcal{F}$ est
(universellement) pur le long de $Z_s$ si et seulement si $i_*\mathcal{F}$
est (universellement) pur le long de $X_s$.
\end{lemma}

\begin{proof}
Cela résulte de
Diviseurs, lemme \ref{divisors-lemma-relative-weak-assassin-finite}.
\end{proof}








\section{Exemples de faisceaux relativement purs}
\label{section-examples-pure-sheaves}

\noindent
Voici quelques exemples où l'on peut expliciter la signification de la pureté.

\begin{lemma}
\label{lemma-proper-pure}
Soit $f : X \to S$ un morphisme de schémas de type fini.
Soit $\mathcal{F}$ un $\mathcal{O}_X$-module quasi-cohérent de type fini.
\begin{enumerate}
\item Si le support de $\mathcal{F}$ est propre sur $S$, alors
$\mathcal{F}$ est universellement pur relativement à $S$.
\item Si $f$ est propre, alors
$\mathcal{F}$ est universellement pur relativement à $S$.
\item Si $f$ est propre, alors $X$ est universellement pur relativement à $S$.
\end{enumerate}
\end{lemma}

\begin{proof}
Ramenons d'abord (1) à (2). Soit $Z \subset X$ le
support schématique de $\mathcal{F}$. Soit $i : Z \to X$
l'immersion fermée correspondante et écrivons
$\mathcal{F} = i_*\mathcal{G}$ pour un
$\mathcal{O}_Z$-module quasi-cohérent de type fini $\mathcal{G}$, voir
Morphismes, section \ref{morphisms-section-support}.
Dans le cas (1), $Z \to S$ est propre par hypothèse.
Ainsi, d'après le lemme \ref{lemma-supported-on-closed}, le cas (1) se ramène au cas (2).

\medskip\noindent
Supposons $f$ propre.
Soit $(g : T \to S, t' \leadsto t, \xi)$ une impureté de $\mathcal{F}$
au-dessus de $s \in S$. Comme $f$ est propre, il est universellement fermé. Ainsi,
$f_T : X_T \to T$ est fermé. Comme $f_T(\xi) = t'$, cela entraîne que
$t \in f(\overline{\{\xi\}})$, d'où une contradiction.
\end{proof}

\begin{lemma}
\label{lemma-quasi-finite-pure}
Soit $f : X \to S$ un morphisme de schémas séparé et de type fini.
Soit $\mathcal{F}$ un $\mathcal{O}_X$-module quasi-cohérent de type fini.
Supposons $\text{Supp}(\mathcal{F}_s)$ fini pour tout $s \in S$.
Alors les conditions suivantes sont équivalentes :
\begin{enumerate}
\item $\mathcal{F}$ est pur relativement à $S$,
\item le support schématique de $\mathcal{F}$ est fini sur $S$, et
\item $\mathcal{F}$ est universellement pur relativement à $S$.
\end{enumerate}
En particulier, pour un morphisme quasi-fini séparé $X \to S$,
$X$ est pur relativement à $S$ si et seulement si $X \to S$ est fini.
\end{lemma}

\begin{proof}
Soit $Z \subset X$ le support schématique de $\mathcal{F}$, voir
Morphismes, définition \ref{morphisms-definition-scheme-theoretic-support}.
Alors $Z \to S$ est un morphisme de schémas séparé de type fini à
fibres finies. Il est donc séparé et quasi-fini, voir
Morphismes, lemme \ref{morphisms-lemma-quasi-finite}.
D'après le
lemme \ref{lemma-supported-on-closed},
il suffit de démontrer le lemme pour $Z \to S$ et le faisceau $\mathcal{F}$
considéré comme module quasi-cohérent de type fini sur $Z$. Nous pouvons donc
supposer que $X \to S$ est séparé et quasi-fini et que
$\text{Supp}(\mathcal{F}) = X$.

\medskip\noindent
Il résulte du
lemme \ref{lemma-proper-pure}
et de
Morphismes, lemme \ref{morphisms-lemma-finite-proper},
que (2) entraîne (3). Il est trivial que (3) entraîne (1). Supposons (1) vérifiée.
Montrons (2). Il est clair que nous pouvons supposer $S$ affine. D'après
Compléments sur les morphismes,
lemme \ref{more-morphisms-lemma-quasi-finite-separated-pass-through-finite},
nous pouvons trouver un diagramme
$$
\xymatrix{
X \ar[rd]_f \ar[rr]_j & & T \ar[ld]^\pi \\
& S &
}
$$
où $\pi$ est fini et $j$ est une immersion ouverte quasi-compacte.
Si nous montrons que $j$ est fermé, alors $j$ est une immersion fermée,
et nous concluons que $f = \pi \circ j$ est fini.
Pour montrer que $j$ est fermé, il suffit de montrer que les spécialisations
se relèvent le long de $j$, voir
Schémas, lemme \ref{schemes-lemma-quasi-compact-closed}.
Soit $x \in X$, posons $t' = j(x)$ et soit $t' \leadsto t$ une spécialisation.
Il faut montrer $t \in j(X)$. Posons $s' = f(x)$ et $s = \pi(t)$, de sorte que
$s' \leadsto s$. D'après
Compléments sur les morphismes, lemme
\ref{more-morphisms-lemma-etale-splits-off-quasi-finite-part-technical},
nous pouvons trouver un voisinage \'etale élémentaire
$(U, u) \to (S, s)$ et une décomposition
$$
T_U = T \times_S U = V \amalg W
$$
en sous-schémas ouverts et fermés, telle que $V \to U$ soit fini,
qu'il existe un unique point $v$ de $V$ d'image $u$, et que
$v$ ait pour image $t$ dans $T$. Comme $V \to T$ est \'etale, nous pouvons relever
les généralisations, voir
Morphismes, lemmes \ref{morphisms-lemma-generalizations-lift-flat} et
\ref{morphisms-lemma-etale-flat}.
Il existe donc une spécialisation $v' \leadsto v$ telle que $v'$
s'envoie sur $t' \in T$. En particulier, $v' \in X_U \subset T_U$.
Notons $u' \in U$ l'image de $t'$. Notons que
$v' \in \text{Ass}_{X_U/U}(\mathcal{F})$, car $X_{u'}$ est un ensemble
fini discret et $X_{u'} = \text{Supp}(\mathcal{F}_{u'})$.
Comme $\mathcal{F}$ est pur relativement à $S$, le point $v'$ doit
se spécialiser en un point de $X_u$. Puisque $v$ est l'unique point de
$V$ au-dessus de $u$ (et qu'aucun point de $W$ ne peut être une spécialisation
de $v'$), nous voyons que $v \in X_u$. Ainsi, $t \in X$.
\end{proof}

\begin{lemma}
\label{lemma-flat-geometrically-integral-fibres-pure}
Soit $f : X \to S$ un morphisme de schémas plat de type fini
à fibres géométriquement intègres. Alors $X$ est universellement pur
sur $S$.
\end{lemma}

\begin{proof}
Soient $\xi \in X$ avec $s' = f(\xi)$ et $s' \leadsto s$ une spécialisation
de $S$. Si $\xi$ est un point associé de $X_{s'}$, alors $\xi$ est
l'unique point générique, car $X_{s'}$ est un schéma intègre. Soit
$\xi_0$ l'unique point générique de $X_s$. Comme $X \to S$ est plat,
nous pouvons relever $s' \leadsto s$ en une spécialisation
$\xi' \leadsto \xi_0$ dans $X$, voir
Morphismes, lemme \ref{morphisms-lemma-generalizations-lift-flat}.
Alors $\xi \leadsto \xi'$, car $\xi$ est le point générique de $X_{s'}$,
donc $\xi \leadsto \xi_0$. Cela signifie que $(\text{id}_S, s' \to s, \xi)$
n'est pas une impureté de $\mathcal{O}_X$ au-dessus de $s$. Comme l'hypothèse
que $f$ est de type fini, plat et à fibres géométriquement intègres
est conservée par changement de base, il n'existe aucune
impureté après un changement de base quelconque. Nous voyons ainsi que $X$ est
universellement $S$-pur.
\end{proof}

\begin{lemma}
\label{lemma-affine-locally-projective-pure}
Soit $f : X \to S$ un morphisme de schémas affine et de type fini.
Soit $\mathcal{F}$ un $\mathcal{O}_X$-module quasi-cohérent de type fini
tel que $f_*\mathcal{F}$ soit localement projectif sur $S$, voir
Propriétés, définition \ref{properties-definition-locally-projective}.
Alors $\mathcal{F}$ est universellement pur sur $S$.
\end{lemma}

\begin{proof}
Après réduction au cas où $S$ est le spectre d'un anneau local
hensélien, cela résulte du
lemme \ref{lemma-explain-why-pure}.
\end{proof}




\section{Un critère de pureté}
\label{section-criterion-purity}

\noindent
Nous démontrons d'abord que, pour une famille plate de faisceaux
quasi-cohérents de type fini, les points de l'assassin relatif
se spécialisent en des points des assassins relatifs des fibres voisines
(lorsqu'ils admettent une spécialisation dans ces fibres).

\begin{lemma}
\label{lemma-associated-point-specializes}
Soit $f : X \to S$ un morphisme de schémas de type fini.
Soit $\mathcal{F}$ un $\mathcal{O}_X$-module quasi-cohérent de type fini.
Soit $s \in S$.
Supposons $\mathcal{F}$ plat sur $S$ en tout point de $X_s$.
Soit $x' \in \text{Ass}_{X/S}(\mathcal{F})$ avec $f(x') = s'$
tel que $s' \leadsto s$ soit une spécialisation dans $S$. Si
$x'$ se spécialise en un point de $X_s$, alors $x' \leadsto x$
pour un $x \in \text{Ass}_{X_s}(\mathcal{F}_s)$.
\end{lemma}

\begin{proof}
Disons que $x' \leadsto t$ avec $t \in X_s$. Nous pouvons alors trouver des spécialisations
$x' \leadsto x \leadsto t$, où $x$ correspond au point générique
d'une composante irréductible de
$\overline{\{x'\}} \cap f^{-1}(\{s\})$. Par hypothèse, $\mathcal{F}$ est plat
sur $S$ en $x$. D'après Compléments sur les morphismes, lemme
\ref{more-morphisms-lemma-associated-point-specializes},
nous avons $x \in \text{Ass}_{X/S}(\mathcal{F})$, comme voulu.
\end{proof}

\begin{lemma}
\label{lemma-criterion}
Soit $f : X \to S$ un morphisme de schémas de type fini.
Soit $\mathcal{F}$ un $\mathcal{O}_X$-module quasi-cohérent
de type fini. Soit $s \in S$. Soit $(S', s') \to (S, s)$ un
voisinage \'etale élémentaire et soit
$$
\xymatrix{
X \ar[d] & X' \ar[l]^g \ar[d] \\
S & S' \ar[l]
}
$$
un diagramme commutatif de morphismes de schémas. Supposons que
\begin{enumerate}
\item $\mathcal{F}$ soit plat sur $S$ en tout point de $X_s$,
\item $X' \to S'$ soit de type fini,
\item $g^*\mathcal{F}$ soit pur le long de $X'_{s'}$,
\item $g : X' \to X$ soit \'etale, et
\item $g(X')$ contienne $\text{Ass}_{X_s}(\mathcal{F}_s)$.
\end{enumerate}
Dans cette situation, $\mathcal{F}$ est pur le long de $X_s$ si et seulement
si l'image de $X' \to X \times_S S'$ contient les points de
$\text{Ass}_{X \times_S S'/S'}(\mathcal{F} \times_S S')$
au-dessus des points de $S'$ qui se spécialisent en $s'$.
\end{lemma}

\begin{proof}
Comme le morphisme $S' \to S$ est \'etale, si $\mathcal{F}$
est pur le long de $X_s$, alors $\mathcal{F} \times_S S'$ est pur le long de
$X_s$, voir le
lemme \ref{lemma-quasi-finite-base-change}.
Comme la pureté satisfait à la descente plate, voir le
lemme \ref{lemma-flat-descend-pure},
si $\mathcal{F} \times_S S'$ est pur le long de $X_{s'}$, alors
$\mathcal{F}$ est pur le long de $X_s$. Nous pouvons donc remplacer $S$ par $S'$
et supposer $S = S'$, de sorte que $g : X' \to X$ est un morphisme \'etale
entre schémas de type fini sur $S$. De plus, nous pouvons remplacer
$S$ par $\Spec(\mathcal{O}_{S, s})$ et supposer $S$ local.

\medskip\noindent
Supposons d'abord $\mathcal{F}$ pur le long de $X_s$.
Dans ce cas, tout point de $\text{Ass}_{X/S}(\mathcal{F})$
se spécialise en un point de $X_s$, par pureté. Ainsi, d'après le
lemme \ref{lemma-associated-point-specializes},
tout point de $\text{Ass}_{X/S}(\mathcal{F})$
se spécialise en un point de $\text{Ass}_{X_s}(\mathcal{F}_s)$.
Tout point de $\text{Ass}_{X/S}(\mathcal{F})$ appartient donc à
l'image de $g$ (car cette image est ouverte et contient
$\text{Ass}_{X_s}(\mathcal{F}_s)$).

\medskip\noindent
Réciproquement, supposons que $g(X')$ contienne $\text{Ass}_{X/S}(\mathcal{F})$.
Soit $S^h = \Spec(\mathcal{O}_{S, s}^h)$ l'hensélisé
de $S$ en $s$. Notons $g^h : (X')^h \to X^h$ le changement de base de $g$
par $S^h \to S$, et $\mathcal{F}^h$ l'image réciproque de $\mathcal{F}$
sur $X^h$. D'après
Diviseurs, lemme \ref{divisors-lemma-base-change-relative-assassin} et
remarque \ref{divisors-remark-base-change-relative-assassin},
l'assassin relatif $\text{Ass}_{X^h/S^h}(\mathcal{F}^h)$
est l'image réciproque de $\text{Ass}_{X/S}(\mathcal{F})$ par la projection
$X^h \to X$. Comme nous avons supposé que $g(X')$ contient
$\text{Ass}_{X/S}(\mathcal{F})$, nous concluons que le changement de base\newline
$g^h((X')^h) = g(X') \times_S S^h$ contient
$\text{Ass}_{X^h/S^h}(\mathcal{F}^h)$. Nous nous ramenons ainsi
au cas où $S$ est le spectre d'un anneau local hensélien.
Soit $x \in \text{Ass}_{X/S}(\mathcal{F})$. Pour achever la démonstration du
lemme, il faut montrer que $x$ se spécialise en un point de $X_s$, voir
le critère (4) de pureté dans la discussion suivant la
définition \ref{definition-pure}.
Par hypothèse, il existe un $x' \in X'$ tel que $g(x') = x$.
Comme $g : X' \to X$ est \'etale, nous avons
$x' \in \text{Ass}_{X'/S}(g^*\mathcal{F})$, voir le
lemme \ref{lemma-etale-weak-assassin-up-down} (appliqué au
morphisme de fibres $X'_w \to X_w$, où $w \in S$ est l'image de $x'$).
Comme $g^*\mathcal{F}$ est pur le long de $X'_s$, nous avons $x' \leadsto y$
pour un $y \in X'_s$. Ainsi, $x = g(x') \leadsto g(y)$ et
$g(y) \in X_s$, comme voulu.
\end{proof}

\begin{lemma}
\label{lemma-finite-type-flat-pure-along-fibre-is-universal}
Soit $f : X \to S$ un morphisme de schémas.
Soit $\mathcal{F}$ un $\mathcal{O}_X$-module quasi-cohérent.
Soit $s \in S$.
Supposons que
\begin{enumerate}
\item $f$ soit de type fini,
\item $\mathcal{F}$ soit de type fini,
\item $\mathcal{F}$ soit plat sur $S$ en tout point de $X_s$, et
\item $\mathcal{F}$ soit pur le long de $X_s$.
\end{enumerate}
Alors $\mathcal{F}$ est universellement pur le long de $X_s$.
\end{lemma}

\begin{proof}
Faisons d'abord une remarque préliminaire.\newline Supposons que $(S', s') \to (S, s)$
soit un voisinage \'etale élémentaire. Notons $\mathcal{F}'$
l'image réciproque de $\mathcal{F}$ sur $X' = X \times_S S'$.
D'après la discussion suivant la
définition \ref{definition-pure},
$\mathcal{F}'$ est pur le long de $X'_{s'}$. De plus, $\mathcal{F}'$
est plat sur $S'$ le long de $X'_{s'}$. Il suffit alors de démontrer
que $\mathcal{F}'$ est universellement pur le long de $X'_{s'}$. En effet, pour
tout morphisme $(T, t) \to (S, s)$ de schémas pointés,
le produit fibré $(T', t') = (T \times_S S', (t, s'))$ est plat sur $(T, t)$ ;
ainsi, si $\mathcal{F}_{T'}$ est pur le long de $X_{t'}$, alors
$\mathcal{F}_T$ est pur le long de $X_t$ d'après le
lemme \ref{lemma-flat-descend-pure}.
Ainsi, au cours de la démonstration, nous pouvons toujours remplacer $(s, S)$ par un voisinage
\'etale élémentaire.
Nous pouvons aussi remplacer $S$ par $\Spec(\mathcal{O}_{S, s})$
en raison du caractère local du problème.

\medskip\noindent
Choisissons un voisinage \'etale élémentaire $(S', s') \to (S, s)$ et
un diagramme commutatif
$$
\xymatrix{
X \ar[d] & X' \ar[l]^g \ar[d] \\
S & \Spec(\mathcal{O}_{S', s'}) \ar[l]
}
$$
tels que $X' \to X \times_S \Spec(\mathcal{O}_{S', s'})$
soit \'etale, que $X_s = g((X')_{s'})$, que le schéma $X'$ soit affine
et que $\Gamma(X', g^*\mathcal{F})$ soit un
$\mathcal{O}_{S', s'}$-module libre, voir le
lemme \ref{lemma-finite-type-flat-along-fibre-free-variant}.
Notons que $X' \to \Spec(\mathcal{O}_{S', s'})$ est de type fini
(c'est un morphisme quasi-compact composé d'un morphisme \'etale
et du changement de base d'un morphisme de type fini).
D'après les remarques préliminaires du premier paragraphe de la démonstration,
nous pouvons remplacer $S$ par $\Spec(\mathcal{O}_{S', s'})$. Nous pouvons donc
supposer qu'il existe un diagramme commutatif
$$
\xymatrix{
X \ar[dr] & & X' \ar[ll]^g \ar[ld] \\
& S &
}
$$
de schémas de type fini sur $S$, où $g$ est \'etale, $X_s \subset g(X')$,
où $S$ est local de point fermé $s$, où $X'$ est affine et où
$\Gamma(X', g^*\mathcal{F})$ est un $\Gamma(S, \mathcal{O}_S)$-module libre.
Notons que, dans ce cas, $g^*\mathcal{F}$ est universellement pur sur $S$, voir le
lemme \ref{lemma-affine-locally-projective-pure}.

\medskip\noindent
Dans cette situation, nous appliquons le
lemme \ref{lemma-criterion}
pour déduire que $\text{Ass}_{X/S}(\mathcal{F}) \subset g(X')$
de l'hypothèse que $\mathcal{F}$ est pur le long de $X_s$
et plat sur $S$ le long de $X_s$. D'après
Diviseurs, lemme \ref{divisors-lemma-base-change-relative-assassin} et
remarque \ref{divisors-remark-base-change-relative-assassin},
pour tout morphisme de schémas pointés
$(T, t) \to (S, s)$, nous avons
$$
\text{Ass}_{X_T/T}(\mathcal{F}_T) \subset
(X_T \to X)^{-1}(\text{Ass}_{X/S}(\mathcal{F})) \subset
g(X') \times_S T = g_T(X'_T).
$$
Ainsi, par le
lemme \ref{lemma-criterion}
appliqué au changement de base à $(T, t)$ du diagramme ci-dessus,
nous concluons que $\mathcal{F}_T$ est pur le long de $X_t$, comme voulu.
\end{proof}

\begin{lemma}
\label{lemma-finite-type-flat-pure-is-universal}
Soit $f : X \to S$ un morphisme de schémas de type fini.
Soit $\mathcal{F}$ un $\mathcal{O}_X$-module quasi-cohérent de type fini.
Supposons $\mathcal{F}$ plat sur $S$. Dans ce cas,
$\mathcal{F}$ est pur relativement à $S$ si et seulement si $\mathcal{F}$
est universellement pur relativement à $S$.
\end{lemma}

\begin{proof}
Conséquence immédiate du
lemme \ref{lemma-finite-type-flat-pure-along-fibre-is-universal}
et des définitions.
\end{proof}

\begin{lemma}
\label{lemma-limit-purity}
Soit $I$ un ensemble ordonné filtrant.
Soit $(S_i, g_{ii'})$ un système projectif de schémas affines indexé par $I$.
Posons $S = \lim_i S_i$ et soit $s \in S$.
Notons $g_i : S \to S_i$ les projections et posons $s_i = g_i(s)$.
Supposons que $f : X \to S$ soit un morphisme de présentation finie,
et $\mathcal{F}$ un $\mathcal{O}_X$-module quasi-cohérent de présentation finie,
pur le long de $X_s$ et plat sur $S$ en tout point de $X_s$.
Alors il existe un $i \in I$, un morphisme de présentation finie
$X_i \to S_i$, un $\mathcal{O}_{X_i}$-module quasi-cohérent $\mathcal{F}_i$
de présentation finie, pur le long de $(X_i)_{s_i}$ et plat sur $S_i$
en tout point de $(X_i)_{s_i}$, tels que $X \cong X_i \times_{S_i} S$
et que l'image réciproque de $\mathcal{F}_i$ sur $X$ soit isomorphe
à $\mathcal{F}$.
\end{lemma}

\begin{proof}
Soit $U \subset X$ l'ensemble des points où $\mathcal{F}$ est
plat sur $S$. D'après
Compléments sur les morphismes, théorème \ref{more-morphisms-theorem-openness-flatness},
c'est un sous-schéma ouvert de $X$. Par hypothèse, $X_s \subset U$.
Comme $X_s$ est quasi-compact, nous pouvons trouver un ouvert quasi-compact
$U' \subset U$ tel que $X_s \subset U'$. D'après
Limites, lemme \ref{limits-lemma-descend-finite-presentation},
nous pouvons trouver un $i \in I$ et un morphisme de présentation finie
$f_i : X_i \to S_i$ dont le changement de base à $S$ est isomorphe à $f_i$.
Fixons un tel choix et posons $X_{i'} = X_i \times_{S_i} S_{i'}$.
Alors $X = \lim_{i'} X_{i'}$, avec des morphismes de transition affines. D'après
Limites, lemme \ref{limits-lemma-descend-modules-finite-presentation},
nous pouvons, quitte à augmenter $i$, supposer qu'il existe un
$\mathcal{O}_{X_i}$-module quasi-cohérent $\mathcal{F}_i$
de présentation finie dont le changement de base à $S$ est isomorphe à
$\mathcal{F}$. D'après
Limites, lemme \ref{limits-lemma-descend-opens},
nous pouvons, quitte à augmenter $i$, supposer qu'il existe un
ouvert $U'_i \subset X_i$ dont l'image réciproque dans $X$ est $U'$.
Notons qu'en particulier $(X_i)_{s_i} \subset U'_i$. D'après
Limites, lemme \ref{limits-lemma-descend-module-flat-finite-presentation},
(en augmentant encore $i$),
nous pouvons supposer $\mathcal{F}_i$ plat sur $U'_i$.
En particulier, $\mathcal{F}_i$ est plat le long de $(X_i)_{s_i}$.

\medskip\noindent
Utilisons ensuite le
lemme \ref{lemma-finite-presentation-flat-along-fibre}
pour choisir un voisinage \'etale élémentaire
$(S_i', s_i') \to (S_i, s_i)$ et un diagramme commutatif de schémas
$$
\xymatrix{
X_i \ar[d] & X_i' \ar[l]^{g_i} \ar[d] \\
S_i & S_i' \ar[l]
}
$$
tels que $g_i$ soit \'etale, que $(X_i)_{s_i} \subset g_i(X_i')$, que les schémas
$X_i'$, $S_i'$ soient affines et que
$\Gamma(X_i', g_i^*\mathcal{F}_i)$ soit un
$\Gamma(S_i', \mathcal{O}_{S_i'})$-module projectif.
Notons que $g_i^*\mathcal{F}_i$ est universellement pur sur $S'_i$, voir le
lemme \ref{lemma-affine-locally-projective-pure}.
Par changement de base, le diagramme ci-dessus donne un diagramme comportant les morphismes
$(S'_{i'}, s'_{i'}) \to (S_{i'}, s_{i'})$ et
$g_{i'} : X'_{i'} \to X_{i'}$ sur $S_{i'}$
pour tout $i' \geq i$, ainsi qu'un diagramme
comportant les morphismes $(S', s') \to (S, s)$ et $g : X' \to X$ sur $S$.

\medskip\noindent
Nous pouvons maintenant utiliser notre critère de pureté.
Soit $W'_i \subset X_i \times_{S_i} S'_i$ l'image du morphisme
\'etale $X'_i \to X_i \times_{S_i} S'_i$. Pour tout $i' \geq i$,
nous avons de même l'image $W'_{i'} \subset X_{i'} \times_{S_{i'}} S'_{i'}$,
ainsi que l'image $W' \subset X \times_S S'$. La formation des images commute
au changement de base ; ainsi, $W'_{i'} = W'_i \times_{S'_i} S'_{i'}$ et
$W' = W_i \times_{S'_i} S'$. Comme
$\mathcal{F}$ est pur le long de $X_s$, le
lemme \ref{lemma-criterion}
entraîne que
\begin{equation}
\label{equation-inclusion}
f^{-1}(\Spec(\mathcal{O}_{S', s'})) \cap
\text{Ass}_{X \times_S S'/S'}(\mathcal{F} \times_S S') \subset W'
\end{equation}
D'après
Compléments sur les morphismes,
lemme \ref{more-morphisms-lemma-relative-assassin-constructible},
les ensembles
$$
E = \{t \in S' \mid \text{Ass}_{X_t}(\mathcal{F}_t) \subset W' \}
\quad\text{and}\quad
E_{i'} = \{t \in S'_{i'} \mid
\text{Ass}_{X_t}(\mathcal{F}_{i', t}) \subset W'_{i'} \}
$$
sont des parties localement constructibles de $S'$ et $S'_{i'}$. D'après
Compléments sur les morphismes,
lemme \ref{more-morphisms-lemma-base-change-assassin-in-U},
$E_{i'}$ est l'image réciproque de $E_i$ par le morphisme
$S'_{i'} \to S'_i$, et $E$ est l'image réciproque de $E_i$ par
le morphisme $S' \to S'_i$. Ainsi,
l'équation (\ref{equation-inclusion})
équivaut à dire que
$\Spec(\mathcal{O}_{S', s'})$ s'envoie dans $E_i$. Comme
$\mathcal{O}_{S', s'} =
\colim_{i' \geq i} \mathcal{O}_{S'_{i'}, s'_{i'}}$,
nous voyons que $\Spec(\mathcal{O}_{S'_{i'}, s'_{i'}})$
s'envoie dans $E_i$ pour un certain $i' \geq i$, voir
Limites, lemme \ref{limits-lemma-limit-contained-in-constructible}.
En appliquant alors le
lemme \ref{lemma-criterion}
à la situation sur $S_{i'}$,
nous concluons que $\mathcal{F}_{i'}$ est pur le long de $(X_{i'})_{s_{i'}}$.
\end{proof}

\begin{lemma}
\label{lemma-flat-finite-presentation-purity-open}
Soit $f : X \to S$ un morphisme de présentation finie.
Soit $\mathcal{F}$ un $\mathcal{O}_X$-module quasi-cohérent
de présentation finie, plat sur $S$. Alors l'ensemble
$$
U = \{s \in S \mid \mathcal{F}\text{ is pure along }X_s\}
$$
est ouvert dans $S$.
\end{lemma}

\begin{proof}
Soit $s \in U$. À l'aide du
lemme \ref{lemma-finite-presentation-flat-along-fibre},
nous pouvons trouver un voisinage \'etale élémentaire
$(S', s') \to (S, s)$ et un diagramme commutatif
$$
\xymatrix{
X \ar[d] & X' \ar[l]^g \ar[d] \\
S & S' \ar[l]
}
$$
tels que $g$ soit \'etale, que $X_s \subset g(X')$, que les schémas
$X'$, $S'$ soient affines et que $\Gamma(X', g^*\mathcal{F})$
soit un $\Gamma(S', \mathcal{O}_{S'})$-module projectif.
Notons que $g^*\mathcal{F}$ est universellement pur sur $S'$, voir le
lemme \ref{lemma-affine-locally-projective-pure}.
Soit $W' \subset X \times_S S'$ l'image du morphisme \'etale
$X' \to X \times_S S'$. Notons que $W$ est ouvert et quasi-compact sur
$S'$. Posons
$$
E = \{t \in S' \mid \text{Ass}_{X_t}(\mathcal{F}_t) \subset W' \}.
$$
D'après
Compléments sur les morphismes,
lemme \ref{more-morphisms-lemma-relative-assassin-constructible},
$E$ est une partie constructible de $S'$. D'après le
lemme \ref{lemma-criterion},
nous avons $\Spec(\mathcal{O}_{S', s'}) \subset E$.
D'après
Morphismes, lemme \ref{morphisms-lemma-constructible-containing-open},
$E$ contient un voisinage ouvert $V'$ de $s'$.
En appliquant le
lemme \ref{lemma-criterion}
une nouvelle fois, nous voyons que, pour tout point $s_1$ de l'image de $V'$ dans $S$,
le faisceau $\mathcal{F}$ est pur le long de $X_{s_1}$. Comme
$S' \to S$ est \'etale, l'image de $V'$ dans $S$ est ouverte, ce qui conclut.
\end{proof}


\section{Applications de la pureté}
\label{section-applications-purity}

\noindent
Voici quelques exemples d'applications de la pureté. Le premier lemme
utilise en fait une forme légèrement plus faible de pureté.

\begin{lemma}
\label{lemma-injectivity-map-source-flat-pure}
Soit $f : X \to S$ un morphisme de type fini.
Soit $\mathcal{F}$ un faisceau quasi-cohérent de type fini sur $X$.
Supposons $S$ local de point fermé $s$.
Supposons $\mathcal{F}$ pur le long de $X_s$ et
$\mathcal{F}$ plat sur $S$.
Soit $\varphi : \mathcal{F} \to \mathcal{G}$ un morphisme de
$\mathcal{O}_X$-modules quasi-cohérents. Alors les conditions suivantes sont équivalentes :
\begin{enumerate}
\item l'application sur les fibres $\varphi_x$ est injective pour tout
$x \in \text{Ass}_{X_s}(\mathcal{F}_s)$, et
\item $\varphi$ est injectif.
\end{enumerate}
\end{lemma}

\begin{proof}
Soit $\mathcal{K} = \Ker(\varphi)$. Notre but est de montrer que
$\mathcal{K} = 0$. Pour cela, il suffit de montrer que
$\text{WeakAss}_X(\mathcal{K}) = \emptyset$, voir
Diviseurs, lemme \ref{divisors-lemma-weakly-ass-zero}.
Nous avons
$\text{WeakAss}_X(\mathcal{K}) \subset \text{WeakAss}_X(\mathcal{F})$, voir
Diviseurs, lemme \ref{divisors-lemma-ses-weakly-ass}.
Comme $\mathcal{F}$ est plat, le
lemme \ref{lemma-bourbaki-finite-type-general-base}
montre que $\text{WeakAss}_X(\mathcal{F}) \subset \text{Ass}_{X/S}(\mathcal{F})$.
Par pureté, tout point $x'$ de $\text{Ass}_{X/S}(\mathcal{F})$
est une généralisation d'un point de $X_s$, et est donc la
spécialisation d'un point $x \in \text{Ass}_{X_s}(\mathcal{F}_s)$, d'après le
lemme \ref{lemma-associated-point-specializes}.
L'injectivité de $\varphi_x$ entraîne donc celle de
$\varphi_{x'}$, d'où $\mathcal{K}_{x'} = 0$.
\end{proof}

\begin{proposition}
\label{proposition-finite-presentation-flat-pure-is-projective}
Soit $f : X \to S$ un morphisme de schémas affine de présentation finie.
Soit $\mathcal{F}$ un $\mathcal{O}_X$-module quasi-cohérent de
présentation finie, plat sur $S$. Alors les conditions suivantes
sont équivalentes :
\begin{enumerate}
\item $f_*\mathcal{F}$ est localement projectif sur $S$, et
\item $\mathcal{F}$ est pur relativement à $S$.
\end{enumerate}
En particulier, pour un homomorphisme d'anneaux $A \to B$ de présentation finie et
un $B$-module de présentation finie $N$, plat sur $A$, nous avons :
$N$ est projectif comme $A$-module si et seulement si $\widetilde{N}$
sur $\Spec(B)$ est pur relativement à $\Spec(A)$.
\end{proposition}

\begin{proof}
L'implication (1) $\Rightarrow$ (2) est le
lemme \ref{lemma-affine-locally-projective-pure}.
Supposons $\mathcal{F}$ pur relativement à $S$.
Notons que, d'après le
lemme \ref{lemma-finite-type-flat-pure-along-fibre-is-universal},
cela entraîne que $\mathcal{F}$ reste pur après tout changement de base. D'après
Descente, lemme \ref{descent-lemma-locally-projective-descends},
il suffit de montrer que $f_*\mathcal{F}$ est localement projectif sur $S$ pour la topologie fpqc.
Choisissons $s \in S$. Nous allons montrer que la restriction de
$f_*\mathcal{F}$ à un voisinage \'etale de $s$ est localement projective.
En effet, d'après le
lemme \ref{lemma-finite-presentation-flat-along-fibre},
après remplacement de $S$ par un voisinage \'etale élémentaire affine
de $s$, nous pouvons supposer qu'il existe un diagramme
$$
\xymatrix{
X \ar[dr] & & X' \ar[ll]^g \ar[ld] \\
& S &
}
$$
de schémas affines et de présentation finie sur $S$,
où $g$ est \'etale, $X_s \subset g(X')$ et
$\Gamma(X', g^*\mathcal{F})$ est un $\Gamma(S, \mathcal{O}_S)$-module projectif.
Notons que, dans ce cas, $g^*\mathcal{F}$ est universellement pur sur $S$, voir le
lemme \ref{lemma-affine-locally-projective-pure}.
Ainsi, d'après le
lemme \ref{lemma-criterion},
l'ouvert $g(X')$ contient les points de
$\text{Ass}_{X/S}(\mathcal{F})$ au-dessus de $\Spec(\mathcal{O}_{S, s})$.
Posons
$$
E = \{t \in S \mid \text{Ass}_{X_t}(\mathcal{F}_t) \subset g(X') \}.
$$
D'après
Compléments sur les morphismes,
lemme \ref{more-morphisms-lemma-relative-assassin-constructible},
$E$ est une partie constructible de $S$. Nous avons vu que
$\Spec(\mathcal{O}_{S, s}) \subset E$. D'après
Morphismes, lemme \ref{morphisms-lemma-constructible-containing-open},
$E$ contient un voisinage ouvert de $s$. Ainsi, après
remplacement de $S$ par un voisinage affine de $s$, nous pouvons supposer que
$\text{Ass}_{X/S}(\mathcal{F}) \subset g(X')$.
D'après le
lemme \ref{lemma-base-change-universally-flat},
cela signifie que
$$
\Gamma(X, \mathcal{F}) \longrightarrow \Gamma(X', g^*\mathcal{F})
$$
est $\Gamma(S, \mathcal{O}_S)$-universellement injective.\newline
D'après Algèbre, lemme \ref{algebra-lemma-pure-submodule-ML},
nous concluons que $\Gamma(X, \mathcal{F})$ est de Mittag-Leffler comme
$\Gamma(S, \mathcal{O}_S)$-module. Comme
$\Gamma(X, \mathcal{F})$ est engendré par un ensemble dénombrable et plat comme
$\Gamma(S, \mathcal{O}_S)$-module, nous concluons
qu'il est projectif d'après
Algèbre, lemme \ref{algebra-lemma-countgen-projective}.
\end{proof}

\noindent
La proposition permet d'améliorer certains résultats précédents.
Le lemme suivant améliore la
proposition \ref{proposition-finite-presentation-flat-at-point}.

\begin{lemma}
\label{lemma-flat-finite-presentation-affine-neighbourhood-projective}
Soit $f : X \to S$ un morphisme localement de présentation finie.
Soit $\mathcal{F}$ un $\mathcal{O}_X$-module quasi-cohérent
de présentation finie. Soit $x \in X$ avec $s = f(x) \in S$.
Si $\mathcal{F}$ est plat en $x$ sur $S$, il existe un voisinage
\'etale élémentaire affine $(S', s') \to (S, s)$ et
un ouvert affine $U' \subset X \times_S S'$ contenant $x' = (x, s')$
tels que $\Gamma(U', \mathcal{F}|_{U'})$ soit un
$\Gamma(S', \mathcal{O}_{S'})$-module projectif.
\end{lemma}

\begin{proof}
Au cours de la démonstration, nous pouvons remplacer $X$ par un voisinage ouvert de $x$
et $S$ par un voisinage \'etale élémentaire de $s$.
Ainsi, par ouverture du lieu de platitude (voir
Compléments sur les morphismes, théorème \ref{more-morphisms-theorem-openness-flatness}),
nous pouvons supposer $\mathcal{F}$ plat sur $S$.
Nous pouvons supposer $S$ et $X$ affines.
Après une nouvelle restriction de $X$, nous pouvons supposer que tout
point de $\text{Ass}_{X_s}(\mathcal{F}_s)$ est une généralisation de $x$.
Cette propriété est conservée lorsqu'on remplace $(S, s)$ par un voisinage
\'etale élémentaire. Nous pouvons donc appliquer le
lemme \ref{lemma-finite-presentation-flat-along-fibre}
pour nous ramener à la situation où il existe un diagramme
$$
\xymatrix{
X \ar[dr] & & X' \ar[ll]^g \ar[ld] \\
& S &
}
$$
de schémas affines et de présentation finie sur $S$,
où $g$ est \'etale, $X_s \subset g(X')$ et
$\Gamma(X', g^*\mathcal{F})$ est un $\Gamma(S, \mathcal{O}_S)$-module projectif.
Notons que, dans ce cas, $g^*\mathcal{F}$ est universellement pur sur $S$, voir le
lemme \ref{lemma-affine-locally-projective-pure}.

\medskip\noindent
Soit $U \subset g(X')$ un voisinage ouvert affine de $x$.
Nous affirmons que $\mathcal{F}|_U$ est pur le long de $U_s$. Si cela est démontré,
le lemme en résulte, car $\mathcal{F}|_U$ sera pur relativement à $S$
après restriction de $S$, voir le
lemme \ref{lemma-flat-finite-presentation-purity-open},
et la projectivité résultera alors de la
proposition \ref{proposition-finite-presentation-flat-pure-is-projective}.
Pour démontrer l'assertion, il faut montrer, après remplacement de $(S, s)$
par un voisinage \'etale élémentaire quelconque, que tout point $\xi$ de
$\text{Ass}_{U/S}(\mathcal{F}|_U)$ au-dessus d'un
$s' \in S$, $s' \leadsto s$, se spécialise en un point de $U_s$.
Comme $U \subset g(X')$, nous pouvons trouver un $\xi' \in X'$ tel que
$g(\xi') = \xi$. Puisque $g^*\mathcal{F}$ est pur sur $S$, le
lemme \ref{lemma-associated-point-specializes}
montre qu'il existe une spécialisation $\xi' \leadsto x'$ avec
$x' \in \text{Ass}_{X'_s}(g^*\mathcal{F}_s)$. Alors
$g(x') \in \text{Ass}_{X_s}(\mathcal{F}_s)$ (voir par exemple le
lemme \ref{lemma-etale-weak-assassin-up-down}
appliqué au morphisme \'etale $X'_s \to X_s$ de schémas noethériens),
donc $g(x') \leadsto x$ par le choix de $X$ ci-dessus ! Puisque
$x \in U$, nous concluons que $g(x') \in U$. Ainsi,
$\xi = g(\xi') \leadsto g(x') \in U_s$, comme voulu.
\end{proof}

\noindent
Le lemme suivant améliore le
lemme \ref{lemma-finite-type-flat-at-point-free-variant}.

\begin{lemma}
\label{lemma-flat-finite-type-affine-neighbourhood-projective}
Soit $f : X \to S$ un morphisme localement de type fini.
Soit $\mathcal{F}$ un $\mathcal{O}_X$-module quasi-cohérent
de type fini. Soit $x \in X$ avec $s = f(x) \in S$.
Si $\mathcal{F}$ est plat en $x$ sur $S$, il existe un voisinage
\'etale élémentaire affine $(S', s') \to (S, s)$ et
un ouvert affine $U' \subset X \times_S \Spec(\mathcal{O}_{S', s'})$
contenant $x' = (x, s')$ tels que
$\Gamma(U', \mathcal{F}|_{U'})$ soit un
$\mathcal{O}_{S', s'}$-module libre.
\end{lemma}

\begin{proof}
La question est locale pour la topologie de Zariski sur $X$ et $S$. Nous pouvons donc supposer
$X$ et $S$ affines. Nous pouvons alors trouver une immersion fermée
$i : X \to \mathbf{A}^n_S$ sur $S$. Il est clair qu'il suffit de
démontrer le lemme pour le faisceau $i_*\mathcal{F}$ sur $\mathbf{A}^n_S$
et le point $i(x)$. Nous nous ramenons ainsi au cas où $X \to S$ est
de présentation finie. Après remplacement de $S$ par
$\Spec(\mathcal{O}_{S', s'})$ et de $X$ par un ouvert de
$X \times_S \Spec(\mathcal{O}_{S', s'})$, nous pouvons supposer que
$\mathcal{F}$ est de présentation finie, voir la
proposition \ref{proposition-finite-type-flat-at-point}.
Dans ce cas, nous pouvons invoquer le
lemme \ref{lemma-flat-finite-presentation-affine-neighbourhood-projective}
et
Algèbre, théorème \ref{algebra-theorem-projective-free-over-local-ring},
pour conclure.
\end{proof}

\begin{lemma}
\label{lemma-flat-finite-type-local-colimit-free}
Soit $A \to B$ un homomorphisme local d'anneaux locaux essentiellement de
type fini. Soit $N$ un $B$-module de type fini, plat comme $A$-module.
Si $A$ est hensélien, alors $N$ est une limite inductive filtrante
$$
N = \colim_i F_i
$$
de $A$-modules libres $F_i$ telle que toutes les applications de transition
$u_i : F_i \to F_{i'}$ du système induisent des applications injectives
$\overline{u}_i : F_i/\mathfrak m_AF_i \to F_{i'}/\mathfrak m_AF_{i'}$.
De plus, $N$ est un $A$-module de Mittag-Leffler.
\end{lemma}

\begin{proof}
Nous pouvons trouver un morphisme de type fini $X \to S = \Spec(A)$,
un point $x \in X$ au-dessus du point fermé $s$ de $S$ et un
$\mathcal{O}_X$-module quasi-cohérent de type fini $\mathcal{F}$ tels que
$\mathcal{F}_x \cong N$ comme $A$-module. Après restriction de $X$,
nous pouvons supposer que tout point de $\text{Ass}_{X_s}(\mathcal{F}_s)$ se spécialise
en $x$. D'après le
lemme \ref{lemma-flat-finite-type-affine-neighbourhood-projective},
il existe un système fondamental de voisinages ouverts affines
$U_i \subset X$ de $x$ tels que $\Gamma(U_i, \mathcal{F})$ soit
un $A$-module libre $F_i$. Notons que, si $U_{i'} \subset U_i$, alors
$$
F_i/\mathfrak m_AF_i = \Gamma(U_{i, s}, \mathcal{F}_s)
\longrightarrow
\Gamma(U_{i', s}, \mathcal{F}_s) = F_{i'}/\mathfrak m_AF_{i'}
$$
est injective, car une section du noyau aurait son support dans
une partie fermée de $X_s$ ne rencontrant pas $x$, ce qui contredit
le choix de $X$ ci-dessus. Comme les applications $F_i \to F_{i'}$ sont
$A$-universellement injectives (lemme \ref{lemma-universally-injective-local}),
il en résulte que $N$ est
de Mittag-Leffler d'après
Algèbre, lemme \ref{algebra-lemma-colimit-universally-injective-ML}.
\end{proof}

\noindent
On pourra omettre le lemme suivant lors d'une première lecture.

\begin{lemma}
\label{lemma-flat-finite-type-local-valuation-ring-has-content}
Soit $A \to B$ un homomorphisme local d'anneaux locaux essentiellement de
type fini. Soit $N$ un $B$-module de type fini, plat comme $A$-module.
Si $A$ est un anneau de valuation, alors tout élément de $N$ possède un idéal de contenu
$I \subset A$ (Compléments d'algèbre, définition
\ref{more-algebra-definition-content-ideal}). De plus, $I$ est un
idéal principal.
\end{lemma}

\begin{proof}
La dernière assertion résulte du fait que $I$ est un idéal de type
fini d'après
Compléments d'algèbre, lemme \ref{more-algebra-lemma-content-finitely-generated},
et Algèbre, lemme \ref{algebra-lemma-characterize-valuation-ring}.

\medskip\noindent
Démontrons l'existence de $I$.
Soit $A \subset A^h$ l'hensélisé. Soit $B'$ le localisé
de $B \otimes_A A^h$ en l'idéal maximal
$\mathfrak m_B \otimes A^h + B \otimes \mathfrak m_{A^h}$.
Alors $B \to B'$ est plat, donc fidèlement plat.
Soit $N' = N \otimes_B B'$.
Soit $x \in N$ et soit $x' \in N'$ son image.
Nous affirmons que, pour un idéal $I \subset A$, nous avons
$x \in IN \Leftrightarrow x' \in IN'$.
En effet, $N/IN \to N'/IN'$ est le produit tensoriel de $B \to B'$
avec $N/IN$, et $B \to B'$ est universellement injectif d'après
Algèbre, lemme \ref{algebra-lemma-faithfully-flat-universally-injective}.
D'après Compléments d'algèbre, lemme \ref{more-algebra-lemma-henselization-valuation-ring},
et Algèbre, lemme \ref{algebra-lemma-ideals-valuation-ring},
l'homomorphisme $A \to A^h$ définit une bijection préservant
l'inclusion $I \mapsto IA^h$ entre les ensembles d'idéaux. Nous concluons que
$x$ possède un idéal de contenu dans $A$ si et seulement si $x'$ possède un idéal de
contenu dans $A^h$. L'assertion pour $x' \in N'$ résulte du
lemme \ref{lemma-flat-finite-type-local-colimit-free} et d'
Algèbre, lemme \ref{algebra-lemma-minimal-contains}.
\end{proof}

\noindent
Voici une application.

\begin{lemma}
\label{lemma-proper-flat-over-dvr-reduced-fibre}
Soit $X \to \Spec(R)$ un morphisme propre et plat, où $R$ est un anneau de valuation.
Si la fibre spéciale est réduite, alors $X$ et toutes les fibres de $X \to \Spec(R)$
sont réduits.
\end{lemma}

\begin{proof}
Supposons la fibre spéciale $X_s$ réduite.
Soit $x \in X$ un point quelconque et montrons que $\mathcal{O}_{X, x}$
est réduit ; cela prouvera que $X$ est réduit.
Soit $x \leadsto x'$ une spécialisation avec $x'$
dans la fibre spéciale ; une telle spécialisation existe,
car un morphisme propre est fermé. Considérons l'anneau
local $A = \mathcal{O}_{X, x'}$. Alors $\mathcal{O}_{X, x}$
est un localisé de $A$ ; il suffit donc de montrer que
$A$ est réduit. Soit $a \in A$ non nul et soit $I = (\pi) \subset R$ son
idéal de contenu, voir le
lemme \ref{lemma-flat-finite-type-local-valuation-ring-has-content}.
Alors $a = \pi a'$ et $a'$ s'envoie sur un élément non nul de
$A/\mathfrak mA$, où $\mathfrak m \subset R$ est l'idéal maximal.
Si $a$ est nilpotent, il en est de même de $a'$, car $\pi$ n'est pas diviseur de zéro
par platitude de $A$ sur $R$.
Mais $a'$ s'envoie sur un élément non nul de l'anneau réduit
$A/\mathfrak m A = \mathcal{O}_{X_s, x'}$.
Cela donne une contradiction à moins que $A$ ne soit réduit, ce
que nous voulions montrer.

\medskip\noindent
Bien entendu, si $X$ est réduit, il en est de même de la fibre générique de $X$ sur $R$.
Si $\mathfrak p \subset R$ est un
idéal premier, alors $R/\mathfrak p$ est un anneau de valuation d'après
Algèbre, lemme \ref{algebra-lemma-make-valuation-rings}.
En reprenant l'argument pour le changement de base de $X$
à $R/\mathfrak p$, on montre donc que la fibre au-dessus de $\mathfrak p$
est réduite.
\end{proof}






\section{Foncteurs d'aplatissement}
\label{section-flattening-functors}

\noindent
Soit $S$ un schéma. Rappelons qu'un foncteur
$F : (\Sch/S)^{opp} \to \textit{Sets}$ est dit compatible aux limites
si, pour tout système projectif filtrant
$\{T_i\}_{i \in I}$ de schémas affines de limite $T$, nous avons
$F(T) = \colim_i F(T_i)$.

\begin{situation}
\label{situation-iso}
Soit $f : X \to S$ un morphisme de schémas.
Soit $u : \mathcal{F} \to \mathcal{G}$ un homomorphisme de
$\mathcal{O}_X$-modules quasi-cohérents. Pour tout schéma $T$ sur
$S$, nous noterons $u_T : \mathcal{F}_T \to \mathcal{G}_T$ le
changement de base de $u$ à $T$ ; autrement dit, $u_T$ est l'image réciproque
de $u$ par la projection $X_T = X \times_S T \to X$.
Dans cette situation, nous pouvons considérer le foncteur
\begin{equation}
\label{equation-iso}
F_{iso} : (\Sch/S)^{opp} \longrightarrow \textit{Sets}, \quad
T \longrightarrow \left\{
\begin{matrix}
\{*\} & \text{if} & u_T \text{ is an isomorphism}, \\
\emptyset & \text{else.} &
\end{matrix}
\right.
\end{equation}
Il existe des variantes $F_{inj}$, $F_{surj}$, $F_{zero}$ où l'on demande que
$u_T$ soit injectif, surjectif ou nul.
\end{situation}

\begin{lemma}
\label{lemma-iso-sheaf}
Plaçons-nous dans la situation \ref{situation-iso}.
\begin{enumerate}
\item Chacun des foncteurs $F_{iso}$, $F_{inj}$, $F_{surj}$, $F_{zero}$
satisfait à la condition de faisceau pour la topologie fpqc.
\item Si $f$ est quasi-compact et $\mathcal{G}$ de type fini,
alors $F_{surj}$ est compatible aux limites.
\item Si $f$ est quasi-compact et $\mathcal{F}$ de type fini, alors
$F_{zero}$ est compatible aux limites.
\item Si $f$ est quasi-compact, $\mathcal{F}$ de type fini et
$\mathcal{G}$ de présentation finie, alors $F_{iso}$ est compatible aux limites.
\end{enumerate}
\end{lemma}

\begin{proof}
Soit $\{T_i \to T\}_{i \in I}$ un recouvrement fpqc de schémas sur $S$.
Posons $X_i = X_{T_i} = X \times_S T_i$ et $u_i = u_{T_i}$.
Notons que $\{X_i \to X_T\}_{i \in I}$ est un recouvrement fpqc de $X_T$, voir
Topologies, lemme \ref{topologies-lemma-fpqc}.
En particulier, pour tout $x \in X_T$, il existe un $i \in I$
et un $x_i \in X_i$ d'image $x$. Comme
$\mathcal{O}_{X_T, x} \to \mathcal{O}_{X_i, x_i}$ est plat, donc
fidèlement plat (voir
Algèbre, lemme \ref{algebra-lemma-local-flat-ff}),
nous concluons que $(u_i)_{x_i}$ est injectif, surjectif, bijectif ou nul
si et seulement si $(u_T)_x$ est injectif, surjectif, bijectif ou nul.
D'où la partie (1) du lemme.

\medskip\noindent
Démonstration de (2). Supposons $f$ quasi-compact et $\mathcal{G}$ de type fini.
Soit $T = \lim_{i \in I} T_i$ une limite dirigée de $S$-schémas affines
et supposons $u_T$ surjectif.
Posons $X_i = X_{T_i} = X \times_S T_i$ et
$u_i = u_{T_i} : \mathcal{F}_i = \mathcal{F}_{T_i}
\to \mathcal{G}_i = \mathcal{G}_{T_i}$.
Pour démontrer (2), il faut montrer que $u_i$ est surjectif pour un certain $i$.
Choisissons $i_0 \in I$ et remplaçons $I$ par $\{i \mid i \geq i_0\}$.
Comme $f$ est quasi-compact, le schéma $X_{i_0}$ est quasi-compact.
Nous pouvons donc choisir des ouverts affines $W_1, \ldots, W_m \subset X$
et un recouvrement ouvert affine
$X_{i_0} = U_{1, i_0} \cup \ldots \cup U_{m, i_0}$ tels que
$U_{j, i_0}$ s'envoie dans $W_j$ par la projection $X_{i_0} \to X$.
Pour tout $i \in I$, soit $U_{j, i}$ l'image réciproque de $U_{j, i_0}$.
En posant $U_j = \lim_i U_{j, i}$, nous voyons que $X_T = U_1 \cup \ldots \cup U_m$
est un recouvrement ouvert affine de $X_T$. Il suffit alors de montrer, pour un
$j \in \{1, \ldots, m\}$ donné, que $u_i|_{U_{j, i}}$ est surjectif pour un certain
$i = i(j) \in I$. À l'aide de
Propriétés, lemme \ref{properties-lemma-finite-type-module},
cela se traduit par le problème algébrique suivant :
soient $A$ un anneau et $u : M \to N$ une application de $A$-modules.
Supposons que $R = \colim_{i \in I} R_i$ soit une limite inductive filtrante
de $A$-algèbres. Si $N$ est un $A$-module de type fini et si
$u \otimes 1 : M \otimes_A R \to N \otimes_A R$ est surjective, alors,
pour un certain $i$, l'application
$u \otimes 1 : M \otimes_A R_i \to N \otimes_A R_i$ est surjective.
C'est
Algèbre, lemme \ref{algebra-lemma-module-map-property-in-colimit}, partie (2).

\medskip\noindent
Démonstration de (3). Les mêmes arguments que dans la démonstration de (2)
ramènent au problème algébrique suivant :
soient $A$ un anneau et $u : M \to N$ une application de $A$-modules.
Supposons que $R = \colim_{i \in I} R_i$ soit une limite inductive filtrante
de $A$-algèbres. Si $M$ est un $A$-module de type fini et si
$u \otimes 1 : M \otimes_A R \to N \otimes_A R$ est nulle, alors,
pour un certain $i$, l'application
$u \otimes 1 : M \otimes_A R_i \to N \otimes_A R_i$ est nulle.
C'est
Algèbre, lemme \ref{algebra-lemma-module-map-property-in-colimit}, partie (1).

\medskip\noindent
Démonstration de (4).
Supposons $f$ quasi-compact et $\mathcal{F}, \mathcal{G}$ de présentation finie.
En raisonnant exactement comme dans le paragraphe précédent
(et en utilisant en outre
Propriétés, lemme \ref{properties-lemma-finite-presentation-module}),
la partie (3) se traduit par l'énoncé algébrique suivant :
soient $A$ un anneau et $u : M \to N$ une application de $A$-modules.
Supposons que $R = \colim_{i \in I} R_i$ soit une limite inductive filtrante
de $A$-algèbres. Supposons $M$ de type fini comme $A$-module, $N$ de présentation
finie comme $A$-module, et
$u \otimes 1 : M \otimes_A R \to N \otimes_A R$ un isomorphisme.
Alors, pour un certain $i$, l'application
$u \otimes 1 : M \otimes_A R_i \to N \otimes_A R_i$ est un isomorphisme.
C'est
Algèbre, lemme \ref{algebra-lemma-module-map-property-in-colimit}, partie (3).
\end{proof}

\begin{situation}
\label{situation-flat-at-point}
Soit $(A, \mathfrak m_A)$ un anneau local.
Notons $\mathcal{C}$ la catégorie dont les objets sont les $A$-algèbres
$A'$ qui sont des anneaux locaux et dont la structure d'algèbre
$A \to A'$ est un homomorphisme local d'anneaux locaux.
Un morphisme entre objets $A', A''$ de $\mathcal{C}$ est un
homomorphisme local $A' \to A''$ de $A$-algèbres.
Soient $A \to B$ un homomorphisme local d'anneaux locaux et $M$ un $B$-module.
Si $A'$ est un objet de $\mathcal{C}$, nous posons $B' = B \otimes_A A'$
et $M' = M \otimes_A A'$, considéré comme $B'$-module.
Pour $A' \in \Ob(\mathcal{C})$, considérons la condition
\begin{equation}
\label{equation-flat-at-primes}
\forall \mathfrak q \in V(\mathfrak m_{A'}B' + \mathfrak m_B B')
\subset \Spec(B') :
M'_{\mathfrak q}\text{ is flat over }A'.
\end{equation}
Notons l'analogie avec
Compléments d'algèbre, équation (\ref{more-algebra-equation-flat-at-primes}).
En particulier, si $A' \to A''$ est un morphisme de $\mathcal{C}$ et si
(\ref{equation-flat-at-primes}) est satisfaite pour $A'$, elle l'est pour $A''$, voir
Compléments d'algèbre,
lemme \ref{more-algebra-lemma-base-change-flat-at-primes}.
Nous obtenons donc un foncteur
\begin{equation}
\label{equation-flat-at-point}
F_{lf} : \mathcal{C} \longrightarrow \textit{Sets}, \quad
A' \longrightarrow \left\{
\begin{matrix}
\{*\} & \text{if }(\ref{equation-flat-at-primes})\text{ holds}, \\
\emptyset & \text{else.} &
\end{matrix}
\right.
\end{equation}
\end{situation}

\begin{lemma}
\label{lemma-flat-at-point}
Plaçons-nous dans la situation \ref{situation-flat-at-point}.
\begin{enumerate}
\item Si $A' \to A''$ est un morphisme plat dans $\mathcal{C}$,
alors $F_{lf}(A') = F_{lf}(A'')$.
\item Si $A \to B$ est essentiellement de présentation finie et
$M$ est un $B$-module de présentation finie, alors $F_{lf}$ est compatible aux
limites : si $\{A_i\}_{i \in I}$ est un
système inductif filtrant d'objets de $\mathcal{C}$, alors
$F_{lf}(\colim_i A_i) = \colim_i F_{lf}(A_i)$.
\end{enumerate}
\end{lemma}

\begin{proof}
La partie (1) est un cas particulier de
Compléments d'algèbre,
lemme \ref{more-algebra-lemma-flat-descent-flat-at-primes}.
La partie (2) est un cas particulier de
Compléments d'algèbre,
lemme \ref{more-algebra-lemma-limit-preserving-flat-at-primes}.
\end{proof}

\begin{lemma}
\label{lemma-flat-at-point-finite}
Dans la situation \ref{situation-flat-at-point}, soient $B \to C$ un homomorphisme local de
$A$-algèbres locales et $N$ un $C$-module. Notons
$F'_{lf} : \mathcal{C} \to \textit{Sets}$ le foncteur associé au couple
$(C, N)$. Si $M \cong N$ comme $B$-modules et si $B \to C$ est fini, alors
$F_{lf} = F'_{lf}$.
\end{lemma}

\begin{proof}
Soit $A'$ un objet de $\mathcal{C}$. Posons $C' = C \otimes_A A'$
et $N' = N \otimes_A A'$, comme dans les définitions de $B'$, $M'$ de la
situation \ref{situation-flat-at-point}.
Notons que $M' \cong N'$ comme $B'$-modules.
L'hypothèse que $B \to C$ est fini a deux conséquences :
(a) $\mathfrak m_C = \sqrt{\mathfrak m_B C}$ et (b)
$B' \to C'$ est fini. La conséquence (a) entraîne que
$$
V(\mathfrak m_{A'}C' + \mathfrak m_C C')
=
\left(
\Spec(C') \to \Spec(B')
\right)^{-1}V(\mathfrak m_{A'}B' + \mathfrak m_B B').
$$
Supposons $\mathfrak q \subset V(\mathfrak m_{A'}B' + \mathfrak m_B B')$.
Alors $M'_{\mathfrak q}$ est plat sur $A'$ si et seulement si
le $C'_{\mathfrak q}$-module $N'_{\mathfrak q}$ est plat sur $A'$
(car ils sont isomorphes comme $A'$-modules), si et seulement si,
pour tout idéal maximal $\mathfrak r$ de $C'_{\mathfrak q}$,
le module $N'_{\mathfrak r}$ est plat sur $A'$ (voir
Algèbre, lemme \ref{algebra-lemma-flat-localization}).
Comme $B'_{\mathfrak q} \to C'_{\mathfrak q}$ est fini d'après (b),
les idéaux maximaux de $C'_{\mathfrak q}$ correspondent exactement
aux idéaux premiers de $C'$ au-dessus de $\mathfrak q$ (voir
Algèbre, lemme \ref{algebra-lemma-integral-going-up}),
et ces idéaux premiers appartiennent tous à
$V(\mathfrak m_{A'}C' + \mathfrak m_C C')$ d'après
l'égalité ci-dessus. Le résultat du lemme en découle.
\end{proof}

\begin{lemma}
\label{lemma-flat-at-point-go-up}
Dans la
situation \ref{situation-flat-at-point},
supposons que $B \to C$ soit un homomorphisme local plat d'anneaux
locaux. Posons $N = M \otimes_B C$. Notons
$F'_{lf} : \mathcal{C} \to \textit{Sets}$ le foncteur associé
au couple $(C, N)$. Alors $F_{lf} = F'_{lf}$.
\end{lemma}

\begin{proof}
Soit $A'$ un objet de $\mathcal{C}$. Posons $C' = C \otimes_A A'$
et $N' = N \otimes_A A' = M' \otimes_{B'} C'$, comme dans les définitions
de $B'$, $M'$ de la
situation \ref{situation-flat-at-point}. Notons que
$$
V(\mathfrak m_{A'}B' + \mathfrak m_B B')
=
\Spec( \kappa(\mathfrak m_B) \otimes_A \kappa(\mathfrak m_{A'}) )
$$
et de même pour $V(\mathfrak m_{A'}C' + \mathfrak m_C C')$.
L'homomorphisme d'anneaux
$$
\kappa(\mathfrak m_B) \otimes_A \kappa(\mathfrak m_{A'})
\longrightarrow
\kappa(\mathfrak m_C) \otimes_A \kappa(\mathfrak m_{A'})
$$
est fidèlement plat, donc
$V(\mathfrak m_{A'}C' + \mathfrak m_C C') \to
V(\mathfrak m_{A'}B' + \mathfrak m_B B')$ est surjectif.
Enfin, si $\mathfrak r \in V(\mathfrak m_{A'}C' + \mathfrak m_C C')$
s'envoie sur $\mathfrak q \in V(\mathfrak m_{A'}B' + \mathfrak m_B B')$, alors
$M'_{\mathfrak q}$ est plat sur $A'$ si et seulement si
$N'_{\mathfrak r}$ est plat sur $A'$, car $B' \to C'$ est plat, voir
Algèbre, lemme \ref{algebra-lemma-flatness-descends-more-general}.
Le lemme résulte formellement de ces remarques.
\end{proof}

\begin{situation}
\label{situation-free-at-generic-points}
Soit $f : X \to S$ un morphisme lisse à fibres géométriquement
irréductibles. Soit $\mathcal{F}$ un $\mathcal{O}_X$-module quasi-cohérent de
type fini. Pour tout schéma $T$ sur $S$, nous noterons
$\mathcal{F}_T$ le changement de base de $\mathcal{F}$ à $T$ ; autrement
dit, $\mathcal{F}_T$ est l'image réciproque de $\mathcal{F}$ par la
projection $X_T = X \times_S T \to X$. Notons que $X_T \to T$
est lisse à fibres géométriquement irréductibles, voir
Morphismes, lemme \ref{morphisms-lemma-base-change-smooth}, et
Compléments sur les morphismes,
lemme \ref{more-morphisms-lemma-base-change-fibres-geometrically-irreducible}.
Soit $p \geq 0$ un entier. Pour un point $t \in T$, considérons la
condition
\begin{equation}
\label{equation-free-at-generic-point-fibre}
\mathcal{F}_T \text{ is free of rank }p\text{ in a neighbourhood of }\xi_t
\end{equation}
où $\xi_t$ est le point générique de la fibre $X_t$. Cette condition
pour tout $t \in T$ est stable par changement de base ; nous obtenons donc un foncteur
\begin{equation}
\label{equation-free-at-generic-points}
H_p : (\Sch/S)^{opp} \longrightarrow \textit{Sets}, \quad
T \longrightarrow \left\{
\begin{matrix}
\{*\} & \text{if }\mathcal{F}_T\text{ satisfies
(\ref{equation-free-at-generic-point-fibre}) }\forall t\in T, \\
\emptyset & \text{else.}
\end{matrix}
\right.
\end{equation}
\end{situation}

\begin{lemma}
\label{lemma-free-at-generic-points}
Plaçons-nous dans la situation \ref{situation-free-at-generic-points}.
\begin{enumerate}
\item Le foncteur $H_p$ satisfait à la condition de faisceau pour la topologie fpqc.
\item Si $\mathcal{F}$ est de présentation finie, alors le foncteur $H_p$ est
compatible aux limites.
\end{enumerate}
\end{lemma}

\begin{proof}
Soit $\{T_i \to T\}_{i \in I}$ un recouvrement fpqc\footnote{Il est assez facile de
montrer que $H_p$
est un faisceau pour la topologie fppf en utilisant le fait que les morphismes plats de présentation
finie sont ouverts. C'est tout ce dont nous aurons besoin par la suite. Mais il est assez
plaisant de démontrer directement qu'il satisfait aussi à la condition de faisceau
pour la topologie fpqc.} de schémas sur $S$.
Posons $X_i = X_{T_i} = X \times_S T_i$ et notons $\mathcal{F}_i$
l'image réciproque de $\mathcal{F}$ sur $X_i$. Supposons que $\mathcal{F}_i$ satisfasse à
(\ref{equation-free-at-generic-point-fibre}) pour tout $i$.
Choisissons $t \in T$ et notons $\xi_t \in X_T$ le point générique de $X_t$.
Il faut montrer que $\mathcal{F}$ est libre dans un voisinage de $\xi_t$.
Pour un certain $i \in I$, nous pouvons trouver un $t_i \in T_i$ d'image $t$.
Notons $\xi_i \in X_i$ le point générique de $X_{t_i}$, de sorte que
$\xi_i$ s'envoie sur $\xi_t$. Le fait que $\mathcal{F}_i$ soit libre de rang
$p$ dans un voisinage de $\xi_i$ entraîne que
$(\mathcal{F}_i)_{x_i} \cong \mathcal{O}_{X_i, x_i}^{\oplus p}$,
ce qui entraîne que
$\mathcal{F}_{T, \xi_t} \cong \mathcal{O}_{X_T, \xi_t}^{\oplus p}$,
puisque $\mathcal{O}_{X_T, \xi_t} \to \mathcal{O}_{X_i, x_i}$ est plat, voir
par exemple
Algèbre, lemme \ref{algebra-lemma-finite-projective-descends}.
Il existe donc un voisinage affine $U$ de $\xi_t$ dans $X_T$
et une surjection
$\mathcal{O}_U^{\oplus p} \to \mathcal{F}_U = \mathcal{F}_T|_U$, voir
Modules, lemme \ref{modules-lemma-finite-type-surjective-on-stalk}.
Après restriction de $T$, nous pouvons supposer $U \to T$ surjectif. Ainsi,
$U \to T$ est un morphisme lisse de schémas affines à fibres géométriquement
irréductibles. De plus, pour tout $t' \in T$, l'application induite
$$
\alpha :
\mathcal{O}_{U, \xi_{t'}}^{\oplus p}
\longrightarrow
\mathcal{F}_{U, \xi_{t'}}
$$
est un isomorphisme (car, par le même argument que précédemment, le module de droite
est libre de rang $p$). Il résulte du
lemme \ref{lemma-induction-step}
que
$$
\Gamma(U, \mathcal{O}_U^{\oplus p})
\otimes_{\Gamma(T, \mathcal{O}_T)} \mathcal{O}_{T, t'}
\longrightarrow
\Gamma(U, \mathcal{F}_U)
\otimes_{\Gamma(T, \mathcal{O}_T)} \mathcal{O}_{T, t'}
$$
est injective pour tout $t' \in T$. La surjection $\alpha$ est donc un
isomorphisme. Cela achève la démonstration de (1).

\medskip\noindent
Supposons $\mathcal{F}$ de présentation finie.
Soit $T = \lim_{i \in I} T_i$ une limite dirigée de $S$-schémas affines
et supposons que $\mathcal{F}_T$ satisfasse à
(\ref{equation-free-at-generic-point-fibre}).
Posons $X_i = X_{T_i} = X \times_S T_i$ et notons $\mathcal{F}_i$
l'image réciproque de $\mathcal{F}$ sur $X_i$.
Notons $U \subset X_T$ le sous-schéma ouvert des points où
$\mathcal{F}_T$ est plat sur $T$, voir
Compléments sur les morphismes, théorème \ref{more-morphisms-theorem-openness-flatness}.
Par hypothèse, tout point générique de toute fibre appartient à $U$, c'est-à-dire que
$U \to T$ est un morphisme lisse surjectif à fibres géométriquement irréductibles.
Nous pouvons restreindre $U$ et supposer $U$ quasi-compact.
À l'aide de
Limites, lemme \ref{limits-lemma-descend-opens},
nous pouvons trouver un $i \in I$ et un ouvert quasi-compact $U_i \subset X_i$
dont l'image réciproque dans $X_T$ est $U$. En augmentant $i$, nous pouvons
supposer $\mathcal{F}_i|_{U_i}$ plat sur $T_i$, voir
Limites, lemme \ref{limits-lemma-descend-module-flat-finite-presentation}.
En particulier, $\mathcal{F}_i|_{U_i}$ est localement libre de type fini,
donc définit une fonction rang
localement constante $\rho : U_i \to \{0, 1, 2, \ldots \}$.
Notons $(U_i)_p \subset U_i$ la partie ouverte et fermée
où $\rho$ a pour valeur $p$. Soit $V_i \subset T_i$
l'image de $(U_i)_p$ ; notons que $V_i$ est ouvert et quasi-compact.
Par hypothèse, l'image de $T \to T_i$ est contenue dans $V_i$.
Il existe donc un $i' \geq i$ tel que $T_{i'} \to T_i$ se factorise
par $V_i$, d'après
Limites, lemme \ref{limits-lemma-descend-opens}.
Alors $\mathcal{F}_{i'}$ satisfait à (\ref{equation-free-at-generic-point-fibre}),
comme voulu. Certains détails sont omis.
\end{proof}

\begin{lemma}
\label{lemma-pre-flat-dimension-n}
Soit $f : X \to S$ un morphisme de schémas localement de type fini.
Soit $\mathcal{F}$ un $\mathcal{O}_X$-module quasi-cohérent de type
fini. Soit $n \geq 0$. Les conditions suivantes sont équivalentes :
\begin{enumerate}
\item pour $s \in S$, la partie fermée $Z \subset X_s$ des points
où $\mathcal{F}$ n'est pas plat sur $S$ (voir le
lemme \ref{lemma-open-in-fibre-where-flat})
vérifie $\dim(Z) < n$, et
\item pour $x \in X$ tel que $\mathcal{F}$ ne soit pas plat en $x$
sur $S$, nous avons\newline $\text{trdeg}_{\kappa(f(x))}(\kappa(x)) < n$.
\end{enumerate}
Si ces conditions sont satisfaites, elles le restent après tout changement de base.
\end{lemma}

\begin{proof}
Soit $x \in X$ un point au-dessus de $s \in S$.
Alors la dimension de l'adhérence de $\{x\}$ dans $X_s$
est $\text{trdeg}_{\kappa(s)}(\kappa(x))$, d'après
Variétés, lemme \ref{varieties-lemma-dimension-locally-algebraic}.
Réciproquement, si $Z \subset X_s$ est une partie fermée
de dimension $d$, il existe un point $x \in Z$
tel que $\text{trdeg}_{\kappa(s)}(\kappa(x)) = d$ (même référence).
Les conditions (1) et (2) sont donc équivalentes (même fibre par fibre).
L'assertion sur le changement de base résulte de
Morphismes, lemmes \ref{morphisms-lemma-base-change-module-flat} et
\ref{morphisms-lemma-dimension-fibre-after-base-change}.
\end{proof}

\begin{definition}
\label{definition-flat-dimension-n}
Soit $f : X \to S$ un morphisme de schémas localement de type fini.
Soit $\mathcal{F}$ un $\mathcal{O}_X$-module quasi-cohérent de type fini.
Soit $n \geq 0$.
Nous disons que {\it $\mathcal{F}$ est plat sur $S$ en dimensions $\geq n$}
si les conditions équivalentes du lemme \ref{lemma-pre-flat-dimension-n}
sont satisfaites.
\end{definition}

\begin{situation}
\label{situation-flat-dimension-n}
Soit $f : X \to S$ un morphisme de schémas localement de type fini.
Soit $\mathcal{F}$ un $\mathcal{O}_X$-module quasi-cohérent de type
fini. Pour tout schéma $T$ sur $S$, nous noterons $\mathcal{F}_T$ le
changement de base de $\mathcal{F}$ à $T$ ; autrement dit, $\mathcal{F}_T$
est l'image réciproque de $\mathcal{F}$ par la projection
$X_T = X \times_S T \to X$. Notons que $X_T \to T$ est de type fini
et que $\mathcal{F}_T$ est un $\mathcal{O}_{X_T}$-module
de type fini (Morphismes, lemme
\ref{morphisms-lemma-base-change-finite-type}, et
Modules, lemme \ref{modules-lemma-pullback-finite-type}).
Soit $n \geq 0$. La définition \ref{definition-flat-dimension-n} et le
lemme \ref{lemma-pre-flat-dimension-n} donnent un foncteur
\begin{equation}
\label{equation-flat-dimension-n}
F_n : (\Sch/S)^{opp} \longrightarrow \textit{Sets}, \quad
T \longrightarrow \left\{
\begin{matrix}
\{*\} & \text{if }\mathcal{F}_T\text{ is flat over }T\text{ in }\dim \geq n, \\
\emptyset & \text{else.}
\end{matrix}
\right.
\end{equation}
\end{situation}

\begin{lemma}
\label{lemma-flat-dimension-n}
Plaçons-nous dans la situation \ref{situation-flat-dimension-n}.
\begin{enumerate}
\item Le foncteur $F_n$ satisfait à la condition de faisceau pour la topologie fpqc.
\item Si $f$ est quasi-compact et localement de présentation finie
et si $\mathcal{F}$ est de présentation finie, alors le foncteur $F_n$ est
compatible aux limites.
\end{enumerate}
\end{lemma}

\begin{proof}
Soit $\{T_i \to T\}_{i \in I}$ un recouvrement fpqc de schémas sur $S$.
Posons $X_i = X_{T_i} = X \times_S T_i$ et notons $\mathcal{F}_i$
l'image réciproque de $\mathcal{F}$ sur $X_i$. Supposons $\mathcal{F}_i$
plat sur $T_i$ en dimensions $\geq n$ pour tout $i$. Soit $t \in T$.
Choisissons un indice $i$ et un point $t_i \in T_i$ d'image $t$.
Considérons le diagramme cartésien
$$
\xymatrix{
X_{\Spec(\mathcal{O}_{T, t})} \ar[d] &
X_{\Spec(\mathcal{O}_{T_i, t_i})} \ar[d] \ar[l] \\
\Spec(\mathcal{O}_{T, t}) &
\Spec(\mathcal{O}_{T_i, t_i}) \ar[l]
}
$$
Comme le morphisme horizontal inférieur est plat, il résulte de
Compléments sur les morphismes, lemme \ref{more-morphisms-lemma-flat-locus-base-change},
que l'ensemble $Z_i \subset X_{t_i}$ où $\mathcal{F}_i$ n'est pas plat
sur $T_i$ et l'ensemble $Z \subset X_t$ où $\mathcal{F}_T$ n'est pas plat
sur $T$ sont liés par la relation $Z_i = Z_{\kappa(t_i)}$. Nous en déduisons
que $\mathcal{F}_T$ est plat sur $T$ en dimensions $\geq n$, d'après
Morphismes, lemme \ref{morphisms-lemma-dimension-fibre-after-base-change}.

\medskip\noindent
Supposons que $f$ soit quasi-compact et localement de présentation finie et
que $\mathcal{F}$ soit de présentation finie.
Dans ce paragraphe, nous ramenons d'abord la démonstration de (2) au cas où
$f$ est de présentation finie.
Soit $T = \lim_{i \in I} T_i$ une limite dirigée de
$S$-schémas affines et supposons $\mathcal{F}_T$ plat en dimensions
$\geq n$. Posons $X_i = X_{T_i} = X \times_S T_i$ et notons $\mathcal{F}_i$
l'image réciproque de $\mathcal{F}$ sur $X_i$. Il faut montrer que
$\mathcal{F}_i$ est plat en dimensions $\geq n$ pour un certain $i$.
Choisissons $i_0 \in I$ et remplaçons $I$ par $\{i \mid i \geq i_0\}$.
Comme $T_{i_0}$ est affine (donc quasi-compact), il existe un nombre fini
d'ouverts affines $W_j \subset S$, $j = 1, \ldots, m$, et un recouvrement ouvert affine
$T_{i_0} = \bigcup_{j = 1, \ldots, m} V_{j, i_0}$
tels que $T_{i_0} \to S$ envoie $V_{j, i_0}$ dans $W_j$.
Pour $i \geq i_0$, notons $V_{j, i}$ l'image réciproque de $V_{j, i_0}$
dans $T_i$. Si nous pouvons montrer, pour chaque $j$, qu'il existe un $i$ tel que
$\mathcal{F}_{V_{j, i_0}}$ soit plat en dimensions $\geq n$,
nous aurons conclu. Nous nous ramenons ainsi au cas où $S$ est affine.
Dans ce cas, $X$ est quasi-compact et admet un recouvrement ouvert
affine fini $X = W_1 \cup \ldots \cup W_m$. Dans ce cas,
le résultat pour $(X \to S, \mathcal{F})$ équivaut au résultat pour
$(\coprod W_j, \coprod \mathcal{F}|_{W_j})$. Nous pouvons donc supposer
$f$ de présentation finie.

\medskip\noindent
Supposons $f$ de présentation finie et $\mathcal{F}$ de présentation
finie. Notons $U \subset X_T$ le sous-schéma ouvert des points où
$\mathcal{F}_T$ est plat sur $T$, voir
Compléments sur les morphismes, théorème \ref{more-morphisms-theorem-openness-flatness}.
Par hypothèse, toute fibre de $Z = X_T \setminus U$ sur
$T$ est de dimension $< n$. D'après
Limites, lemme \ref{limits-lemma-approximate-given-relative-dimension},
nous pouvons trouver un sous-schéma fermé $Z \subset Z' \subset X_T$
tel que $\dim(Z'_t) < n$ pour tout $t \in T$ et que
$Z' \to X_T$ soit de présentation finie. D'après
Limites, lemmes \ref{limits-lemma-descend-finite-presentation} et
\ref{limits-lemma-descend-closed-immersion-finite-presentation},
il existe un $i \in I$ et un sous-schéma fermé $Z'_i \subset X_i$
de présentation finie dont le changement de base à $T$ est $Z'$. D'après
Limites, lemme \ref{limits-lemma-limit-dimension},
nous pouvons supposer toutes les fibres de $Z'_i \to T_i$ de dimension $< n$. D'après
Limites, lemme \ref{limits-lemma-descend-module-flat-finite-presentation},
nous pouvons supposer $\mathcal{F}_i|_{X_i \setminus T'_i}$
plat sur $T_i$. Cela entraîne que $\mathcal{F}_i$ est
plat en dimensions $\geq n$ ; nous utilisons ici que $Z' \to X_T$ est de présentation
finie, donc que le complémentaire $X_T \setminus Z'$ est quasi-compact !
La partie (2) est ainsi démontrée et la démonstration du lemme est achevée.
\end{proof}

\begin{situation}
\label{situation-flat}
Soit $f : X \to S$ un morphisme de schémas.
Soit $\mathcal{F}$ un $\mathcal{O}_X$-module quasi-cohérent.
Pour tout schéma $T$ sur $S$, nous noterons $\mathcal{F}_T$ le
changement de base de $\mathcal{F}$ à $T$ ; autrement dit, $\mathcal{F}_T$
est l'image réciproque de $\mathcal{F}$ par la projection
$X_T = X \times_S T \to X$. Comme le changement de base d'un module plat
est plat, nous obtenons un foncteur
\begin{equation}
\label{equation-flat}
F_{flat} : (\Sch/S)^{opp} \longrightarrow \textit{Sets}, \quad
T \longrightarrow \left\{
\begin{matrix}
\{*\} & \text{if } \mathcal{F}_T \text{ is flat over }T, \\
\emptyset & \text{else.}
\end{matrix}
\right.
\end{equation}
\end{situation}

\begin{lemma}
\label{lemma-flat}
Plaçons-nous dans la situation \ref{situation-flat}.
\begin{enumerate}
\item Le foncteur $F_{flat}$ satisfait à la condition de faisceau pour la topologie fpqc.
\item Si $f$ est quasi-compact et localement de présentation finie
et si $\mathcal{F}$ est de présentation finie, alors le foncteur
$F_{flat}$ est compatible aux limites.
\end{enumerate}
\end{lemma}

\begin{proof}
La partie (1) résulte de l'énoncé suivant : si $T' \to T$ est un morphisme plat
surjectif de schémas sur $S$, alors $\mathcal{F}_{T'}$ est plat sur $T'$
si et seulement si $\mathcal{F}_T$ est plat sur $T$, voir
Compléments sur les morphismes, lemme \ref{more-morphisms-lemma-flat-locus-base-change}.
La partie (2) résulte de
Limites, lemme \ref{limits-lemma-descend-module-flat-finite-presentation},
après réduction au cas où $X$ et $S$ sont affines (comparer avec
la démonstration du
lemme \ref{lemma-flat-dimension-n}).
\end{proof}




\section{Stratifications d'aplatissement}
\label{section-flattening}

\noindent
Seulement les définitions. Le lecteur qui cherche une
``stratification de platitude générique'' consultera
Compléments sur les morphismes, Section
\ref{more-morphisms-section-generic-flatness-stratification}.

\begin{definition}
\label{definition-flattening}
Soit $X \to S$ un morphisme de schémas.
Soit $\mathcal{F}$ un $\mathcal{O}_X$-module quasi-cohérent.
Nous disons que {\it l'aplatissement universel de $\mathcal{F}$ existe}
si le foncteur $F_{flat}$ défini dans la situation \ref{situation-flat}
est représentable par un schéma $S'$ sur $S$.
Nous disons que {\it l'aplatissement universel de $X$ existe}
si l'aplatissement universel de $\mathcal{O}_X$ existe.
\end{definition}

\noindent
Notons que si l'aplatissement universel $S'$\footnote{Le schéma $S'$ est parfois
appelé {\it aplatisseur universel}. Dans \cite{GruRay}, il est appelé
le {\it platificateur universel}. L'existence de l'aplatissement universel
ne doit pas être confondue avec le type de résultats étudiés dans
Compléments d'algèbre, Section \ref{more-algebra-section-blowup-flat}.} de
$\mathcal{F}$ existe, alors le morphisme $S' \to S$ est un monomorphisme
surjectif de schémas tel que $\mathcal{F}_{S'}$ soit plat sur $S'$
et tel qu'un morphisme $T \to S$ se factorise par $S'$ si et seulement si
$\mathcal{F}_T$ est plat sur $T$.

\begin{example}
\label{example-no-universal-flattening}
Soit $X = S = \Spec(k[x, y])$, où $k$ est un corps. Soit
$\mathcal{F} = \widetilde{M}$, où $M = k[x, x^{-1}, y]/(y)$.
Pour une $k[x, y]$-algèbre $A$, posons $F_{flat}(A) = F_{flat}(\Spec(A))$.
Alors $F_{flat}(k[x, y]/(x, y)^n) = \{*\}$ pour tout $n$, tandis que
$F_{flat}(k[[x, y]]) = \emptyset$. Cela signifie que $F_{flat}$ n'est pas
représentable (même par un espace algébrique, voir
Espaces formels, Lemme
\ref{formal-spaces-lemma-map-into-algebraic-space}).
Ainsi, l'aplatissement universel n'existe pas dans ce cas.
\end{example}

\noindent
Nous définissons (comparer à
Topologie, Remarque \ref{topology-remark-locally-finite-stratification})
une {\it stratification} (localement finie, au sens schématique) d'un schéma $S$
comme la donnée de sous-schémas fermés $Z_i \subset S$ indexés par un
ensemble partiellement ordonné $I$ tels que
$S = \bigcup Z_i$ (ensemblistement), que tout point de $S$ possède
un voisinage ne rencontrant qu'un nombre fini de $Z_i$, et que
$$
Z_i \cap Z_j = \bigcup\nolimits_{k \leq i, j} Z_k.
$$
En posant $S_i = Z_i \setminus \bigcup_{j < i} Z_j$, la stratification
proprement dite est la décomposition $S = \coprod S_i$ en
sous-schémas localement fermés. Souvent, nous indiquons seulement les strates
$S_i$ et laissons au lecteur la construction des sous-schémas fermés $Z_i$.
Étant donnée une stratification, nous obtenons un monomorphisme
$$
S' = \coprod\nolimits_{i \in I} S_i \longrightarrow S.
$$
Nous l'appellerons le {\it monomorphisme associé à la stratification}.
Avec cette terminologie, nous pouvons définir ce que signifie posséder une
stratification d'aplatissement.

\begin{definition}
\label{definition-flattening-stratification}
Soit $X \to S$ un morphisme de schémas.
Soit $\mathcal{F}$ un $\mathcal{O}_X$-module quasi-cohérent.
Nous disons que $\mathcal{F}$ admet une {\it stratification d'aplatissement}
si le foncteur $F_{flat}$ défini dans la situation \ref{situation-flat}
est représentable par un monomorphisme $S' \to S$ associé
à une stratification de $S$ par des sous-schémas localement fermés.
Nous disons que $X$ admet une {\it stratification d'aplatissement}
si $\mathcal{O}_X$ admet une stratification d'aplatissement.
\end{definition}

\noindent
Lorsqu'une stratification d'aplatissement existe, il importe souvent
de comprendre l'ensemble d'indices qui étiquette les strates et son ordre partiel.
Cela est souvent lié aux rangs de modules. Par exemple, si
$X = S$ et si $\mathcal{F}$ est un $\mathcal{O}_S$-module de présentation finie,
alors la stratification d'aplatissement existe et est donnée par les idéaux de Fitting
de $\mathcal{F}$, voir
Diviseurs, Lemme \ref{divisors-lemma-finite-presentation-module}.




\section{Stratification d'aplatissement sur un anneau artinien}
\label{section-flattening-artinian}

\noindent
Une stratification d'aplatissement existe lorsque le schéma de base est le spectre
d'un anneau artinien.

\begin{lemma}
\label{lemma-flattening-stratification-artinian}
Soit $S$ le spectre d'un anneau artinien.
Pour tout schéma $X$ sur $S$ et tout $\mathcal{O}_X$-module quasi-cohérent,
il existe un aplatissement universel. En fait, l'aplatissement universel
est donné par une immersion fermée $S' \to S$ et constitue donc aussi une
stratification d'aplatissement pour $\mathcal{F}$.
\end{lemma}

\begin{proof}
Choisissons un recouvrement ouvert affine $X = \bigcup U_i$.
Alors $F_{flat}$ est le produit des foncteurs associés à
chacun des couples $(U_i, \mathcal{F}|_{U_i})$.
Il suffit donc de démontrer le résultat pour chaque
$(U_i, \mathcal{F}|_{U_i})$.
Dans le cas affine, le lemme résulte immédiatement de
Compléments d'algèbre,
Lemme \ref{more-algebra-lemma-flattening-artinian-universal-property}.
\end{proof}






\section{Aplatissement d'une application}
\label{section-flattening-map}

\noindent
Le théorème \ref{theorem-flattening-map}
est la clef des énoncés d'aplatissement qui suivent.

\begin{lemma}
\label{lemma-universally-separating}
Soit $S$ un schéma.
Soit $g : X' \to X$ un morphisme plat de schémas sur $S$,
où $X$ est localement de type fini sur $S$.
Soit $\mathcal{F}$ un $\mathcal{O}_X$-module quasi-cohérent de type fini
qui est plat sur $S$. Si $\text{Ass}_{X/S}(\mathcal{F}) \subset g(X')$,
alors l'application canonique
$$
\mathcal{F} \longrightarrow g_*g^*\mathcal{F}
$$
est injective et le reste après tout changement de base.
\end{lemma}

\begin{proof}
La dernière assertion signifie que $\mathcal{F}_T \to (g_T)_*g_T^*\mathcal{F}_T$
est injective pour tout morphisme $T \to S$. L'hypothèse
$\text{Ass}_{X/S}(\mathcal{F}) \subset g(X')$ est préservée par changement de base, voir
Diviseurs, Lemme \ref{divisors-lemma-base-change-relative-assassin} et
Remarque \ref{divisors-remark-base-change-relative-assassin}.
Il en va de même de l'hypothèse de platitude et de type fini.
Il suffit donc de démontrer l'injectivité de la flèche affichée.
Soit $\mathcal{K} = \Ker(\mathcal{F} \to g_*g^*\mathcal{F})$.
Notre but est de montrer que $\mathcal{K} = 0$.
Pour ce faire, il suffit de montrer que
$\text{WeakAss}_X(\mathcal{K}) = \emptyset$, voir
Diviseurs, Lemme \ref{divisors-lemma-weakly-ass-zero}.
Nous avons
$\text{WeakAss}_X(\mathcal{K}) \subset \text{WeakAss}_X(\mathcal{F})$, voir
Diviseurs, Lemme \ref{divisors-lemma-ses-weakly-ass}.
Puisque $\mathcal{F}$ est plat, il résulte du
lemme \ref{lemma-bourbaki-finite-type-general-base}
que $\text{WeakAss}_X(\mathcal{F}) \subset \text{Ass}_{X/S}(\mathcal{F})$.
Par hypothèse, tout point $x$ de $\text{Ass}_{X/S}(\mathcal{F})$
est l'image d'un certain $x' \in X'$. Comme $g$ est plat, l'homomorphisme
d'anneaux locaux $\mathcal{O}_{X, x} \to \mathcal{O}_{X', x'}$
est fidèlement plat ; par conséquent, l'application
$$
\mathcal{F}_x
\longrightarrow
g^*\mathcal{F}_{x'} =
\mathcal{F}_x \otimes_{\mathcal{O}_{X, x}} \mathcal{O}_{X', x'}
$$
est injective (voir
Algèbre, Lemme \ref{algebra-lemma-faithfully-flat-universally-injective}).
Cela implique que $\mathcal{K}_x = 0$, comme souhaité.
\end{proof}

\begin{lemma}
\label{lemma-flattening-module-map}
Soit $A$ un anneau. Soit $u : M \to N$ une application surjective de $A$-modules.
Si $M$ est projectif comme $A$-module, alors il existe un idéal
$I \subset A$ tel que, pour tout homomorphisme d'anneaux $\varphi : A \to B$,
les conditions suivantes soient équivalentes :
\begin{enumerate}
\item $u \otimes 1 : M \otimes_A B \to N \otimes_A B$ est un
isomorphisme ;
\item $\varphi(I) = 0$.
\end{enumerate}
\end{lemma}

\begin{proof}
Comme $M$ est projectif, nous pouvons trouver un $A$-module projectif $C$
tel que $F = M \oplus C$ soit un $A$-module libre.
En remplaçant $u$ par $u \oplus 1 : F = M \oplus C \to N \oplus C$,
nous voyons que nous pouvons supposer $M$ libre. Dans ce cas, soit $I$
l'idéal de $A$ engendré par les coefficients de tous les éléments de
$\Ker(u)$ relativement à une base (fixée) de $M$.
Cela fonctionne parce que, puisque $u$ est surjective et
$\otimes_A B$ exact à droite, $\Ker(u \otimes 1)$ est
l'image de $\Ker(u) \otimes_A B$ dans $M \otimes_A B$.
\end{proof}

\begin{theorem}
\label{theorem-flattening-map}
Dans la
situation \ref{situation-iso},
supposons que
\begin{enumerate}
\item $f$ soit de présentation finie ;
\item $\mathcal{F}$ soit de présentation finie, plat sur $S$ et
pur relativement à $S$ ;
\item $u$ soit surjective.
\end{enumerate}
Alors $F_{iso}$ est représentable par une immersion fermée $Z \to S$.
De plus, $Z \to S$ est de présentation finie si $\mathcal{G}$ est
de présentation finie.
\end{theorem}

\begin{proof}
Nous utiliserons sans autre mention le fait que $\mathcal{F}$ est universellement pur
sur $S$, voir
le lemme \ref{lemma-finite-type-flat-pure-along-fibre-is-universal}.
D'après
le lemme \ref{lemma-iso-sheaf}
et
Descente, Lemmes \ref{descent-lemma-closed-immersion} et
\ref{descent-lemma-descent-data-sheaves},
la question est locale pour la topologie étale sur $S$.
Il suffit donc de montrer, étant donné $s \in S$, qu'il existe
un voisinage étale de $(S, s)$ sur lequel le théorème soit vrai.

\medskip\noindent
En utilisant
le lemme \ref{lemma-finite-presentation-flat-along-fibre}
et après avoir remplacé $S$ par un voisinage étale élémentaire de $s$,
nous pouvons supposer qu'il existe un diagramme commutatif
$$
\xymatrix{
X \ar[dr] & & X' \ar[ll]^g \ar[ld] \\
& S &
}
$$
de schémas de présentation finie sur $S$,
où $g$ est étale, $X_s \subset g(X')$, les schémas $X'$ et $S$ sont affines,
et $\Gamma(X', g^*\mathcal{F})$ est un $\Gamma(S, \mathcal{O}_S)$-module projectif.
Notons que $g^*\mathcal{F}$ est universellement pur sur $S$, voir
le lemme \ref{lemma-affine-locally-projective-pure}.
Ainsi, d'après
le lemme \ref{lemma-criterion},
nous voyons que l'ouvert $g(X')$ contient les points de
$\text{Ass}_{X/S}(\mathcal{F})$ situés au-dessus de $\Spec(\mathcal{O}_{S, s})$.
Posons
$$
E = \{t \in S \mid \text{Ass}_{X_t}(\mathcal{F}_t) \subset g(X') \}.
$$
D'après
Compléments sur les morphismes,
Lemme \ref{more-morphisms-lemma-relative-assassin-constructible},
$E$ est une partie constructible de $S$. Nous avons vu que
$\Spec(\mathcal{O}_{S, s}) \subset E$. D'après
Morphismes, Lemme \ref{morphisms-lemma-constructible-containing-open},
nous voyons que $E$ contient un voisinage ouvert de $s$. Ainsi, après avoir
remplacé $S$ par un voisinage affine plus petit de $s$, nous pouvons supposer que
$\text{Ass}_{X/S}(\mathcal{F}) \subset g(X')$.

\medskip\noindent
Puisque nous avons supposé $u$ surjective, nous avons $F_{iso} = F_{inj}$. D'après
le lemme \ref{lemma-universally-separating},
il s'ensuit que $u : \mathcal{F} \to \mathcal{G}$ est injective si et seulement si
$g^*u : g^*\mathcal{F} \to g^*\mathcal{G}$ est injective, et cela reste vrai
après tout changement de base. Nous nous sommes donc ramenés au cas où,
outre les hypothèses du théorème, $X \to S$ est un morphisme de
schémas affines et $\Gamma(X, \mathcal{F})$ est un
$\Gamma(S, \mathcal{O}_S)$-module projectif. Ce cas résulte immédiatement du
lemme \ref{lemma-flattening-module-map}.

\medskip\noindent
Pour voir que $Z$ est de présentation finie si $\mathcal{G}$ est de présentation
finie, combinons la partie (4) du
lemme \ref{lemma-iso-sheaf}
avec
Limites, Remarque \ref{limits-remark-limit-preserving}.
\end{proof}

\begin{lemma}
\label{lemma-Weil-restriction-closed-subschemes}
Soit $f:X\to S$ un morphisme de schémas de présentation finie,
plat et pur. Soit $Y$ un sous-schéma fermé de $X$. Soit $F=f_*Y$ le
foncteur de restriction de Weil de $Y$ le long de $f$, défini par
$$
F : (\Sch/S)^{opp} \to \textit{Sets}, \quad
T \mapsto
\left\{
\begin{matrix}
\{*\} & \text{si} & Y_T\to X_T \text{ est un isomorphisme, }\\
\emptyset & \text{sinon.} &
\end{matrix}
\right.
$$
Alors $F$ est représentable par une immersion fermée $Z\to S$. De plus,
$Z\to S$ est de présentation finie si $Y\to S$ l'est.
\end{lemma}

\begin{proof}
Soit $\mathcal{I}$ le faisceau d'idéaux définissant $Y$ dans $X$ et soit
$u:\mathcal{O}_X\to\mathcal{O}_X/\mathcal{I}$ la surjection.
Alors, pour un $S$-schéma $T$, l'immersion fermée $Y_T\to X_T$
est un isomorphisme si et seulement si $u_T$ est un isomorphisme. Le résultat
découle donc du
théorème \ref{theorem-flattening-map}.
\end{proof}



\section{Aplatissement dans le cas local}
\label{section-flattening-local}

\noindent
Dans cette section, nous commençons à appliquer les résultats précédents afin d'obtenir une
ébauche de la stratification d'aplatissement.

\begin{theorem}
\label{theorem-flattening-local}
Dans la
situation \ref{situation-flat-at-point},
supposons $A$ hensélien, $B$ essentiellement de type fini sur $A$, et
$M$ un $B$-module fini. Il existe alors un idéal
$I \subset A$ tel que $A/I$ coreprésente le foncteur $F_{lf}$ sur la catégorie
$\mathcal{C}$. Autrement dit, étant donné un homomorphisme local d'anneaux locaux
$\varphi : A \to A'$, avec $B' = B \otimes_A A'$ et $M' = M \otimes_A A'$,
les conditions suivantes sont équivalentes :
\begin{enumerate}
\item $\forall \mathfrak q \in V(\mathfrak m_{A'}B' + \mathfrak m_B B')
\subset \Spec(B') :
M'_{\mathfrak q}\text{ est plat sur }A'$ ;
\item $\varphi(I) = 0$.
\end{enumerate}
Si $B$ est essentiellement de présentation finie sur $A$ et $M$
de présentation finie sur $B$, alors $I$ est un idéal de type fini.
\end{theorem}

\begin{proof}
Choisissons un homomorphisme d'anneaux de type fini $A \to C$, un $C$-module fini $N$
et un idéal premier $\mathfrak q$ de $C$ tels que $B = C_{\mathfrak q}$
et $M = N_{\mathfrak q}$. Dans la suite, lorsque nous disons
``le théorème vaut pour $(N/C/A, \mathfrak q)$'', nous entendons qu'il
vaut pour $(A \to B, M)$, où $B = C_{\mathfrak q}$ et
$M = N_{\mathfrak q}$. D'après le
lemme \ref{lemma-flat-at-point-go-up},
le foncteur $F_{lf}$ ne change pas si nous remplaçons $B$ par un anneau local
plat sur $B$. Ainsi, puisque $A$ est hensélien, nous pouvons appliquer le
lemme \ref{lemma-existence-algebra}
et supposer qu'il existe un dévissage complet de
$N/C/A$ en $\mathfrak q$.

\medskip\noindent
Soit $(A_i, B_i, M_i, \alpha_i, \mathfrak q_i)_{i = 1, \ldots, n}$
un tel dévissage complet de $N/C/A$ en $\mathfrak q$. Soit
$\mathfrak q'_i \subset A_i$ l'unique idéal premier situé au-dessus de
$\mathfrak q_i \subset B_i$, comme dans la
définition \ref{definition-complete-devissage-at-x-algebra}.
Puisque $C \to A_1$ est surjectif et $N \cong M_1$ comme $C$-modules,
nous voyons d'après le
lemme \ref{lemma-flat-at-point-finite}
qu'il suffit de montrer que le théorème vaut pour $(M_1/A_1/A, \mathfrak q'_1)$.
Puisque $B_1 \to A_1$ est fini et que $\mathfrak q_1$ est l'unique idéal premier
de $B_1$ au-dessus de $\mathfrak q'_1$, nous voyons que
$(A_1)_{\mathfrak q'_1} \to (B_1)_{\mathfrak q_1}$ est fini (voir
Algèbre, Lemme \ref{algebra-lemma-unique-prime-over-localize-below}, ou
Compléments sur les morphismes,
Lemme \ref{more-morphisms-lemma-finite-morphism-single-point-in-fibre}).
Ainsi, d'après le
lemme \ref{lemma-flat-at-point-finite},
il suffit de montrer que le théorème vaut pour $(M_1/B_1/A, \mathfrak q_1)$.

\medskip\noindent
À ce stade, nous pouvons supposer, par récurrence sur la longueur $n$ du
dévissage, que le théorème vaut pour $(M_2/B_2/A, \mathfrak q_2)$.
(Si $n = 1$, alors $M_2 = 0$, qui est plat sur $A$.)
En inversant les deux dernières étapes du paragraphe précédent et en utilisant
$M_2 \cong \Coker(\alpha_2)$ comme $B_1$-modules, nous voyons
que le théorème vaut pour $(\Coker(\alpha_1)/B_1/A, \mathfrak q_1)$.

\medskip\noindent
Soit $A'$ un objet de $\mathcal{C}$. À ce stade, nous utilisons le
lemme \ref{lemma-induction-step}
pour voir que si $(M_1 \otimes_A A')_{\mathfrak q'}$ est plat
sur $A'$ pour un idéal premier $\mathfrak q'$ de $B_1 \otimes_A A'$
situé au-dessus de $\mathfrak m_{A'}$, alors
$(\Coker(\alpha_1) \otimes_A A')_{\mathfrak q'}$ est plat sur $A'$.
Nous en concluons que $F_{lf}$ est un sous-foncteur du
foncteur $F'_{lf}$ associé au module
$\Coker(\alpha_1)_{\mathfrak q_1}$ sur $(B_1)_{\mathfrak q_1}$.
D'après le paragraphe précédent, $F'_{lf}$ est coreprésenté par
$A/J$ pour un certain idéal $J \subset A$. Nous pouvons donc remplacer $A$ par
$A/J$ et supposer que $\Coker(\alpha_1)_{\mathfrak q_1}$ est
plat sur $A$.

\medskip\noindent
Puisque $\Coker(\alpha_1)$ est un $B_1$-module pour lequel
il existe un dévissage complet de $N_1/B_1/A$ en $\mathfrak q_1$
et puisque $\Coker(\alpha_1)_{\mathfrak q_1}$ est
plat sur $A$, d'après le
lemme \ref{lemma-complete-devissage-flat-finite-type-module},
nous voyons que $\Coker(\alpha_1)$ est libre comme $A$-module, donc en particulier
plat comme $A$-module. Par conséquent, le
lemme \ref{lemma-induction-step}
implique que $F_{lf}(A')$ est non vide si et seulement si $\alpha \otimes 1_{A'}$
est injective. Soit $N_1 = \Im(\alpha_1) \subset M_1$, de sorte que
nous avons les suites exactes
$$
0 \to N_1 \to M_1 \to \Coker(\alpha_1) \to 0
\quad\text{et}\quad
B_1^{\oplus r_1} \to N_1 \to 0
$$
La platitude de $\Coker(\alpha_1)$ implique que la première suite
est universellement exacte (voir
Algèbre, Lemme \ref{algebra-lemma-flat-universally-injective}).
Ainsi, $\alpha \otimes 1_{A'}$ est injective si et seulement si
$B_1^{\oplus r_1} \otimes_A A' \to N_1 \otimes_A A'$
est un isomorphisme. Enfin, le
théorème \ref{theorem-flattening-map}
montre que ce foncteur est coreprésentable par $A/I$ pour un certain idéal $I$,
et nous concluons que $F_{lf}$ est lui aussi coreprésentable par $A/I$.

\medskip\noindent
Pour démontrer la dernière assertion, supposons que $A \to B$ soit essentiellement de
présentation finie et $M$ de présentation finie sur $B$. Soit $I \subset A$
l'idéal tel que $F_{lf}$ soit coreprésenté par $A/I$.
Écrivons $I = \bigcup I_\lambda$, où $I_\lambda$ parcourt les idéaux de type fini
contenus dans $I$. Comme $F_{lf}(A/I) = \{*\}$,
nous voyons alors que $F_{lf}(A/I_\lambda) = \{*\}$ pour un certain $\lambda$, voir
la partie (2) du lemme \ref{lemma-flat-at-point}.
Cela implique clairement que $I = I_\lambda$.
\end{proof}

\begin{remark}
\label{remark-flattening-local-scheme-theoretic}
Voici une reformulation schématique du
théorème \ref{theorem-flattening-local}.
Soit $(X, x) \to (S, s)$ un morphisme de schémas pointés
qui est localement de type fini. Soit $\mathcal{F}$ un
$\mathcal{O}_X$-module quasi-cohérent de type fini.
Supposons $S$ local hensélien, de point fermé $s$.
Il existe un sous-schéma fermé $Z \subset S$ ayant la propriété suivante :
pour tout morphisme de schémas pointés $(T, t) \to (S, s)$, les conditions
suivantes sont équivalentes :
\begin{enumerate}
\item $\mathcal{F}_T$ est plat sur $T$ en tous les points de la fibre
$X_t$ qui s'envoient sur $x \in X_s$ ;
\item $\Spec(\mathcal{O}_{T, t}) \to S$ se factorise par $Z$.
\end{enumerate}
De plus, si $X \to S$ est de présentation finie en $x$ et $\mathcal{F}_x$
de présentation finie sur $\mathcal{O}_{X, x}$, alors $Z \to S$
est de présentation finie.
\end{remark}

\noindent
À ce stade, nous pouvons déduire gratuitement du résultat précédent des énoncés
très généraux. Notons que le cas le plus intéressant est peut-être
celui où $E = X_s$ !

\begin{lemma}
\label{lemma-freebie}
Soit $S$ le spectre d'un anneau local hensélien, de point fermé $s$.
Soit $X \to S$ un morphisme de schémas localement de type fini.
Soit $\mathcal{F}$ un $\mathcal{O}_X$-module quasi-cohérent de type fini.
Soit $E \subset X_s$ une partie. Il existe un sous-schéma fermé
$Z \subset S$ ayant la propriété suivante : pour tout morphisme de schémas pointés
$(T, t) \to (S, s)$, les conditions suivantes sont équivalentes :
\begin{enumerate}
\item $\mathcal{F}_T$ est plat sur $T$ en tous les points de la fibre
$X_t$ qui s'envoient sur un point de $E \subset X_s$ ;
\item $\Spec(\mathcal{O}_{T, t}) \to S$ se factorise par $Z$.
\end{enumerate}
De plus, si $X \to S$ est localement de présentation finie,
si $\mathcal{F}$ est de présentation finie et si $E \subset X_s$ est
fermé et quasi-compact, alors $Z \to S$ est de présentation finie.
\end{lemma}

\begin{proof}
Pour $x \in X_s$, notons $Z_x \subset S$ le sous-schéma fermé obtenu dans la
remarque \ref{remark-flattening-local-scheme-theoretic}.
Il est alors clair que $Z = \bigcap_{x \in E} Z_x$ convient !

\medskip\noindent
Pour démontrer la dernière assertion, supposons $X$ localement de présentation finie,
$\mathcal{F}$ de présentation finie et $Z$ fermé et quasi-compact.
Choisissons d'abord un nombre fini d'ouverts affines $W_j \subset X$ tels que
$E \subset \bigcup W_j$. Il suffit clairement de démontrer le
résultat pour chaque morphisme $W_j \to S$, muni du faisceau $\mathcal{F}|_{X_j}$
et de la partie fermée $E \cap W_j$. Nous pouvons donc supposer $X$ affine.
Dans ce cas,
Compléments d'algèbre, Lemme \ref{more-algebra-lemma-limit-preserving-flat-at-primes},
montre que le foncteur défini par (1) est ``compatible aux limites''.
Nous pouvons donc montrer que $Z \to S$ est de présentation finie exactement
comme dans la dernière partie de la démonstration du
théorème \ref{theorem-flattening-local}.
\end{proof}

\begin{remark}
\label{remark-flattening-complete-noetherian}
En remontant la démonstration du
lemme \ref{lemma-freebie}
jusqu'à ses origines, nous découvrons un chemin long et sinueux. Mais si nous supposons que
\begin{enumerate}
\item $f$ est de type fini ;
\item $\mathcal{F}$ est un $\mathcal{O}_X$-module de type fini ;
\item $E = X_s$ ;
\item $S$ est le spectre d'un anneau local noethérien complet,
\end{enumerate}
alors il existe une démonstration reposant entièrement sur une algèbre plus élémentaire,
comme suit : nous nous ramenons d'abord au cas où $X$ est affine en choisissant
un recouvrement ouvert affine fini. Dans ce cas, $Z$ existe d'après
Compléments d'algèbre,
Lemme \ref{more-algebra-lemma-flattening-complete-local-universal-property}.
L'étape essentielle de cette démonstration consiste à construire le sous-schéma fermé $Z$
pas à pas à l'intérieur des troncations
$\Spec(\mathcal{O}_{S, s}/\mathfrak m_s^n)$.
Cela repose sur le fait que les stratifications d'aplatissement existent toujours
lorsque la base est artinienne, et sur le fait que
$\mathcal{O}_{S, s} = \lim \mathcal{O}_{S, s}/\mathfrak m_s^n$.
\end{remark}




\section{Variantes d'un lemme}
\label{section-variants-mod-injective}

\noindent
Dans cette section, nous étudions des variantes des
lemmes d'Algèbre \ref{algebra-lemma-mod-injective-general} et
\ref{algebra-lemma-mod-injective}.
La version la plus générale est la
proposition \ref{proposition-finite-type-injective-into-flat-mod-m} ;
elle est énoncée comme \cite[Lemme 4.2.2]{GruRay}, mais la démonstration de
loc.cit.\ ne donne que le résultat plus faible énoncé dans le
lemme \ref{lemma-upstairs-finite-type-injective-into-flat-mod-m}.
La démonstration délicate de la
proposition \ref{proposition-finite-type-injective-into-flat-mod-m}
est due à Ofer Gabber. Comme nous n'avons actuellement aucune application de
cette proposition, nous encourageons le lecteur à passer à la section suivante
après avoir lu la démonstration du
lemme \ref{lemma-upstairs-finite-type-injective-into-flat-mod-m} ;
ce lemme sera utilisé dans la section suivante pour démontrer le
théorème \ref{theorem-check-flatness-at-associated-points}.

\begin{situation}
\label{situation-mod-injective}
Soit $\varphi : A \to B$ un homomorphisme local d'anneaux locaux
qui est essentiellement de type fini. Soient $M$ un $A$-module plat,
$N$ un $B$-module fini et $u : N \to M$ une application de $A$-modules telle que
$\overline{u} : N/\mathfrak m_AN \to M/\mathfrak m_AM$ soit injective.
\end{situation}

\noindent
Dans cette situation, notre but est de montrer que $u$ est $A$-universellement injective,
que $N$ est de présentation finie sur $B$ et que $N$ est plat comme $A$-module.
Si cela est vrai, nous dirons que {\it le lemme vaut} dans la situation donnée.

\begin{lemma}
\label{lemma-Noetherian-finite-type-injective-into-flat-mod-m}
Si, dans la situation \ref{situation-mod-injective}, l'anneau $A$ est noethérien,
alors le lemme vaut.
\end{lemma}

\begin{proof}
En appliquant Algèbre, Lemme \ref{algebra-lemma-mod-injective}, nous voyons que
$u$ est injective et que $N/u(M)$ est plat sur $A$. Alors $u$ est
$A$-universellement injective
(Algèbre, Lemme \ref{algebra-lemma-flat-tor-zero}) et $N$ est $A$-plat
(Algèbre, Lemme \ref{algebra-lemma-flat-ses}). Puisque $B$ est noethérien
dans ce cas, nous voyons que $N$ est de présentation finie.
\end{proof}

\begin{lemma}
\label{lemma-reduce-finite-type-injective-into-flat-mod-m}
Soit $A_0$ un anneau local. Si le lemme vaut pour toute
situation \ref{situation-mod-injective} avec $A = A_0$, où $B$ est une
localisation d'une algèbre de polynômes sur $A$ et $N$ est de présentation finie
sur $B$, alors le lemme vaut pour toute
situation \ref{situation-mod-injective} avec $A = A_0$.
\end{lemma}

\begin{proof}
Soient $A \to B$ et $u : N \to M$ comme dans la situation \ref{situation-mod-injective}.
Écrivons $B = C/I$, où $C$ est la localisation d'une algèbre de polynômes
sur $A$ en un idéal premier. Si nous pouvons montrer que $N$ est de présentation finie comme
$C$-module, cela montre a fortiori que $N$ est de présentation finie
comme $B$-module (voir
Algèbre, Lemme \ref{algebra-lemma-finitely-presented-over-subring}).
Nous pouvons donc supposer que $B$ est la localisation d'une algèbre de polynômes.
Écrivons ensuite $N = B^{\oplus n}/K$ pour un sous-module $K \subset B^{\oplus n}$.
Comme $B/\mathfrak m_AB$ est noethérien (étant essentiellement de type fini
sur un corps), il existe un nombre fini d'éléments $k_1, \ldots, k_s \in K$
tels que, pour $K' = \sum Bk_i$ et $N' = B^{\oplus n}/K'$, la
surjection canonique $N' \to N$ induise un isomorphisme
$N'/\mathfrak m_AN' \cong N/\mathfrak m_AN$.
Maintenant, si le lemme vaut pour la composée $u' : N' \to M$,
alors $u'$ est injective, donc $N' = N$ et $u' = u$. Le lemme vaut donc pour
la situation initiale.
\end{proof}

\begin{lemma}
\label{lemma-henselian-finite-type-injective-into-flat-mod-m}
Si, dans la situation \ref{situation-mod-injective}, l'anneau $A$ est hensélien,
alors le lemme vaut.
\end{lemma}

\begin{proof}
Il suffit de le démontrer lorsque $B$ est essentiellement de présentation finie
sur $A$ et $N$ est de présentation finie sur $B$, voir le
lemme \ref{lemma-reduce-finite-type-injective-into-flat-mod-m}.
Faisons temporairement l'hypothèse supplémentaire que $N$ est plat sur $A$.
Alors $N$ est une colimite filtrante $N = \colim_i F_i$
de $A$-modules libres $F_i$ telle que les applications de transition
$u_{ii'} : F_i \to F_{i'}$ soient injectives modulo $\mathfrak m_A$, voir le
lemme \ref{lemma-flat-finite-type-local-colimit-free}.
Chacune des composées $u_i : F_i \to M$ est $A$-universellement
injective d'après le
lemme \ref{lemma-universally-injective-local} ;
par conséquent, $u = \colim u_i$ est $A$-universellement injective, comme souhaité.

\medskip\noindent
Supposons $A$ local hensélien, $B$ essentiellement
de présentation finie sur $A$ et $N$ de présentation finie sur $B$. D'après le
théorème \ref{theorem-flattening-local},
il existe un idéal de type fini $I \subset A$ tel que
$N/IN$ soit plat sur $A/I$ et que $N/I^2N$ ne soit pas plat sur
$A/I^2$, sauf si $I = 0$. Le résultat du paragraphe précédent montre que
le lemme vaut pour $u \bmod I : N/IN \to M/IM$ sur $A/I$.
Considérons le diagramme commutatif
$$
\xymatrix{
0 \ar[r] &
M \otimes_A I/I^2 \ar[r] &
M/I^2M \ar[r] &
M/IM \ar[r] & 0 \\
&
N \otimes_A I/I^2 \ar[r] \ar[u]^u &
N/I^2N \ar[r] \ar[u]^u &
N/IN \ar[r] \ar[u]^u & 0
}
$$
dont les lignes sont exactes par exactitude à droite de $\otimes$ et parce que
$M$ est plat sur $A$. Notons que la flèche verticale de gauche est l'application
$N/IN \otimes_{A/I} I/I^2 \to M/IM \otimes_{A/I} I/I^2$ et qu'elle est donc
injective. Une chasse au diagramme montre que la flèche inférieure gauche est injective,
c'est-à-dire $\text{Tor}^1_{A/I^2}(I/I^2, M/I^2) = 0$, voir
Algèbre, Remarque \ref{algebra-remark-Tor-ring-mod-ideal}.
Ainsi, $N/I^2N$ est plat sur $A/I^2$ d'après
Algèbre, Lemme \ref{algebra-lemma-what-does-it-mean},
ce qui est une contradiction, sauf si $I = 0$.
\end{proof}

\noindent
Le lemme suivant traite le cas particulier de la
situation \ref{situation-mod-injective} où $M$ est muni
d'une structure de $B$-module et $u$ est $B$-linéaire. C'est le cas le plus
souvent utilisé en pratique, et sa démonstration est nettement plus facile
que celle du cas général.

\begin{lemma}
\label{lemma-upstairs-finite-type-injective-into-flat-mod-m}
Soit $A \to B$ un homomorphisme local d'anneaux locaux qui est
essentiellement de type fini. Soit $u : N \to M$ une application de $B$-modules.
Si $N$ est un $B$-module fini, $M$ est plat sur $A$ et
$\overline{u} : N/\mathfrak m_A N \to M/\mathfrak m_A M$ est injective,
alors $u$ est $A$-universellement injective, $N$ est de présentation finie sur
$B$ et $N$ est plat sur $A$.
\end{lemma}

\begin{proof}
Soit $A \to A^h$ l'hensélisé de $A$. Soit $B'$ la localisation
de $B \otimes_A A^h$ en l'idéal maximal
$\mathfrak m_B \otimes A^h + B \otimes \mathfrak m_{A^h}$.
Comme $B \to B'$ est plat (donc fidèlement plat, voir
Algèbre, Lemme \ref{algebra-lemma-local-flat-ff}),
nous pouvons remplacer $A \to B$ par $A^h \to B'$,
le module $M$ par $M \otimes_B B'$, le module $N$ par $N \otimes_B B'$
et $u$ par $u \otimes \text{id}_{B'}$, voir
Algèbre, Lemmes \ref{algebra-lemma-descend-properties-modules} et
\ref{algebra-lemma-flatness-descends-more-general}.
Nous pouvons ainsi supposer que $A$ est un anneau local hensélien.
Dans ce cas, notre lemme résulte du lemme plus général
\ref{lemma-henselian-finite-type-injective-into-flat-mod-m}.
\end{proof}

\begin{lemma}
\label{lemma-valuation-ring-finite-type-injective-into-flat-mod-m}
Si, dans la situation \ref{situation-mod-injective}, l'anneau $A$ est un
anneau de valuation, alors le lemme vaut.
\end{lemma}

\begin{proof}
Rappelons qu'un $A$-module est plat si et seulement s'il est sans torsion, voir
Compléments d'algèbre, Lemme
\ref{more-algebra-lemma-valuation-ring-torsion-free-flat}.
Soit $T \subset N$ la $A$-torsion. Alors $u(T) = 0$ et
$N/T$ est $A$-plat. Ainsi, $N/T$ est de présentation finie sur $B$, voir
Compléments d'algèbre, Lemme
\ref{more-algebra-lemma-flat-finite-type-valuation-ring-finite-presentation}.
Par conséquent, $T$ est un $B$-module fini, voir
Algèbre, Lemme \ref{algebra-lemma-extension}.
Comme $N/T$ est $A$-plat, nous voyons que
$T/\mathfrak m_A T \subset N/\mathfrak m_A N$, voir
Algèbre, Lemme \ref{algebra-lemma-flat-tor-zero}.
Puisque $\overline{u}$ est injective mais $u(T) = 0$, nous concluons que
$T/\mathfrak m_A T = 0$. Ainsi, $T = 0$ par le lemme de Nakayama, voir
Algèbre, Lemme \ref{algebra-lemma-NAK}. À ce stade, nous avons
démontré deux des trois assertions ($N$ est $A$-plat et
de présentation finie sur $B$) ; il reste à montrer que
$u$ est universellement injective.

\medskip\noindent
D'après Algèbre, Théorème \ref{algebra-theorem-universally-exact-criteria},
il suffit de montrer que $N \otimes_A Q \to M \otimes_A Q$ est injective
pour tout $A$-module de présentation finie $Q$. D'après
Compléments d'algèbre, Lemme
\ref{more-algebra-lemma-generalized-valuation-ring-modules},
nous pouvons supposer $Q = A/(a)$, avec $a \in \mathfrak m_A$ non nul.
Il suffit donc de montrer que $N/aN \to M/aM$ est injective.
Soit $x \in N$ tel que $u(x) \in aM$. D'après le
lemme \ref{lemma-flat-finite-type-local-valuation-ring-has-content},
nous savons que $x$ possède un idéal de contenu $I \subset A$. Comme
$I$ est de type fini
(Compléments d'algèbre, Lemme \ref{more-algebra-lemma-content-finitely-generated})
et $A$ est un anneau de valuation, nous avons $I = (b)$ pour un certain $b$
(par Algèbre, Lemme \ref{algebra-lemma-characterize-valuation-ring}).
D'après Compléments d'algèbre, Lemme \ref{more-algebra-lemma-equal-content},
l'élément $u(x)$ a lui aussi pour idéal de contenu $I$.
Puisque $u(x) \in aM$, nous voyons que $(b) \subset (a)$
d'après Compléments d'algèbre, Définition \ref{more-algebra-definition-content-ideal}.
Comme $x \in bN$, nous concluons $x \in aN$, comme souhaité.
\end{proof}

\noindent
Considérons la situation suivante :
\begin{equation}
\label{equation-star}
\begin{matrix}
A \to B\text{ de présentation finie, }S \subset B
\text{ une partie multiplicative, et }\\
N\text{ un }S^{-1}B\text{-module de présentation finie}
\end{matrix}
\end{equation}
Dans cette situation, une {\it extension pure} est la donnée d'un ouvert affine
$U \subset \Spec(B)$ tel que $\Spec(S^{-1}B) \subset U$ et d'un
$\mathcal{O}(U)$-module de présentation finie $N'$ prolongeant $N$
tel que $N'$ soit $A$-projectif et que $N' \to N = S^{-1}N'$
soit $A$-universellement injective.

\medskip\noindent
Dans (\ref{equation-star}), si $A \to A_1$ est un homomorphisme d'anneaux, nous
pouvons effectuer un changement de base : prenons $B_1 = B \otimes_A A_1$, soit
$S_1 \subset B_1$ l'image de $S$, et soit $N_1 = N \otimes_A A_1$. Cela fonctionne
parce que $S_1^{-1}B_1 = S^{-1}B \otimes_A A_1$. Nous l'utiliserons
sans autre mention dans la suite.

\begin{lemma}
\label{lemma-properties-pure-spreadout}
Dans (\ref{equation-star}), s'il existe une extension pure, alors
\begin{enumerate}
\item les éléments de $N$ possèdent des idéaux de contenu dans $A$ ;
\item si $u : N \to M$ est un morphisme vers un $A$-module plat $M$
tel que $N/\mathfrak m N \to M/\mathfrak m M$ soit injectif
pour tout idéal maximal $\mathfrak m$ de $A$,\newline alors $u$ est
$A$-universellement injectif.
\end{enumerate}
\end{lemma}

\begin{proof}
Choisissons $U$, $N'$ comme dans la définition d'une extension pure.
Tout élément $x' \in N'$ possède un idéal de contenu dans $A$ parce que
$N'$ est $A$-projectif (cela se voit facilement directement, mais
résulte aussi de Compléments d'algèbre, Lemme
\ref{more-algebra-lemma-content-exists-flat-Mittag-Leffler}, et
Algèbre, Exemple \ref{algebra-example-ML}).
Puisque $N' \to N$ est $A$-universellement injective, nous voyons que
l'image $x \in N$ de tout $x' \in N'$ possède un idéal de contenu dans $A$
(c'est le même que l'idéal de contenu de $x'$).
Pour un élément général $x \in N$, nous choisissons $s \in S$ tel que
$s x$ appartienne à l'image de $N' \to N$ et utilisons le fait que $x$ et $sx$
ont le même idéal de contenu.

\medskip\noindent
Soit $u : N \to M$ comme en (2). Pour montrer que $u$ est $A$-universellement
injective, nous pouvons remplacer $A$ par une localisation en un idéal maximal
(nous omettons ce petit détail). Supposons $A$ local, d'idéal maximal $\mathfrak m$.
Choisissons $s \in S$ et considérons la composée
$$
N' \to N \xrightarrow{1/s} N \xrightarrow{u} M
$$
Chacune de ces applications est injective modulo $\mathfrak m$ ; la composée
est donc $A$-universellement injective d'après le
lemme \ref{lemma-universally-injective-local}.
Puisque $N = \colim_{s \in S} (1/s)N'$, nous concluons que $u$
est $A$-universellement injective comme colimite d'applications universellement injectives.
\end{proof}

\begin{lemma}
\label{lemma-find-pure-spreadout}
Dans (\ref{equation-star}), pour tout $\mathfrak p \in \Spec(A)$,
il existe un idéal de type fini $I \subset \mathfrak pA_\mathfrak p$
tel que nous ayons une extension pure sur $A_\mathfrak p/I$.
\end{lemma}

\begin{proof}
Nous pouvons remplacer $A$ par $A_\mathfrak p$. Ainsi, nous pouvons supposer $A$
local et $\mathfrak p$ égal à l'idéal maximal $\mathfrak m$ de $A$.
Nous pouvons écrire $N = S^{-1}N'$ pour un $B$-module de présentation finie $N'$,
en chassant les dénominateurs dans une présentation de $N$ sur $S^{-1}B$.
Puisque $B/\mathfrak m B$ est noethérien, le noyau $K$ de
$N'/\mathfrak m N' \to N/\mathfrak m N$ est de type fini.
Nous pouvons donc choisir $s \in S$ tel que $K$ soit annulé par $s$.
Après avoir remplacé $B$ par $B_s$, ce qui est permis puisque cela revient simplement à passer
à un sous-schéma ouvert affine de $\Spec(B)$, nous constatons que les éléments de $S$
agissent injectivement sur $N'/\mathfrak m N'$. Nous choisissons alors
un sous-anneau local $A_0 \subset A$ essentiellement de type fini sur $\mathbf{Z}$,
un homomorphisme d'anneaux de type fini $A_0 \to B_0$ tel que $B = A \otimes_{A_0} B_0$,
et un $B_0$-module fini $N'_0$ tel que
$N' = B \otimes_{B_0} N'_0 = A \otimes_{A_0} N'_0$.
Nous affirmons que $I = \mathfrak m_{A_0} A$ convient.
En effet, nous avons
$$
N'/IN' = N'_0/\mathfrak m_{A_0} N'_0 \otimes_{\kappa_{A_0}} A/I
$$
qui est libre sur $A/I$. La multiplication par les éléments de $S$
est injective après passage au quotient par l'idéal maximal ; par conséquent,
$N'/IN' \to N/IN$ est universellement injective, par exemple d'après le
lemme \ref{lemma-invert-universally-injective}.
\end{proof}

\begin{lemma}
\label{lemma-universally-injective-if-flat}
Dans (\ref{equation-star}), supposons que $N$ soit $A$-plat, que $M$ soit un $A$-module plat,
et que $u : N \to M$ soit un homomorphisme de $A$-modules tel que
$u \otimes \text{id}_{\kappa(\mathfrak p)}$ soit injectif pour tout
$\mathfrak p \in \Spec(A)$. Alors $u$ est $A$-universellement injectif.
\end{lemma}

\begin{proof}
D'après Algèbre, Lemme \ref{algebra-lemma-check-universally-injective-into-flat},
il suffit de vérifier que $N/IN \to M/IM$ est injectif pour tout
idéal $I \subset A$. Après avoir remplacé $A$ par $A/I$, nous voyons qu'il suffit
de démontrer que $u$ est injectif.

\medskip\noindent
Démonstration de l'injectivité de $u$. Soit $x \in N$ un élément non nul du
noyau de $u$. Il existe alors un idéal premier faiblement associé $\mathfrak p$
au module $Ax$ ; voir Algèbre, Lemme \ref{algebra-lemma-weakly-ass-zero}.
En remplaçant $A$ par $A_\mathfrak p$, nous pouvons supposer que $A$ est local et
nous trouvons un élément non nul $y \in Ax$ dont l'annulateur a pour radical
$\mathfrak m_A$ ; voir
Algèbre, Lemme \ref{algebra-lemma-weakly-ass-local}.
Ainsi, $\text{Supp}(y) \subset \Spec(S^{-1}B)$ est non vide et
contenu dans la fibre fermée de $\Spec(S^{-1}B) \to \Spec(A)$.
Soit $I \subset \mathfrak m_A$ un
idéal de type fini tel que nous ayons une extension pure sur $A/I$ ; voir le
lemme \ref{lemma-find-pure-spreadout}. Alors $I^n y = 0$ pour un certain $n$. Or,
$y \in \text{Ann}_M(I^n) = \text{Ann}_A(I^n) \otimes_R N$ par platitude.
Ainsi, pour obtenir la contradiction souhaitée, il suffit de montrer que
$$
\text{Ann}_A(I^n) \otimes_R N
\longrightarrow
\text{Ann}_A(I^n) \otimes_R M
$$
est injectif. Puisque $N$ et $M$ sont plats et que $\text{Ann}_A(I^n)$
est annulé par $I^n$, il suffit de montrer que
$Q \otimes_A N \to Q \otimes_A M$ est injectif pour tout $A$-module
$Q$ annulé par $I$. Cela découle de notre choix de $I$ et du
lemme \ref{lemma-properties-pure-spreadout}, partie (2).
\end{proof}

\begin{lemma}
\label{lemma-big-intersection-is-zero}
Soit $A$ un anneau intègre local qui n'est pas un corps.
Soit $S$ un ensemble d'idéaux de type fini de $A$.
Supposons que $S$ soit stable par produits et que
$\bigcup_{I \in S} V(I)$ soit le complémentaire du point générique de $\Spec(A)$.
Alors $\bigcap_{I \in S} I = (0)$.
\end{lemma}

\begin{proof}
Puisque $\mathfrak m_A \subset A$ n'est pas le point générique de $\Spec(A)$,
nous voyons que $I \subset \mathfrak m_A$ pour au moins un $I \in S$.
Par conséquent, $\bigcap_{I \in S} I \subset \mathfrak m_A$.
Soit $f \in \mathfrak m_A$ non nul. Alors
$V(f) \subset \bigcup_{I \in S} V(I)$.
Comme la topologie constructible sur $V(f)$ est quasi-compacte
(Topologie, Lemme \ref{topology-lemma-constructible-hausdorff-quasi-compact}
et
Algèbre, Lemme \ref{algebra-lemma-spec-spectral}),
nous trouvons que $V(f) \subset V(I_1) \cup \ldots \cup V(I_n)$
pour certains $I_j \in S$. Comme $I_1 \ldots I_n \in S$, nous voyons que
$V(f) \subset V(I)$ pour un certain $I$. Puisque $I$ est de type fini,
cela implique que $I^m \subset (f)$ pour un certain $m$ et, puisque
$S$ est stable par produits, nous voyons que $I \subset (f^2)$ pour
un certain $I \in S$. Il est alors impossible que $f \in I$.
\end{proof}

\begin{lemma}
\label{lemma-closed-points-complement}
Soit $A$ un anneau local. Soient $I, J \subset A$ des idéaux.
Si $J$ est de type fini et $I \subset J^n$ pour tout $n \geq 1$,
alors $V(I)$ contient les points fermés de $\Spec(A) \setminus V(J)$.
\end{lemma}

\begin{proof}
Soit $\mathfrak p \subset A$ un point fermé de $\Spec(A) \setminus V(J)$.
Nous voulons montrer que $I \subset \mathfrak p$. Sinon, un certain $f \in I$
s'envoie sur un élément non nul de $A/\mathfrak p$. Remarquons que
$V(J) \cap \Spec(A/\mathfrak p)$ est l'ensemble des points non génériques.
Le lemme \ref{lemma-big-intersection-is-zero}, appliqué
à la famille d'idéaux $J^nA/\mathfrak p$, entraîne donc que
l'image de $f$ est nulle dans $A/\mathfrak p$.
\end{proof}

\begin{lemma}
\label{lemma-make-smaller-flatness-ideal}
Soit $A$ un anneau local. Soit $I \subset A$ un idéal.
Soit $U \subset \Spec(A)$ un ouvert quasi-compact.
Soit $M$ un $A$-module. Supposons que
\begin{enumerate}
\item $M/IM$ soit plat sur $A/I$,
\item $M$ soit plat sur $U$.
\end{enumerate}
Alors $M/I_2M$ est plat sur $A/I_2$, où
$I_2 = \Ker(I \to \Gamma(U, I/I^2))$.
\end{lemma}

\begin{proof}
Il suffit de montrer que
$M \otimes_A I/I_2 \to IM/I_2M$ est injectif ; voir
Algèbre, Lemme \ref{algebra-lemma-what-does-it-mean-again}.
C'est vrai sur $U$ par l'hypothèse (2). Il suffit donc de montrer
que $M \otimes_A I/I_2$ s'injecte dans ses sections sur $U$.
Nous avons $M \otimes_A I/I_2 = M/IM \otimes_A I/I_2$ et
$M/IM$ est une colimite filtrante de $A/I$-modules libres finis
(Algèbre, Théorème \ref{algebra-theorem-lazard}).
Il suffit donc de montrer que $I/I_2$ s'injecte dans ses sections
sur $U$, ce qui résulte de la construction de $I_2$.
\end{proof}

\begin{proposition}
\label{proposition-finite-type-injective-into-flat-mod-m}
Soit $A \to B$ un homomorphisme local d'anneaux locaux
qui est essentiellement de type fini. Soient $M$ un $A$-module plat,
$N$ un $B$-module fini et $u : N \to M$ un homomorphisme de $A$-modules tel que
$\overline{u} : N/\mathfrak m_AN \to M/\mathfrak m_AM$ soit injectif.
Alors $u$ est $A$-universellement injectif, $N$ est de présentation finie sur
$B$, et $N$ est plat sur $A$.
\end{proposition}

\begin{proof}
Nous pouvons supposer que $B$ est la localisation d'une
$A$-algèbre de présentation finie $B_0$ et que $N$ est la localisation d'un
$B_0$-module de présentation finie $M_0$ ; voir le
lemme \ref{lemma-reduce-finite-type-injective-into-flat-mod-m}.
D'après Compléments sur les morphismes, Lemme
\ref{more-morphisms-lemma-generic-flatness-stratification},
il existe une « stratification de platitude générique »
pour $\widetilde{M_0}$ sur $\Spec(B_0)$ au-dessus de $\Spec(A)$.
En revenant à $N$, nous obtenons une suite de sous-schémas fermés
$$
S = \Spec(A) \supset S_0 \supset S_1 \supset \ldots \supset S_t = \emptyset
$$
où $S_i \subset S$ est défini par un idéal de type fini de $A$
et telle que l'image inverse de $\widetilde{N}$ sur
$\Spec(B) \times_S (S_i \setminus S_{i + 1})$ soit plate sur
$S_i \setminus S_{i + 1}$. Nous allons démontrer la proposition par
récurrence sur $t$ (le cas initial $t = 1$ sera démontré en même temps
que les autres étapes). Soit $\Spec(A/J_i)$ l'adhérence schématique
de $S_i \setminus S_{i + 1}$.

\medskip\noindent
{\bf Assertion 1.} $N/J_iN$ est plat sur $A/J_i$. C'est immédiat pour
$i = t - 1$ et cela résulte de l'hypothèse de récurrence pour $i > 0$.
Nous pouvons donc supposer $t > 1$, $S_{t - 1} \not = \emptyset$, et
$J_0 = 0$, et nous devons démontrer que $N$ est plat. Soit $J \subset A$
l'idéal définissant $S_1$. En raisonnant de nouveau par récurrence sur $t$, nous avons aussi
la platitude modulo les puissances de $J$. Soit $A^h$ l'hensélisé de $A$
et soit $B'$ la localisation de $B \otimes_A A^h$ en l'idéal maximal
$\mathfrak m_B \otimes A^h + B \otimes \mathfrak m_{A^h}$. Alors $B \to B'$
est fidèlement plat. Posons $N' = N \otimes_B B'$. Remarquons que $N'$
est $A^h$-plat si et seulement si $N$ est $A$-plat. D'après le
théorème \ref{theorem-flattening-local}, il existe un plus petit idéal
$I \subset A^h$ tel que $N'/IN'$ soit plat sur $A^h/I$, et
$I$ est de type fini. D'après ce qui précède, $I \subset J^nA^h$ pour
tout $n \geq 1$. Soit $S_i^h \subset \Spec(A^h)$ l'image inverse
de $S_i \subset \Spec(A)$. D'après le lemme \ref{lemma-closed-points-complement},
nous voyons que $V(I)$ contient les points fermés de $U = \Spec(A^h) - S_1^h$.
Par construction, $N'$ est $A^h$-plat sur $U$.
D'après le lemme \ref{lemma-make-smaller-flatness-ideal}, nous voyons que $N'/I_2N'$
est plat sur $A/I_2$, où
$I_2 = \Ker(I \to \Gamma(U, I/I^2))$. Ainsi $I = I_2$ par minimalité
de $I$. Cela implique que $I = I^2$ localement sur $U$, c'est-à-dire
que nous avons $I\mathcal{O}_{U, u} = (0)$ ou $I\mathcal{O}_{U, u} = (1)$
pour tout $u \in U$. Comme $V(I)$ contient les points fermés de $U$,
nous voyons que $I = 0$ sur $U$. Puisque $U \subset \Spec(A^h)$ est schématiquement
dense (parce que nous avons remplacé $A$ par $A/J_0$ au début
de ce paragraphe), nous voyons que $I = 0$. Ainsi $N'$ est $A^h$-plat
et, par conséquent, l'assertion 1 est vraie.

\medskip\noindent
Revenons à la situation décrite avant l'assertion 1. Avec
$A^h$ l'hensélisé de $A$, avec $B'$ la localisation
de $B \otimes_A A^h$ en l'idéal maximal
$\mathfrak m_B \otimes A^h + B \otimes \mathfrak m_{A^h}$, et avec
$N' = N \otimes_B B'$, nous voyons maintenant que l'idéal d'aplatissement $I \subset A^h$
du théorème \ref{theorem-flattening-local} est nilpotent.
Si $nil(A^h)$ désigne l'idéal des éléments nilpotents, alors
$nil(A^h) = nil(A) A^h$
(Compléments d'algèbre, Lemme \ref{more-algebra-lemma-henselization-nil}).
Il existe donc un idéal de type fini
nilpotent $I_0 \subset A$ tel que $N/I_0N$ soit plat sur $A/I_0$.

\medskip\noindent
{\bf Assertion 2.} Pour tout idéal premier $\mathfrak p \subset A$,
l'application
$\kappa(\mathfrak p) \otimes_A N \to \kappa(\mathfrak p) \otimes_A M$
est injective. Nous disons que $\mathfrak p$ est mauvais si cela est faux.
Supposons que $C$ soit une chaîne non vide de premiers mauvais et posons
$\mathfrak p^* = \bigcup_{\mathfrak p \in C} \mathfrak p$.
D'après le lemme \ref{lemma-find-pure-spreadout},
il existe un idéal de type fini
$\mathfrak a \subset \mathfrak p^*A_{\mathfrak p^*}$
tel qu'il existe une extension pure sur $V(\mathfrak a)$.
Si $\mathfrak p^*$ était bon, il résulterait du
lemme \ref{lemma-properties-pure-spreadout}
que les points de $V(\mathfrak a)$ sont bons.
Cependant, comme $\mathfrak a$ est de type fini et que
$\mathfrak p^*A_{\mathfrak p^*} = \bigcup_{\mathfrak p \in C}A_{\mathfrak p^*}$,
nous voyons que $V(\mathfrak a)$ contient un $\mathfrak p \in C$, contradiction.
Ainsi $\mathfrak p^*$ est mauvais. D'après le lemme de Zorn, s'il existe un
premier mauvais, il en existe un maximal, disons $\mathfrak p$.
Autrement dit, nous pouvons supposer que tout $\mathfrak p' \supset \mathfrak p$,
$\mathfrak p' \not = \mathfrak p$, est bon.
Dans ce cas, nous voyons que, pour tout $f \in A$, $f \not \in \mathfrak p$,
l'application $u \otimes \text{id}_{A/(\mathfrak p + f)}$ est universellement
injective ; voir le lemme \ref{lemma-universally-injective-if-flat}.
Il suffit donc de montrer que $N/\mathfrak p N$ est séparé
pour la topologie définie par les sous-modules $f(N/\mathfrak pN)$.
Puisque $B \to B'$ est fidèlement plat, il suffit de démontrer le
même énoncé pour le module $N'/\mathfrak p N'$.
D'après le lemme \ref{lemma-flat-finite-type-local-colimit-free} et
Compléments d'algèbre, Lemme
\ref{more-algebra-lemma-content-exists-flat-Mittag-Leffler},
les éléments de $N'/\mathfrak pN'$ possèdent des idéaux de contenu dans $A^h/\mathfrak pA^h$.
Il suffit donc de montrer que
$\bigcap_{f \in A, f \not \in \mathfrak p} f(A^h/\mathfrak p A^h) = 0$.
Il suffit ensuite de montrer la même chose pour
$A^h/\mathfrak q A^h$ pour tout idéal premier $\mathfrak q \subset A^h$
minimal au-dessus de $\mathfrak p A^h$. Comme $A \to A^h$ est l'hensélisation,
tout $\mathfrak q$ se contracte en $\mathfrak p$ et tout
$\mathfrak q' \supset \mathfrak q$, $\mathfrak q' \not = \mathfrak q$,
se contracte en un idéal premier $\mathfrak p'$ contenant strictement $\mathfrak p$.
L'annulation des intersections découle donc du
lemme \ref{lemma-big-intersection-is-zero}.

\medskip\noindent
Nous pouvons maintenant tout rassembler. En effet, en utilisant
les assertions 1 et 2, nous voyons que $N/I_0 N \to M/I_0M$ est
$A/I_0$-universellement injectif d'après le
lemme \ref{lemma-universally-injective-if-flat}.
Alors les diagrammes
$$
\xymatrix{
N \otimes_A (I_0^n/I_0^{n + 1}) \ar[r] \ar[d] &
M \otimes_A (I_0^n/I_0^{n + 1}) \ar@{=}[d] \\
I_0^n N /I_0^{n + 1} N \ar[r] &
I_0^n M /I_0^{n + 1} M
}
$$
montrent que les flèches verticales de gauche sont injectives. Ainsi, d'après
Algèbre, Lemme \ref{algebra-lemma-what-does-it-mean-again},
nous voyons que $N$ est plat. De manière analogue, l'injectivité
universelle de $u$ peut être ramenée (même sans
démontrer d'abord la platitude de $N$) à celle modulo $I_0$. Ceci achève
la démonstration.
\end{proof}


\section{Modules plats de type fini, IIIe partie}
\label{section-finite-type-flat-III}

\noindent
Le résultat suivant est l'un des résultats principaux de ce chapitre.

\begin{theorem}
\label{theorem-check-flatness-at-associated-points}
Soit $f : X \to S$ localement de type fini.
Soit $\mathcal{F}$ un $\mathcal{O}_X$-module quasi-cohérent de type fini.
Soit $x \in X$ d'image $s \in S$.
Les conditions suivantes sont équivalentes :
\begin{enumerate}
\item $\mathcal{F}$ est plat en $x$ sur $S$ ;
\item pour tout $x' \in \text{Ass}_{X_s}(\mathcal{F}_s)$ qui
se spécialise en $x$, le module $\mathcal{F}$ est plat en $x'$ sur $S$.
\end{enumerate}
\end{theorem}

\begin{proof}
Il est clair que (1) implique (2), puisque $\mathcal{F}_{x'}$ est une localisation
de $\mathcal{F}_x$ pour tout point qui se spécialise en $x$.
Posons $A = \mathcal{O}_{S, s}$, $B = \mathcal{O}_{X, x}$ et
$N = \mathcal{F}_x$. Soit $\Sigma \subset B$ la partie multiplicative
de $B$ formée des éléments qui agissent comme non-diviseurs de zéro sur $N/\mathfrak m_AN$.
L'hypothèse (2) implique que $\Sigma^{-1}N$ est $A$-plat d'après la description
de $\Spec(\Sigma^{-1}N)$ donnée dans le
lemme \ref{lemma-homothety-spectrum}.
D'autre part, l'application $N \to \Sigma^{-1}N$ est injective modulo
$\mathfrak m_A$ par construction. En appliquant donc le
lemme \ref{lemma-upstairs-finite-type-injective-into-flat-mod-m},
nous concluons.
\end{proof}

\noindent
Nous appliquons maintenant directement ce résultat afin d'obtenir les résultats utiles suivants.

\begin{lemma}
\label{lemma-check-along-closed-fibre}
Soit $S$ un schéma local de point fermé $s$.
Soit $f : X \to S$ localement de type fini.
Soit $\mathcal{F}$ un $\mathcal{O}_X$-module quasi-cohérent de type fini.
Supposons que
\begin{enumerate}
\item tout point de $\text{Ass}_{X/S}(\mathcal{F})$ se spécialise
en un point de la fibre fermée $X_s$\footnote{Cela vaut par exemple si
$f$ est de type fini et $\mathcal{F}$ est pur le long de $X_s$, ou
si $f$ est propre.},
\item $\mathcal{F}$ est plat sur $S$ en tout point de $X_s$.
\end{enumerate}
Alors $\mathcal{F}$ est plat sur $S$.
\end{lemma}

\begin{proof}
Cela résulte immédiatement du fait qu'il suffit de vérifier la
platitude aux points de l'assassin relatif de $\mathcal{F}$
sur $S$, d'après le
théorème \ref{theorem-check-flatness-at-associated-points}.
\end{proof}











\section{Aplatissement universel}
\label{section-flattening-final}

\noindent
Si $f : X \to S$ est un morphisme propre de présentation finie
de schémas, alors on peut trouver un aplatissement universel de $f$.
Dans cette section, nous étudions ce résultat ainsi que certaines de ses variantes.

\begin{lemma}
\label{lemma-free-at-generic-points-representable}
Dans la situation \ref{situation-free-at-generic-points}.
Pour chaque $p \geq 0$, le foncteur $H_p$
(\ref{equation-free-at-generic-points}) est représentable
par une immersion localement fermée $S_p \to S$. Si $\mathcal{F}$
est de présentation finie, alors $S_p \to S$ est de présentation finie.
\end{lemma}

\begin{proof}
Pour chaque $S$, nous allons démontrer simultanément l'énoncé pour tout $p \geq 0$.
Le foncteur $H_p$ est un faisceau pour la topologie fppf d'après le
lemme \ref{lemma-free-at-generic-points}.
En combinant donc
Descente, Lemme \ref{descent-lemma-descent-data-sheaves},
Compléments sur les morphismes, Lemme
\ref{more-morphisms-lemma-separated-locally-quasi-finite-morphisms-fppf-descend}
, et
Descente, Lemme \ref{descent-lemma-descending-fppf-property-immersion},
nous voyons que la question est locale pour la topologie \'etale sur $S$.
En particulier, la question est locale pour la topologie de Zariski sur $S$.

\medskip\noindent
Pour $s \in S$, notons $\xi_s$ l'unique point générique de la fibre $X_s$.
Remarquons que, pour tout $s \in S$, la restriction $\mathcal{F}_s$ de
$\mathcal{F}$ est localement libre d'un certain rang $p(s) \geq 0$ dans un
voisinage de $\xi_s$. (Comme $X_s$ est irréductible
et lisse, cela résulte de la platitude générique de $\mathcal{F}_s$ sur
$X_s$ ; voir
Algèbre, Lemme \ref{algebra-lemma-generic-flatness-Noetherian},
bien que ce soit excessif.) Pour référence ultérieure, remarquons que
$$
p(s) =
\dim_{\kappa(\xi_s)}(
\mathcal{F}_{\xi_s} \otimes_{\mathcal{O}_{X, \xi_s}} \kappa(\xi_s)
).
$$
En particulier, $H_{p(s)}(s)$ est non vide et $H_q(s)$ est vide
si $q \not = p(s)$.

\medskip\noindent
Soit $U \subset X$ un sous-schéma ouvert.
Comme $f : X \to S$ est lisse, il est ouvert.
Il résulte immédiatement de (\ref{equation-free-at-generic-points})
que le foncteur $H_p$ pour la paire $(f|_U : U \to f(U), \mathcal{F}|_U)$
et le foncteur $H_p$ pour la paire
$(f|_{f^{-1}(f(U))}, \mathcal{F}|_{f^{-1}(f(U))})$
sont les mêmes. Ainsi, pour démontrer l'existence de $S_p$ sur $f(U)$, nous pouvons
toujours remplacer $X$ par $U$.

\medskip\noindent
Choisissons $s \in S$. Il existe un voisinage ouvert affine $U$
de $\xi_s$ tel que $\mathcal{F}|_U$ puisse être engendré par au plus
$p(s)$ éléments. D'après les arguments précédents, pour démontrer
l'énoncé concernant $H_{p(s)}$ dans un voisinage de $s$, nous pouvons supposer
que $\mathcal{F}$ est engendré par $p(s)$ éléments, c'est-à-dire qu'il existe
une surjection
$$
u : \mathcal{O}_X^{\oplus p(s)} \longrightarrow \mathcal{F}
$$
Dans ce cas, il est clair que $H_{p(s)}$ est égal à $F_{iso}$
(\ref{equation-iso}) pour l'application $u$ (cela résulte immédiatement du
lemme \ref{lemma-injectivity-map-source-flat-pure},
mais aussi du
lemme \ref{lemma-induction-step-fp},
après avoir encore un peu rétréci de sorte que $S$ et $X$ soient tous deux affines.)
Nous pouvons donc appliquer le
théorème \ref{theorem-flattening-map}
pour voir que $H_{p(s)}$ est représentable par une immersion fermée dans un
voisinage de $s$.

\medskip\noindent
Le résultat découle formellement de ce qui précède.
En effet, les arguments ci-dessus montrent que, localement sur $S$, la fonction
$s \mapsto p(s)$ est bornée. Nous pouvons donc raisonner par récurrence
sur $p = \max_{s \in S} p(s)$. Le foncteur $H_p$ est représentable
par une immersion fermée $S_p \to S$ d'après ce qui précède. Remplaçons $S$ par
$S \setminus S_p$, ce qui diminue le maximum d'au moins un, et
nous concluons par l'hypothèse de récurrence.

\medskip\noindent
Supposons que $\mathcal{F}$ soit de présentation finie.
Alors $S_p \to S$ est localement de présentation finie d'après le
lemme \ref{lemma-free-at-generic-points}, partie (2), combiné à
Limites, Remarque \ref{limits-remark-limit-preserving}.
Nous reprenons alors l'argument de récurrence du paragraphe précédent pour
voir que chaque $S_p$ est quasi-compact lorsque $S$ est affine :
d'abord, si $p = \max_{s \in S} p(s)$, alors $S_p \subset S$
est fermé (voir ci-dessus), donc quasi-compact. Ensuite, $U = S \setminus S_p$
est un ouvert quasi-compact de $S$, car $S_p \to S$ est une immersion
fermée de présentation finie (voir par exemple la discussion dans Morphismes, Section
\ref{morphisms-section-constructible}). Puis $S_{p - 1} \to U$
est une immersion fermée de présentation finie, de sorte que $S_{p - 1}$
est quasi-compact et $U' = S \setminus (S_p \cup S_{p - 1})$
est quasi-compact. Et ainsi de suite.
\end{proof}

\begin{lemma}
\label{lemma-localize-flat-dimension-n}
Dans la situation \ref{situation-flat-dimension-n}.
Soit $h : X' \to X$ un morphisme \'etale.
Posons $\mathcal{F}' = h^*\mathcal{F}$ et $f' = f \circ h$.
Soit $F_n'$ le foncteur (\ref{equation-flat-dimension-n})
associé à $(f' : X' \to S, \mathcal{F}')$.
Alors $F_n$ est un sous-foncteur de $F_n'$ et, si
$h(X') \supset \text{Ass}_{X/S}(\mathcal{F})$, alors $F_n = F'_n$.
\end{lemma}

\begin{proof}
Soit $T \to S$ un morphisme quelconque. Alors $h_T : X'_T \to X_T$ est \'etale comme
changement de base du morphisme \'etale $g$. Pour $t \in T$, notons
$Z \subset X_t$ l'ensemble des points où $\mathcal{F}_T$ n'est pas
plat sur $T$, et notons de même $Z' \subset X'_t$ l'ensemble des
points où $\mathcal{F}'_T$ n'est pas plat sur $T$. Comme
$\mathcal{F}'_T = h_T^*\mathcal{F}_T$, nous voyons que
$Z' = h_t^{-1}(Z)$ ; voir
Morphismes, Lemme \ref{morphisms-lemma-flat-permanence}.
Ainsi, $Z' \to Z$ est un morphisme \'etale, donc $\dim(Z') \leq \dim(Z)$
(par exemple d'après
Descente, Lemme \ref{descent-lemma-dimension-at-point-local},
ou simplement parce qu'un morphisme \'etale est lisse de dimension relative $0$).
Cela implique que $F_n \subset F_n'$.

\medskip\noindent
Enfin, supposons que $h(X') \supset \text{Ass}_{X/S}(\mathcal{F})$
et que $T \to S$ soit un morphisme tel que $F_n'(T)$ soit non vide, c'est-à-dire
tel que $\mathcal{F}'_T$ soit plat en dimensions $\geq n$ sur $T$.
Choisissons un point $t \in T$ et soient $Z \subset X_t$ et $Z' \subset X'_t$
comme ci-dessus. Pour obtenir une contradiction, supposons que $\dim(Z) \geq n$.
Choisissons un point générique $\xi \in Z$ correspondant à une composante
de dimension $\geq n$. Soit $x \in \text{Ass}_{X_t}(\mathcal{F}_t)$
une généralisation de $\xi$. Alors $x$ s'envoie sur un point de
$\text{Ass}_{X/S}(\mathcal{F})$ d'après
Diviseurs, Lemme \ref{divisors-lemma-base-change-relative-assassin} et
Remarque \ref{divisors-remark-base-change-relative-assassin}.
Nous voyons donc que $x$ appartient à l'image de $h_T$, disons
$x = h_T(x')$ pour un certain $x' \in X'_T$. Mais $x' \not \in Z'$
puisque $x \leadsto \xi$ et $\dim(Z') < n$. Ainsi $\mathcal{F}'_T$
est plat sur $T$ en $x'$, ce qui implique que $\mathcal{F}_T$ est plat
en $x$ sur $T$ (d'après
Morphismes, Lemme \ref{morphisms-lemma-flat-permanence}).
Comme cela vaut pour tout $x$ de cette forme, nous concluons
que $\mathcal{F}_T$ est plat sur $T$ en $\xi$ d'après le
théorème \ref{theorem-check-flatness-at-associated-points},
ce qui est la contradiction souhaitée.
\end{proof}

\begin{lemma}
\label{lemma-compare-H-F}
Supposons que $X \to S$ soit un morphisme lisse de schémas affines
à fibres géométriquement irréductibles de dimension $d$ et que
$\mathcal{F}$ soit un $\mathcal{O}_X$-module quasi-cohérent de
présentation finie. Alors $F_d = \coprod_{p = 0, \ldots, c} H_p$
pour un certain $c \geq 0$, où $F_d$ est comme dans
(\ref{equation-flat-dimension-n}) et $H_p$ comme dans
(\ref{equation-free-at-generic-points}).
\end{lemma}

\begin{proof}
Comme $X$ est affine et $\mathcal{F}$ est quasi-cohérent de présentation finie,
nous savons que $\mathcal{F}$ peut être engendré par $c \geq 0$ éléments.
Alors $\dim_{\kappa(x)}(\mathcal{F}_x \otimes \kappa(x))$
en tout point $x \in X$ ne dépasse jamais $c$. En particulier, $H_p = \emptyset$
pour $p > c$. De plus, remarquons qu'il existe certainement une inclusion
$\coprod H_p \to F_d$. Cela étant dit, le contenu
du lemme est que, si un changement de base $\mathcal{F}_T$ est plat en
dimensions $\geq d$ sur $T$ et si $t \in T$, alors $\mathcal{F}_T$ est
libre d'un certain rang $r$ dans un voisinage ouvert $U \subset X_T$
de l'unique point générique $\xi$ de $X_t$. En effet, $H_r$
contient alors l'image de $U$, qui est un voisinage ouvert de $t$.
L'existence de $U$ résulte de
Compléments sur les morphismes, Lemme
\ref{more-morphisms-lemma-flat-and-free-at-point-fibre}.
\end{proof}

\begin{lemma}
\label{lemma-flat-dimension-n-representable}
Dans la situation \ref{situation-flat-dimension-n}.
Soient $s \in S$ et $d \geq 0$. Supposons que
\begin{enumerate}
\item il existe un d\'evissage complet
de $\mathcal{F}/X/S$ au-dessus d'un certain point $s \in S$,
\item $X$ soit de présentation finie sur $S$,
\item $\mathcal{F}$ soit un $\mathcal{O}_X$-module de présentation finie, et
\item $\mathcal{F}$ soit plat en dimensions $\geq d + 1$ sur $S$.
\end{enumerate}
Alors, après avoir éventuellement remplacé $S$ par un voisinage ouvert
de $s$, le foncteur $F_d$ (\ref{equation-flat-dimension-n})
est représentable par un monomorphisme $Z_d \to S$ de présentation finie.
\end{lemma}

\begin{proof}
Une remarque préliminaire est que $X$ et $S$ sont des schémas affines et qu'il
suffit de démontrer que $F_d$ est représentable par un monomorphisme de présentation
finie $Z_d \to S$ dans la catégorie des schémas affines sur $S$.
(Bien entendu, nous n'exigeons pas que $Z_d$ soit affine.)
Ainsi, tout au long de la démonstration du
lemme, nous travaillons dans la catégorie des schémas affines sur $S$.

\medskip\noindent
Soit $(Z_k, Y_k, i_k, \pi_k, \mathcal{G}_k, \alpha_k)_{k = 1, \ldots, n}$
un d\'evissage complet de $\mathcal{F}/X/S$ au-dessus de $s$ ; voir la
définition \ref{definition-complete-devissage}.
Nous allons raisonner par récurrence sur la longueur $n$ du d\'evissage.
Rappelons que $Y_k \to S$ est lisse à fibres géométriquement irréductibles ; voir la
définition \ref{definition-one-step-devissage}.
Soit $d_k$ la dimension relative de $Y_k$ sur $S$.
Rappelons que $i_{k, *}\mathcal{G}_k = \Coker(\alpha_k)$ et
que $i_k$ est une immersion fermée.
D'après les définitions citées ci-dessus, nous avons
$d_1 = \dim(\text{Supp}(\mathcal{F}_s))$ et
$$
d_k = \dim(\text{Supp}(\Coker(\alpha_{k - 1})_s))
= \dim(\text{Supp}(\mathcal{G}_{k, s}))
$$
pour $k = 2, \ldots, n$. Il s'ensuit que $d_1 > d_2 > \ldots > d_n \geq 0$
parce que $\alpha_k$ est un isomorphisme au point générique de $(Y_k)_s$.

\medskip\noindent
Remarquons que $i_1$ est une immersion fermée et que
$\mathcal{F} = i_{1, *}\mathcal{G}_1$.
Ainsi, pour tout morphisme de schémas $T \to S$ avec $T$ affine,
nous avons $\mathcal{F}_T = i_{1, T, *}\mathcal{G}_{1, T}$ et
$i_{1, T}$ est encore une immersion fermée de schémas sur $T$.
Ainsi, $\mathcal{F}_T$ est plat en dimensions $\geq d$ sur $T$
si et seulement si $\mathcal{G}_{1, T}$ est plat en dimensions $\geq d$ sur $T$.
Comme $\pi_1 : Z_1 \to Y_1$ est fini, nous voyons de même que
$\mathcal{G}_{1, T}$ est plat en dimensions $\geq d$ sur $T$
si et seulement si $\pi_{1, T, *}\mathcal{G}_{1, T}$ est plat en dimensions
$\geq d$ sur $T$. Les mêmes arguments s'appliquent à
« plat en dimensions $\geq d + 1$ » et nous concluons en particulier que
$\pi_{1, *}\mathcal{G}_1$ est plat sur $S$ en dimensions $\geq d + 1$
d'après notre hypothèse sur $\mathcal{F}$.

\medskip\noindent
Supposons que $d_1 > d$. Il résulte de la discussion précédente qu'en
particulier $\pi_{1, *}\mathcal{G}_1$ est plat sur $S$ au
point générique de $(Y_1)_s$. D'après le
lemme \ref{lemma-induction-step-fp},
nous pouvons remplacer $S$ par un voisinage affine de $s$ et supposer que
$\alpha_1$ est $S$-universellement injectif. Puisque $\alpha_1$ est
$S$-universellement injectif, pour tout morphisme $T \to S$ avec $T$ affine,
nous avons une suite exacte courte
$$
0 \to \mathcal{O}_{Y_{1, T}}^{\oplus r_1}
\to \pi_{1, T, *}\mathcal{G}_{1, T} \to \Coker(\alpha_1)_T \to 0
$$
et la première flèche reste $T$-universellement injective. Ainsi, l'ensemble
des points de $(Y_1)_T$ où $\pi_{1, T, *}\mathcal{G}_{1, T}$ est plat sur
$T$ est le même que l'ensemble des points de $(Y_1)_T$ où
$\Coker(\alpha_1)_T$ est plat sur $S$. De cette manière, la question
se ramène au faisceau $\Coker(\alpha_1)$, qui possède un d\'evissage
complet de longueur $n - 1$, et nous concluons par récurrence.

\medskip\noindent
Si $d_1 < d$, alors $F_d$ est représenté par $S$ et nous concluons.

\medskip\noindent
Le dernier cas est celui où $d_1 = d$. Ce cas résulte d'une combinaison du
lemme \ref{lemma-compare-H-F}
et du
lemme \ref{lemma-free-at-generic-points-representable}.
\end{proof}

\begin{theorem}
\label{theorem-flat-dimension-n-representable}
Dans la situation \ref{situation-flat-dimension-n}.
Supposons en outre que $f$ soit de présentation finie, que
$\mathcal{F}$ soit un $\mathcal{O}_X$-module de présentation finie,
et que $\mathcal{F}$ soit pur relativement à $S$.
Alors $F_n$ est représentable par un monomorphisme
$Z_n \to S$ de présentation finie.
\end{theorem}

\begin{proof}
Le foncteur $F_n$ est un faisceau pour la topologie fppf d'après le
lemme \ref{lemma-flat-dimension-n}.
Observons qu'un monomorphisme de présentation finie est
séparé et quasi-fini (Morphismes, Lemme
\ref{morphisms-lemma-monomorphism-loc-finite-type-loc-quasi-finite}).
En combinant donc
Descente, Lemme \ref{descent-lemma-descent-data-sheaves},
Compléments sur les morphismes, Lemme
\ref{more-morphisms-lemma-separated-locally-quasi-finite-morphisms-fppf-descend}
, et
Descente, Lemmes \ref{descent-lemma-descending-property-monomorphism} et
\ref{descent-lemma-descending-property-finite-presentation},
nous voyons que la question est locale pour la topologie \'etale sur $S$.

\medskip\noindent
En particulier, la situation est locale pour la topologie de Zariski sur $S$
et nous pouvons supposer que $S$ est affine. Dans ce cas, la dimension des
fibres de $f$ est majorée ; nous voyons donc que $F_n$ est représentable
pour $n$ assez grand. Nous pouvons ainsi raisonner par récurrence descendante sur $n$.
Supposons que nous sachions que $F_{n + 1}$ est représentable par un monomorphisme
$Z_{n + 1} \to S$ de présentation finie. Considérons le changement de base
$X_{n + 1} = Z_{n + 1} \times_S X$ et l'image inverse $\mathcal{F}_{n + 1}$
de $\mathcal{F}$ sur $X_{n + 1}$. Le morphisme $Z_{n + 1} \to S$ est
quasi-fini, puisqu'il s'agit d'un monomorphisme de présentation finie ; le
lemme \ref{lemma-quasi-finite-base-change}
implique donc que $\mathcal{F}_{n + 1}$ est pur relativement à $Z_{n + 1}$.
Puisque $F_n$ est un sous-foncteur de $F_{n + 1}$, nous concluons que, pour
démontrer le résultat pour $F_n$, il suffit de le démontrer pour le
foncteur correspondant à la situation
$\mathcal{F}_{n + 1}/X_{n + 1}/Z_{n + 1}$.
Nous nous ramenons ainsi à démontrer le résultat pour $F_n$ dans le cas où
$S_{n + 1} = S$, c'est-à-dire que nous pouvons supposer que $\mathcal{F}$ est plat
en dimensions $\geq n + 1$ sur $S$.

\medskip\noindent
Fixons $n$ et supposons que $\mathcal{F}$ soit plat en dimensions $\geq n + 1$
sur $S$. Pour achever la démonstration, nous devons montrer que $F_n$ est représentable
par un monomorphisme $Z_n \to S$ de présentation finie.
Puisque la question est locale pour la topologie \'etale sur $S$, il suffit de
montrer que, pour tout $s \in S$, il existe un voisinage \'etale élémentaire
$(S', s') \to (S, s)$ tel que le résultat soit vrai après changement de base à $S'$.
Ainsi, d'après le
lemme \ref{lemma-existence-complete},
nous pouvons supposer qu'il existe des morphismes \'etales $h_j : Y_j \to X$,
$j = 1, \ldots, m$, tels que, pour chaque $j$, il existe un d\'evissage complet
de $\mathcal{F}_j/Y_j/S$ au-dessus de $s$,
où $\mathcal{F}_j$ est l'image inverse de $\mathcal{F}$ sur $Y_j$,
et tels que $X_s \subset \bigcup h_j(Y_j)$. Remarquons que, d'après le
lemme \ref{lemma-localize-flat-dimension-n},
les faisceaux $\mathcal{F}_j$ sont encore plats
en dimensions $\geq n + 1$ sur $S$.
Posons $W = \bigcup h_j(Y_j)$, qui est un ouvert quasi-compact de $X$.
Comme $\mathcal{F}$ est pur le long de $X_s$, nous voyons que
$$
E = \{t \in S \mid \text{Ass}_{X_t}(\mathcal{F}_t) \subset W \}.
$$
contient toutes les généralisations de $s$. D'après
Compléments sur les morphismes,
Lemme \ref{more-morphisms-lemma-relative-assassin-constructible},
$E$ est une partie constructible de $S$. Nous avons vu que
$\Spec(\mathcal{O}_{S, s}) \subset E$. D'après
Morphismes, Lemme \ref{morphisms-lemma-constructible-containing-open},
nous voyons que $E$ contient un voisinage ouvert de $s$.
Ainsi, après avoir rétréci $S$, nous pouvons supposer que $E = S$.
Il résulte du
lemme \ref{lemma-localize-flat-dimension-n}
qu'il suffit de démontrer le lemme pour le foncteur $F_n$ associé à
$X = \coprod Y_j$ et $\mathcal{F} = \coprod \mathcal{F}_j$.
Si $F_{j, n}$ désigne le foncteur pour $Y_j \to S$ et le faisceau
$\mathcal{F}_i$, nous voyons que $F_n = \prod F_{j, n}$. Il suffit donc
de démontrer que chaque $F_{j, n}$ est représentable par un monomorphisme
$Z_{j, n} \to S$ de présentation finie, puisque alors
$$
Z_n = Z_{1, n} \times_S \ldots \times_S Z_{m, n}
$$
Nous avons ainsi ramené le théorème au cas particulier traité dans le
lemme \ref{lemma-flat-dimension-n-representable}.
\end{proof}

\noindent
Nous explicitons dans le lemme suivant la signification du théorème en termes
d'aplatissements universels.

\begin{lemma}
\label{lemma-when-universal-flattening}
Soit $f : X \to S$ un morphisme de schémas.
Soit $\mathcal{F}$ un $\mathcal{O}_X$-module quasi-cohérent.
\begin{enumerate}
\item Si $f$ est de présentation finie, si $\mathcal{F}$ est un
$\mathcal{O}_X$-module de présentation finie et si $\mathcal{F}$ est
pur relativement à $S$, alors il existe un aplatissement universel
$S' \to S$ de $\mathcal{F}$. De plus, $S' \to S$ est un monomorphisme
de présentation finie.
\item Si $f$ est de présentation finie et $X$ est pur relativement à $S$,
alors il existe un aplatissement universel $S' \to S$ de $X$.
De plus, $S' \to S$ est un monomorphisme de présentation finie.
\item Si $f$ est propre et de présentation finie et si $\mathcal{F}$ est un
$\mathcal{O}_X$-module de présentation finie, alors il existe un
aplatissement universel $S' \to S$ de $\mathcal{F}$. De plus, $S' \to S$ est
un monomorphisme de présentation finie.
\item Si $f$ est propre et de présentation finie,
alors il existe un aplatissement universel $S' \to S$ de $X$.
\end{enumerate}
\end{lemma}

\begin{proof}
Ces énoncés résultent immédiatement du
théorème \ref{theorem-flat-dimension-n-representable}
appliqué à $F_0 = F_{flat}$,
et du fait que, si $f$ est propre, alors $\mathcal{F}$ est automatiquement
pur sur la base ; voir le
lemme \ref{lemma-proper-pure}.
\end{proof}







\section{Théorème d'existence de Grothendieck, IV}
\label{section-existence}

\noindent
Cette section poursuit l'étude entreprise dans
Cohomologie des schémas, sections \ref{coherent-section-existence},
\ref{coherent-section-existence-proper}, et
\ref{coherent-section-existence-proper-support}.
Nous nous plaçons dans la situation suivante.

\begin{situation}
\label{situation-existence}
On dispose ici d'un système projectif d'anneaux $(A_n)$ dont les morphismes de transition
sont surjectifs et dont les noyaux sont localement nilpotents. Posons $A = \lim A_n$. On dispose d'un
schéma $X$ séparé et de présentation finie sur $A$.
Posons $X_n = X \times_{\Spec(A)} \Spec(A_n)$ et
considérons-le comme un sous-schéma fermé de $X$. On suppose en outre donné un système
$(\mathcal{F}_n, \varphi_n)$ où $\mathcal{F}_n$ est un
$\mathcal{O}_{X_n}$-module de présentation finie, plat sur $A_n$, à support propre sur $A_n$,
et
$$
\varphi_n :
\mathcal{F}_n \otimes_{\mathcal{O}_{X_n}} \mathcal{O}_{X_{n - 1}}
\longrightarrow
\mathcal{F}_{n - 1}
$$
est un isomorphisme (notation fondée sur l'équivalence de
Morphismes, lemme \ref{morphisms-lemma-i-star-equivalence}).
\end{situation}

\noindent
Notre but est de déterminer s'il existe un faisceau quasi-cohérent $\mathcal{F}$
sur $X$ tel que
$\mathcal{F}_n = \mathcal{F} \otimes_{\mathcal{O}_X} \mathcal{O}_{X_n}$
pour tout $n$.

\begin{lemma}
\label{lemma-compute-what-it-should-be}
Dans la situation \ref{situation-existence}, considérons
$$
K = R\lim_{D_\QCoh(\mathcal{O}_X)}(\mathcal{F}_n) =
DQ_X(R\lim_{D(\mathcal{O}_X)}\mathcal{F}_n)
$$
Alors $K$ appartient à $D^b_{\QCoh}(\mathcal{O}_X)$ et, en fait,
$K$ n'a de faisceaux de cohomologie non nuls qu'en degrés $\geq 0$.
\end{lemma}

\begin{proof}
C'est un cas particulier de
Catégories dérivées de schémas,\newline exemple
\ref{perfect-example-inverse-limit-quasi-coherent}.
\end{proof}

\begin{lemma}
\label{lemma-compute-against-perfect}
Dans la situation \ref{situation-existence}, soit $K$ comme dans le
lemme \ref{lemma-compute-what-it-should-be}. Pour tout
objet parfait $E$ de $D(\mathcal{O}_X)$, on a
\begin{enumerate}
\item $M = R\Gamma(X, K \otimes^\mathbf{L} E)$ est un objet parfait de $D(A)$
et il existe un isomorphisme canonique
$R\Gamma(X_n, \mathcal{F}_n \otimes^\mathbf{L} E|_{X_n}) =
M \otimes_A^\mathbf{L} A_n$
dans $D(A_n)$,
\item $N = R\Hom_X(E, K)$ est un objet parfait de $D(A)$
et il existe un isomorphisme canonique
$R\Hom_{X_n}(E|_{X_n}, \mathcal{F}_n) = N \otimes_A^\mathbf{L} A_n$
dans $D(A_n)$.
\end{enumerate}
Dans les deux assertions, $E|_{X_n}$ désigne l'image inverse dérivée
de $E$ sur $X_n$.
\end{lemma}

\begin{proof}
Démontrons (2). Posons $E_n = E|_{X_n}$ et
$N_n = R\Hom_{X_n}(E_n, \mathcal{F}_n)$.
Rappelons que $R\Hom_{X_n}(-, -)$ est égal à
$R\Gamma(X_n, R\SheafHom(-, -))$, voir
Cohomologie, section \ref{cohomology-section-global-RHom}.
Il résulte alors de Catégories dérivées de schémas, lemme
\ref{perfect-lemma-base-change-RHom-perfect}
que $N_n$ est un objet parfait de $D(A_n)$
dont la formation commute au changement de base. Ainsi, les morphismes
$N_n \otimes_{A_n}^\mathbf{L} A_{n - 1} \to N_{n - 1}$
induits par $\varphi_n$ sont des isomorphismes.
D'après Compléments d'algèbre, lemme \ref{more-algebra-lemma-Rlim-perfect-gives-perfect}, nous constatons que
$R\lim N_n$ est parfait et
que son changement de base à $A_n$ redonne $N_n$.
D'autre part, le foncteur exact
$R\Hom_X(E, -) : D_\QCoh(\mathcal{O}_X) \to D(A)$
entre catégories triangulées commute aux produits
et donc aux limites dérivées, d'où
$$
R\Hom_X(E, K) =
R\lim R\Hom_X(E, \mathcal{F}_n) =
R\lim R\Hom_X(E_n, \mathcal{F}_n) =
R\lim N_n
$$
Ceci prouve (2). Pour établir (1), on se ramène à (2)
à l'aide de Cohomologie, lemme \ref{cohomology-lemma-dual-perfect-complex}.
\end{proof}

\begin{lemma}
\label{lemma-relative-pseudo-coherence}
Dans la situation \ref{situation-existence}, soit $K$ comme dans le
lemme \ref{lemma-compute-what-it-should-be}. Alors $K$
est pseudo-cohérent relativement à $A$.
\end{lemma}

\begin{proof}
En combinant le lemme \ref{lemma-compute-against-perfect} et
Catégories dérivées de schémas, lemme \ref{perfect-lemma-perfect-enough},
on voit que $R\Gamma(X, K \otimes^\mathbf{L} E)$
est pseudo-cohérent dans $D(A)$ pour tout objet pseudo-cohérent
$E$ de $D(\mathcal{O}_X)$. Le lemme découle donc de
Compléments sur les morphismes, lemme
\ref{more-morphisms-lemma-characterize-pseudo-coherent}.
\end{proof}

\begin{lemma}
\label{lemma-compute-over-affine}
Dans la situation \ref{situation-existence}, soit $K$ comme dans le
lemme \ref{lemma-compute-what-it-should-be}. Pour tout
ouvert quasi-compact $U \subset X$, on a
$$
R\Gamma(U, K) \otimes_A^\mathbf{L} A_n =
R\Gamma(U_n, \mathcal{F}_n)
$$
dans $D(A_n)$, où $U_n = U \cap X_n$.
\end{lemma}

\begin{proof}
Fixons $n$. D'après Catégories dérivées de schémas, lemme
\ref{perfect-lemma-computing-sections-as-colim}
il existe un système de complexes parfaits $E_m$
sur $X$ tel que\newline
$R\Gamma(U, K) = \text{hocolim} R\Gamma(X, K \otimes^\mathbf{L} E_m)$.
En fait, cette formule vaut non seulement pour $K$, mais pour tout objet de
$D_\QCoh(\mathcal{O}_X)$.
En l'appliquant à $\mathcal{F}_n$
nous obtenons
\begin{align*}
R\Gamma(U_n, \mathcal{F}_n)
& =
R\Gamma(U, \mathcal{F}_n) \\
& =
\text{hocolim}_m R\Gamma(X, \mathcal{F}_n \otimes^\mathbf{L} E_m) \\
& =
\text{hocolim}_m R\Gamma(X_n, \mathcal{F}_n \otimes^\mathbf{L} E_m|_{X_n})
\end{align*}
En utilisant le lemme \ref{lemma-compute-against-perfect}
et le fait que $- \otimes_A^\mathbf{L} A_n$
commute aux colimites homotopiques, nous obtenons le résultat.
\end{proof}

\begin{lemma}
\label{lemma-finitely-presented}
Dans la situation \ref{situation-existence}, soit $K$ comme dans le
lemme \ref{lemma-compute-what-it-should-be}.
Notons $X_0 \subset X$ le sous-ensemble fermé
constitué des points situés au-dessus du sous-ensemble fermé
$\Spec(A_1) = \Spec(A_2) = \ldots$ de $\Spec(A)$.
Il existe un ouvert $W \subset X$ contenant $X_0$
tel que
\begin{enumerate}
\item $H^i(K)|_W$ est nul sauf si $i = 0$,
\item $\mathcal{F} = H^0(K)|_W$ est de présentation finie, et
\item $\mathcal{F}_n = \mathcal{F} \otimes_{\mathcal{O}_X} \mathcal{O}_{X_n}$.
\end{enumerate}
\end{lemma}

\begin{proof}
Fixons $n \geq 1$. Par construction, il existe un morphisme canonique
$K \to \mathcal{F}_n$ dans $D_\QCoh(\mathcal{O}_X)$
et donc un morphisme canonique $H^0(K) \to \mathcal{F}_n$
de faisceaux quasi-cohérents. Ceci explique le sens du point (3).

\medskip\noindent
Soit $x \in X_0$ un point.
Nous allons trouver un voisinage ouvert $W$
de $x$ tel que (1), (2) et (3) soient vrais. Puisque $X_0$ est quasi-compact,
ceci démontrera le lemme. Soit $U \subset X$ un voisinage ouvert affine
de $x$. Écrivons $U = \Spec(B)$.
Choisissons une surjection $P \to B$ avec $P$ lisse sur $A$.
D'après le lemme \ref{lemma-relative-pseudo-coherence}
et la définition de la pseudo-cohérence relative,
il existe un complexe $F^\bullet$ borné supérieurement
de $P$-modules libres de rang fini représentant
$Ri_*K$, où $i : U \to \Spec(P)$ est l'immersion
fermée induite par la présentation.
Soit $M_n$ le $B$-module correspondant à $\mathcal{F}_n|_U$.
D'après le lemme \ref{lemma-compute-over-affine},
$$
H^i(F^\bullet \otimes_A A_n) =
\left\{
\begin{matrix}
0 & \text{si} & i \not = 0 \\
M_n & \text{si} & i = 0
\end{matrix}
\right.
$$
Soit $i$ l'indice maximal tel que $F^i$ soit non nul.
Si $i \leq 0$, alors (1), (2) et (3) sont vrais.
Sinon, $i > 0$ et nous voyons que le rang du morphisme
$$
F^{i - 1} \to F^i
$$
au point $x$ est maximal. Il est donc maximal dans un voisinage ouvert
de $x$ dans $\Spec(P)$. Ainsi, après avoir remplacé
$P$ par une localisation principale, nous pouvons supposer que le morphisme
affiché est surjectif. Puisque $F^i$ est libre de rang fini, nous pouvons choisir
une décomposition $F^{i - 1} = F' \oplus F^i$. Nous pouvons alors
remplacer $F^\bullet$ par le complexe
$$
\ldots \to F^{i - 2} \to F' \to 0 \to \ldots
$$
et conclure par récurrence sur $i$.
\end{proof}

\begin{lemma}
\label{lemma-proper-support}
Dans la situation \ref{situation-existence}, soit $K$ comme dans le
lemme \ref{lemma-compute-what-it-should-be}. Soit $W \subset X$
comme dans le lemme \ref{lemma-finitely-presented}.
Posons $\mathcal{F} = H^0(K)|_W$. Alors, après avoir éventuellement restreint l'ouvert $W$,
le support de $\mathcal{F}$ est propre sur $A$.
\end{lemma}

\begin{proof}
Fixons $n \geq 1$. Soit $I_n = \Ker(A \to A_n)$.
D'après le lemme \ref{more-algebra-lemma-limit-henselian} de Compléments d'algèbre,
le couple $(A, I_n)$ est hensélien.
Soit $Z \subset W$ le support de $\mathcal{F}$.
C'est un sous-ensemble fermé puisque $\mathcal{F}$ est de présentation finie.
D'après le point (3) du lemme \ref{lemma-finitely-presented},
nous voyons que $Z \times_{\Spec(A)} \Spec(A_n)$
est égal au support de $\mathcal{F}_n$ et est donc
propre sur $\Spec(A/I)$.
D'après le lemme de Compléments sur les morphismes
\ref{more-morphisms-lemma-split-off-proper-part-henselian}
nous pouvons écrire $Z = Z_1 \amalg Z_2$, où $Z_1, Z_2$ sont ouverts et
fermés dans $Z$, où $Z_1$ est propre
sur $A$ et où $Z_1 \times_{\Spec(A)} \Spec(A/I_n)$
est égal au support de $\mathcal{F}_n$.
Autrement dit, $Z_2$ ne rencontre pas $X_0$.
Ainsi, après avoir remplacé $W$ par $W \setminus Z_2$, nous obtenons le lemme.
\end{proof}

\begin{lemma}
\label{lemma-monomorphism-isomorphism}
Soit $A = \lim A_n$ la limite d'un système d'anneaux
dont les morphismes de transition sont surjectifs et les noyaux
localement nilpotents. Soit $S = \Spec(A)$. Soit $T \to S$ un monomorphisme
localement de type fini. Si $\Spec(A_n) \to S$
se factorise par $T$ pour tout $n$, alors $T = S$.
\end{lemma}

\begin{proof}
Posons $S_n = \Spec(A_n)$. Soit $T_0 \subset T$ l'image
commune des factorisations $S_n \to T$. Alors $T_0$ est quasi-compact.
Soit $T' \subset T$ un ouvert quasi-compact contenant $T_0$.
Alors $S_n \to T$ se factorise par $T'$.
Si nous pouvons montrer que $T' = S$, alors $T' = T = S$.
Nous pouvons donc supposer que $T$ est quasi-compact.

\medskip\noindent
Supposons que $T$ soit quasi-compact.
Dans ce cas, $T \to S$ est séparé et quasi-fini
(lemme de Morphismes
\ref{morphisms-lemma-monomorphism-loc-finite-type-loc-quasi-finite}).
En utilisant le théorème principal de Zariski
(sous la forme du lemme de Compléments sur les morphismes
\ref{more-morphisms-lemma-quasi-finite-separated-pass-through-finite})
nous choisissons une factorisation $T \to W \to S$ avec $W \to S$ fini
et $T \to W$ une immersion ouverte. Écrivons $W = \Spec(B)$.
Les factorisations (uniques) $S_n \to T$ peuvent être vues
comme des morphismes vers $W$ et nous obtenons
$$
A \longrightarrow B \longrightarrow \lim A_n = A
$$
Considérons le morphisme $h : S = \Spec(A) \to \Spec(B) = W$ provenant
de la flèche de droite. Alors
$$
T \times_{W, h} S
$$
est un sous-schéma ouvert de $S$ contenant l'image de $S_n \to S$ pour tout $n$.
Pour achever la démonstration, il suffit de montrer que tout ouvert $U \subset S$
contenant l'image de $S_n \to S$ pour un certain $n \geq 1$ est égal à $S$.
C'est vrai parce que $(A, \Ker(A \to A_n))$ est un couple hensélien
(lemme \ref{more-algebra-lemma-limit-henselian} de Compléments d'algèbre)
et que, par conséquent, tout point fermé de $S$ appartient à l'image de $S_n \to S$.
\end{proof}

\begin{theorem}[Théorème d'existence de Grothendieck]
\label{theorem-existence}
Dans la situation \ref{situation-existence},
il existe un $\mathcal{O}_X$-module de présentation finie
$\mathcal{F}$, plat sur $A$, dont le support est propre sur $A$,
tel que
$\mathcal{F}_n = \mathcal{F} \otimes_{\mathcal{O}_X} \mathcal{O}_{X_n}$
pour tout $n$, de manière compatible avec les morphismes $\varphi_n$.
\end{theorem}

\begin{proof}
Appliquons les lemmes \ref{lemma-compute-what-it-should-be},
\ref{lemma-compute-against-perfect},
\ref{lemma-relative-pseudo-coherence},
\ref{lemma-compute-over-affine},
\ref{lemma-finitely-presented}, et
\ref{lemma-proper-support}
pour obtenir un sous-schéma ouvert $W \subset X$ contenant tous les points
situés au-dessus de $\Spec(A_n)$
et un $\mathcal{O}_W$-module de présentation finie $\mathcal{F}$
dont le support est propre sur $A$, avec
$\mathcal{F}_n = \mathcal{F} \otimes_{\mathcal{O}_W} \mathcal{O}_{X_n}$
pour tout $n \geq 1$. (Ceci a un sens puisque $X_n \subset W$.)
D'après le lemme \ref{lemma-proper-pure}, nous voyons que $\mathcal{F}$
est universellement pur relativement à $\Spec(A)$.
D'après le théorème \ref{theorem-flat-dimension-n-representable}
(pour une explication, voir le lemme \ref{lemma-when-universal-flattening}),
il existe un aplatissement universel $S' \to \Spec(A)$
de $\mathcal{F}$ et, de plus, le morphisme $S' \to \Spec(A)$
est un monomorphisme de présentation finie.
Puisque le changement de base de $\mathcal{F}$ à $\Spec(A_n)$
est $\mathcal{F}_n$, nous trouvons que $\Spec(A_n) \to \Spec(A)$
se factorise (de manière unique) par $S'$ pour tout $n$.
D'après le lemme \ref{lemma-monomorphism-isomorphism},
nous voyons que $S' = \Spec(A)$.
Cela signifie que $\mathcal{F}$ est plat sur $A$.
Enfin, puisque le support schématique $Z$ de $\mathcal{F}$
est propre sur $\Spec(A)$, le morphisme $Z \to X$ est fermé.
Par conséquent, l'image directe $(W \to X)_*\mathcal{F}$ a son support
dans $W$ et possède toutes les propriétés voulues.
\end{proof}






\section{Théorème d'existence de Grothendieck, V}
\label{section-existence-derived}

\noindent
Dans cette section, nous démontrons un analogue du théorème d'existence de Grothendieck
dans la catégorie dérivée, en suivant la méthode utilisée dans la
section \ref{section-existence} pour les modules quasi-cohérents.
Le cas classique est étudié dans
Cohomologie des schémas, sections \ref{coherent-section-existence},
\ref{coherent-section-existence-proper}, et
\ref{coherent-section-existence-proper-support}.
Nous travaillerons dans la situation suivante.

\begin{situation}
\label{situation-existence-derived}
Nous avons ici un système projectif d'anneaux $(A_n)$ dont les morphismes de transition
sont surjectifs et dont les noyaux sont localement nilpotents. Posons $A = \lim A_n$. Nous avons un
schéma $X$ propre, plat et de présentation finie sur $A$.
Nous posons $X_n = X \times_{\Spec(A)} \Spec(A_n)$ et le
considérons comme un sous-schéma fermé de $X$. Supposons en outre donné un système
$(K_n, \varphi_n)$ où $K_n$ est un objet pseudo-cohérent de
$D(\mathcal{O}_{X_n})$ et
$$
\varphi_n : K_n \longrightarrow K_{n - 1}
$$
est un morphisme dans $D(\mathcal{O}_{X_n})$ qui induit un isomorphisme
$K_n \otimes_{\mathcal{O}_{X_n}}^\mathbf{L} \mathcal{O}_{X_{n - 1}}
\to K_{n - 1}$ dans $D(\mathcal{O}_{X_{n - 1}})$.
\end{situation}

\noindent
Plus précisément, nous devrions écrire
$\varphi_n : K_n \to Ri_{n - 1, *}K_{n - 1}$
où $i_{n - 1} : X_{n - 1} \to X_n$ est le morphisme d'inclusion,
et, dans cette notation, la condition est que le morphisme
adjoint $Li_{n - 1}^*K_n \to K_{n - 1}$ soit un isomorphisme.
Notre but est de trouver un $K \in D(\mathcal{O}_X)$ pseudo-cohérent
tel que $K_n = K \otimes_{\mathcal{O}_X}^\mathbf{L} \mathcal{O}_{X_n}$
pour tout $n$ (avec le même abus de notation).

\begin{lemma}
\label{lemma-compute-what-it-should-be-derived}
Dans la situation \ref{situation-existence-derived}, considérons
$$
K = R\lim_{D_\QCoh(\mathcal{O}_X)}(K_n) =
DQ_X(R\lim_{D(\mathcal{O}_X)} K_n)
$$
Alors $K$ appartient à $D^-_{\QCoh}(\mathcal{O}_X)$.
\end{lemma}

\begin{proof}
Le foncteur $DQ_X$ existe parce que $X$ est quasi-compact et
quasi-séparé, voir le lemme de Catégories dérivées de schémas
\ref{perfect-lemma-better-coherator}.
Puisque $DQ_X$ est un adjoint à droite, il commute aux produits
et donc aux limites dérivées. D'où l'égalité
dans l'énoncé du lemme.

\medskip\noindent
D'après le lemme de Catégories dérivées de schémas
\ref{perfect-lemma-boundedness-better-coherator},
le foncteur $DQ_X$ est de dimension cohomologique bornée.
Il suffit donc de montrer que $R\lim K_n \in D^-(\mathcal{O}_X)$.
Pour le voir, soit $U \subset X$ un ouvert affine.
Il existe alors une suite exacte canonique
$$
0 \to
R^1\lim H^{m - 1}(U, K_n) \to H^m(U, R\lim K_n) \to
\lim H^m(U, K_n) \to 0
$$
d'après le lemme \ref{cohomology-lemma-RGamma-commutes-with-Rlim} de Cohomologie.
Puisque $U$ est affine et $K_n$ est pseudo-cohérent (et possède donc des
faisceaux de cohomologie quasi-cohérents d'après le lemme
\ref{perfect-lemma-pseudo-coherent} de Catégories dérivées de schémas),
nous voyons que $H^m(U, K_n) = H^m(K_n)(U)$\newline d'après le
lemme de Catégories dérivées de schémas
\ref{perfect-lemma-affine-compare-bounded}.
Nous concluons ainsi qu'il suffit de montrer que $K_n$
est borné supérieurement indépendamment de $n$.

\medskip\noindent
Puisque $K_n$ est pseudo-cohérent, nous avons $K_n \in D^-(\mathcal{O}_{X_n})$.
Supposons que $a_n$ soit maximal tel que $H^{a_n}(K_n)$ soit non nul.
Bien entendu, $a_1 \leq a_2 \leq a_3 \leq \ldots$.
Notons que $H^{a_n}(K_n)$ est un
$\mathcal{O}_{X_n}$-module de présentation finie
(lemme \ref{cohomology-lemma-finite-cohomology} de Cohomologie).
Nous avons $H^{a_n}(K_{n - 1}) =
H^{a_n}(K_n) \otimes_{\mathcal{O}_{X_n}} \mathcal{O}_{X_{n - 1}}$.
Puisque $X_{n - 1} \to X_n$ est un épaississement, il résulte du
lemme de Nakayama (lemme \ref{algebra-lemma-NAK} d'Algèbre) que si
$H^{a_n}(K_n) \otimes_{\mathcal{O}_{X_n}} \mathcal{O}_{X_{n - 1}}$
est nul, alors $H^{a_n}(K_n)$ est lui aussi nul. Ainsi $a_n = a_{n - 1}$
pour tout $n$, et nous concluons.
\end{proof}

\begin{lemma}
\label{lemma-compute-against-perfect-derived}
Dans la situation \ref{situation-existence-derived}, soit $K$ comme dans le
lemme \ref{lemma-compute-what-it-should-be-derived}. Pour tout objet parfait
$E$ de $D(\mathcal{O}_X)$, la cohomologie
$$
M = R\Gamma(X, K \otimes^\mathbf{L} E)
$$
est un objet pseudo-cohérent de $D(A)$ et il existe un isomorphisme canonique
$$
R\Gamma(X_n, K_n \otimes^\mathbf{L} E|_{X_n}) = M \otimes_A^\mathbf{L} A_n
$$
dans $D(A_n)$. Ici, $E|_{X_n}$ désigne l'image inverse dérivée de $E$ sur $X_n$.
\end{lemma}

\begin{proof}
Écrivons $E_n = E|_{X_n}$ et
$M_n = R\Gamma(X_n, K_n \otimes^\mathbf{L} E|_{X_n})$.
D'après le lemme de Catégories dérivées de schémas
\ref{perfect-lemma-flat-proper-pseudo-coherent-direct-image-general}
nous voyons que $M_n$ est un objet pseudo-cohérent de $D(A_n)$
dont la formation commute au changement de base. Ainsi, les morphismes
$M_n \otimes_{A_n}^\mathbf{L} A_{n - 1} \to M_{n - 1}$
provenant de $\varphi_n$ sont des isomorphismes. D'après le
lemme de Compléments d'algèbre
\ref{more-algebra-lemma-Rlim-pseudo-coherent-gives-pseudo-coherent}
nous trouvons que $R\lim M_n$ est pseudo-cohérent et
que son changement de base à $A_n$ redonne $M_n$.
D'autre part, le foncteur exact
$R\Gamma(X, -) : D_\QCoh(\mathcal{O}_X) \to D(A)$
entre catégories triangulées commute aux produits
et donc aux limites dérivées, d'où
$$
R\Gamma(X, E \otimes^\mathbf{L} K) =
R\lim R\Gamma(X,  E \otimes^\mathbf{L} K_n) =
R\lim R\Gamma(X_n, E_n \otimes^\mathbf{L} K_n) =
R\lim M_n
$$
comme souhaité.
\end{proof}

\begin{lemma}
\label{lemma-relative-pseudo-coherence-derived}
Dans la situation \ref{situation-existence-derived}, soit $K$ comme dans le
lemme \ref{lemma-compute-what-it-should-be-derived}. Alors $K$
est pseudo-cohérent sur $X$.
\end{lemma}

\begin{proof}
En combinant le lemme \ref{lemma-compute-against-perfect-derived} et le
lemme de Catégories dérivées de schémas
\ref{perfect-lemma-perfect-enough}
nous voyons que $R\Gamma(X, K \otimes^\mathbf{L} E)$
est pseudo-cohérent dans $D(A)$ pour tout objet pseudo-cohérent
$E$ de $D(\mathcal{O}_X)$. Il résulte donc du
lemme de Compléments sur les morphismes
\ref{more-morphisms-lemma-characterize-pseudo-coherent}
que $K$ est pseudo-cohérent relativement à $A$.
Puisque $X$ est plat et de présentation finie
sur $A$, cela équivaut à être pseudo-cohérent sur $X$, voir le
lemme de Compléments sur les morphismes
\ref{more-morphisms-lemma-check-relative-pseudo-coherence-on-charts}.
\end{proof}

\begin{lemma}
\label{lemma-compute-over-affine-derived}
Dans la situation \ref{situation-existence-derived}, soit $K$ comme dans le
lemme \ref{lemma-compute-what-it-should-be-derived}. Pour tout ouvert quasi-compact
$U \subset X$, nous avons
$$
R\Gamma(U, K) \otimes_A^\mathbf{L} A_n =
R\Gamma(U_n, K_n)
$$
dans $D(A_n)$, où $U_n = U \cap X_n$.
\end{lemma}

\begin{proof}
Fixons $n$. D'après le lemme de Catégories dérivées de schémas
\ref{perfect-lemma-computing-sections-as-colim}
il existe un système de complexes parfaits $E_m$
sur $X$ tel que
$R\Gamma(U, K) = \text{hocolim} R\Gamma(X, K \otimes^\mathbf{L} E_m)$.
En fait, cette formule vaut non seulement pour $K$, mais pour tout objet de
$D_\QCoh(\mathcal{O}_X)$.
En l'appliquant à $K_n$,
nous obtenons
\begin{align*}
R\Gamma(U_n, K_n)
& =
R\Gamma(U, K_n) \\
& =
\text{hocolim}_m R\Gamma(X, K_n \otimes^\mathbf{L} E_m) \\
& =
\text{hocolim}_m R\Gamma(X_n, K_n \otimes^\mathbf{L} E_m|_{X_n})
\end{align*}
En utilisant le lemme \ref{lemma-compute-against-perfect-derived}
et le fait que $- \otimes_A^\mathbf{L} A_n$
commute aux colimites homotopiques, nous obtenons le résultat.
\end{proof}

\begin{theorem}[Théorème d'existence de Grothendieck dérivé]
\label{theorem-existence-derived}
Dans la situation \ref{situation-existence-derived},
il existe un $K$ pseudo-cohérent dans $D(\mathcal{O}_X)$
tel que $K_n = K \otimes_{\mathcal{O}_X}^\mathbf{L} \mathcal{O}_{X_n}$
pour tout $n$, de manière compatible avec les morphismes $\varphi_n$.
\end{theorem}

\begin{proof}
Appliquons les lemmes \ref{lemma-compute-what-it-should-be-derived},
\ref{lemma-compute-against-perfect-derived},
\ref{lemma-relative-pseudo-coherence-derived}
pour obtenir un objet pseudo-cohérent $K$ de $D(\mathcal{O}_X)$.
En choisissant des ouverts affines dans le lemme
\ref{lemma-compute-over-affine-derived}
il s'ensuit immédiatement que $K$ a pour restriction $K_n$ sur $X_n$.
\end{proof}

\begin{remark}
\label{remark-correct-generality}
Le résultat de cette section peut être généralisé. Il est probablement valable
si nous supposons seulement que $X \to \Spec(A)$ est séparé et de présentation finie,
et que $K_n$ est pseudo-cohérent relativement à $A_n$ et à support dans un sous-ensemble
fermé de $X_n$ propre sur $A_n$. Le résultat sera un $K$ qui
est pseudo-cohérent relativement à $A$ et à support dans un sous-ensemble fermé
propre sur $A$. Si nous en avons un jour besoin, nous
formulerons un énoncé précis et le démontrerons ici.
\end{remark}






\section{Éclatements et platitude}
\label{section-blowup-flat}

\noindent
Dans cette section, nous poursuivons l'étude de résultats de la forme : « Après un
éclatement, le transformé strict devient plat », voir la
section \ref{more-algebra-section-blowup-flat} de Compléments d'algèbre et la
section \ref{divisors-section-blowup-flat} de Diviseurs.
Nous utiliserons dans cette section la notation suivante
(plus ou moins standard). Si $X \to S$ est
un morphisme de schémas, si $\mathcal{F}$ est un module quasi-cohérent
sur $X$ et si $T \to S$ est un morphisme de schémas, alors nous notons
$\mathcal{F}_T$ l'image inverse de $\mathcal{F}$ sur le changement de base
$X_T = X \times_S T$.

\begin{remark}
\label{remark-successive-blowups}
Soit $S$ un schéma quasi-compact et quasi-séparé. Soit $f : X \to S$
un morphisme de schémas. Soit $\mathcal{F}$ un module quasi-cohérent sur $X$.
Soit $U \subset S$ un sous-schéma ouvert quasi-compact. Étant donné un éclatement $U$-admissible
$S' \to S$, nous notons $X'$ le transformé strict de $X$ et $\mathcal{F}'$
le transformé strict de $\mathcal{F}$, considéré comme un module quasi-cohérent
sur $X'$ (via le lemme \ref{divisors-lemma-strict-transform} de Diviseurs).
Soit $P$ une propriété de $\mathcal{F}/X/S$ qui est stable par passage au transformé
strict (comme ci-dessus) pour les éclatements $U$-admissibles. Le problème général de
cette section est de montrer (sous des conditions auxiliaires sur $\mathcal{F}/X/S$)
qu'il existe un éclatement $U$-admissible $S' \to S$
tel que le transformé strict $\mathcal{F}'/X'/S'$ possède $P$.

\medskip\noindent
La stratégie générale consistera à utiliser le fait qu'une composée
d'éclatements $U$-admissibles est un éclatement $U$-admissible, voir le
lemme \ref{divisors-lemma-composition-admissible-blowups} de Diviseurs.
En fait, nous utiliserons le lemme plus précis
\ref{divisors-lemma-composition-finite-type-blowups} de Diviseurs
et le combinerons avec le
lemme \ref{divisors-lemma-strict-transform-composition-blowups} de Diviseurs.
Il en résulte qu'il suffit de trouver une suite d'éclatements $U$-admissibles
successifs
$$
S = S_0 \leftarrow S_1 \leftarrow \ldots \leftarrow S_n
$$
telle que, en posant $\mathcal{F}_0 = \mathcal{F}$ et $X_0 = X$, puis en définissant
$\mathcal{F}_i/X_i$ comme le transformé strict de
$\mathcal{F}_{i - 1}/X_{i - 1}$, nous
aboutissions à $\mathcal{F}_n/X_n/S_n$ possédant la propriété $P$.

\medskip\noindent
En particulier, choisissons un faisceau quasi-cohérent d'idéaux de type fini
$\mathcal{I} \subset \mathcal{O}_S$ tel que $V(\mathcal{I}) = S \setminus U$,
voir le lemme \ref{properties-lemma-quasi-coherent-finite-type-ideals} de Propriétés.
Soit $S' \to S$ l'éclatement le long de $\mathcal{I}$ et soit $E \subset S'$
le diviseur exceptionnel (lemme de Diviseurs
\ref{divisors-lemma-blowing-up-gives-effective-Cartier-divisor}).
Nous voyons alors que nous avons ramené le
problème au cas où il existe un diviseur de Cartier effectif
$D \subset S$ dont le support est $X \setminus U$. En particulier, nous pouvons
supposer que $U$ est schématiquement dense dans $S$
(lemme \ref{divisors-lemma-complement-effective-Cartier-divisor} de Diviseurs).

\medskip\noindent
Supposons que $P$ soit locale sur $S$ : si $S = \bigcup S_i$ est un recouvrement ouvert
fini par des ouverts quasi-compacts et si $P$ vaut pour
$\mathcal{F}_{S_i}/X_{S_i}/S_i$, alors $P$ vaut pour $\mathcal{F}/X/S$.
Dans ce cas, le problème général ci-dessus est lui aussi local sur $S$, c'est-à-dire que,
si, pour tout $s \in S$, nous pouvons trouver un voisinage ouvert quasi-compact $W$ de $s$
tel que le problème pour $\mathcal{F}_W/X_W/W$ soit résoluble, alors le
problème est résoluble pour $\mathcal{F}/X/S$. Cela résulte des
lemmes \ref{divisors-lemma-extend-admissible-blowups} et
\ref{divisors-lemma-dominate-admissible-blowups}.
\end{remark}

\begin{lemma}
\label{lemma-helper-blowup-affine-space}
Soit $R$ un anneau et soit $f \in R$. Soit $r\geq 0$ un entier.
Soit $R \to S$ un morphisme d'anneaux et soit $M$ un $S$-module. Supposons que
\begin{enumerate}
\item $R \to S$ est de présentation finie et plat,
\item tout anneau fibre $S \otimes_R \kappa(\mathfrak p)$ est
géométriquement intègre sur $R$,
\item $M$ est un $S$-module de type fini,
\item $M_f$ est un $S_f$-module de présentation finie,
\item pour tout $\mathfrak p \in R$, $f \not \in \mathfrak p$, en posant
$\mathfrak q = \mathfrak pS$, le module $M_{\mathfrak q}$ est libre
de rang $r$ sur $S_\mathfrak q$.
\end{enumerate}
Alors il existe un idéal de type fini $I \subset R$ avec
$V(f) = V(I)$ tel que, pour tout $a \in I$, en posant $R' = R[\frac{I}{a}]$,
le quotient
$$
M' = (M \otimes_R R')/a\text{-torsion}
$$
sur $S' = S \otimes_R R'$ possède la propriété suivante : pour tout idéal premier
$\mathfrak p' \subset R'$, il existe $g \in S'$,
$g \not \in \mathfrak p'S'$, tel que $M'_g$ soit un $S'_g$-module libre
de rang $r$.
\end{lemma}

\begin{proof}
Ce lemme est une généralisation du
lemme \ref{more-algebra-lemma-blowup-module} de Compléments d'algèbre ;
nous recommandons au lecteur de lire d'abord cette démonstration.
Choisissons une surjection $S^{\oplus n} \to M$, ce qui est possible par (1).
Choisissons un sous-module de type fini $K \subset \Ker(S^{\oplus n} \to M)$
tel que $S^{\oplus n}/K \to M$ devienne un isomorphisme après inversion de $f$.
Ceci est possible par (4). Posons $M_1 = S^{\oplus n}/K$ et supposons que nous puissions
démontrer le lemme pour $M_1$. Soit $I \subset R$ l'idéal correspondant.
Alors, pour $a \in I$, le morphisme
$$
M_1' = (M_1 \otimes_R R')/a\text{-torsion}
\longrightarrow
M' = (M \otimes_R R')/a\text{-torsion}
$$
est surjectif. C'est aussi un isomorphisme après inversion de $a$ dans $R'$,
car $R'_a = R_f$, voir le lemme \ref{algebra-lemma-blowup-in-principal} d'Algèbre.
Mais $a$ est un non-diviseur de zéro dans $M'_1$, de sorte que le morphisme affiché est un
isomorphisme. Il suffit donc de démontrer le lemme lorsque $M$ est un module
de présentation finie sur $S$.

\medskip\noindent
Supposons que $M$ soit un $S$-module de présentation finie vérifiant (3).
Alors $J = \text{Fit}_r(M) \subset S$ est un idéal de type fini.
D'après le lemme \ref{lemma-fibres-irreducible-flat-projective-nonnoetherian},
nous pouvons écrire $S$ comme facteur direct d'un
$R$-module libre : $\bigoplus_{\alpha \in A} R = S \oplus C$.
Pour tout élément $h \in S$, en écrivant $h = \sum a_\alpha$ dans la
décomposition ci-dessus, nous disons que les $a_\alpha$ sont les coefficients de $h$.
Soit $I' \subset R$ l'idéal des coefficients
des éléments de $J$. La multiplication par un élément de $S$ définit
un morphisme $R$-linéaire $S \to S$ ; par conséquent, $I'$ est engendré par les coefficients
des générateurs de $J$, c'est-à-dire que $I'$ est un idéal de type fini.
Nous affirmons que $I = fI'$ convient.

\medskip\noindent
Vérifions d'abord que $V(f) = V(I)$. L'inclusion $V(f) \subset V(I)$ est
claire. Réciproquement, si $f \not \in \mathfrak p$, alors
$\mathfrak q =  \mathfrak p S$ n'est pas un élément de $V(J)$ d'après la
propriété (5) et le lemme de Compléments d'algèbre
\ref{more-algebra-lemma-fitting-ideal-generate-locally}.
Il existe donc un
élément de $J$ dont l'image dans $S \otimes_R \kappa(\mathfrak p)$ n'est pas nulle.
Ainsi, il existe un élément de $I'$ qui n'appartient pas à
$\mathfrak p$, donc $\mathfrak p \not \in V(fI') = V(I)$.

\medskip\noindent
Soit $a \in I$ et posons $R' = R[\frac{I}{a}]$. Nous pouvons écrire $a = fa'$
pour un certain $a' \in I'$. D'après les lemmes \ref{algebra-lemma-affine-blowup} et
\ref{algebra-lemma-blowup-add-principal} d'Algèbre, nous voyons que $I' R' = a'R'$
et que $a'$ est un non-diviseur de zéro dans $R'$. Posons $S' = S \otimes_S R'$.
Tout élément $g$ de $JS' = \text{Fit}_r(M \otimes_S S')$ peut
s'écrire $g = \sum_\alpha c_\alpha$ avec $c_\alpha \in I'R'$.
Puisque $I'R' = a'R'$, nous pouvons écrire $c_\alpha = a'c'_\alpha$ pour un certain
$c'_\alpha \in R'$ et $g = (\sum c'_\alpha)a' = g' a'$ dans $S'$.
De plus, il existe $g_0 \in J$ tel que $a' = c_\alpha$
pour un certain $\alpha$. Pour cet élément, nous avons $g_0 = g'_0 a'$ dans $S'$,
où $g'_0$ est une unité de $S'$.
Soit $\mathfrak p' \subset R'$
un idéal premier et $\mathfrak q' = \mathfrak p'S'$.
D'après ce qui précède, nous voyons que $JS'_{\mathfrak q'}$ est l'idéal
principal engendré par le non-diviseur de zéro $a'$.
Il résulte du lemme de Compléments d'algèbre
\ref{more-algebra-lemma-principal-fitting-ideal}
que $M'_{\mathfrak q'}$ peut être engendré par $r$ éléments.
Puisque $M'$ est de type fini, il existe $m_1, \ldots, m_r \in M'$ et
$g \in S'$, $g \not \in \mathfrak q'$, tels que le morphisme correspondant
$(S')^{\oplus r} \to M'$ devienne surjectif après inversion de $g$.

\medskip\noindent
Enfin, considérons l'idéal $J' = \text{Fit}_{r - 1}(M')$.
Notons que $J'S'_g$ est engendré par les coefficients des relations entre
$m_1, \ldots, m_r$ (compatibilité de l'idéal de Fitting avec le changement de base).
Il suffit donc de montrer que $J' = 0$, voir le
lemme de Compléments d'algèbre
\ref{more-algebra-lemma-fitting-ideal-finite-locally-free}.
Puisque $R'_a = R_f$ (lemme \ref{algebra-lemma-blowup-in-principal} d'Algèbre)
et $M'_a = M_f$, nous voyons d'après (5)
que l'image de $J'_a$ est nulle dans $S_{\mathfrak q''}$ pour tout idéal premier
$\mathfrak q'' \subset S'$ de la forme $\mathfrak q'' = \mathfrak p''S'$,
où $\mathfrak p'' \subset R'_a$. Puisque
$S'_a \subset \prod_{\mathfrak q''\text{ comme ci-dessus}} S'_{\mathfrak q''}$
(car $(S'_a)_{\mathfrak p''} \subset S'_{\mathfrak q''}$ d'après le
lemme \ref{lemma-base-change-universally-flat}),
nous voyons que $J'R'_a = 0$. Puisque $a$ est un non-diviseur de zéro dans $R'$, nous
concluons que $J' = 0$, ce qui achève la démonstration.
\end{proof}

\begin{lemma}
\label{lemma-flatten-module-pre}
Soit $S$ un schéma quasi-compact et quasi-séparé.\newline
Soit $X \to S$ un morphisme de schémas.
Soit $\mathcal{F}$ un module quasi-cohérent sur $X$.
Soit $U \subset S$ un ouvert quasi-compact. Supposons que
\begin{enumerate}
\item $X \to S$ est affine, de présentation finie, plat, à
fibres géométriquement intègres,
\item $\mathcal{F}$ est un module de type fini,
\item $\mathcal{F}_U$ est de présentation finie,
\item $\mathcal{F}$ est plat sur $S$ en tous les points génériques des
fibres situées au-dessus de points de $U$.
\end{enumerate}
Alors il existe un éclatement $U$-admissible $S' \to S$
et un sous-schéma ouvert $V \subset X_{S'}$
tels que (a) le transformé strict $\mathcal{F}'$ de $\mathcal{F}$
se restreigne en un $\mathcal{O}_V$-module localement libre de rang fini et que
(b) $V \to S'$ soit surjectif.
\end{lemma}

\begin{proof}
Étant donnés $\mathcal{F}/X/S$ et $U \subset S$ sous les hypothèses du
lemme, notons $P$ la propriété « $\mathcal{F}$ est plat sur $S$ en
tous les points génériques des fibres ». Il est clair que $P$ est préservée par passage au
transformé strict, voir le
lemme \ref{divisors-lemma-strict-transform-flat} de Diviseurs
et le lemme \ref{morphisms-lemma-base-change-module-flat} de Morphismes.
Il est également clair que $P$ est locale sur $S$. Par conséquent, toutes les observations
de la remarque \ref{remark-successive-blowups} s'appliquent au problème posé par
le lemme.

\medskip\noindent
Considérons la fonction $r : U \to \mathbf{Z}_{\geq 0}$ qui associe à
$u \in U$ l'entier
$$
r(u) = \dim_{\kappa(\xi_u)}(\mathcal{F}_{\xi_u} \otimes \kappa(\xi_u))
$$
où $\xi_u$ est le point générique de la fibre $X_u$.
D'après le lemme de Compléments sur les morphismes
\ref{more-morphisms-lemma-flat-and-free-at-point-fibre}
et le fait que l'image d'un ouvert de $X_S$ dans $S$ est ouverte,
nous voyons que $r(u)$ est localement constant. Par conséquent,
$U = U_0 \amalg U_1 \amalg \ldots \amalg U_c$ est une réunion disjointe finie
de sous-schémas ouverts et fermés sur lesquels $r$ est constante de valeur
$i$ sur $U_i$. D'après le
lemme \ref{divisors-lemma-separate-disjoint-opens-by-blowing-up} de Diviseurs,
nous pouvons trouver un éclatement $U$-admissible qui décompose $S$ en
la réunion disjointe de deux schémas, le premier contenant $U_0$ et
le second $U_1 \cup \ldots \cup U_c$. En répétant ceci encore $c - 1$
fois, nous pouvons supposer que $S$ est une réunion disjointe
$S = S_0 \amalg S_1 \amalg \ldots \amalg S_c$ avec $U_i \subset S_i$.
Nous pouvons donc supposer que la fonction $r$ définie ci-dessus est constante, disons
de valeur $r$.

\medskip\noindent
D'après la remarque \ref{remark-successive-blowups}, nous voyons que nous pouvons supposer qu'il
existe un diviseur de Cartier effectif $D \subset S$ dont le support est
$S \setminus U$. Une autre application de la remarque \ref{remark-successive-blowups},
combinée avec le
lemme \ref{divisors-lemma-characterize-effective-Cartier-divisor} de Diviseurs,
nous indique que nous pouvons supposer que
$S = \Spec(R)$ et $D = \Spec(R/(f))$ pour un certain non-diviseur de zéro
$f \in R$. Ce cas est traité par le
lemme \ref{lemma-helper-blowup-affine-space}.
\end{proof}

\begin{lemma}
\label{lemma-trick-fitting-ideal}
Soit $A \to C$ un morphisme d'anneaux fini et localement libre de rang $d$.
Soit $h \in C$ un élément tel que $C_h$ soit \'etale sur $A$.
Soit $J \subset C$ un idéal. Posons $I = \text{Fit}_0(C/J)$, où nous
considérons $C/J$ comme un $A$-module de type fini. Alors $IC_h = JJ'$ pour un certain idéal
$J' \subset C_h$. Si $J$ est de type fini, $I$ et $J'$ le sont aussi.
\end{lemma}

\begin{proof}
Nous utiliserons les propriétés élémentaires des idéaux de Fitting, voir le
lemme \ref{more-algebra-lemma-fitting-ideal-basics} de Compléments d'algèbre.
Alors $IC$ est l'idéal de Fitting de $C/J \otimes_A C$.
Notons que $C \to C \otimes_A C$, $c \mapsto 1 \otimes c$, possède une section
(le morphisme de multiplication). Par hypothèse, $C \to C \otimes_A C$ est \'etale
en tout idéal premier de l'image de $\Spec(C_h)$ par cette section.
Par conséquent, le morphisme de multiplication $C \otimes_A C_h \to C_h$ est \'etale
donc en particulier plat, voir le lemme \ref{algebra-lemma-map-between-etale} d'Algèbre.
Il existe donc une $C_h$-algèbre telle que
$C \otimes_A C_h \cong C_h \oplus C'$ comme $C_h$-algèbres, voir le
lemme \ref{algebra-lemma-surjective-flat-finitely-presented} d'Algèbre.
Ainsi $(C/J) \otimes_A C_h \cong (C_h/J_h) \oplus C'/I'$ comme
$C_h$-modules pour un certain idéal $I' \subset C'$.
Par conséquent, $IC_h = JJ'$ avec $J' = \text{Fit}_0(C'/I')$, où nous considérons
$C'/J'$ comme un $C_h$-module.
\end{proof}

\begin{lemma}
\label{lemma-push-ideal}
Soit $A \to B$ un morphisme d'anneaux \'etale. Soit $a \in A$ un non-diviseur de zéro.
Soit $J \subset B$ un idéal de type fini avec $V(J) \subset V(aB)$.
Pour tout idéal premier $\mathfrak q \subset B$, il existe un idéal de type fini
$I \subset A$ avec $V(I) \subset V(a)$ et un élément
$g \in B$, $g \not \in \mathfrak q$, tels que
$IB_g = JJ'$ pour un certain idéal de type fini $J' \subset B_g$.
\end{lemma}

\begin{proof}
Nous pouvons remplacer $B$ par une localisation principale en
un élément $g \in B$, $g \not \in \mathfrak q$. Nous pouvons donc supposer
que $B$ est \'etale standard, voir la
proposition \ref{algebra-proposition-etale-locally-standard} d'Algèbre.
Nous pouvons donc supposer que $B$ est une localisation
de $C = A[x]/(f)$ pour un certain polynôme unitaire $f \in A[x]$ de degré $d$.
Écrivons $B = C_h$ pour un certain $h \in C$. Choisissons des éléments $h_1, \ldots, h_n \in C$
qui engendrent $J$
sur $B$. La condition $V(J) \subset V(aB)$ signifie que
$a^m = \sum b_i h_i$ dans $B$ pour un certain entier $m$ suffisamment grand. Posons $h_{n + 1} = a^m$.
Comme dans le lemme \ref{lemma-trick-fitting-ideal}, prenons
$I = \text{Fit}_0(C/(h_1, \ldots, h_{r + 1}))$.
Puisque le module $C/(h_1, \ldots, h_{r + 1})$
est annulé par $a^m$, nous voyons que $a^{dm} \in I$, ce qui implique
que $V(I) \subset V(a)$.
\end{proof}

\begin{lemma}
\label{lemma-flatten-module-etale-localize}
Soit $S$ un schéma quasi-compact et quasi-séparé.
Soit $X \to S$ un morphisme de schémas.
Soit $\mathcal{F}$ un module quasi-cohérent sur $X$.
Soit $U \subset S$ un ouvert quasi-compact.
Supposons qu'il existe un nombre fini de diagrammes commutatifs
$$
\xymatrix{
& X_i \ar[r]_{j_i} \ar[d] & X \ar[d] \\
S_i^* \ar[r] & S_i \ar[r]^{e_i} & S
}
$$
où
\begin{enumerate}
\item $e_i : S_i \to S$ sont des morphismes \'etales quasi-compacts et
$S = \bigcup e_i(S_i)$,
\item $j_i : X_i \to X$ sont des morphismes \'etales et
$X = \bigcup j_i(X_i)$,
\item $S^*_i \to S_i$ est un éclatement $e_i^{-1}(U)$-admissible
tel que le transformé strict $\mathcal{F}_i^*$ de $j_i^*\mathcal{F}$
soit plate sur $S^*_i$.
\end{enumerate}
Alors il existe un éclatement $U$-admissible $S' \to S$ tel que
le transformé strict de $\mathcal{F}$ soit plat sur $S'$.
\end{lemma}

\begin{proof}
Nous affirmons que les hypothèses du lemme sont préservées par les éclatements $U$-admissibles.
En effet, supposons que $b : S' \to S$ soit un éclatement $U$-admissible le long du
faisceau quasi-cohérent d'idéaux $\mathcal{I}$. De plus, soit $S^*_i \to S_i$
l'éclatement du faisceau quasi-cohérent d'idéaux $\mathcal{J}_i$.
Alors la famille de morphismes $e'_i : S'_i = S_i \times_S S' \to S'$
et $j'_i : X_i' = X_i \times_S S' \to X \times_S S'$ satisfait aux conditions
(1), (2), (3) pour le transformé strict $\mathcal{F}'$ de $\mathcal{F}$
relativement à l'éclatement $S' \to S$. Premièrement, observons que $S_i'$ est
l'éclatement de $S_i$ le long de l'image inverse de $\mathcal{I}$, voir le
lemme \ref{divisors-lemma-flat-base-change-blowing-up} de Diviseurs.
Deuxièmement, considérons l'éclatement
$S_i^{\prime *} \to S_i'$ de $S_i'$ le long de l'image inverse de l'idéal
$\mathcal{J}_i$.
D'après le lemme \ref{divisors-lemma-blowing-up-two-ideals} de Diviseurs,
nous obtenons un diagramme commutatif
$$
\xymatrix{
S_i^{\prime *} \ar[r] \ar[rd] \ar[d] & S'_i \ar[d] \\
S_i^* \ar[r] & S_i
}
$$
et tous les morphismes du diagramme ci-dessus sont des éclatements. Par conséquent, d'après les
lemmes \ref{divisors-lemma-strict-transform-flat} et
\ref{divisors-lemma-strict-transform-composition-blowups} de Diviseurs,
nous voyons
\begin{align*}
& \text{ le transformé strict de }(j'_i)^*\mathcal{F}'\text{ par }
S_i^{\prime *} \to S_i' \\
= &
\text{ le transformé strict de }j_i^*\mathcal{F}\text{ par }
S_i^{\prime *} \to S_i \\
= &
\text{ le transformé strict de }\mathcal{F}_i'\text{ par }
S_i^{\prime *} \to S_i' \\
= &
\text{ l'image inverse de }\mathcal{F}_i^*\text{ par }
X_i \times_{S_i} S_i^{\prime *} \to X_i
\end{align*}
qui est donc plate sur $S_i^{\prime *}$
(lemme \ref{morphisms-lemma-base-change-module-flat} de Morphismes).
Ceci étant dit, nous voyons que toutes les observations de la
remarque \ref{remark-successive-blowups} s'appliquent au
problème consistant à trouver un éclatement $U$-admissible tel que le
transformé strict de $\mathcal{F}$ devienne plat sur la base
sous les hypothèses du lemme. En particulier, nous pouvons supposer
que $S \setminus U$ est le support d'un diviseur de Cartier effectif
$D \subset S$. Une autre application de la remarque \ref{remark-successive-blowups},
combinée avec le
lemme \ref{divisors-lemma-characterize-effective-Cartier-divisor} de Diviseurs,
montre que nous pouvons supposer $S = \Spec(A)$ et $D = \Spec(A/(a))$
pour un certain non-diviseur de zéro $a \in A$.

\medskip\noindent
Choisissons un $i$ et $s \in S_i$. Le lemme \ref{lemma-push-ideal}
implique que nous pouvons trouver un voisinage ouvert $s \in W_i \subset S_i$
et un idéal quasi-cohérent de type fini $\mathcal{I} \subset \mathcal{O}_S$
tel que $\mathcal{I} \cdot \mathcal{O}_{W_i} = \mathcal{J}_i \mathcal{J}'_i$
pour un certain idéal quasi-cohérent de type fini
$\mathcal{J}'_i \subset \mathcal{O}_{W_i}$
et tel que $V(\mathcal{I}) \subset V(a) = S \setminus U$.
Puisque $S_i$ est quasi-compact, nous pouvons remplacer $S_i$ par une famille finie
$W_1, \ldots, W_n$ de tels ouverts et supposer que, pour tout $i$, il existe
un faisceau quasi-cohérent d'idéaux $\mathcal{I}_i \subset \mathcal{O}_S$ tel
que $\mathcal{I}_i \cdot \mathcal{O}_{S_i} = \mathcal{J}_i \mathcal{J}'_i$
pour un certain idéal quasi-cohérent de type fini
$\mathcal{J}'_i \subset \mathcal{O}_{S_i}$.
Comme dans la discussion du premier paragraphe de la démonstration, considérons
l'éclatement $S'$ de $S$ le long du produit $\mathcal{I}_1 \ldots \mathcal{I}_n$
(cet éclatement est $U$-admissible par construction). Le changement de base de $S' \to S$
à $S_i$ est l'éclatement le long de
$$
\mathcal{J}_i \cdot
\mathcal{J}'_i \mathcal{I}_1 \ldots \hat{\mathcal{I}_i} \ldots \mathcal{I}_n
$$
qui se factorise par l'éclatement donné $S_i^* \to S_i$
(lemme \ref{divisors-lemma-blowing-up-two-ideals} de Diviseurs). Dans la notation
du diagramme ci-dessus, cela signifie que $S_i^{\prime *} = S_i'$. Par conséquent,
après avoir remplacé $S$ par $S'$, nous arrivons à la situation où
$j_i^*\mathcal{F}$ est plat sur $S_i$. Par conséquent, $j_i^*\mathcal{F}$ est plat
sur $S$, voir le
lemme \ref{lemma-etale-flat-up-down}.
D'après le lemme \ref{morphisms-lemma-flat-permanence} de Morphismes,
nous voyons que $\mathcal{F}$ est plat sur $S$.
\end{proof}

\begin{theorem}
\label{theorem-flatten-module}
Soit $S$ un schéma quasi-compact et quasi-séparé.
Soit $X$ un schéma sur $S$.
Soit $\mathcal{F}$ un module quasi-cohérent sur $X$.
Soit $U \subset S$ un ouvert quasi-compact. Supposons que
\begin{enumerate}
\item $X$ est quasi-compact,
\item $X$ est localement de présentation finie sur $S$,
\item $\mathcal{F}$ est un module de type fini,
\item $\mathcal{F}_U$ est de présentation finie, et
\item $\mathcal{F}_U$ est plat sur $U$.
\end{enumerate}
Alors il existe un éclatement $U$-admissible $S' \to S$ tel que le
transformé strict $\mathcal{F}'$ de $\mathcal{F}$ soit un
$\mathcal{O}_{X \times_S S'}$-module de présentation finie et
plat sur $S'$.
\end{theorem}

\begin{proof}
Montrons d'abord que nous pouvons trouver un éclatement $U$-admissible tel que le
transformé strict soit plat. La question est locale pour la topologie \'etale sur la source
et sur le but, voir le
lemme \ref{lemma-flatten-module-etale-localize} pour un énoncé précis.
En particulier, nous pouvons supposer que $S = \Spec(R)$ et $X = \Spec(A)$
sont affines. Pour $s \in S$, écrivons $\mathcal{F}_s = \mathcal{F}|_{X_s}$
(image inverse de $\mathcal{F}$ sur la fibre). Comme $X \to S$ est de type fini,
$d = \max_{s \in S} \dim(\text{Supp}(\mathcal{F}_s))$
est un entier. Nous allons raisonner par récurrence sur $d$.

\medskip\noindent
Soit $x \in X$ un point de $X$ situé au-dessus de $s \in S$ avec
$\dim_x(\text{Supp}(\mathcal{F}_s)) = d$.
Appliquons le lemme \ref{lemma-elementary-devissage} pour obtenir
$g : X' \to X$, $e : S' \to S$, $i : Z' \to X'$, et $\pi : Z' \to Y'$.
Observons que $Y' \to S'$ est
un morphisme lisse de schémas affines à fibres géométriquement irréductibles
de dimension $d$. Puisque le problème est local pour la topologie \'etale, il suffit
de démontrer le théorème pour $g^*\mathcal{F}/X'/S'$. Puisque $i : Z' \to X'$
est une immersion fermée de présentation finie (et puisque le transformé strict
commute à l'image directe affine, voir le
lemme \ref{divisors-lemma-strict-transform-affine} de Diviseurs),
il suffit de démontrer le résultat d'aplatissement pour $\mathcal{G}$.
Puisque $\pi$ est fini (donc également affine), il suffit de démontrer le
résultat d'aplatissement pour $\pi_*\mathcal{G}/Y'/S'$. Nous pouvons donc supposer
que $X \to S$ est un morphisme lisse de schémas affines à fibres
géométriquement irréductibles de dimension $d$.

\medskip\noindent
Ensuite, nous appliquons un éclatement comme dans le lemme \ref{lemma-flatten-module-pre}.
Nous arrivons ainsi à la situation où il existe un
ouvert $V \subset X$ surjectif sur $S$ tel que $\mathcal{F}|_V$
soit localement libre de rang fini. Soit $\xi \in X$ le point générique de $X_s$. Soit
$r = \dim_{\kappa(\xi)} \mathcal{F}_\xi \otimes \kappa(\xi)$.
Choisissons un morphisme $\alpha : \mathcal{O}_X^{\oplus r} \to \mathcal{F}$
qui induit un isomorphisme
$\kappa(\xi)^{\oplus r} \to \mathcal{F}_\xi \otimes \kappa(\xi)$.
Puisque $\mathcal{F}$ est localement libre sur $V$, nous trouvons un voisinage ouvert
$W$ de $\xi$ où $\alpha$ est un isomorphisme. Restreignons $S$ à un
voisinage ouvert affine de $s$ tel que $W \to S$ soit surjectif. Disons que $\mathcal{F}$
est le module quasi-cohérent associé au $A$-module $N$. Puisque
$\mathcal{F}$ est plat sur $S$ en tous les points génériques des fibres
(en fait en tous les points de $W$), nous voyons que
$$
\alpha_\mathfrak p : A_\mathfrak p^{\oplus r} \to N_\mathfrak p
$$
est universellement injectif pour tout idéal premier $\mathfrak p$ de $R$, voir le
lemme \ref{lemma-induction-step}. Par conséquent, $\alpha$ est universellement injectif,
voir le lemme \ref{algebra-lemma-universally-injective-check-stalks} d'Algèbre.
Posons $\mathcal{H} = \Coker(\alpha)$.
D'après le lemme \ref{divisors-lemma-strict-transform-universally-injective} de Diviseurs,
nous voyons que, pour tout éclatement $U$-admissible $S' \to S$,
les transformés stricts $\mathcal{F}'$ et $\mathcal{H}'$
s'insèrent dans une suite exacte
$$
0 \to \mathcal{O}_{X \times_S S'}^{\oplus r} \to \mathcal{F}'
\to \mathcal{H}' \to 0
$$
Le lemme \ref{lemma-induction-step} montre donc aussi que $\mathcal{F}'$
est plat en un point $x'$ si et seulement si
$\mathcal{H}'$ est plat en ce point. En particulier, $\mathcal{H}_U$ est
plat sur $U$ et $\mathcal{H}_U$ est un module de présentation finie.
L'hypothèse de récurrence appliquée à $\mathcal{H}$ montre l'existence
d'un éclatement $U$-admissible tel que le transformé strict
$\mathcal{H}'$ soit plate, comme souhaité.

\medskip\noindent
Pour achever la démonstration du théorème, il nous reste à montrer que $\mathcal{F}'$
est un module de présentation finie (après éventuellement un autre
éclatement $U$-admissible). Cela résulte du
lemme \ref{lemma-flat-finite-type-finitely-presented-over-dense-open},
car nous pouvons supposer que $U \subset S$ est schématiquement dense (voir le
troisième paragraphe de la remarque \ref{remark-successive-blowups}).
Ceci achève la démonstration du théorème.
\end{proof}





\section{Applications}
\label{section-applications}

\noindent
Dans cette section, nous appliquons certains des résultats ci-dessus.

\begin{lemma}
\label{lemma-flat-after-blowing-up}
Soit $S$ un schéma quasi-compact et quasi-séparé.
Soit $X$ un schéma sur $S$.
Soit $U \subset S$ un ouvert quasi-compact.
Supposons que
\begin{enumerate}
\item $X \to S$ soit de type fini et quasi-séparé, et que
\item $X_U \to U$ soit plat et localement de présentation finie.
\end{enumerate}
Alors il existe un éclatement $U$-admissible $S' \to S$ tel que
le transformé strict de $X$ soit plat et de présentation finie
sur $S'$.
\end{lemma}

\begin{proof}
Puisque $X \to S$ est quasi-compact et quasi-séparé par hypothèse,
le transformé strict de $X$ relativement à un éclatement $S' \to S$
est lui aussi quasi-compact et quasi-séparé. Ainsi, pour démontrer le lemme,
il suffit de trouver un éclatement $U$-admissible tel que le transformé
strict soit plat et localement de présentation finie.
Soit $X = W_1 \cup \ldots \cup W_n$ un recouvrement ouvert affine fini.
Si nous pouvons trouver un éclatement $U$-admissible $S_i \to S$ tel que le
transformé strict de $W_i$ soit plat et localement de présentation finie,
alors il existe un éclatement $U$-admissible $S' \to S$ dominant
tous les $S_i \to S$ et qui convient (voir
Diviseurs, lemme \ref{divisors-lemma-dominate-admissible-blowups} ;
voir aussi la remarque \ref{remark-successive-blowups}).
Nous pouvons donc supposer que $X$ est affine.

\medskip\noindent
Supposons que $X$ soit affine. D'après
Morphismes, lemme \ref{morphisms-lemma-quasi-affine-finite-type-over-S},
nous pouvons choisir une immersion $j : X \to \mathbf{A}^n_S$ sur $S$.
Soit $V \subset \mathbf{A}^n_S$ un sous-schéma ouvert quasi-compact
tel que $j$ induise une immersion fermée $i : X \to V$ sur $S$. Appliquons
le théorème \ref{theorem-flatten-module}
à $V \to S$ et au module quasi-cohérent $i_*\mathcal{O}_X$
pour obtenir un éclatement $U$-admissible $S' \to S$ tel que le transformé
strict de $i_*\mathcal{O}_X$ soit plat sur $S'$ et de présentation finie
comme $\mathcal{O}_{V \times_S S'}$-module. Soit $X'$ le transformé strict
de $X$ relativement à $S' \to S$. Soit $i' : X' \to V \times_S S'$
le morphisme induit.
Comme la formation du transformé strict commute à l'image directe par les
morphismes affines (Diviseurs, lemme \ref{divisors-lemma-strict-transform-affine}),
nous voyons que $i'_*\mathcal{O}_{X'}$ est plat sur $S$ et de
présentation finie comme $\mathcal{O}_{V \times_S S'}$-module.
Cela implique le lemme.
\end{proof}

\begin{lemma}
\label{lemma-finite-after-blowing-up}
Soit $S$ un schéma quasi-compact et quasi-séparé.
Soit $X$ un schéma sur $S$.
Soit $U \subset S$ un ouvert quasi-compact.
Supposons que
\begin{enumerate}
\item $X \to S$ soit propre, et que
\item $X_U \to U$ soit fini localement libre.
\end{enumerate}
Alors il existe un éclatement $U$-admissible $S' \to S$ tel que
le transformé strict de $X$ soit fini localement libre sur $S'$.
\end{lemma}

\begin{proof}
D'après le lemme \ref{lemma-flat-after-blowing-up}, nous pouvons supposer que
$X \to S$ est plat et de présentation finie. Après avoir remplacé
$S$ par un éclatement $U$-admissible si nécessaire, nous pouvons supposer
que $U \subset S$ est schématiquement dense. Alors $f$ est
fini d'après le lemme \ref{lemma-proper-flat-finite-over-dense-open}.
Ainsi $f$ est fini localement libre d'après
Morphismes, lemme \ref{morphisms-lemma-finite-flat}.
\end{proof}

\begin{lemma}
\label{lemma-zariski-after-blowup}
Soit $\varphi : X \to S$ un morphisme séparé de type fini, avec
$S$ quasi-compact et quasi-séparé. Soit $U \subset S$ un
ouvert quasi-compact tel que $\varphi^{-1}U \to U$ soit un isomorphisme.
Alors il existe un éclatement $U$-admissible $S' \to S$ tel que
le transformé strict $X'$ de $X$ soit isomorphe à un sous-schéma ouvert
de $S'$.
\end{lemma}

\begin{proof}
La discussion de la remarque \ref{remark-successive-blowups} s'applique.
Nous pouvons donc effectuer un premier éclatement $U$-admissible et supposer que le complémentaire
$S \setminus U$ est le support d'un diviseur de Cartier effectif $D$.
En particulier, $U$ est schématiquement dense dans $S$.
Ensuite, nous effectuons un autre éclatement $U$-admissible pour nous ramener au cas où
$X \to S$ est plat et de présentation finie, voir
le lemme \ref{lemma-flat-after-blowing-up}.
Dans ce cas, le résultat découle du lemme \ref{lemma-zariski}.
\end{proof}

\noindent
Le lemme suivant affirme qu'une modification propre peut être dominée
par un éclatement.

\begin{lemma}
\label{lemma-dominate-modification-by-blowup}
Soit $\varphi : X \to S$ un morphisme propre, où
$S$ est quasi-compact et quasi-séparé. Soit $U \subset S$ un
ouvert quasi-compact tel que $\varphi^{-1}U \to U$ soit un isomorphisme.
Alors il existe un éclatement $U$-admissible $S' \to S$
qui domine $X$, c'est-à-dire tel qu'il existe une factorisation
$S' \to X \to S$ du morphisme d'éclatement.
\end{lemma}

\begin{proof}
La discussion de la remarque \ref{remark-successive-blowups} s'applique.
Nous pouvons donc effectuer un premier éclatement $U$-admissible et supposer que le complémentaire
$S \setminus U$ est le support d'un diviseur de Cartier effectif $D$.
En particulier, $U$ est schématiquement dense dans $S$.
Choisissons un autre éclatement $U$-admissible $S' \to S$ tel que le transformé
strict $X'$ de $X$ soit un sous-schéma ouvert de $S'$, voir
le lemme \ref{lemma-zariski-after-blowup}.
Puisque $X' \to S'$ est propre et que
$U \subset S'$ est dense, nous voyons que $X' = S'$. Quelques détails sont omis.
\end{proof}

\begin{lemma}
\label{lemma-get-section-after-blowup}
Soit $S$ un schéma. Soient $U \subset W \subset S$ des sous-schémas ouverts.
Soit $f : X \to W$ un morphisme et soit $s : U \to X$ un
morphisme tel que $f \circ s = \text{id}_U$. Supposons que
\begin{enumerate}
\item $f$ soit propre,
\item $S$ soit quasi-compact et quasi-séparé, et que
\item $U$ et $W$ soient quasi-compacts.
\end{enumerate}
Alors il existe un éclatement $U$-admissible $b : S' \to S$ et un morphisme
$s' : b^{-1}(W) \to X$ prolongeant $s$, avec $f \circ s' = b|_{b^{-1}(W)}$.
\end{lemma}

\begin{proof}
Nous pouvons, et nous allons, remplacer $X$ par l'image schématique de $s$.
Alors $X \to W$ est un isomorphisme au-dessus de $U$, voir
Morphismes, lemme
\ref{morphisms-lemma-scheme-theoretic-image-of-partial-section}.
D'après le lemme \ref{lemma-dominate-modification-by-blowup},
il existe un éclatement $U$-admissible $W' \to W$ et un
prolongement $W' \to X$ de $s$.
Nous achevons la démonstration en appliquant
Diviseurs, lemme \ref{divisors-lemma-extend-admissible-blowups},
afin de prolonger $W' \to W$ en un éclatement $U$-admissible de $S$.
\end{proof}









\section{Compactifications}
\label{section-compactify}

\noindent
Soit $S$ un schéma quasi-compact et quasi-séparé. Nous dirons qu'un
schéma $X$ sur $S$ {\it admet une compactification sur $S$}
ou {\it est compactifiable sur $S$} s'il existe
une immersion ouverte quasi-compacte $X \to \overline{X}$ dans un schéma
$\overline{X}$ propre sur $S$. Si $X$ admet une compactification sur
$S$, alors $X \to S$ est séparé et de type fini. Un théorème de
Nagata, voir \cite{Lutkebohmert}, \cite{Conrad-Nagata}, \cite{Nagata-1},
\cite{Nagata-2}, \cite{Nagata-3} et \cite{Nagata-4},
affirme que la réciproque est également vraie. Nous démontrerons ce théorème
dans la section suivante, voir le théorème \ref{theorem-nagata}.

\medskip\noindent
Soit $S$ un schéma quasi-compact et quasi-séparé.
Soit $X \to S$ un morphisme de schémas séparé et de type fini. La catégorie
des {\it compactifications de $X$ sur $S$} est la catégorie définie
comme suit :
\begin{enumerate}
\item Ses objets sont les immersions ouvertes $j : X \to \overline{X}$ sur $S$ telles que
$\overline{X} \to S$ soit propre.
\item Ses morphismes $(j' : X \to \overline{X}') \to (j : X \to \overline{X})$
sont les morphismes $f : \overline{X}' \to \overline{X}$ de schémas sur $S$
tels que $f \circ j' = j$.
\end{enumerate}
Si $j : X \to \overline{X}$ est une compactification, alors $j$ est une
immersion ouverte quasi-compacte, voir
Schémas, remarque \ref{schemes-remark-quasi-compact-and-quasi-separated}.

\medskip\noindent
{\bf Attention.} Nous ne supposons {\it pas} que les compactifications $j : X \to \overline{X}$
aient une image dense. Par conséquent, si $f : \overline{X}' \to \overline{X}$
est un morphisme de compactifications, il se peut que
$f^{-1}(j(X)) = j'(X)$ ne soit pas vérifiée.

\begin{lemma}
\label{lemma-compactifications-cofiltered}
Soit $S$ un schéma quasi-compact et quasi-séparé.
Soit $X$ un schéma compactifiable sur $S$.
\begin{enumerate}
\item[(a)] La catégorie des compactifications de $X$ sur $S$ est
cofiltrante.
\item[(b)] La sous-catégorie pleine formée des compactifications
$j : X \to \overline{X}$ telles que $j(X)$ soit dense et
schématiquement dense dans $\overline{X}$ est initiale
(Catégories, définition \ref{categories-definition-initial}).
\item[(c)] Si $f : \overline{X}' \to \overline{X}$ est un morphisme
de compactifications de $X$ tel que $j'(X)$ soit dense dans $\overline{X}'$,
alors $f^{-1}(j(X)) = j'(X)$.
\end{enumerate}
\end{lemma}

\begin{proof}
Pour démontrer (a), nous devons vérifier les conditions (1), (2), (3) de
Catégories, définition \ref{categories-definition-codirected}.
La condition (1) est satisfaite précisément parce que nous avons supposé $X$
compactifiable.
Soient $j_i : X \to \overline{X}_i$, $i = 1, 2$, deux compactifications.
Nous pouvons alors considérer l'image schématique $\overline{X}$
de $(j_1, j_2) : X \to \overline{X}_1 \times_S \overline{X}_2$.
Cela détermine une troisième compactification $j : X \to \overline{X}$
qui domine les deux $j_i$ :
$$
\xymatrix{
(X, \overline{X}_1) & (X, \overline{X}) \ar[l] \ar[r] & (X, \overline{X}_2)
}
$$
Ainsi (2) est satisfaite. Soient $f_1, f_2 : \overline{X}_1 \to \overline{X}_2$
deux morphismes entre les compactifications
$j_i : X \to \overline{X}_i$, $i = 1, 2$.
Soit $\overline{X} \subset \overline{X}_1$ l'égalisateur de
$f_1$ et $f_2$. Comme $\overline{X}_2 \to S$ est séparé, nous voyons
que $\overline{X}$ est un sous-schéma fermé de $\overline{X}_1$
et est donc propre sur $S$. De plus, nous obtenons une
immersion ouverte $X \to \overline{X}$ puisque $f_1|_X = f_2|_X = \text{id}_X$.
Le morphisme $(X \to \overline{X}) \to (j_1 : X \to \overline{X}_1)$
donné par l'immersion fermée $\overline{X} \to \overline{X}_1$
égalise $f_1$ et $f_2$, ce qui démontre la condition (3).

\medskip\noindent
Démonstration de (b). Soit $j : X \to \overline{X}$ une compactification.
Si $\overline{X}'$ désigne l'adhérence schématique de $X$
dans $\overline{X}$, alors $X$ est dense et schématiquement dense
dans $\overline{X}'$ d'après
Morphismes, lemme \ref{morphisms-lemma-quasi-compact-immersion}.
Cela démontre la première condition de
Catégories, définition \ref{categories-definition-initial}.
Puisque nous avons déjà montré que la catégorie des compactifications
de $X$ est cofiltrante, la seconde condition de
Catégories, définition \ref{categories-definition-initial},
découle de la première (nous omettons la résolution de cet
exercice catégorique).

\medskip\noindent
Démonstration de (c). Après avoir remplacé $\overline{X}'$ par l'adhérence schématique
de $j'(X)$ (ce qui ne modifie pas l'espace topologique sous-jacent),
cela découle de Morphismes, lemme
\ref{morphisms-lemma-scheme-theoretic-image-of-partial-section}.
\end{proof}

\noindent
Nous pouvons également considérer la catégorie de toutes les compactifications (pour $X$ variable).
Il se trouve que cette catégorie, localisée par rapport à l'ensemble des morphismes
qui induisent un isomorphisme sur l'intérieur,
est équivalente à la catégorie des schémas compactifiables sur $S$.

\begin{lemma}
\label{lemma-compactifyable}
Soit $S$ un schéma quasi-compact et quasi-séparé. Soit $f : X \to Y$
un morphisme de schémas sur $S$, avec $Y$ séparé et de type fini
sur $S$ et $X$ compactifiable sur $S$. Alors $X$ admet une compactification
sur $Y$.
\end{lemma}

\begin{proof}
Soit $j : X \to \overline{X}$ une compactification de $X$ sur $S$.
Nous définissons alors $\overline{X}'$ comme
l'image schématique de $(j, f) : X \to \overline{X} \times_S Y$.
Le morphisme $\overline{X}' \to Y$ est propre, car
$\overline{X} \times_S Y \to Y$ est propre comme changement de base de
$\overline{X} \to S$. D'autre part, puisque $Y$ est séparé
sur $S$, le morphisme $(1, f) : X \to X \times_S Y$ est une
immersion fermée (Schémas, lemme \ref{schemes-lemma-semi-diagonal})
et donc $X \to \overline{X}'$ est une immersion ouverte d'après Morphismes, lemme
\ref{morphisms-lemma-scheme-theoretic-image-of-partial-section}, appliqué
à la ``section partielle'' $s = (j, f)$ de la projection
$\overline{X} \times_S Y \to \overline{X}$.
\end{proof}

\noindent
Soit $S$ un schéma quasi-compact et quasi-séparé.
Nous définissons la {\it catégorie des compactifications} comme la catégorie
dont les objets sont les paires $(X, \overline{X})$ où $\overline{X}$
est un schéma propre sur $S$ et $X \subset \overline{X}$ est un
ouvert quasi-compact, et dont les morphismes
sont les diagrammes commutatifs
$$
\xymatrix{
X \ar[d] \ar[r]_f & Y \ar[d] \\
\overline{X} \ar[r]^{\overline{f}} & \overline{Y}
}
$$
de morphismes de schémas sur $S$.

\begin{lemma}
\label{lemma-right-multiplicative-system}
Soit $S$ un schéma quasi-compact et quasi-séparé.
La famille des morphismes
$(u, \overline{u}) : (X', \overline{X}') \to (X, \overline{X})$
tels que $u$ soit un isomorphisme forme un système multiplicatif à droite
(Catégories, définition \ref{categories-definition-multiplicative-system})
de flèches dans la catégorie des compactifications.
\end{lemma}

\begin{proof}
L'axiome RMS1 se vérifie immédiatement. Vérifions que RMS2 est satisfait.
Supposons donné un diagramme
$$
\xymatrix{
& (X', \overline{X}') \ar[d]_{(u, \overline{u})} \\
(Y, \overline{Y}) \ar[r]^{(f, \overline{f})} & (X, \overline{X})
}
$$
où $u : X' \to X$ est un isomorphisme. Posons alors $Y' = Y \times_X X'$
et notons $v : Y' \to Y$ la projection (qui est un isomorphisme). Posons aussi
$\overline{Y}' = \overline{Y} \times_{\overline{X}} \overline{X}'$
et notons $\overline{v} : \overline{Y}' \to \overline{Y}$ la projection.
Il est clair que $Y' \to \overline{Y}'$ est une immersion ouverte.
Le diagramme
$$
\xymatrix{
(Y', \overline{Y}') \ar[r]_{(g, \overline{g})} \ar[d]_{(v, \overline{v})} &
(X', \overline{X}') \ar[d]_{(u, \overline{u})} \\
(Y, \overline{Y}) \ar[r]^{(f, \overline{f})} & (X, \overline{X})
}
$$
montre que l'axiome RMS2 est satisfait.

\medskip\noindent
Vérifions que RMS3 est satisfait. Supposons donnés deux morphismes
$(f, \overline{f}), (g, \overline{g}) :
(X, \overline{X}) \to (Y, \overline{Y})$
de compactifications et un morphisme
$(v, \overline{v}) : (Y, \overline{Y}) \to (Y', \overline{Y}')$
tel que $v$ soit un isomorphisme et que
$(v, \overline{v}) \circ (f, \overline{f}) =
(v, \overline{v}) \circ (g, \overline{g})$. Alors $f = g$.
Ainsi, si nous notons $\overline{X}' \subset \overline{X}$
l'égalisateur de $\overline{f}$ et $\overline{g}$,
alors $(u, \overline{u}) : (X, \overline{X}') \to (X, \overline{X})$
sera un morphisme de la catégorie des compactifications
tel que $(f, \overline{f}) \circ (u, \overline{u}) =
(g, \overline{g}) \circ (u, \overline{u})$, comme voulu.
\end{proof}

\begin{lemma}
\label{lemma-invert-right-multiplicative-system}
Soit $S$ un schéma quasi-compact et quasi-séparé.\newline
Le foncteur $(X, \overline{X}) \mapsto X$ définit une
équivalence de la catégorie des compactifications localisée
(Catégories, lemme \ref{categories-lemma-right-localization})
par rapport au système
multiplicatif à droite du lemme \ref{lemma-right-multiplicative-system}
vers la catégorie des schémas compactifiables sur $S$.
\end{lemma}

\begin{proof}
Notons $\mathcal{C}$ la catégorie des compactifications et
notons $Q : \mathcal{C} \to \mathcal{C}'$ le foncteur de
localisation de Catégories, lemme
\ref{categories-lemma-properties-right-localization}.
Notons $\mathcal{D}$ la catégorie des schémas compactifiables
sur $S$. Il est clair, d'après le lemme que nous venons de citer et notre
choix du système multiplicatif, que nous
obtenons un foncteur $\mathcal{C}' \to \mathcal{D}$.
Ce foncteur est manifestement essentiellement surjectif.
Si $f : X \to Y$ est un morphisme de schémas
compactifiables, choisissons une immersion ouverte $Y \to \overline{Y}$
dans un schéma propre sur $S$, puis choisissons un plongement
$X \to \overline{X}$ dans un schéma $\overline{X}$ propre sur
$\overline{Y}$ (ce qui est possible d'après le lemme \ref{lemma-compactifyable}
appliqué à $X \to \overline{Y}$). Nous obtenons ainsi un morphisme
$(X, \overline{X}) \to (Y, \overline{Y})$ de compactifications
qui induit le morphisme donné $X \to Y$.
Enfin, supposons donnés deux morphismes de la
catégorie localisée ayant même source et même but, disons
$$
a = ((f, \overline{f}) : (X', \overline{X}') \to (Y, \overline{Y}),
(u, \overline{u}) : (X', \overline{X}') \to (X, \overline{X}))
$$
et
$$
b = ((g, \overline{g}) : (X'', \overline{X}'') \to (Y, \overline{Y}),
(v, \overline{v}) : (X'', \overline{X}'') \to (X, \overline{X}))
$$
qui induisent le même morphisme $X \to Y$ sur $S$, autrement dit
$f \circ u^{-1} = g \circ v^{-1}$. D'après
Catégories, lemme \ref{categories-lemma-morphisms-right-localization},
nous pouvons supposer que $(X', \overline{X}') = (X'', \overline{X}'')$
et que $(u, \overline{u}) = (v, \overline{v})$. Dans ce cas, nous
pouvons considérer l'égalisateur $\overline{X}''' \subset \overline{X}'$
de $\overline{f}$ et $\overline{g}$. Le morphisme
$(w, \overline{w}) : (X', \overline{X}''') \to (X', \overline{X}')$ appartient
au système multiplicatif et nous voyons que $a = b$ dans la catégorie
localisée en précomposant avec $(w, \overline{w})$.
\end{proof}






\section{Compactification de Nagata}
\label{section-compactifications}

\noindent
Dans cette section, nous démontrons le théorème annoncé à la
section \ref{section-compactify}.

\begin{lemma}
\label{lemma-check-separated}
Soit $X \to S$ un morphisme de schémas. Si $X = U \cup V$
est un recouvrement ouvert tel que $U \to S$ et $V \to S$ soient séparés
et que $U \cap V \to U \times_S V$ soit fermé, alors
$X \to S$ est séparé.
\end{lemma}

\begin{proof}
Omis. Indication : vérifier que $\Delta : X \to X \times_S X$ est
fermé en utilisant le recouvrement ouvert de $X \times_S X$ donné par
$U \times_S U$, $U \times_S V$, $V \times_S U$ et $V \times_S V$.
\end{proof}

\begin{lemma}
\label{lemma-separate-disjoint-locally-closed-by-blowing-up}
Soit $X$ un schéma quasi-compact et quasi-séparé.
Soit $U \subset X$ un ouvert quasi-compact.
\begin{enumerate}
\item Si $Z_1, Z_2 \subset X$ sont des sous-schémas fermés de
présentation finie tels que $Z_1 \cap Z_2 \cap U = \emptyset$, alors
il existe un éclatement $U$-admissible $X' \to X$
tel que les transformés stricts de $Z_1$ et $Z_2$ soient disjoints.
\item Si $T_1, T_2 \subset U$ sont des parties disjointes fermées pour la topologie constructible, alors
il existe un éclatement $U$-admissible $X' \to X$ tel que les adhérences de
$T_1$ et $T_2$ soient disjointes.
\end{enumerate}
\end{lemma}

\begin{proof}
Démonstration de (1). L'hypothèse que $Z_i \to X$ est de présentation finie
signifie que le faisceau quasi-cohérent d'idéaux $\mathcal{I}_i$ de $Z_i$
est de type fini, voir 
Morphismes, lemme \ref{morphisms-lemma-closed-immersion-finite-presentation}.
Notons $Z \subset X$ le sous-schéma fermé
découpé par le produit $\mathcal{I}_1 \mathcal{I}_2$.
Observons que $Z \cap U$ est la réunion disjointe
de $Z_1 \cap U$ et $Z_2 \cap U$. D'après Diviseurs, lemme
\ref{divisors-lemma-separate-disjoint-opens-by-blowing-up},
il existe un éclatement $U \cap Z$-admissible $Z' \to Z$ tel que
les transformés stricts de $Z_1$ et $Z_2$ soient disjoints.
Notons $Y \subset Z$ le centre de cet éclatement.
Alors $Y \to X$ est une immersion fermée de présentation finie comme composée
de $Y \to Z$ et $Z \to X$ (Diviseurs, définition
\ref{divisors-definition-admissible-blowup}
et Morphismes, lemme \ref{morphisms-lemma-composition-finite-presentation}).
Ainsi, l'éclatement $X' \to X$ de centre $Y$ est un éclatement $U$-admissible.
D'après les propriétés générales des transformés stricts, les
transformés stricts de $Z_1, Z_2$ relativement à $X' \to X$
sont les mêmes que les transformés stricts de $Z_1, Z_2$ relativement
à $Z' \to Z$, voir
Diviseurs, lemme \ref{divisors-lemma-strict-transform}.
Cela démontre (1).

\medskip\noindent
Démonstration de (2). D'après Propriétés, lemme
\ref{properties-lemma-quasi-coherent-finite-type-ideals},
il existe un faisceau quasi-cohérent d'idéaux de type fini
$\mathcal{J}_i \subset \mathcal{O}_U$ tel que
$T_i = V(\mathcal{J}_i)$ (ensemblistement).
D'après Propriétés, lemme \ref{properties-lemma-extend},
il existe un faisceau quasi-cohérent de type fini
d'idéaux $\mathcal{I}_i \subset \mathcal{O}_X$
dont la restriction à $U$ est $\mathcal{J}_i$.
Appliquons le résultat de la partie (1) aux
sous-schémas fermés $Z_i = V(\mathcal{I}_i)$ pour conclure.
\end{proof}

\begin{lemma}
\label{lemma-blowup-iso-along}
Soit $f : X \to Y$ un morphisme propre de schémas quasi-compacts et
quasi-séparés. Soit $V \subset Y$ un ouvert quasi-compact,
et soit $U = f^{-1}(V)$. Soit $T \subset V$ une partie fermée telle que
$f|_U : U \to V$ soit un isomorphisme au-dessus d'un voisinage ouvert de $T$
dans $V$. Alors il existe un éclatement $V$-admissible $Y' \to Y$
tel que le transformé strict $f' : X' \to Y'$ de $f$
soit un isomorphisme au-dessus d'un voisinage ouvert de l'adhérence
de $T$ dans $Y'$.
\end{lemma}

\begin{proof}
Soit $T' \subset V$ le complémentaire du plus grand ouvert au-dessus duquel
$f|_U$ est un isomorphisme. Alors $T', T$ sont fermés dans $V$ et
$T \cap T' = \emptyset$. Puisque $V$ est un espace topologique
spectral, nous pouvons trouver des parties $T_c, T'_c$ fermées pour la topologie constructible
telles que $T \subset T_c$, $T' \subset T'_c$ et
$T_c \cap T'_c = \emptyset$ (choisir un ouvert quasi-compact
$W$ de $V$ contenant $T'$ et ne rencontrant pas $T$,
et poser $T_c = V \setminus W$, puis choisir un ouvert quasi-compact
$W'$ de $V$ contenant $T_c$ et ne rencontrant pas $T'$,
et poser $T'_c = V \setminus W'$).
D'après le lemme \ref{lemma-separate-disjoint-locally-closed-by-blowing-up},
nous pouvons, après avoir remplacé $Y$ par un éclatement $V$-admissible,
supposer que $T_c$ et $T'_c$ ont des adhérences disjointes dans $Y$.
Posons $Y_0 = Y \setminus \overline{T}'_c$, $V_0 = V \setminus T'_c$,
$U_0 = U \times_V V_0$ et $X_0 = X \times_Y Y_0$.
Puisque $U_0 \to V_0$ est un isomorphisme, nous pouvons trouver un
éclatement $V_0$-admissible $Y'_0 \to Y_0$ tel que le
transformé strict $X'_0$ de $X_0$ s'envoie isomorphiquement sur $Y'_0$, voir
le lemme \ref{lemma-zariski-after-blowup}.
D'après Diviseurs, lemme \ref{divisors-lemma-extend-admissible-blowups},
il existe un éclatement $V$-admissible $Y' \to Y$ dont la restriction
à $Y_0$ est $Y'_0 \to Y_0$. Si $f' : X' \to Y'$ désigne le
transformé strict de $f$, alors nous obtenons le résultat voulu, car
$f'$ se restreint en un isomorphisme au-dessus de $Y'_0$.
\end{proof}

\begin{lemma}
\label{lemma-find-common-blowups}
Soit $S$ un schéma quasi-compact et quasi-séparé.
Soient $U \to X_1$ et $U \to X_2$ des immersions ouvertes
de schémas sur $S$, et supposons que $U$, $X_1$, $X_2$ soient de type
fini et séparés sur $S$. Alors il existe un diagramme commutatif
$$
\xymatrix{
X_1' \ar[d] \ar[r] & X & X_2' \ar[l] \ar[d] \\
X_1 & U \ar[l] \ar[lu] \ar[u] \ar[ru] \ar[r] & X_2
}
$$
de schémas sur $S$, où $X_i' \to X_i$ est un éclatement $U$-admissible,
$X_i' \to X$ est une immersion ouverte et $X$ est séparé et de
type fini sur $S$.
\end{lemma}

\begin{proof}
Tout au long de la démonstration, tous les schémas seront séparés et de type fini sur $S$.
Cela implique en particulier que ces schémas sont quasi-compacts et quasi-séparés,
et que les morphismes entre eux sont quasi-compacts et séparés.
Voir Schémas, sections \ref{schemes-section-quasi-compact} et
\ref{schemes-section-separation-axioms}.
Nous utiliserons le fait que, si $U \to W$ est une immersion de tels schémas sur $S$,
alors l'image schématique $Z$ de $U$ dans $W$ est un sous-schéma fermé
de $W$, que $U \to Z$ est une immersion ouverte, que $U \subset Z$ est
schématiquement dense et que $U \subset Z$ est topologiquement dense. Voir
Morphismes, lemme
\ref{morphisms-lemma-quasi-compact-immersion}.

\medskip\noindent
Soit $X_{12} \subset X_1 \times_S X_2$ l'image schématique
de $U \to X_1 \times_S X_2$. Les projections $p_i : X_{12} \to X_i$
induisent des isomorphismes $p_i^{-1}(U) \to U$ d'après
Morphismes, lemme
\ref{morphisms-lemma-scheme-theoretic-image-of-partial-section}.
Choisissons un éclatement $U$-admissible $X_i^i \to X_i$ tel que
le transformé strict $X_{12}^i$ de $X_{12}$ soit isomorphe à un
sous-schéma ouvert de $X_i^i$, voir
le lemme \ref{lemma-zariski-after-blowup}.
Soit $\mathcal{I}_i \subset \mathcal{O}_{X_i}$ le faisceau quasi-cohérent
d'idéaux de type fini correspondant.
Rappelons que $X_{12}^i \to X_{12}$ est l'éclatement de
$p_i^{-1}\mathcal{I}_i \mathcal{O}_{X_{12}}$, voir
Diviseurs, lemme \ref{divisors-lemma-strict-transform}.
Soit $X_{12}'$ l'éclatement de $X_{12}$ de centre
$p_1^{-1}\mathcal{I}_1 p_2^{-1}\mathcal{I}_2 \mathcal{O}_{X_{12}}$, voir
Diviseurs, lemme \ref{divisors-lemma-blowing-up-two-ideals}
pour la signification de cette construction. Nous obtenons en particulier un diagramme commutatif
$$
\xymatrix{
X_{12}' \ar[d] \ar[r] & X_{12}^2 \ar[d] \\
X_{12}^1 \ar[r] & X_{12}
}
$$
où tous les morphismes sont des éclatements $U$-admissibles.
Puisque $X_{12}^i \subset X_i^i$ est un ouvert, nous pouvons choisir un éclatement $U$-admissible
$X_i' \to X_i^i$ dont la restriction est $X_{12}' \to X_{12}^i$, voir
Diviseurs, lemme \ref{divisors-lemma-extend-admissible-blowups}.
Alors $X_{12}' \subset X_i'$ est un sous-schéma ouvert et le diagramme
$$
\xymatrix{
X_{12}' \ar[d] \ar[r] & X_i' \ar[d] \\
X_{12}^i \ar[r] & X_i^i
}
$$
est commutatif, les flèches verticales étant des éclatements et les flèches horizontales
des immersions ouvertes. Remarquons que $X'_{12} \to X_1' \times_S X_2'$ est
une immersion propre (utiliser que $X'_{12} \to X_{12}$ est propre,
que $X_{12} \to X_1 \times_S X_2$ est fermé, que $X_1' \times_S X_2' \to
X_1 \times_S X_2$ est séparé et appliquer Morphismes, lemme
\ref{morphisms-lemma-image-proper-scheme-closed}).
Ainsi $X'_{12} \to  X_1' \times_S X_2'$ est une immersion fermée.
Il s'ensuit que, si nous définissons $X$ en recollant $X_1'$ et $X_2'$
le long du sous-schéma ouvert commun $X_{12}'$, alors $X \to S$ est de type fini
et séparé (lemme \ref{lemma-check-separated}).
Comme les composées d'éclatements $U$-admissibles sont des éclatements $U$-admissibles
(Diviseurs, lemme \ref{divisors-lemma-composition-admissible-blowups}),
le lemme est démontré.
\end{proof}

\begin{lemma}
\label{lemma-replaced-by-strict-transform}
Soient $X \to S$ et $Y \to S$ des morphismes de schémas.
Soit $U \subset X$ un sous-schéma ouvert.
Soit $V \to X \times_S Y$ un morphisme quasi-compact
dont la composée avec la première projection est à valeurs dans $U$.
Soit $Z \subset X \times_S Y$ l'image schématique de
$V \to X \times_S Y$. Soit $X' \to X$ un éclatement $U$-admissible.
Alors l'image schématique de $V \to X' \times_S Y$ est le
transformé strict de $Z$ relativement à l'éclatement.
\end{lemma}

\begin{proof}
Notons $Z' \to Z$ le transformé strict. Le morphisme $Z' \to X'$
induit un morphisme $Z' \to X' \times_S Y$ qui est une immersion fermée
(car $Z'$ est par définition un sous-schéma fermé de $X' \times_X Z$).
Ainsi, pour achever la démonstration, il suffit de montrer que l'image schématique
$Z''$ de $V \to Z'$ est $Z'$. Observons que $Z'' \subset Z'$
est un sous-schéma fermé tel que $V \to Z'$ se factorise par $Z''$.
Puisque $V \to X \times_S Y$ et $V \to X' \times_S Y$ sont tous deux
quasi-compacts (pour le second, cela découle de Schémas, lemme
\ref{schemes-lemma-quasi-compact-permanence},
et du fait que $X' \times_S Y \to X \times_S Y$ est séparé
comme changement de base d'un morphisme propre), d'après Morphismes, lemme
\ref{morphisms-lemma-quasi-compact-scheme-theoretic-image},
nous voyons que $Z \cap (U \times_S Y) = Z'' \cap (U \times_S Y)$.
Ainsi, le morphisme d'inclusion $Z'' \to Z'$ est un isomorphisme
hors du diviseur exceptionnel $E$ de $Z' \to Z$. Cependant, le
faisceau structural de $Z'$ n'a aucune section non nulle à support
dans $E$ (par définition des transformés stricts), et nous concluons que
la surjection $\mathcal{O}_{Z'} \to \mathcal{O}_{Z''}$
doit être un isomorphisme.
\end{proof}

\begin{lemma}
\label{lemma-compactification-dominates}
Soit $S$ un schéma quasi-compact et quasi-séparé. Soit $U$ un
schéma de type fini et séparé sur $S$. Soit $V \subset U$ un
ouvert quasi-compact. Si $V$ admet une compactification $V \subset Y$
sur $S$, alors il existe un éclatement $V$-admissible $Y' \to Y$ et un
ouvert $V \subset V' \subset Y'$ tel que $V \to U$
se prolonge en un morphisme propre $V' \to U$.
\end{lemma}

\begin{proof}
Considérons l'image schématique $Z \subset Y \times_S U$
du morphisme ``diagonal'' $V \to Y \times_S U$. Si nous remplaçons
$Y$ par un éclatement $V$-admissible, alors $Z$ est remplacé par
le transformé strict relativement à cet éclatement, voir
le lemme \ref{lemma-replaced-by-strict-transform}. Ainsi, d'après
le lemme \ref{lemma-zariski-after-blowup}, nous pouvons supposer que $Z \to Y$
est une immersion ouverte. Si $V' \subset Y$ désigne son image, alors nous
voyons que le morphisme induit $V' \to U$ est propre, car la
projection $Y \times_S U \to U$ est propre et $V' \cong Z$
est un sous-schéma fermé de $Y \times_S U$.
\end{proof}

\noindent
Le lemme suivant n'est formulé que dans le cas noethérien.
La version pour les schémas quasi-compacts et quasi-séparés est
également vraie, mais elle découlera trivialement du théorème principal
de cette section.

\begin{lemma}
\label{lemma-two-compactifications}
Soit $S$ un schéma noethérien. Soit $U$ un schéma de type fini
et séparé sur $S$. Soit $U = U_1 \cup U_2$ avec des ouverts tels que
$U_1$ et $U_2$ admettent des compactifications sur $S$ et que
$U_1 \cap U_2$ soit dense dans $U$. Alors $U$ admet une compactification sur $S$.
\end{lemma}

\begin{proof}
Choisissons une compactification $U_i \subset X_i$ pour $i = 1, 2$. Nous pouvons
supposer que $U_i$ est schématiquement dense dans $X_i$. Nous pouvons supposer qu'il
existe un ouvert $V_i \subset X_i$ et un morphisme propre
$\psi_i : V_i \to U$ prolongeant $\text{id} : U_i \to U_i$, voir
le lemme \ref{lemma-compactification-dominates}. Diagramme :
$$
\xymatrix{
U_i \ar[r] \ar[d] & V_i \ar[r] \ar[dl]^{\psi_i} & X_i \\
U
}
$$
Pour $\{i, j\} = \{1, 2\}$, posons
$Z_i = U \setminus U_j = U_i \setminus (U_1 \cap U_2)$
et
$Z_j = U \setminus U_i = U_j \setminus (U_1 \cap U_2)$.
Ainsi, nous avons
$$
U = U_1 \amalg Z_2 = Z_1 \amalg U_2 = Z_1 \amalg (U_1 \cap U_2) \amalg Z_2
$$
Notons $Z_{i, i} \subset V_i$ l'image réciproque de $Z_i$ par $\psi_i$.
Observons que $\psi_i$ est un isomorphisme au-dessus d'un voisinage ouvert de $Z_i$.
Notons $Z_{i, j} \subset V_i$ l'image réciproque de $Z_j$ par $\psi_i$.
Observons que $\psi_i : Z_{i, j} \to Z_j$ est un morphisme propre.
Puisque $Z_i$ et $Z_j$ sont des parties fermées disjointes de
$U$, nous voyons que $Z_{i, i}$ et $Z_{i, j}$ sont des parties fermées disjointes
de $V_i$.

\medskip\noindent
Notons $\overline{Z}_{i, i}$ et $\overline{Z}_{i, j}$ les adhérences de
$Z_{i, i}$ et $Z_{i, j}$ dans $X_i$. Après avoir remplacé $X_i$ par un
éclatement $V_i$-admissible, nous pouvons supposer que
$\overline{Z}_{i, i}$ et $\overline{Z}_{i, j}$ sont disjointes, voir
le lemme \ref{lemma-separate-disjoint-locally-closed-by-blowing-up}.
Nous supposons que cela vaut à la fois pour $X_1$ et $X_2$.
Observons que cette propriété est préservée si nous remplaçons $X_i$
par un nouvel éclatement $V_i$-admissible.

\medskip\noindent
Posons $V_{12} = V_1 \times_U V_2$. Nous avons une immersion
$V_{12} \to X_1 \times_S X_2$ qui est la composée de l'immersion fermée
$V_{12} = V_1 \times_U V_2 \to V_1 \times_S V_2$
(Schémas, lemme \ref{schemes-lemma-fibre-product-after-map})
et de l'immersion ouverte $V_1 \times_S V_2 \to X_1 \times_S X_2$.
Soit $X_{12} \subset X_1 \times_S X_2$ l'image schématique
de $V_{12} \to X_1 \times_S X_2$. Les morphismes de projection
$$
p_1 : X_{12} \to X_1
\quad\text{et}\quad
p_2 : X_{12} \to X_2
$$
sont propres, car $X_1$ et $X_2$ sont propres sur $S$. Si nous remplaçons $X_1$ par un
éclatement $V_1$-admissible, alors $X_{12}$ est remplacé par
le transformé strict relativement à cet éclatement, voir
le lemme \ref{lemma-replaced-by-strict-transform}.

\medskip\noindent
Notons $\psi : V_{12} \to U$ le morphisme donné par les composées
$\psi = \psi_1 \circ p_1|_{V_{12}} = \psi_2 \circ p_2|_{V_{12}}$.
Considérons le sous-schéma fermé
$$
Z_{12, 2} =
(p_1|_{V_{12}})^{-1}(Z_{1, 2}) =
(p_2|_{V_{12}})^{-1}(Z_{2, 2}) =
\psi^{-1}(Z_2) \subset V_{12}
$$
Le morphisme $p_1|_{V_{12}} : V_{12} \to V_1$ est un isomorphisme
au-dessus d'un voisinage ouvert de $Z_{1, 2}$, car $\psi_2 : V_2 \to U$
est un isomorphisme au-dessus d'un voisinage ouvert de $Z_2$ et
$V_{12} = V_1 \times_U V_2$.
D'après le lemme \ref{lemma-blowup-iso-along},
il existe un éclatement $V_1$-admissible $X_1' \to X_1$
tel que le transformé strict $p'_1 : X'_{12} \to X'_1$
de $p_1$ soit un isomorphisme au-dessus d'un voisinage ouvert de
l'adhérence de $Z_{1, 2}$ dans $X'_1$.
Après avoir remplacé $X_1$ par $X'_1$ et $X_{12}$ par $X'_{12}$,
nous pouvons supposer que $p_1$ est un isomorphisme au-dessus d'un voisinage
ouvert de $\overline{Z}_{1, 2}$.

\medskip\noindent
La réduction du paragraphe précédent nous dit que
$$
X_{12} \cap (\overline{Z}_{1, 2} \times_S \overline{Z}_{2, 1}) = \emptyset
$$
où l'intersection est prise dans $X_1 \times_S X_2$. En effet, l'image
réciproque $p_1^{-1}(\overline{Z}_{1, 2})$ dans $X_{12}$ est envoyée isomorphiquement
sur $\overline{Z}_{1, 2}$. En particulier, nous voyons que $Z_{12, 2}$
est dense dans $p_1^{-1}(\overline{Z}_{1, 2})$. Ainsi $p_2$ envoie
$p_1^{-1}(\overline{Z}_{1, 2})$ dans $\overline{Z}_{2, 2}$.
Puisque $\overline{Z}_{2, 2} \cap \overline{Z}_{2, 1} = \emptyset$,
nous concluons.

\medskip\noindent
Considérons les schémas
$$
W_i = U \coprod\nolimits_{U_i} (X_i \setminus \overline{Z}_{i, j}),
\quad i = 1, 2
$$
obtenus par recollement. Appliquons le lemme \ref{lemma-check-separated}
pour voir que $W_i \to S$ est séparé. Tout d'abord,
$U \to S$ et $X_i \to S$ sont séparés. L'immersion
$U_i \to U \times_S (X_i \setminus \overline{Z}_{i, j})$
est fermée, car toute spécialisation $u_i \leadsto u$
avec $u_i \in U_i$ et $u \in U \setminus U_i$
se relève de manière unique en une spécialisation
$u_i \leadsto v_i$ dans $V_i$ le long du morphisme propre
$\psi_i : V_i \to U$, et alors $v_i$ appartient nécessairement à $Z_{i, j}$.
Ainsi, l'image de l'immersion est fermée, et l'immersion
est donc une immersion fermée.

\medskip\noindent
D'autre part, pour tout anneau de valuation $A$ sur $S$, de corps des fractions $K$,
et tout morphisme $\gamma : \Spec(K) \to (U_1 \cap U_2)$ sur $S$, il
existe un $i$ et un prolongement de $\gamma$ en un morphisme $h_i : \Spec(A) \to W_i$.
En effet, pour chacun des $i = 1, 2$, il existe un morphisme
$g_i : \Spec(A) \to X_i$ prolongeant $\gamma$ d'après le
critère valuatif de propreté de $X_i$ sur $S$, voir
Morphismes, lemme \ref{morphisms-lemma-characterize-proper}.
Ainsi, nous ne rencontrons une difficulté
que si $g_i(\mathfrak m_A) \in \overline{Z}_{i, j}$ pour $i = 1, 2$. Cela est
impossible, car l'intersection de $X_{12}$ et
$\overline{Z}_{1, 2} \times_S \overline{Z}_{2, 1}$ est vide, comme démontré ci-dessus.

\medskip\noindent
Considérons un diagramme
$$
\xymatrix{
W_1' \ar[d] \ar[r] & W & W_2' \ar[l] \ar[d] \\
W_1 & U \ar[l] \ar[lu] \ar[u] \ar[ru] \ar[r] & W_2
}
$$
comme dans le lemme \ref{lemma-find-common-blowups}. D'après le paragraphe précédent,
pour tout diagramme aux flèches pleines
$$
\xymatrix{
\Spec(K) \ar[r]_\gamma  \ar[d] & W \ar[d] \\
\Spec(A) \ar@{..>}[ru] \ar[r] & S
}
$$
où $\Im(\gamma) \subset U_1 \cap U_2$, il existe un $i$ et
un prolongement $h_i : \Spec(A) \to W_i$ de $\gamma$.
En utilisant le critère valuatif de propreté pour $W'_i \to W_i$,
nous pouvons alors relever $h_i$ en $h'_i : \Spec(A) \to W'_i$.
La flèche pointillée du diagramme existe donc. Comme $W$
est séparé sur $S$, nous voyons que cette flèche est également unique.
Cela implique que $W \to S$ est universellement fermé d'après
Morphismes, lemme
\ref{morphisms-lemma-refined-valuative-criterion-universally-closed}.
Comme $W \to S$ est déjà de type fini et séparé, on conclut.
\end{proof}

\begin{theorem}
\label{theorem-nagata}
\begin{reference}
Voir \cite{Lutkebohmert}, \cite{Conrad-Nagata}, \cite{Nagata-1},
\cite{Nagata-2}, \cite{Nagata-3} et \cite{Nagata-4}
\end{reference}
Soit $S$ un schéma quasi-compact et quasi-séparé. Soit
$X \to S$ un morphisme séparé de type fini.
Alors $X$ admet une compactification sur $S$.
\end{theorem}

\begin{proof}
Nous nous ramenons d'abord au cas noethérien. Nous recommandons vivement au lecteur
de sauter ce paragraphe. Il existe une immersion fermée
$X \to X'$ telle que $X' \to S$ soit de présentation finie et séparé.
Voir Limites, proposition
\ref{limits-proposition-separated-closed-in-finite-presentation}.
Si nous trouvons une compactification de $X'$ sur $S$, alors
prendre l'image schématique de $X$ dans celle-ci donnera
une compactification de $X$ sur $S$. Nous pouvons donc supposer que
$X \to S$ est séparé et de présentation finie.
Nous pouvons écrire $S = \lim S_i$ comme limite d'un
système filtrant de schémas noethériens à morphismes de transition affines.
Voir Limites, proposition \ref{limits-proposition-approximate}.
Nous pouvons choisir un $i$ et un morphisme $X_i \to S_i$ de présentation
finie dont le changement de base à $S$ est $X \to S$, voir
Limites, lemme \ref{limits-lemma-descend-finite-presentation}.
Quitte à augmenter $i$, nous pouvons supposer que $X_i \to S_i$ est séparé, voir
Limites, lemme \ref{limits-lemma-descend-separated-finite-presentation}.
Si nous pouvons trouver une compactification de $X_i$ sur $S_i$, alors le
changement de base de celle-ci à $S$ sera une compactification de $X$ sur $S$.
Cela nous ramène au cas traité dans le paragraphe suivant.

\medskip\noindent
Supposons que $S$ soit noethérien. Nous pouvons choisir un recouvrement ouvert affine fini
$X = \bigcup_{i = 1, \ldots, n} U_i$ tel que $U_1 \cap \ldots \cap U_n$
soit dense dans $X$. Cela découle de
Propriétés, lemme \ref{properties-lemma-point-and-maximal-points-affine},
et du fait que $X$ est quasi-compact et possède un nombre fini de
composantes irréductibles. Pour chaque $i$, nous pouvons choisir un $n_i \geq 0$ et une
immersion $U_i \to \mathbf{A}^{n_i}_S$ d'après
Morphismes, lemme \ref{morphisms-lemma-quasi-affine-finite-type-over-S}.
Ainsi, $U_i$ admet une compactification sur $S$ pour $i = 1, \ldots, n$
en prenant l'image schématique dans $\mathbf{P}^{n_i}_S$.
En appliquant le lemme \ref{lemma-two-compactifications}
$(n - 1)$ fois, nous concluons que le théorème est vrai.
\end{proof}






\section{La topologie h}
\label{section-h}

\noindent
Pour nous, en termes informels, un faisceau h est un faisceau pour la topologie de Zariski
qui satisfait à la propriété de faisceau pour les morphismes propres surjectifs
de présentation finie ; voir le lemme \ref{lemma-characterize-sheaf-h}.
Il peut toutefois être utile de souligner que la définition de la topologie h
sur la catégorie des schémas dépend de la référence choisie.

\medskip\noindent
Voevodsky a initialement défini un
recouvrement h comme une famille finie de morphismes de type fini qui sont
conjointement universellement submersifs (Morphismes, définition
\ref{morphisms-definition-submersive}). Voir \cite[Definition 3.1.2]{Voevodsky}.
Cette définition est particulièrement adaptée lorsque la catégorie de schémas sous-jacente est
restreinte aux schémas de type fini sur un schéma de base noethérien fixé.
Dans ce cadre, Voevodsky relie les recouvrements h aux recouvrements ph.
La topologie ph est engendrée par les recouvrements de Zariski et les morphismes propres
surjectifs. Voir Topologies, section \ref{topologies-section-ph}
pour plus de détails.

\medskip\noindent
Dans Topologies, section \ref{topologies-section-V}, nous étudions la topologie V.
Un morphisme quasi-compact $X \to Y$ définit un recouvrement V si toute spécialisation
de points de $Y$ est l'image d'une spécialisation de points de $X$
et si cette propriété subsiste après tout changement de base
(Topologies, lemme \ref{topologies-lemma-refine-qcqs-V}).
Dans ce cas, $X \to Y$ est universellement submersif
(Topologies, lemme \ref{topologies-lemma-V-covering-universally-submersive}).
Il s'avère que la notion de recouvrement V remplace avantageusement les
``familles de morphismes de but fixé qui sont
conjointement universellement submersifs'' lorsqu'on travaille avec des schémas non noethériens.

\medskip\noindent
Notre démarche consistera d'abord à démontrer l'équivalence entre les recouvrements ph et
les recouvrements V pour les familles (éventuellement infinies) de morphismes qui sont
localement de présentation finie. Nous utiliserons ensuite ces familles
comme notion de recouvrement h dans le projet Stacks.
Pour les schémas noethériens et les familles finies, ces recouvrements
coïncident avec ceux de la définition de Voevodsky ; voir le
lemme \ref{lemma-Noetherian-h-covering}.
Sur la catégorie des schémas de présentation finie sur un schéma de base
fixé $S$, quasi-compact et quasi-séparé, ces recouvrements
définissent la même topologie que celle de \cite[Definition 2.7]{Witt-Grass}.

\begin{lemma}
\label{lemma-equivalence-h-v-locally-finite-presentation}
Soit $\{f_i : X_i \to X\}_{i \in I}$ une famille de morphismes
de schémas de but fixé, où $f_i$ est localement de
présentation finie pour tout $i$. Les assertions suivantes sont équivalentes :
\begin{enumerate}
\item $\{X_i \to X\}$ est un recouvrement ph, et
\item $\{X_i \to X\}$ est un recouvrement V.
\end{enumerate}
\end{lemma}

\begin{proof}
Soit $U \subset X$ un ouvert affine. D'après
Topologies, définitions \ref{topologies-definition-ph-covering}
et \ref{topologies-definition-V-covering}, il suffit de montrer que
la famille obtenue par changement de base $\{X_i \times_X U \to U\}$ admet un raffinement
par un recouvrement ph standard si et seulement si elle admet un raffinement par
un recouvrement V standard. On peut donc supposer $X$ affine et il faut montrer que
$\{X_i \to X\}$ admet un raffinement par un recouvrement ph standard
si et seulement si elle admet un raffinement par un recouvrement V standard.
Puisqu'un recouvrement ph standard est un recouvrement V standard, voir
Topologies, lemme \ref{topologies-lemma-standard-ph-standard-V},
il suffit de démontrer la réciproque.

\medskip\noindent
Supposons $X$ affine et supposons que $\{f_i : X_i \to X\}_{i \in I}$
admette un raffinement par un recouvrement V standard
$\{g_j : Y_j \to X\}_{j = 1, \ldots, m}$.
Pour chaque $j$, choisissons un indice $i_j$ et un morphisme
$h_j : Y_j \to X_{i_j}$ tels que $g_j = f_{i_j} \circ h_j$.
Puisque $Y_j$ est affine, donc quasi-compact,
nous pouvons, pour chaque $j$, trouver un nombre fini d'ouverts affines
$U_{j, k} \subset X_{i_j}$ tels que $\Im(h_j) \subset \bigcup U_{j, k}$.
Alors $\{U_{j, k} \to X\}_{j, k}$ raffine $\{X_i \to X\}$
et constitue un recouvrement V standard (car c'est une famille finie de morphismes
entre schémas affines et que sa propriété de relèvement pour les anneaux de valuation
découle de la propriété correspondante de $\{Y_j \to X\}$).
Nous sommes ainsi ramenés au cas traité au paragraphe suivant.

\medskip\noindent
Supposons que $\{f_i : X_i \to X\}_{i = 1, \ldots, n}$
soit un recouvrement V standard, les $f_i$ étant de présentation finie.
Il faut montrer que $\{X_i \to X\}$ admet un raffinement par un recouvrement ph standard.
Choisissons une stratification de platitude générique
$$
X = S \supset S_0 \supset S_1 \supset \ldots \supset S_t = \emptyset
$$
comme dans Compléments sur les morphismes, lemme
\ref{more-morphisms-lemma-generic-flatness-stratification-scheme}
pour le morphisme de présentation finie
$$
\coprod\nolimits_{i = 1, \ldots, n} f_i :
\coprod\nolimits_{i = 1, \ldots, n} X_i
\longrightarrow
X
$$
de schémas affines. Nous utiliserons toutes les propriétés de la stratification
sans plus le préciser. Par construction, le changement de base de chaque $f_i$ à
$U_k = S_k \setminus S_{k + 1}$ est plat.
Notons $Y_k$ l'adhérence schématique de $U_k$ dans $S_k$. Comme
$U_k \to S_k$ est une immersion ouverte quasi-compacte (voir
Propriétés, lemme \ref{properties-lemma-quasi-coherent-finite-type-ideals}),
on voit que $U_k \subset Y_k$ est une immersion ouverte quasi-compacte et dense
(et schématiquement dense) ; voir
Morphismes, lemme \ref{morphisms-lemma-quasi-compact-scheme-theoretic-image}.
Le morphisme $\coprod_{k = 0, \ldots, t - 1} Y_k \to X$
est fini et surjectif ; $\{Y_k \to X\}$ est donc un recouvrement ph standard
et, par conséquent, un recouvrement V standard (voir ci-dessus). Par la propriété de transitivité
des recouvrements V standards
(Topologies, lemme \ref{topologies-lemma-composition-standard-V}),
il suffit de montrer que le tiré en arrière
du recouvrement $\{X_i \to X\}$ sur chaque $Y_k$ admet un raffinement par un
recouvrement V standard. Cela nous ramène au cas décrit dans le
paragraphe suivant.

\medskip\noindent
Supposons que $\{f_i : X_i \to X\}_{i = 1, \ldots, n}$ soit un recouvrement V standard
où chaque $f_i$ est de présentation finie et qu'il existe un ouvert quasi-compact dense
$U \subset X$ tel que $X_i \times_X U \to U$ soit plat.
D'après le théorème \ref{theorem-flatten-module}
il existe un éclatement $U$-admissible $X' \to X$ tel que
le transformé strict $f'_i : X'_i \to X'$ de $f_i$ soit plat.
Observons que le morphisme projectif (donc fermé) $X' \to X$
est surjectif puisque $U \subset X$ est dense et que $U$ s'identifie
à un ouvert de $X'$. Après avoir remplacé $X'$ par un nouvel
éclatement $U$-admissible si nécessaire,
nous pouvons également supposer que $U \subset X'$ est schématiquement dense
(voir la remarque \ref{remark-successive-blowups}).
Ainsi, pour tout point $x \in X'$, il existe un anneau de valuation $V$
et un morphisme $g : \Spec(V) \to X'$ tels que le point générique
de $\Spec(V)$ s'envoie dans $U$ et que le point fermé de
$\Spec(V)$ s'envoie sur $x$, voir Morphismes, lemme
\ref{morphisms-lemma-reach-points-scheme-theoretic-image}.
Puisque $\{X_i \to X\}$ est un recouvrement V standard, nous pouvons choisir
une extension d'anneaux de valuation $V \subset W$, un indice $i$ et un morphisme
$\Spec(W) \to X_i$ tels que le diagramme
$$
\xymatrix{
\Spec(W) \ar[d] \ar[rr] & & X_i \ar[d] \\
\Spec(V) \ar[r] & X' \ar[r] & X
}
$$
soit commutatif. Puisque $X'_i \subset X' \times_X X_i$ est un sous-schéma fermé
contenant l'ouvert $U \times_X X_i$, puisque $\Spec(W)$ est un schéma intègre,
et puisque le morphisme induit $h : \Spec(W) \to X' \times_X X_i$ envoie
le point générique de $\Spec(W)$ dans $U \times_X X_i$, nous en déduisons
que $h$ se factorise par le sous-schéma fermé $X'_i \subset X' \times_X X_i$.
Nous en déduisons que $\{f'_i : X'_i \to X'\}$ est un recouvrement V.
En particulier, $\coprod f'_i$ est surjectif. En particulier,
$\{X'_i \to X'\}$ est un recouvrement fppf. Puisqu'un recouvrement fppf est un recouvrement ph
(Compléments sur les morphismes, lemme \ref{more-morphisms-lemma-fppf-ph}),
il existe un recouvrement ph standard $\{Y_j \to X'\}$ qui raffine
$\{X'_i \to X\}$. Disons que ce recouvrement est donné par un morphisme propre surjectif
$Y \to X'$ et un recouvrement fini par des ouverts affines
$Y = \bigcup Y_j$. Alors la composée $Y \to X$ est propre et surjective,
et nous en déduisons que $\{Y_j \to X\}$ est un recouvrement ph standard.
Cela achève la démonstration.
\end{proof}

\noindent
Voici notre définition.

\begin{definition}
\label{definition-h-covering}
Soit $T$ un schéma. Un {\it recouvrement h de $T$} est une famille de morphismes
$\{f_i : T_i \to T\}_{i \in I}$ telle que chaque $f_i$ soit
localement de présentation finie et que l'une des conditions équivalentes du
lemme \ref{lemma-equivalence-h-v-locally-finite-presentation} soit satisfaite.
\end{definition}

\noindent
Pour les schémas noethériens, cette notion coïncide avec celle de recouvrement ph
(nous le consignons dans le lemme \ref{lemma-Noetherian-h-ph} ci-dessous), et
nous retrouvons la notion de Voevodsky.

\begin{lemma}
\label{lemma-Noetherian-h-covering}
Soit $X$ un schéma noethérien. Soit $\{X_i \to X\}_{i \in I}$
une famille finie de morphismes de type fini. Les assertions suivantes sont équivalentes :
\begin{enumerate}
\item $\coprod_{i \in I} X_i \to X$ est universellement submersif
(Morphismes, définition
\ref{morphisms-definition-submersive}), et
\item $\{X_i \to X\}_{i \in I}$ est un recouvrement h.
\end{enumerate}
\end{lemma}

\begin{proof}
L'implication (2) $\Rightarrow$ (1) découle du résultat plus général
Topologies, lemme \ref{topologies-lemma-V-covering-universally-submersive}
et de notre définition des recouvrements h. Supposons que $\coprod X_i \to X$
soit universellement submersif. Nous allons montrer que $\{X_i \to X\}$
admet un raffinement qui est un recouvrement ph ; cela suffira d'après
Topologies, lemme \ref{topologies-lemma-refine-by-ph} et
notre définition des recouvrements h.
L'argument sera le même que celui employé dans la démonstration du
lemme \ref{lemma-equivalence-h-v-locally-finite-presentation}.

\medskip\noindent
Choisissons une stratification de platitude générique
$$
X = S \supset S_0 \supset S_1 \supset \ldots \supset S_t = \emptyset
$$
comme dans Compléments sur les morphismes, lemme
\ref{more-morphisms-lemma-generic-flatness-stratification-scheme}
pour le morphisme de présentation finie
$$
\coprod\nolimits_{i = 1, \ldots, n} f_i :
\coprod\nolimits_{i = 1, \ldots, n} X_i
\longrightarrow
X
$$
Nous allons utiliser toutes les propriétés de la stratification
sans autre mention. Par construction, le changement de base de chaque
$f_i$ à $U_k = S_k \setminus S_{k + 1}$ est plat.
Notons $Y_k$ l'adhérence schématique de $U_k$ dans $S_k$. Puisque
$U_k \to S_k$ est une immersion ouverte quasi-compacte (tous les schémas de
ce paragraphe sont noethériens),
nous voyons que $U_k \subset Y_k$ est une immersion ouverte quasi-compacte et dense
(et schématiquement dense), voir
Morphismes, lemme \ref{morphisms-lemma-quasi-compact-scheme-theoretic-image}.
Le morphisme $\coprod_{k = 0, \ldots, t - 1} Y_k \to X$
est fini et surjectif, donc $\{Y_k \to X\}$ est un recouvrement ph.
Par la propriété de transitivité des recouvrements ph
(Topologies, lemme \ref{topologies-lemma-ph}),
il suffit de montrer que le recouvrement obtenu par changement de base
de $\{X_i \to X\}$ à chaque $Y_k$ admet un raffinement qui est un
recouvrement ph. Cela nous ramène au cas décrit dans le
paragraphe suivant.

\medskip\noindent
Supposons que $\coprod X_i \to X$ soit universellement submersif et
qu'il existe un ouvert dense $U \subset X$ tel que
$X_i \times_X U \to U$ soit plat pour tout $i$.
D'après le théorème \ref{theorem-flatten-module}
il existe un éclatement $U$-admissible $X' \to X$ tel que
la transformée stricte $f'_i : X'_i \to X'$ de $f_i$ soit plate pour tout $i$.
Observons que le morphisme projectif (donc fermé) $X' \to X$
est surjectif puisque $U \subset X$ est dense et que $U$ s'identifie
à un ouvert de $X'$. En remplaçant au besoin $X'$ par un nouvel
éclatement $U$-admissible, nous pouvons aussi supposer que $U \subset X'$ est dense
(voir la remarque \ref{remark-successive-blowups}).
Ainsi, pour tout point $x \in X'$, il existe un anneau de valuation discrète $A$
et un morphisme $g : \Spec(A) \to X'$ tel que le point
générique de $\Spec(A)$ soit envoyé dans $U$ et que le point fermé de
$\Spec(A)$ soit envoyé sur $x$; voir Limites, lemme
\ref{limits-lemma-reach-point-closure-Noetherian}.
Posons
$$
W = \Spec(A) \times_X \coprod X_i = \coprod \Spec(A) \times_X X_i
$$
Puisque $\coprod X_i \to X$ est universellement submersif,
il existe une spécialisation $w' \leadsto w$ dans $W$
telle que $w'$ ait pour image le point générique de $\Spec(A)$
et que $w$ ait pour image le point fermé de $\Spec(A)$.
(Sinon, la fibre fermée de $W \to \Spec(A)$
serait stable par générisation, donc ouverte, ce qui
contredirait le fait que $W \to \Spec(A)$ est submersif.)
Disons que $w' \in \Spec(A) \times_X X_i$; alors, bien sûr,
$w \in \Spec(A) \times_X X_i$ également. Soit
$x'_i \leadsto x_i$ l'image de $w' \leadsto w$ dans
$X' \times_X X_i$. Puisque $x'_i \in X'_i$ et que
$X'_i \subset X' \times_X X_i$ est un sous-schéma fermé,
nous voyons que $x_i \in X'_i$. Puisque $x_i$ a pour image $x \in X'$,
nous concluons que
$\coprod X'_i \to X'$ est surjectif ! En particulier,
$\{X'_i \to X'\}$ est un recouvrement fppf. Or, tout recouvrement fppf est un recouvrement ph
(Compléments sur les morphismes, lemme \ref{more-morphisms-lemma-fppf-ph}).
Comme $X' \to X$ est propre et surjectif, nous concluons
que $\{X'_i \to X\}$ est un recouvrement ph, ce qui achève la démonstration.
\end{proof}

\begin{lemma}
\label{lemma-Noetherian-h-ph}
Soit $X$ un schéma localement noethérien. Une famille de morphismes
$\{f_i : X_i \to X\}_{i \in I}$ de cible $X$ est un recouvrement h
si et seulement si elle est un recouvrement ph.
\end{lemma}

\begin{proof}
D'après la définition \ref{definition-h-covering}, tout recouvrement h est un recouvrement ph.
Réciproquement, si $\{f_i : X_i \to X\}$ est un recouvrement ph, alors les morphismes
$f_i$ sont localement de type fini (Topologies, définition
\ref{topologies-definition-ph-covering}). Puisque $X$ est localement noethérien,
chaque $f_i$ est localement de présentation finie, et nous voyons qu'il s'agit
d'un recouvrement h par définition.
\end{proof}

\noindent
Le lemme suivant et \cite[Theorem 8.4]{rydh_descent}
montrent que notre définition coïncide avec (ou du moins est
étroitement liée à) celle de l'article \cite{rydh_descent}
de David Rydh. Pour simplifier, nous nous limitons à une base affine.

\begin{lemma}
\label{lemma-approximate-h-cover}
Soit $X$ un schéma affine. Soit $\{X_i \to X\}_{i \in I}$
un recouvrement h. Alors il existe un morphisme propre surjectif
$$
Y \longrightarrow X
$$
de présentation finie (!) et un recouvrement ouvert affine fini
$Y = \bigcup_{j = 1, \ldots, m} Y_j$ tel que
$\{Y_j \to X\}_{j = 1, \ldots, m}$ raffine $\{X_i \to X\}_{i \in I}$.
\end{lemma}

\begin{proof}
Par hypothèse, il existe un morphisme propre surjectif
$Y \to X$ et un recouvrement ouvert affine fini
$Y = \bigcup_{j = 1, \ldots, m} Y_j$ tel que
$\{Y_j \to X\}_{j = 1, \ldots, m}$ raffine $\{X_i \to X\}_{i \in I}$.
Cela signifie que, pour tout $j$, il existe un indice $i_j \in I$
et un morphisme $h_j : Y_j \to X_{i_j}$ au-dessus de $X$.
Voir la définition \ref{definition-h-covering} et
Topologies, définition \ref{topologies-definition-ph-covering}.
Le problème est que nous ne savons pas si $Y \to X$ est de
présentation finie.
D'après
Limites, lemme \ref{limits-lemma-proper-limit-of-proper-finite-presentation},
nous pouvons écrire
$$
Y = \lim Y_\lambda
$$
comme limite d'un système projectif filtrant de schémas $Y_\lambda$ propres et de présentation finie
sur $X$, tels que les morphismes $Y \to Y_\lambda$ et les
morphismes de transition soient des immersions fermées. Observons que
chaque $Y_\lambda \to X$ est surjectif.
D'après Limites, lemme \ref{limits-lemma-descend-opens},
nous pouvons trouver un $\lambda$ et des ouverts quasi-compacts
$Y_{\lambda, j} \subset Y_\lambda$, $j = 1, \ldots, m$,
qui recouvrent $Y_\lambda$ et induisent $Y_j$ sur $Y$.
Alors $Y_j = \lim Y_{\lambda, j}$.
Quitte à augmenter $\lambda$, nous pouvons supposer $Y_{\lambda, j}$
affine pour tout $j$; voir
Limites, lemme \ref{limits-lemma-limit-affine}.
Enfin, puisque $X_i \to X$ est localement de présentation finie,
nous pouvons utiliser la caractérisation fonctorielle des morphismes
localement de présentation finie
(Limites, proposition
\ref{limits-proposition-characterize-locally-finite-presentation})
pour trouver un $\lambda$ tel que, pour tout $j$, il existe
un morphisme $h_{\lambda, j} : Y_{\lambda, j} \to X_{i_j}$
dont la restriction à $Y_j$ est le morphisme $h_j$ choisi ci-dessus.
Ainsi $\{Y_{\lambda, j} \to X\}$ raffine
$\{X_i \to X\}$, ce qui achève la démonstration.
\end{proof}

\noindent
Nous revenons au développement de la théorie générale des recouvrements h.

\begin{lemma}
\label{lemma-zariski-h}
Un recouvrement fppf est un recouvrement h. Par conséquent, les recouvrements syntomiques, lisses, \'etales,
et de Zariski sont également des recouvrements h.
\end{lemma}

\begin{proof}
Cela est vrai car, dans un recouvrement fppf, on exige que les morphismes soient
localement de présentation finie, et car les recouvrements fppf sont des
recouvrements ph, voir Compléments sur les morphismes,
lemme \ref{more-morphisms-lemma-fppf-ph}.
La seconde assertion découle de la première et de
Topologies, lemme \ref{topologies-lemma-zariski-etale-smooth-syntomic-fppf}.
\end{proof}

\begin{lemma}
\label{lemma-surjective-proper-finite-presentation-h}
Soit $f : Y \to X$ un morphisme surjectif et propre de schémas
qui est de présentation finie. Alors $\{Y \to X\}$ est un recouvrement h.
\end{lemma}

\begin{proof}
Combiner les lemmes de Topologies
\ref{topologies-lemma-zariski-etale-smooth-syntomic-fppf-fpqc-ph-V} et
\ref{topologies-lemma-surjective-proper-ph}.
\end{proof}

\begin{lemma}
\label{lemma-refine-by-h}
Soit $T$ un schéma. Soit $\{f_i : T_i \to T\}_{i \in I}$ une famille
de morphismes telle que $f_i$ soit localement de présentation finie pour tout $i$.
Les assertions suivantes sont équivalentes :
\begin{enumerate}
\item $\{T_i \to T\}_{i \in I}$ est un recouvrement h,
\item il existe un recouvrement h qui raffine $\{T_i \to T\}_{i \in I}$, et
\item $\{\coprod_{i \in I} T_i \to T\}$ est un recouvrement h.
\end{enumerate}
\end{lemma}

\begin{proof}
Cela découle de l'énoncé analogue pour les recouvrements ph
(Topologies, lemme \ref{topologies-lemma-refine-by-ph})
ou de l'énoncé analogue pour les recouvrements V
(Topologies, lemme \ref{topologies-lemma-refine-by-V}).
\end{proof}

\noindent
Ensuite, nous montrons que notre notion de recouvrement h satisfait aux conditions de
Sites, définition \ref{sites-definition-site}.

\begin{lemma}
\label{lemma-h}
Soit $T$ un schéma.
\begin{enumerate}
\item Si $T' \to T$ est un isomorphisme, alors $\{T' \to T\}$
est un recouvrement h de $T$.
\item Si $\{T_i \to T\}_{i\in I}$ est un recouvrement h et si, pour chaque
$i$, nous avons un recouvrement h $\{T_{ij} \to T_i\}_{j\in J_i}$, alors
$\{T_{ij} \to T\}_{i \in I, j\in J_i}$ est un recouvrement h.
\item Si $\{T_i \to T\}_{i\in I}$ est un recouvrement h
et si $T' \to T$ est un morphisme de schémas, alors
$\{T' \times_T T_i \to T'\}_{i\in I}$ est un recouvrement h.
\end{enumerate}
\end{lemma}

\begin{proof}
Cela découle immédiatement de l'énoncé correspondant pour
les recouvrements ph ou les recouvrements V
(Topologies, lemme \ref{topologies-lemma-ph} ou
\ref{topologies-lemma-V})
et du fait que la classe des morphismes localement
de présentation finie est stable par changement de base et par
composition.
\end{proof}

\noindent
Ensuite, nous définissons les grands sites h avec lesquels nous travaillerons dans le
projet Stacks. Il convient de lire la discussion générale
dans Topologies, section \ref{topologies-section-procedure}
avant de poursuivre.

\begin{definition}
\label{definition-big-h-site}
Un {\it grand site h} est un site $\Sch_h$ quelconque comme dans
Sites, définition \ref{sites-definition-site}, construit comme suit :
\begin{enumerate}
\item Choisir un ensemble quelconque de schémas $S_0$ et un ensemble quelconque de recouvrements h
$\text{Cov}_0$ parmi ces schémas.
\item Prendre comme catégorie sous-jacente une catégorie quelconque\newline $\Sch_\alpha$
construite comme dans Ensembles, lemme \ref{sets-lemma-construct-category}
à partir de l'ensemble $S_0$.
\item Choisir un ensemble quelconque de recouvrements comme dans
Ensembles, lemme \ref{sets-lemma-coverings-site}, à partir de la
catégorie $\Sch_\alpha$, de la classe des recouvrements h
et de l'ensemble $\text{Cov}_0$ choisi ci-dessus.
\end{enumerate}
\end{definition}

\noindent
Voir les remarques qui suivent
Topologies, définition \ref{topologies-definition-big-zariski-site},
pour la motivation et les explications concernant la définition des grands sites.

\begin{definition}
\label{definition-standard-h}
Soit $T$ un schéma affine. Un {\it recouvrement h standard} de $T$
est une famille $\{f_i : T_i \to T\}_{i = 1, \ldots, n}$ dans laquelle chaque $T_i$ est
affine et chaque $f_i$ est de présentation finie, et qui satisfait à l'une des
conditions équivalentes suivantes : (1) $\{T_i \to T\}$ admet un raffinement par
un recouvrement ph standard ou (2) $\{T_i \to T\}$ est un recouvrement V.
\end{definition}

\noindent
L'équivalence de ces conditions résulte du
lemme \ref{lemma-equivalence-h-v-locally-finite-presentation}, de
Topologies, définition \ref{topologies-definition-ph-covering}, et du
lemme \ref{topologies-lemma-refine-by-ph}.

\medskip\noindent
Avant de poursuivre l'introduction du grand site h d'un
schéma $S$, signalons que la topologie sur un grand site h
$\Sch_h$ est, en un certain sens, induite par la topologie h
sur la catégorie de tous les schémas.

\begin{lemma}
\label{lemma-h-induced}
Soit $\Sch_h$ un grand site h comme dans la
définition \ref{definition-big-h-site}.
Soit $T \in \Ob(\Sch_h)$.
Soit $\{T_i \to T\}_{i \in I}$ un recouvrement h arbitraire de $T$.
\begin{enumerate}
\item Il existe un recouvrement $\{U_j \to T\}_{j \in J}$ de $T$ dans le site
$\Sch_h$ qui raffine $\{T_i \to T\}_{i \in I}$.
\item Si $\{T_i \to T\}_{i \in I}$ est un recouvrement h standard, alors
il est tautologiquement équivalent à un recouvrement de $\Sch_h$.
\item Si $\{T_i \to T\}_{i \in I}$ est un recouvrement de Zariski, alors
il est tautologiquement équivalent à un recouvrement de $\Sch_h$.
\end{enumerate}
\end{lemma}

\begin{proof}
Omis. Indication : il s'agit exactement de la même démonstration que celle de
Topologies, lemme \ref{topologies-lemma-ph-induced}.
\end{proof}

\begin{definition}
\label{definition-big-small-h}
Soit $S$ un schéma. Soit $\Sch_h$ un grand site h contenant $S$.
\begin{enumerate}
\item Le {\it grand site h de $S$}, noté
$(\Sch/S)_h$, est le site $\Sch_h/S$
introduit dans Sites, section \ref{sites-section-localize}.
\item Le {\it grand site h affine de $S$}, noté
$(\textit{Aff}/S)_h$, est la sous-catégorie pleine de
$(\Sch/S)_h$ dont les objets sont les $U/S$ affines.
Un recouvrement de $(\textit{Aff}/S)_h$ est tout recouvrement
$\{U_i \to U\}$ de $(\Sch/S)_h$ qui est un recouvrement h standard.
\end{enumerate}
\end{definition}

\noindent
Précisons explicitement que le grand site h affine est un site.

\begin{lemma}
\label{lemma-verify-site-h}
Soit $S$ un schéma. Soit $\Sch_h$ un grand
site h contenant $S$. Alors $(\textit{Aff}/S)_h$ est un site.
\end{lemma}

\begin{proof}
En raisonnant comme dans la démonstration de\newline
Topologies, lemme \ref{topologies-lemma-verify-site-etale},
il suffit de montrer que la collection des recouvrements h standards
satisfait aux propriétés (1), (2) et (3) de
Sites, définition \ref{sites-definition-site}.
Cela est clair puisque, par exemple, étant donné un recouvrement h standard
$\{T_i \to T\}_{i\in I}$ et, pour chaque
$i$, un recouvrement h standard $\{T_{ij} \to T_i\}_{j \in J_i}$, alors
$\{T_{ij} \to T\}_{i \in I, j\in J_i}$ est un recouvrement h
(lemme \ref{lemma-h}), $\bigcup_{i\in I} J_i$ est fini et
chaque $T_{ij}$ est affine. Ainsi $\{T_{ij} \to T\}_{i \in I, j\in J_i}$
est un recouvrement h standard.
\end{proof}

\begin{lemma}
\label{lemma-fibre-products-h}
Soit $S$ un schéma. Soit $\Sch_h$ un grand
site h contenant $S$. Les catégories sous-jacentes aux sites
$\Sch_h$, $(\Sch/S)_h$ et $(\textit{Aff}/S)_h$ possèdent des produits fibrés.
Dans chaque cas, le foncteur évident vers la catégorie $\Sch$ de
tous les schémas commute aux produits fibrés. La catégorie
$(\Sch/S)_h$ possède un objet final, à savoir $S/S$.
\end{lemma}

\begin{proof}
Pour $\Sch_h$, cela est vrai par construction, voir
Ensembles, lemme \ref{sets-lemma-what-is-in-it}.
Supposons donnés des morphismes de schémas $U \to S$, $V \to U$, $W \to U$
avec $U, V, W \in \Ob(\Sch_h)$.
Le produit fibré $V \times_U W$ dans $\Sch_h$
est un produit fibré dans $\Sch$ et
est le produit fibré de $V/S$ et de $W/S$ au-dessus de $U/S$ dans
la catégorie de tous les schémas au-dessus de $S$, et donc également un
produit fibré dans $(\Sch/S)_h$.
Cela démontre le résultat pour $(\Sch/S)_h$.
Si $U, V, W$ sont affines, alors $V \times_U W$ l'est aussi, d'où le
résultat pour $(\textit{Aff}/S)_h$.
\end{proof}

\noindent
Ensuite, vérifions que le grand site affine définit le même
topos que le grand site.

\begin{lemma}
\label{lemma-affine-big-site-h}
Soit $S$ un schéma. Soit $\Sch_h$ un grand
site h contenant $S$.
Le foncteur $(\textit{Aff}/S)_h \to (\Sch/S)_h$
est cocontinu et induit une équivalence de topos de
$\Sh((\textit{Aff}/S)_h)$ vers
$\Sh((\Sch/S)_h)$.
\end{lemma}

\begin{proof}
La notion de foncteur cocontinu spécial est introduite dans
Sites, définition \ref{sites-definition-special-cocontinuous-functor}.
Nous devons donc vérifier les hypothèses (1) -- (5) de
Sites, lemme \ref{sites-lemma-equivalence}.
Notons le foncteur d'inclusion
$u : (\textit{Aff}/S)_h \to (\Sch/S)_h$.
La cocontinuité résulte du fait que tout recouvrement h de
$T/S$, avec $T$ affine, peut être raffiné par un recouvrement h standard,
par exemple d'après le lemme \ref{lemma-approximate-h-cover}. Ainsi (1) est satisfaite.
Nous voyons que $u$ est continu simplement parce qu'un recouvrement h standard
est un recouvrement h.
Ainsi, (2) est vérifiée. Les assertions (3) et (4) découlent immédiatement du fait
que $u$ est pleinement fidèle. Enfin, la condition (5) découle du
fait que tout schéma possède un recouvrement ouvert affine (qui est
un recouvrement h).
\end{proof}

\begin{lemma}
\label{lemma-characterize-sheaf-h}
Soit $\mathcal{F}$ un préfaisceau sur $(\Sch/S)_h$.
Alors $\mathcal{F}$ est un faisceau si et seulement si
\begin{enumerate}
\item $\mathcal{F}$ satisfait la condition de faisceau pour les
recouvrements de Zariski, et
\item si $f : V \to U$ est propre, surjectif et de présentation finie, alors
l'application de $\mathcal{F}(U)$ vers l'égalisateur
des deux applications $\mathcal{F}(V) \to \mathcal{F}(V \times_U V)$ est bijective.
\end{enumerate}
De plus, en présence de (1), la propriété (2) équivaut à la
propriété
\begin{enumerate}
\item[(2')] la propriété de faisceau pour $\{V \to U\}$ comme en (2), avec $U$ affine.
\end{enumerate}
\end{lemma}

\begin{proof}
Nous allons montrer que, si (1) et (2) sont vérifiées, alors $\mathcal{F}$ est un faisceau.
Soit $\{T_i \to T\}$ un recouvrement dans $(\Sch/S)_h$.
Vérifions la condition de faisceau pour ce recouvrement.
Soient $s_i \in \mathcal{F}(T_i)$ des sections dont les restrictions donnent la même
section sur $T_i \times_T T_{i'}$. Montrons qu'il existe une
unique section $s \in \mathcal{F}(T)$ dont la restriction à $T_i$ est $s_i$.
Soit $T = \bigcup U_j$ un recouvrement ouvert affine.
D'après la propriété (1), il suffit de construire des sections $s_j \in \mathcal{F}(U_j)$
qui coïncident sur $U_j \cap U_{j'}$ pour obtenir $s$.
Considérons les recouvrements $\{T_i \times_T U_j \to U_j\}$.
Alors les $s_{ji} = s_i|_{T_i \times_T U_j}$ sont des sections qui coïncident
sur $(T_i \times_T U_j) \times_{U_j} (T_{i'} \times_T U_j)$.
Choisissons un morphisme propre surjectif $V_j \to U_j$ de présentation finie
et un recouvrement ouvert affine fini $V_j = \bigcup V_{jk}$
tel que $\{V_{jk} \to U_j\}$ raffine $\{T_i \times_T U_j \to U_j\}$.
Voir le lemme \ref{lemma-approximate-h-cover}.
Si $s_{jk} \in \mathcal{F}(V_{jk})$
désigne le tiré en arrière de $s_{ji}$ sur $V_{jk}$ par les
morphismes sous-entendus, alors les $s_{jk}$ se recollent en une section
$s'_j \in \mathcal{F}(V_j)$. En utilisant encore la coïncidence sur les intersections,
nous voyons que $s'_j$ appartient à l'égalisateur des deux
applications $\mathcal{F}(V_j) \to \mathcal{F}(V_j \times_{U_j} V_j)$.
Ainsi, d'après (2), $s'_j$ provient d'une unique section
$s_j \in \mathcal{F}(U_j)$. Nous omettons de vérifier que ces
sections $s_j$ possèdent toutes les propriétés voulues.

\medskip\noindent
Démonstration de l'équivalence de (2) et (2') en présence de (1).
Supposons que $V \to U$ soit un morphisme de $(\Sch/S)_h$,
propre, surjectif et de présentation finie. Choisissons un
recouvrement ouvert affine $U = \bigcup U_i$ et posons $V_i = V \times_U U_i$.
Alors $\mathcal{F}(U) \to \mathcal{F}(V)$
est injective, car $\mathcal{F}(U_i) \to \mathcal{F}(V_i)$
est injective d'après (2') et $\mathcal{F}(U) \to \prod \mathcal{F}(U_i)$
est injective d'après (1). Enfin, supposons que l'on se donne un élément
$t \in \mathcal{F}(V)$ appartenant à l'égalisateur des deux applications
$\mathcal{F}(V) \to \mathcal{F}(V \times_U V)$.
Alors $t|_{V_i}$ appartient à l'égalisateur des deux applications
$\mathcal{F}(V_i) \to \mathcal{F}(V_i \times_{U_i} V_i)$
pour tout $i$. Nous obtenons donc une unique section $s_i \in \mathcal{F}(U_i)$
qui est envoyée sur $t|_{V_i}$ pour tout $i$, d'après (2').
Nous omettons de vérifier que $s_i|_{U_i \cap U_j} = s_j|_{U_i \cap U_j}$
pour tous $i, j$; cela utilise la propriété d'unicité que nous venons d'établir.
La propriété de faisceau pour le recouvrement $U = \bigcup U_i$ donne
une section $s \in \mathcal{F}(U)$. Nous omettons de démontrer que $s$
a pour image $t$ dans $\mathcal{F}(V)$.
\end{proof}

\noindent
Établissons maintenant quelques relations entre les topos
associés à ces sites.

\begin{lemma}
\label{lemma-morphism-big-h}
Soit $\Sch_h$ un grand site h.
Soit $f : T \to S$ un morphisme dans $\Sch_h$.
Le foncteur
$$
u : (\Sch/T)_h \longrightarrow (\Sch/S)_h,
\quad
V/T \longmapsto V/S
$$
est cocontinu et admet un adjoint à droite continu
$$
v : (\Sch/S)_h \longrightarrow (\Sch/T)_h,
\quad
(U \to S) \longmapsto (U \times_S T \to T).
$$
Ces deux foncteurs induisent le même morphisme de topos
$$
f_{big} :
\Sh((\Sch/T)_h)
\longrightarrow
\Sh((\Sch/S)_h)
$$
Nous avons $f_{big}^{-1}(\mathcal{G})(U/T) = \mathcal{G}(U/S)$.
Nous avons $f_{big, *}(\mathcal{F})(U/S) = \mathcal{F}(U \times_S T/T)$.
De plus, $f_{big}^{-1}$ admet un adjoint à gauche $f_{big!}$ qui commute aux
produits fibrés et aux égalisateurs.
\end{lemma}

\begin{proof}
Le foncteur $u$ est cocontinu, continu et commute aux produits fibrés
et aux égalisateurs. Par conséquent,
Sites, lemmes \ref{sites-lemma-when-shriek} et
\ref{sites-lemma-preserve-equalizers}
s'appliquent, et nous en déduisons la formule
donnant $f_{big}^{-1}$ ainsi que l'existence de $f_{big!}$. De plus,
le foncteur $v$ est adjoint à droite car, étant donnés $U/T$ et $V/S$,
nous avons $\Mor_S(u(U), V) = \Mor_T(U, V \times_S T)$,
comme voulu. Nous pouvons donc appliquer
Sites, lemmes \ref{sites-lemma-have-functor-other-way} et
\ref{sites-lemma-have-functor-other-way-morphism} pour obtenir la
formule pour $f_{big, *}$.
\end{proof}

\begin{lemma}
\label{lemma-composition-h}
Étant donnés des schémas $X$, $Y$, $Y$ dans $(\Sch/S)_h$
et des morphismes $f : X \to Y$, $g : Y \to Z$, nous avons
$g_{big} \circ f_{big} = (g \circ f)_{big}$.
\end{lemma}

\begin{proof}
Cela résulte de la description élémentaire des images directe
et inverse pour les foncteurs entre les grands sites donnée dans le
lemme \ref{lemma-morphism-big-h}.
\end{proof}










\section{Compléments sur la topologie h}
\label{section-h-more}

\noindent
Dans cette section, nous démontrons quelques résultats supplémentaires sur la topologie h.
Commençons par quelques contre-exemples.

\begin{example}
\label{example-structure-sheaf-not-h-sheaf}
Le ``faisceau structural'' $\mathcal{O}$ n'est pas un faisceau pour la topologie h.
Considérons par exemple une immersion fermée surjective de
présentation finie $X \to Y$. Alors $\{X \to Y\}$ est un recouvrement h,
par exemple d'après le lemme \ref{lemma-surjective-proper-finite-presentation-h}.
De plus, remarquons que $X \times_Y X = X$.
Ainsi, si $\mathcal{O}$ était un faisceau pour la topologie h, alors
$\mathcal{O}_Y(Y) \to \mathcal{O}_X(X)$ serait bijective.
Ce n'est pas le cas dès que $X$ et $Y$ sont affines et que le morphisme
$X \to Y$ n'est pas un isomorphisme.
\end{example}

\begin{example}
\label{example-representable-sheaf-not-h-sheaf}
Sur chacun des sites $(\Sch/S)_h$,\newline la topologie n'est pas sous-canonique,
autrement dit, les préfaisceaux représentables ne sont pas tous des faisceaux.
En effet, le ``faisceau structural'' $\mathcal{O}$ est représentable
puisque $\mathcal{O}(X) = \Mor_S(X, \mathbf{A}^1_S)$
dans $(\Sch/S)_h$, et nous avons vu dans
l'exemple \ref{example-structure-sheaf-not-h-sheaf}
que $\mathcal{O}$ n'est pas un faisceau.
\end{example}

\begin{lemma}
\label{lemma-limit-h-topology}
Soit $T$ un schéma affine qui s'écrit comme une limite
$T = \lim_{i \in I} T_i$ d'un système projectif filtrant de schémas affines.
\begin{enumerate}
\item Soit $\mathcal{V} = \{V_j \to T\}_{j = 1, \ldots, m}$ un
recouvrement h standard de $T$, voir la définition
\ref{definition-standard-h}.
Il existe alors un indice $i$ et un recouvrement h standard
$\mathcal{V}_i = \{V_{i, j} \to T_i\}_{j = 1, \ldots, m}$
dont le changement de base $T \times_{T_i} \mathcal{V}_i$ à $T$
est isomorphe à $\mathcal{V}$.
\item Soient $\mathcal{V}_i$, $\mathcal{V}'_i$ deux recouvrements h standards
de $T_i$. Si
$f : T \times_{T_i} \mathcal{V}_i \to T \times_{T_i} \mathcal{V}'_i$ est
un morphisme de recouvrements de $T$, alors il existe un indice
$i' \geq i$ et un morphisme
$f_{i'} : T_{i'} \times_{T_i} \mathcal{V} \to
T_{i'} \times_{T_i} \mathcal{V}'_i$
dont le changement de base à $T$ est $f$.
\item Si
$f, g : \mathcal{V} \to \mathcal{V}'_i$
sont des morphismes de recouvrements h standards de $T_i$ dont les
changements de base $f_T, g_T$ à $T$ sont égaux, alors il existe un
indice $i' \geq i$ tel que $f_{T_{i'}} = g_{T_{i'}}$.
\end{enumerate}
Autrement dit, la catégorie des recouvrements h standards de $T$ est
la colimite sur $I$ des catégories de recouvrements h standards de $T_i$.
\end{lemma}

\begin{proof}
D'après Limites, lemme \ref{limits-lemma-descend-finite-presentation},
la catégorie des schémas de présentation finie sur $T$ est la
colimite sur $I$ des catégories de schémas de présentation finie sur $T_i$. D'après
Limites, lemme \ref{limits-lemma-descend-affine-finite-presentation},
il en va de même pour la catégorie des schémas qui sont affines et
de présentation finie sur $T$.
Pour achever la démonstration du lemme, il suffit de montrer que si
$\{V_{j, i} \to T_i\}_{j = 1, \ldots, m}$ est une famille finie de
morphismes de présentation finie, les $V_{j, i}$ étant affines, et que la
famille obtenue par changement de base $\{T \times_{T_i} V_{j, i} \to T\}$ est
un recouvrement h, alors, pour un certain $i' \geq i$, la famille
$\{T_{i'} \times_{T_i} V_{j, i} \to T_{i'}\}$ est un recouvrement h.
Pour le voir, appliquons le lemme \ref{lemma-approximate-h-cover} afin de
choisir un morphisme de présentation finie, propre et surjectif
$Y \to T$ ainsi qu'un recouvrement fini par des ouverts affines
$Y = \bigcup_{k = 1, \ldots, n} Y_k$ de sorte que
$\{Y_k \to T\}_{k = 1, \ldots, n}$ raffine
$\{T \times_{T_i} V_{j, i} \to T\}$.
En utilisant les arguments ci-dessus ainsi que
Limites, lemmes \ref{limits-lemma-eventually-proper},
\ref{limits-lemma-descend-surjective}, et
\ref{limits-lemma-descend-opens},
nous pouvons trouver un $i' \geq i$ et un morphisme de présentation finie, surjectif et propre
$Y_{i'} \to T_{i'}$ ainsi qu'un recouvrement par des ouverts affines
$Y_{i'} = \bigcup_{k = 1, \ldots, n} Y_{i', k}$
de sorte que, de plus, $\{Y_{i', k} \to Y_{i'}\}$ raffine
$\{T_{i'} \times_{T_i} V_{j, i} \to T_{i'}\}$.
Il s'ensuit que cette dernière famille est
un recouvrement h, ce qui achève la démonstration.
\end{proof}

\begin{lemma}
\label{lemma-extend-sheaf-h}
Soit $S$ un schéma contenu dans un grand site $\Sch_h$.
Soit $F : (\Sch/S)_h^{opp} \to \textit{Sets}$ un faisceau h satisfaisant à
la propriété (b) de Topologies, lemme \ref{topologies-lemma-extend},
avec $\mathcal{C} = (\Sch/S)_h$.
Alors l'extension $F'$ de $F$ à la catégorie de tous les
schémas sur $S$ satisfait à la condition de faisceau pour tout recouvrement h
et préserve les limites (Limites, remarque \ref{limits-remark-limit-preserving}).
\end{lemma}

\begin{proof}
Cela résulte des arguments donnés dans les démonstrations de
Topologies, lemmes \ref{topologies-lemma-extend-sheaf-general} et
\ref{topologies-lemma-extend-sheaf}, en utilisant les
lemmes \ref{lemma-limit-h-topology} et \ref{lemma-h-induced}.
Détails omis.
\end{proof}








\section{Carrés d'éclatement et topologie ph}
\label{section-blow-up-ph}

\noindent
Soit $X$ un schéma. Soit $Z \subset X$ un sous-schéma fermé
tel que le morphisme d'inclusion soit de présentation finie, c'est-à-dire que
le faisceau quasi-cohérent d'idéaux correspondant à $Z$ soit de
type fini. Soit $b : X' \to X$ l'éclatement de $X$ le long de $Z$
et soit $E = b^{-1}(Z)$ le diviseur exceptionnel. Voir
Diviseurs, section \ref{divisors-section-blowing-up}.
Dans cette situation et dans toute cette section, nous dirons que
\begin{equation}
\label{equation-blow-up-square}
\vcenter{
\xymatrix{
E \ar[d] \ar[r] & X' \ar[d]^b \\
Z \ar[r] & X
}
}
\end{equation}
est un {\it carré d'éclatement}.

\begin{lemma}
\label{lemma-blow-up-square-ph}
Soit $\mathcal{F}$ un faisceau sur un site $(\Sch/S)_{ph}$, voir
Topologies, définition \ref{topologies-definition-big-small-ph}.
Alors, pour tout carré d'éclatement (\ref{equation-blow-up-square})
dans la catégorie $(\Sch/S)_{ph}$, le diagramme
$$
\xymatrix{
\mathcal{F}(E) & \mathcal{F}(X') \ar[l] \\
\mathcal{F}(Z) \ar[u] & \mathcal{F}(X) \ar[u] \ar[l]
}
$$
est cartésien dans la catégorie des ensembles.
\end{lemma}

\begin{proof}
Puisque $Z \amalg X' \to X$ est un morphisme propre surjectif,
on voit que $\{Z \amalg X' \to X\}$ est un recouvrement ph
(Topologies, lemme \ref{topologies-lemma-surjective-proper-ph}).
Nous avons
$$
(Z \amalg X') \times_X (Z \amalg X') =
Z \amalg E \amalg E \amalg X' \times_X X'
$$
Puisque $\mathcal{F}$ est un faisceau de Zariski, on voit que
$\mathcal{F}$ transforme les unions disjointes en produits.
La condition de faisceau pour le recouvrement
$\{Z \amalg X' \to X\}$ affirme donc que
$\mathcal{F}(X) \to \mathcal{F}(Z) \times \mathcal{F}(X')$
est injective et a pour image l'ensemble des couples $(t, s')$ tels que
(a) $t|_E = s'|_E$ et (b) $s'$ appartient à l'égalisateur des deux applications
$\mathcal{F}(X') \to \mathcal{F}(X' \times_X X')$.
Observons ensuite que le morphisme évident
$$
E \times_Z E \amalg X' \longrightarrow X' \times_X X'
$$
est un morphisme propre surjectif, puisque $b$ induit
un isomorphisme $X' \setminus E \to X \setminus Z$. Nous en déduisons que
$\mathcal{F}(X' \times_X X') \to
\mathcal{F}(E \times_Z E) \times \mathcal{F}(X')$ est injective.
Il s'ensuit que (a) $\Rightarrow$ (b), ce qui prouve le lemme.
\end{proof}

\begin{lemma}
\label{lemma-thickening-ph}
Soit $\mathcal{F}$ un faisceau sur un site $(\Sch/S)_{ph}$
comme dans Topologies, définition \ref{topologies-definition-big-small-ph}.
Soit $X \to X'$ un morphisme de $(\Sch/S)_{ph}$ qui est
un épaississement. Alors
$\mathcal{F}(X') \to \mathcal{F}(X)$ est bijective.
\end{lemma}

\begin{proof}
Remarquons que $X \to X'$ est un morphisme propre surjectif et que
$X \times_{X'} X = X$.
La propriété de faisceau pour le recouvrement ph $\{X \to X'\}$
(Topologies, lemme
\ref{topologies-lemma-surjective-proper-ph})
permet de conclure.
\end{proof}








\section{Carrés de quasi-éclatement et topologie h}
\label{section-blow-up-h}

\noindent
Considérons un carré d'éclatement (\ref{equation-blow-up-square}).
Bien que le morphisme $b : X' \to X$
soit projectif (Diviseurs, lemme \ref{divisors-lemma-blowing-up-projective}),
il n'existe en général aucun moyen simple de garantir que
$b$ soit de présentation finie. Comme les recouvrements h sont construits
à l'aide de morphismes de présentation finie, il nous faut une variante.
Précisément, nous dirons qu'un diagramme commutatif
\begin{equation}
\label{equation-almost-blow-up-square}
\vcenter{
\xymatrix{
E \ar[d] \ar[r] & X' \ar[d]^b \\
Z \ar[r] & X
}
}
\end{equation}
de schémas est un {\it carré de quasi-éclatement} si les conditions suivantes
sont satisfaites
\begin{enumerate}
\item $Z \to X$ est une immersion fermée de présentation finie,
\item $E = b^{-1}(Z)$ est un sous-schéma fermé localement principal de $X'$,
\item $b$ est propre et de présentation finie,
\item le sous-schéma fermé $X'' \subset X'$ défini par l'idéal quasi-cohérent
des sections de $\mathcal{O}_{X'}$ à support dans $E$
(Propriétés, lemme \ref{properties-lemma-sections-supported-on-closed-subset})
est l'éclatement de $X$ le long de $Z$.
\end{enumerate}
Il s'ensuit que le morphisme $b$ induit un isomorphisme
$X' \setminus E \to X \setminus Z$.
Pour quelques exemples très simples de carrés de quasi-éclatement, voir les
exemples \ref{example-one-generator} et
\ref{example-two-generators}.

\medskip\noindent
Le changement de base d'un éclatement n'est en général pas un éclatement,
mais les quasi-éclatements sont compatibles avec le changement de base.

\begin{lemma}
\label{lemma-base-change-almost-blow-up}
Considérons un carré de quasi-éclatement (\ref{equation-almost-blow-up-square}).
Soit $Y \to X$ un morphisme quelconque. Alors le changement de base
$$
\xymatrix{
Y \times_X E \ar[d] \ar[r] & Y \times_X X' \ar[d] \\
Y \times_X Z \ar[r] & Y
}
$$
est lui aussi un carré de quasi-éclatement.
\end{lemma}

\begin{proof}
Le morphisme $Y \times_X X' \to Y$ est propre et de présentation finie
d'après Morphismes, lemmes \ref{morphisms-lemma-base-change-proper} et
\ref{morphisms-lemma-base-change-finite-presentation}.
Le morphisme $Y \times_X Z \to Y$ est une immersion fermée
(Morphismes, lemme \ref{morphisms-lemma-base-change-closed-immersion}) de
présentation finie. L'image inverse de $Y \times_X Z$ dans $Y \times_X X'$
est égale à l'image inverse de $E$ dans $Y \times_X X'$ et est donc
localement principale
(Diviseurs, lemme \ref{divisors-lemma-pullback-locally-principal}).
Soient $X'' \subset X'$, resp.\ $Y'' \subset Y \times_X X'$, les sous-schémas fermés
correspondant aux idéaux quasi-cohérents de sections de
$\mathcal{O}_{X'}$, resp.\ $\mathcal{O}_{Y \times_Y X'}$
à support dans $E$, resp.\ $Y \times_X E$.
Manifestement, $Y'' \subset Y \times_X X''$ est le sous-schéma fermé
correspondant à l'idéal quasi-cohérent des sections de
$\mathcal{O}_{Y \times_Y X''}$ à support dans $Y \times_X (E \cap X'')$.
Ainsi, $Y''$ est la transformée stricte de $Y$ relativement à l'éclatement
$X'' \to X$, voir
Diviseurs, définition \ref{divisors-definition-strict-transform}.
Ainsi, d'après Diviseurs, lemme \ref{divisors-lemma-strict-transform},
nous voyons que $Y''$ est l'éclatement de centre $Y \times_X Z$ sur $Y$.
\end{proof}

\noindent
On peut rétrécir les carrés de quasi-éclatement.

\begin{lemma}
\label{lemma-shrink-almost-blow-up}
Considérons un carré de quasi-éclatement (\ref{equation-almost-blow-up-square}).
Soit $W \to X'$ une immersion fermée de présentation finie.
Les assertions suivantes sont équivalentes
\begin{enumerate}
\item $X' \setminus E$ est schématiquement contenu dans $W$,
\item l'éclatement $X''$ de $X$ le long de $Z$ est schématiquement contenu dans $W$,
\item le diagramme
$$
\xymatrix{
E \cap W \ar[d] \ar[r] & W \ar[d] \\
Z \ar[r] & X
}
$$
est un carré de quasi-éclatement. Ici, $E \cap W$ est l'intersection
schématique.
\end{enumerate}
\end{lemma}

\begin{proof}
Supposons (1). Alors la surjection $\mathcal{O}_{X'} \to \mathcal{O}_W$
est un isomorphisme au-dessus de l'ouvert $X' \subset E$. Comme le faisceau d'idéaux
de $X'' \subset X'$ est constitué des sections de $\mathcal{O}_{X'}$ à support dans $E$
(par notre définition des carrés de quasi-éclatement),
nous concluons que (2) est vraie. Si (2) est vraie, alors (3) l'est.
Si (3) est vraie, alors (1) l'est parce que $X'' \cap (X' \setminus E)$
est isomorphe à $X \setminus Z$, qui est à son tour isomorphe
à $X' \setminus E$.
\end{proof}

\noindent
L'éclatement proprement dit est la limite\newline des rétrécissements de tout quasi-éclatement donné.

\begin{lemma}
\label{lemma-blow-up-limit-almost-blow-up}
Considérons un carré de quasi-éclatement (\ref{equation-almost-blow-up-square})
où $X$ est quasi-compact et quasi-séparé. Alors l'éclatement $X''$ de $X$
le long de $Z$ peut s'écrire
$$
X'' = \lim X'_i
$$
où la limite porte sur le système filtrant des sous-schémas fermés
$X'_i \subset X'$ de présentation finie satisfaisant aux conditions équivalentes
du lemme \ref{lemma-shrink-almost-blow-up}.
\end{lemma}

\begin{proof}
Soit $\mathcal{I} \subset \mathcal{O}_{X'}$ le faisceau quasi-cohérent
d'idéaux correspondant à $X''$. D'après
Propriétés, lemme \ref{properties-lemma-quasi-coherent-colimit-finite-type},
nous pouvons écrire $\mathcal{I}$ comme la colimite filtrante
$\mathcal{I} = \colim \mathcal{I}_i$ de ses sous-modules quasi-cohérents
de type fini. Comme ces modules correspondent
$1$-à-$1$ aux sous-schémas fermés $X'_i$, la démonstration est terminée.
\end{proof}

\noindent
Il existe des carrés de quasi-éclatement.

\begin{lemma}
\label{lemma-almost-blow-up-square}
Soit $X$ un schéma quasi-compact et quasi-séparé. Soit $Z \subset X$
un sous-schéma fermé défini par un faisceau quasi-cohérent d'idéaux
de type fini. Alors il existe un carré de quasi-éclatement comme dans
(\ref{equation-almost-blow-up-square}).
\end{lemma}

\begin{proof}
Nous pouvons écrire $X = \lim X_i$ comme la limite d'un système projectif filtrant
de schémas noethériens dont les morphismes de transition sont affines, voir
Limites, proposition \ref{limits-proposition-approximate}.
Nous pouvons trouver un indice $i$ et une immersion fermée $Z_i \to X_i$
dont le changement de base à $X$ est l'immersion fermée $Z \to X$.
Voir Limites, lemmes \ref{limits-lemma-descend-finite-presentation} et
\ref{limits-lemma-descend-closed-immersion-finite-presentation}.
Soit $b_i : X'_i \to X_i$ l'éclatement de centre $Z_i$.
Cela donne un carré d'éclatement
$$
\xymatrix{
E_i \ar[r] \ar[d] & X'_i \ar[d]^{b_i} \\
Z_i \ar[r] & X_i
}
$$
où tous les morphismes sont des morphismes de type fini entre schémas noethériens
et sont donc de présentation finie. Il s'agit donc d'un carré de quasi-éclatement.
D'après le lemme \ref{lemma-base-change-almost-blow-up},
le changement de base de ce diagramme à $X$ fournit le
carré de quasi-éclatement recherché.
\end{proof}

\noindent
Les carrés de quasi-éclatement sont uniques à rétrécissement près, comme dans le
lemme \ref{lemma-shrink-almost-blow-up}.

\begin{lemma}
\label{lemma-almost-blow-up-unique}
Soit $X$ un schéma quasi-compact et quasi-séparé et soit $Z \subset X$
un sous-schéma fermé défini par un faisceau quasi-cohérent d'idéaux de type fini.
Supposons donnés des carrés de quasi-éclatement (\ref{equation-almost-blow-up-square})
$$
\xymatrix{
E_k \ar[r] \ar[d] & X_k' \ar[d] \\
Z \ar[r] & X
}
$$
pour $k = 1, 2$ ; alors il existe un carré de quasi-éclatement
$$
\xymatrix{
E \ar[r] \ar[d] & X' \ar[d] \\
Z \ar[r] & X
}
$$
et des immersions fermées $i_k : X' \to X'_k$ au-dessus de $X$
telles que $E = i_k^{-1}(E_k)$.
\end{lemma}

\begin{proof}
Notons $X'' \to X$ l'éclatement de centre $Z$ dans $X$.
Nous considérons $X''$ comme un sous-schéma fermé de $X'_1$ et de $X'_2$.
Écrivons $X'' = \lim X'_{1, i}$ comme dans le
lemme \ref{lemma-blow-up-limit-almost-blow-up}.
D'après Limites, proposition
\ref{limits-proposition-characterize-locally-finite-presentation},
il existe un $i$ et un morphisme $h : X'_{1, i} \to X'_2$ qui coïncide
avec les inclusions $X'' \subset X'_{1, i}$ et $X'' \subset X'_2$.
D'après Limites, le lemme \ref{limits-lemma-eventually-closed-immersion},
la restriction de $h$ à $X'_{1, i'}$ est une immersion fermée
pour un certain $i' \geq i$. Ceci achève la démonstration.
\end{proof}

\noindent
Nos techniques d'aplatissement par éclatement s'étendent aux
quasi-éclatements dans les situations favorables.

\begin{lemma}
\label{lemma-flat-after-almost-blowing-up}
Soit $Y$ un schéma quasi-compact et quasi-séparé.
Soit $X$ un schéma de présentation finie sur $Y$.
Soit $V \subset Y$ un ouvert quasi-compact tel que
$X_V \to V$ soit plat. Alors il existe un diagramme commutatif
$$
\xymatrix{
E \ar[ddd] \ar[rd] & & & D \ar[lll] \ar[ddd] \ar[ld] \\
& Y' \ar[d] & X' \ar[l] \ar[d] \\
& Y & X \ar[l] \\
Z \ar[ru] & & & T \ar[lll] \ar[lu]
}
$$
dont les carrés de droite et de gauche sont des carrés de quasi-éclatement,
dont les carrés inférieur et supérieur sont cartésiens, et tel que
$Z \cap V = \emptyset$, et
tel que $X' \to Y'$ soit plat (et de présentation finie).
\end{lemma}

\begin{proof}
Si $Y$ est un schéma noethérien, ce lemme découle immédiatement
du lemme \ref{lemma-flat-after-blowing-up},
car, dans ce cas, les carrés d'éclatement sont des carrés de quasi-éclatement
(nous utilisons aussi le fait que les transformées strictes sont des éclatements).
Le cas général se ramène au cas noethérien par
approximation noethérienne absolue.

\medskip\noindent
Nous pouvons écrire $Y = \lim Y_i$ comme la limite d'un système projectif filtrant
de schémas noethériens dont les morphismes de transition sont affines, voir
Limites, proposition \ref{limits-proposition-approximate}.
Nous pouvons trouver un indice $i$ et un morphisme $X_i \to Y_i$
de présentation finie dont le changement de base à $Y$ est $X \to Y$.
Voir Limites, lemme \ref{limits-lemma-descend-finite-presentation}.
Quitte à augmenter $i$, nous pouvons supposer que $V$ est l'image réciproque
d'un sous-schéma ouvert $V_i \subset Y_i$, voir
Limites, lemme \ref{limits-lemma-descend-opens}.
Enfin, quitte à augmenter $i$, nous pouvons supposer que $X_{i, V_i} \to V_i$
est plat, voir
Limites, lemme \ref{limits-lemma-descend-flat-finite-presentation}.
Dans le cas noethérien, nous pouvons construire un diagramme comme dans le
lemme pour $X_i \to Y_i \supset V_i$. Le changement de base de ce
diagramme par $Y \to Y_i$ donne la solution. Utilisons le fait que le changement de base
préserve les propriétés des morphismes, voir Morphismes, lemmes
\ref{morphisms-lemma-base-change-proper},
\ref{morphisms-lemma-base-change-finite-presentation},
\ref{morphisms-lemma-base-change-closed-immersion}, et
\ref{morphisms-lemma-base-change-flat},
ainsi que le fait que le changement de base d'un carré de
quasi-éclatement est un carré de quasi-éclatement, voir
le lemme \ref{lemma-base-change-almost-blow-up}.
\end{proof}

\begin{lemma}
\label{lemma-blow-up-square-h}
Soit $\mathcal{F}$ un faisceau sur l'un des sites $(\Sch/S)_h$
construits dans la définition \ref{definition-big-small-h}.
Alors, pour tout carré de quasi-éclatement (\ref{equation-almost-blow-up-square})
dans la catégorie $(\Sch/S)_h$, le diagramme
$$
\xymatrix{
\mathcal{F}(E) & \mathcal{F}(X') \ar[l] \\
\mathcal{F}(Z) \ar[u] & \mathcal{F}(X) \ar[u] \ar[l]
}
$$
est cartésien dans la catégorie des ensembles.
\end{lemma}

\begin{proof}
Puisque $Z \amalg X' \to X$ est un morphisme propre surjectif
de présentation finie, on voit que $\{Z \amalg X' \to X\}$ est un recouvrement h
(lemme \ref{lemma-surjective-proper-finite-presentation-h}).
Nous avons
$$
(Z \amalg X') \times_X (Z \amalg X') =
Z \amalg E \amalg E \amalg X' \times_X X'
$$
Puisque $\mathcal{F}$ est un faisceau pour la topologie de Zariski, on voit que
$\mathcal{F}$ envoie les unions disjointes sur des produits.
Ainsi, la condition de faisceau pour le recouvrement
$\{Z \amalg X' \to X\}$ dit que
$\mathcal{F}(X) \to \mathcal{F}(Z) \times \mathcal{F}(X')$
est injective et a pour image l'ensemble des couples $(t, s')$ tels que
(a) $t|_E = s'|_E$ et (b) $s'$ appartient à l'égalisateur des deux applications
$\mathcal{F}(X') \to \mathcal{F}(X' \times_X X')$.
Observons ensuite que le morphisme évident
$$
E \times_Z E \amalg X' \longrightarrow X' \times_X X'
$$
est un morphisme propre surjectif de présentation finie puisque $b$ induit
un isomorphisme $X' \setminus E \to X \setminus Z$. Nous en déduisons que
$\mathcal{F}(X' \times_X X') \to
\mathcal{F}(E \times_Z E) \times \mathcal{F}(X')$ est injective.
Il s'ensuit que (a) $\Rightarrow$ (b), ce qui prouve le lemme.
\end{proof}

\begin{lemma}
\label{lemma-thickening-h}
Soit $\mathcal{F}$ un faisceau sur l'un des sites $(\Sch/S)_h$
construits dans la définition \ref{definition-big-small-h}.
Soit $X \to X'$ un morphisme de $(\Sch/S)_h$ qui est
un épaississement et qui est de présentation finie. Alors
$\mathcal{F}(X') \to \mathcal{F}(X)$ est bijective.
\end{lemma}

\begin{proof}
Première démonstration. Remarquons que $X \to X'$ est un morphisme propre surjectif
de présentation finie et que $X \times_{X'} X = X$.
Par la propriété de faisceau pour le recouvrement h $\{X \to X'\}$
(lemme \ref{lemma-surjective-proper-finite-presentation-h}),
nous pouvons conclure.

\medskip\noindent
Deuxième démonstration (sotte). L'éclatement de $X'$ le long de $X$ est le schéma vide.
Cela tient au fait que l'algèbre affine d'éclatement $A[\frac{I}{a}]$
(Algèbre, section \ref{algebra-section-blow-up})
est nulle si $a$ est un élément nilpotent de $A$. Les détails sont omis.
On obtient donc un carré de quasi-éclatement de la forme
$$
\xymatrix{
\emptyset \ar[r] \ar[d] & \emptyset \ar[d] \\
X \ar[r] & X'
}
$$
Puisque $\mathcal{F}$ est un faisceau, $\mathcal{F}(\emptyset)$
est un singleton.
En appliquant le lemme \ref{lemma-blow-up-square-h}, on obtient la conclusion.
\end{proof}

\begin{proposition}
\label{proposition-check-h}
Soit $\mathcal{F}$ un préfaisceau sur l'un des sites $(\Sch/S)_h$
construits dans la définition \ref{definition-big-small-h}.
Alors $\mathcal{F}$ est un faisceau si et seulement si les
conditions suivantes sont satisfaites
\begin{enumerate}
\item $\mathcal{F}$ est un faisceau pour la topologie de Zariski,
\item pour tout morphisme $f : X \to Y$ de $(\Sch/S)_h$ tel que $Y$ soit affine
et que $f$ soit surjectif, plat, propre et de présentation finie,
$\mathcal{F}(Y)$ est l'égalisateur des deux applications
$\mathcal{F}(X) \to \mathcal{F}(X \times_Y X)$,
\item pour tout carré de quasi-éclatement (\ref{equation-almost-blow-up-square})
avec $X$ affine dans la catégorie $(\Sch/S)_h$, le diagramme
$$
\xymatrix{
\mathcal{F}(E) & \mathcal{F}(X') \ar[l] \\
\mathcal{F}(Z) \ar[u] & \mathcal{F}(X) \ar[u] \ar[l]
}
$$
est cartésien dans la catégorie des ensembles.
\end{enumerate}
\end{proposition}

\begin{proof}
Supposons que $\mathcal{F}$ soit un faisceau. La condition (1) est vérifiée car
un recouvrement de Zariski est un recouvrement h, voir
le lemme \ref{lemma-zariski-h}.
La condition (2) est vérifiée car, pour $f$ comme en (2),
$\{X \to Y\}$ est un recouvrement fppf (c'est clair)
et donc un recouvrement h, voir le lemme \ref{lemma-zariski-h}.
La condition (3) résulte du lemme \ref{lemma-blow-up-square-h}.

\medskip\noindent
Réciproquement, supposons que $\mathcal{F}$ satisfasse (1), (2) et (3).
Nous allons montrer que $\mathcal{F}$ est un faisceau en appliquant
le lemme \ref{lemma-characterize-sheaf-h}. Considérons
un morphisme surjectif, de présentation finie et propre
$f : X \to Y$ dans $(\Sch/S)_h$, avec $Y$ affine. Il suffit de montrer
que $\mathcal{F}(Y)$ est l'égalisateur des deux applications
$\mathcal{F}(X) \to \mathcal{F}(X \times_Y X)$.

\medskip\noindent
Supposons d'abord qu'en outre $f : X \to Y$ soit une immersion fermée
(autrement dit, $f$ est un épaississement). Alors l'éclatement de $Y$ le long de $X$
est le schéma vide, ce qui fournit un carré de quasi-éclatement
dont les sommets sont $\emptyset, \emptyset, X, Y$
(comparer avec la deuxième démonstration du lemme \ref{lemma-thickening-h}).
La condition (3) nous dit donc que
$$
\xymatrix{
\mathcal{F}(\emptyset) & \mathcal{F}(\emptyset) \ar[l] \\
\mathcal{F}(X) \ar[u] & \mathcal{F}(Y) \ar[u] \ar[l]
}
$$
est cartésien dans la catégorie des ensembles. Puisque $\mathcal{F}$ est un faisceau
pour la topologie de Zariski, on voit que $\mathcal{F}(\emptyset)$
est un singleton. On en déduit que $\mathcal{F}(X) = \mathcal{F}(Y)$.

\medskip\noindent
Intermède A : soit $T \to T'$ un morphisme de $(\Sch/S)_h$ qui est
un épaississement et qui est de présentation finie. Alors
$\mathcal{F}(T') \to \mathcal{F}(T)$ est bijective.
En effet, choisissons un recouvrement ouvert affine $T' = \bigcup T'_i$
et posons $T_i = T \times_{T'} T'_i$ ; ce sont les ouverts affines correspondants
de $T$. Alors $\mathcal{F}(T'_i) \to \mathcal{F}(T_i)$
est bijective pour tout $i$ d'après le résultat du paragraphe précédent.
La propriété de faisceau pour la topologie de Zariski montre que
$\mathcal{F}(T') \to \mathcal{F}(T)$ est injective. En répétant le
raisonnement, on constate qu'elle est bijective. Nous omettons les détails mineurs.

\medskip\noindent
Intermède B : considérons un carré de quasi-éclatement
(\ref{equation-almost-blow-up-square}) dans la catégorie $(\Sch/S)_h$.
Nous affirmons alors que le diagramme
$$
\xymatrix{
\mathcal{F}(E) & \mathcal{F}(X') \ar[l] \\
\mathcal{F}(Z) \ar[u] & \mathcal{F}(X) \ar[u] \ar[l]
}
$$
est cartésien dans la catégorie des ensembles. Cela résulte de la condition
(3) de la manière suivante : on choisit un recouvrement ouvert affine de $X$ et on raisonne
comme dans l'intermède A. Nous omettons les détails.

\medskip\noindent
Ensuite, soit $f : X \to Y$ un morphisme surjectif, propre et de présentation finie
dans $(\Sch/S)_h$, avec $Y$ affine. Choisissons une stratification par platitude générique
$$
Y \supset Y_0 \supset Y_1 \supset \ldots \supset Y_t = \emptyset
$$
comme dans Compléments sur les morphismes, lemme
\ref{more-morphisms-lemma-generic-flatness-stratification-scheme},
pour $f : X \to Y$.
Nous allons utiliser toutes les propriétés de cette stratification
sans les mentionner davantage. Posons $X_0 = X \times_Y Y_0$.
D'après l'intermède B, nous avons $\mathcal{F}(Y_0) = \mathcal{F}(Y)$,
$\mathcal{F}(X_0) = \mathcal{F}(X)$, et
$\mathcal{F}(X_0 \times_{Y_0} X_0) = \mathcal{F}(X \times_Y X)$.

\medskip\noindent
Nous allons démontrer le résultat par récurrence sur $t$. Si $t = 1$,
alors $X_0 \to Y_0$ est surjectif, propre, plat et de présentation finie,
et le résultat découle de la propriété (2).
Pour $t > 1$, nous pouvons remplacer $Y$ par $Y_0$ et $X$ par $X_0$
(voir ci-dessus) et supposer que $Y = Y_0$.

\medskip\noindent
Considérons le sous-schéma ouvert quasi-compact
$V = Y \setminus Y_1 = Y_0 \setminus Y_1$.
Choisissons un diagramme
$$
\xymatrix{
E \ar[ddd] \ar[rd] & & & D \ar[lll] \ar[ddd] \ar[ld] \\
& Y' \ar[d] & X' \ar[l] \ar[d] \\
& Y & X \ar[l] \\
Z \ar[ru] & & & T \ar[lll] \ar[lu]
}
$$
comme dans le lemme \ref{lemma-flat-after-almost-blowing-up}
pour $f : X \to Y \supset V$. Alors $f' : X' \to Y'$ est plat et de
présentation finie. De plus, $f'$ est propre (utiliser
Morphismes, lemmes \ref{morphisms-lemma-composition-proper} et
\ref{morphisms-lemma-image-proper-scheme-closed} pour le voir).
Ainsi, l'image $W = f'(X') \subset Y'$ est un sous-schéma ouvert
(Morphismes, lemme \ref{morphisms-lemma-fppf-open}) et fermé
de $Y'$. Observons que $Y' \setminus E$ est contenu
dans $W$. D'après le lemme \ref{lemma-shrink-almost-blow-up},
cela signifie que nous pouvons remplacer $Y'$ par $W$
dans le diagramme ci-dessus. Autrement dit, nous pouvons supposer, et supposons,
que $f'$ est surjectif. À ce stade, nous savons que
$$
\vcenter{
\xymatrix{
\mathcal{F}(E) & \mathcal{F}(Y') \ar[l] \\
\mathcal{F}(Z) \ar[u] & \mathcal{F}(Y) \ar[u] \ar[l]
}
}
\quad\text{et}\quad
\vcenter{
\xymatrix{
\mathcal{F}(D) & \mathcal{F}(X') \ar[l] \\
\mathcal{F}(T) \ar[u] & \mathcal{F}(X) \ar[u] \ar[l]
}
}
$$
sont cartésiens d'après l'intermède B. Remarquons que
$Z \cap Y_1 \to Z$ est un épaississement de
présentation finie (puisque $Z$ est ensemblistement contenu dans $Y_1$
comme sous-schéma fermé de $Y$ disjoint de $V$).
Nous obtenons donc une filtration
$$
Z \supset Z \cap Y_1 \supset Z \cap Y_2 \subset \ldots \subset
Z \cap Y_t = \emptyset
$$
comme ci-dessus pour la restriction $T = Z \times_Y X \to Z$ de $f$ à $T$.
Par l'hypothèse de récurrence, nous trouvons donc que
$\mathcal{F}(Z) \to \mathcal{F}(T)$
est une application injective d'ensembles dont l'image est l'égalisateur
des deux applications $\mathcal{F}(T) \to
\mathcal{F}(T \times_Z T)$.

\medskip\noindent
Soit $s \in \mathcal{F}(X)$ un élément de l'égalisateur des deux applications
$\mathcal{F}(X) \to \mathcal{F}(X \times_Y X)$.
D'après ce qui précède, la restriction $s|_T$
provient d'un unique élément $t \in \mathcal{F}(Z)$,
et, de même, la restriction $s|_{X'}$
provient d'un unique élément $t' \in \mathcal{F}(Y')$.
En chassant les sections au moyen des applications de restriction de $\mathcal{F}$
correspondant aux flèches du grand diagramme commutatif ci-dessus,
le lecteur constate que $t$ et $t'$ se restreignent au même élément de
$\mathcal{F}(E)$, car ils se restreignent au même élément de
$\mathcal{F}(D)$ et que (2) s'applique ; nous utilisons ici le fait que $D \to E$ est
surjectif, plat, propre et de présentation finie en tant que restriction
de $X' \to Y'$. Ainsi, le premier des deux
carrés cartésiens affichés ci-dessus nous donne une unique section
$u \in \mathcal{F}(Y)$
dont les restrictions valent respectivement $t$ et $t'$ sur $Z$ et $Y'$.
Pour voir que $u$ se restreint à $s$ sur $X$, utiliser le second diagramme.
\end{proof}

\begin{example}
\label{example-one-generator}
Soit $A$ un anneau. Soit $f \in A$ un élément. Soit $J \subset A$
un idéal de type fini annulé par une puissance de $f$.
Alors
$$
\xymatrix{
E = \Spec(A/fA + J) \ar[r] \ar[d] & \Spec(A/J) = X' \ar[d] \\
Z = \Spec(A/fA) \ar[r] & \Spec(A) = X
}
$$
est un carré de quasi-éclatement.
\end{example}

\begin{example}
\label{example-two-generators}
Soit $A$ un anneau. Soient $f_1, f_2 \in A$ des éléments.
$$
\xymatrix{
E = \text{Proj}(A/(f_1, f_2)[T_0, T_1]) \ar[r] \ar[d] &
\text{Proj}(A[T_0, T_1]/(f_2 T_0 - f_1 T_1) = X' \ar[d] \\
Z = \Spec(A/(f_1, f_2)) \ar[r] & \Spec(A) = X
}
$$
est un carré de quasi-éclatement.
\end{example}

\begin{lemma}
\label{lemma-refine-check-h}
Soit $\mathcal{F}$ un préfaisceau sur l'un des sites $(\Sch/S)_h$
construits dans la définition \ref{definition-big-small-h}.
Alors $\mathcal{F}$ est un faisceau si et seulement si les
conditions suivantes sont satisfaites
\begin{enumerate}
\item $\mathcal{F}$ est un faisceau pour la topologie de Zariski,
\item étant donné un morphisme $f : X \to Y$ de $(\Sch/S)_h$ avec $Y$ affine
et $f$ surjectif, plat, propre et de présentation finie, alors
$\mathcal{F}(Y)$ est l'égalisateur des deux applications
$\mathcal{F}(X) \to \mathcal{F}(X \times_Y X)$,
\item $\mathcal{F}$ transforme un carré de quasi-éclatement comme dans
l'exemple \ref{example-one-generator} dans la catégorie $(\Sch/S)_h$
en un diagramme cartésien d'ensembles, et
\item $\mathcal{F}$ transforme un carré de quasi-éclatement comme dans
l'exemple \ref{example-two-generators} dans la catégorie $(\Sch/S)_h$
en un diagramme cartésien d'ensembles.
\end{enumerate}
\end{lemma}

\begin{proof}
D'après la proposition \ref{proposition-check-h}, il suffit de montrer qu'étant donné
un carré de quasi-éclatement (\ref{equation-almost-blow-up-square})
avec $X$ affine dans la catégorie $(\Sch/S)_h$, le diagramme
$$
\xymatrix{
\mathcal{F}(E) & \mathcal{F}(X') \ar[l] \\
\mathcal{F}(Z) \ar[u] & \mathcal{F}(X) \ar[u] \ar[l]
}
$$
est cartésien dans la catégorie des ensembles. L'idée générale de la démonstration
consiste à dominer le morphisme par d'autres carrés de quasi-éclatement auxquels
nous pouvons appliquer localement les hypothèses (3) et (4).

\medskip\noindent
Supposons que nous ayons un carré de quasi-éclatement
(\ref{equation-almost-blow-up-square}) dans la catégorie $(\Sch/S)_h$,
un recouvrement ouvert $X = \bigcup U_i$ et des recouvrements ouverts
$U_i \cap U_j = \bigcup U_{ijk}$ tels que les diagrammes
$$
\vcenter{
\xymatrix{
\mathcal{F}(E \cap b^{-1}(U_i)) & \mathcal{F}(b^{-1}(U_i)) \ar[l] \\
\mathcal{F}(Z \cap U_i) \ar[u] & \mathcal{F}(U_i) \ar[u] \ar[l]
}
}
\quad\text{et}\quad
\vcenter{
\xymatrix{
\mathcal{F}(E \cap b^{-1}(U_{ijk})) &
\mathcal{F}(b^{-1}(U_{ijk})) \ar[l] \\
\mathcal{F}(Z \cap U_{ijk}) \ar[u] &
\mathcal{F}(U_{ijk}) \ar[u] \ar[l]
}
}
$$
sont cartésiens ; il en va alors de même pour
$$
\xymatrix{
\mathcal{F}(E) & \mathcal{F}(X') \ar[l] \\
\mathcal{F}(Z) \ar[u] & \mathcal{F}(X) \ar[u] \ar[l]
}
$$
Cela découle du fait que $\mathcal{F}$ est un faisceau pour la topologie de Zariski.

\medskip\noindent
En particulier, si nous avons un carré d'éclatement
(\ref{equation-almost-blow-up-square}) tel que $b : X' \to X$
soit une immersion fermée et que $Z$ soit un sous-schéma fermé
localement principal, alors nous voyons que
$\mathcal{F}(X) = \mathcal{F}(X') \times_{\mathcal{F}(E)} \mathcal{F}(Z)$.
En effet, sur un recouvrement de $X$ par des ouverts affines, nous obtenons un carré de quasi-éclatement comme en (3).

\medskip\noindent
Soient $Z \subset X$, $E_k \subset X'_k \to X$, $E \subset X' \to X$, et
$i_k : X' \to X'_k$ comme dans l'énoncé du
lemme \ref{lemma-almost-blow-up-unique}. Alors
$$
\xymatrix{
E \ar[d] \ar[r] & X' \ar[d] \\
E_k \ar[r] & X'_k
}
$$
est un carré de quasi-éclatement du type considéré au paragraphe
précédent. Ainsi
$$
\mathcal{F}(X'_k) = \mathcal{F}(X') \times_{\mathcal{F}(E)} \mathcal{F}(E_k)
$$
pour $k = 1, 2$, d'après le résultat du paragraphe précédent. Il s'ensuit que
$$
\mathcal{F}(X) \longrightarrow
\mathcal{F}(X'_k) \times_{\mathcal{F}(E_k)} \mathcal{F}(Z)
$$
est bijective pour $k = 1$ si et seulement si elle est bijective pour $k = 2$.
Ainsi, étant donnée une immersion fermée $Z \to X$ de présentation finie,
avec $X$ quasi-compact et quasi-séparé, le fait que
$\mathcal{F}(X) = \mathcal{F}(X') \times_{\mathcal{F}(E)} \mathcal{F}(Z)$
soit vérifiée ou non est indépendant du choix du carré de quasi-éclatement
(\ref{equation-almost-blow-up-square}) effectué.
(De plus, d'après le lemme \ref{lemma-almost-blow-up-square},
il existe effectivement un carré de
quasi-éclatement pour $Z \subset X$.)

\medskip\noindent
Enfin, considérons un objet affine $X$ de $(\Sch/S)_h$
et une immersion fermée $Z \to X$ de présentation finie.
Nous allons démontrer la propriété voulue pour la paire $(X, Z)$
par récurrence sur le nombre $r$ de générateurs de l'idéal
définissant $Z$ dans $X$. Si le nombre de générateurs est $\leq 2$,
nous pouvons choisir notre carré de quasi-éclatement comme
dans l'exemple \ref{example-two-generators}
et conclure par l'hypothèse (4).

\medskip\noindent
Étape de récurrence. Supposons $X = \Spec(A)$ et $Z = \Spec(A/(f_1, \ldots, f_r))$
avec $r > 2$. Choisissons un carré d'éclatement
(\ref{equation-almost-blow-up-square})
pour la paire $(X, Z)$. Posons $Z_1 = \Spec(A/(f_1, f_2))$
et soit
$$
\xymatrix{
E_1 \ar[d] \ar[r] & Y \ar[d] \\
Z_1 \ar[r] & X
}
$$
le carré de quasi-éclatement construit dans
l'exemple \ref{example-two-generators}.
D'après le lemme \ref{lemma-base-change-almost-blow-up}, les changements de base
$$
(I)
\vcenter{
\xymatrix{
Y \times_X E \ar[r] \ar[d] & Y \times_X X' \ar[d] \\
Y \times_X Z \ar[r] & Y
}
}
\quad\text{et}\quad
(II)
\vcenter{
\xymatrix{
E \ar[r] \ar[d] & Z_1 \times_X X' \ar[d] \\
Z \ar[r] & Z_1
}
}
$$
sont des carrés de quasi-éclatement. L'idéal de $Z$ dans $Z_1$ est engendré
par $r - 2$ éléments. L'idéal de $Y \times_X Z$
est engendré par les images réciproques de $f_1, \ldots, f_r$ dans $Y$. Localement sur $Y$,
l'idéal engendré par $f_1, f_2$ peut être engendré par
un élément ; ainsi, $Y \times_X Z$ est, sur des ouverts affines de $Y$,
défini par au plus $r - 1$ éléments.
Par les hypothèses de récurrence et la discussion ci-dessus,
$$
\mathcal{F}(Y) =
\mathcal{F}(Y \times_X X') \times_{\mathcal{F}(Y \times_X E)}
\mathcal{F}(Y \times_X Z)
$$
et
$$
\mathcal{F}(Z_1) =
\mathcal{F}(Z_1 \times_X X') \times_{\mathcal{F}(E)}
\mathcal{F}(Z)
$$
Par l'hypothèse (4), nous avons
$$
\mathcal{F}(X) =
\mathcal{F}(Y) \times_{\mathcal{F}(E_1)} \mathcal{F}(Z_1)
$$
Supposons maintenant que nous ayons une paire $(s', t)$ avec $s' \in \mathcal{F}(X')$
et $t \in \mathcal{F}(Z)$ ayant la même restriction dans $\mathcal{F}(E)$.
Alors $(s'|{Z_1 \times_X X'}, t)$ est l'image d'un unique élément
$t_1 \in \mathcal{F}(Z_1)$. De même,
$(s'|_{Y \times_X X'}, t|_{Y \times_X Z})$ est l'image d'un
unique élément $s_Y \in \mathcal{F}(Y)$.
Nous affirmons que $s_Y$ et $t_1$ se restreignent au même élément
de $\mathcal{F}(E_1)$. Cela est vrai parce que le
carré de quasi-éclatement
$$
\xymatrix{
E_1 \times_X E \ar[r] \ar[d] & E_1 \times_X X' \ar[d] \\
E_1 \times_X Z \ar[r] & E_1
}
$$
est le changement de base du carré de quasi-éclatement (I) par $E_1 \to Y$ et
le changement de base du carré de quasi-éclatement (II) par $E_1 \to Z_1$, et
parce que les paires de sections utilisées pour construire $s_Y$ et $t_1$ coïncident.
Ainsi, par la troisième égalité de produits fibrés, nous voyons qu'il existe
un unique $s \in \mathcal{F}(X)$ dont l'image est $s_Y$ dans $\mathcal{F}(Y)$
et $t_1$ dans $\mathcal{F}(Z)$.
Nous omettons de vérifier que l'image de $s$ est $s'$ dans $\mathcal{F}(X')$
et $t$ dans $\mathcal{F}(Z)$ ; indication : utiliser l'unicité de $s$ que nous venons
de construire et raisonner localement sur des ouverts affines.
\end{proof}

\begin{lemma}
\label{lemma-refine-check-h-stack}
Soit $p : \mathcal{S} \to (\Sch/S)_h$ une catégorie fibrée en groupoïdes.
Alors $\mathcal{S}$ est un champ en groupoïdes si et seulement si les conditions
suivantes sont satisfaites
\begin{enumerate}
\item $\mathcal{S}$ est un champ en groupoïdes pour la topologie de Zariski,
\item étant donné un morphisme $f : X \to Y$ de $(\Sch/S)_h$ avec $Y$ affine
et $f$ surjectif, plat, propre et de présentation finie, alors
$$
\mathcal{S}_Y \longrightarrow
\mathcal{S}_X \times_{\mathcal{S}_{X \times_Y X}} \mathcal{S}_X
$$
est une équivalence de catégories,
\item pour un carré de quasi-éclatement comme dans
l'exemple \ref{example-one-generator} ou \ref{example-two-generators}
dans la catégorie $(\Sch/S)_h$, le foncteur
$$
\mathcal{S}_X \longrightarrow
\mathcal{S}_Z \times_{\mathcal{S}_E} \mathcal{S}_{X'}
$$
est une équivalence de catégories.
\end{enumerate}
\end{lemma}

\begin{proof}
Ce lemme est une conséquence formelle du lemme \ref{lemma-refine-check-h}
et de notre définition des champs en groupoïdes.
Par exemple, supposons (1), (2), (3). Pour montrer que
$\mathcal{S}$ est un champ, nous devons
démontrer la descente des morphismes et des objets, voir
Champs, définition \ref{stacks-definition-stack-in-groupoids}.

\medskip\noindent
Si $x, y$ sont des objets de $\mathcal{S}$ au-dessus d'un objet $U$ de
$(\Sch/S)_h$, alors nos hypothèses impliquent que $\mathit{Isom}(x, y)$
est un préfaisceau sur $(\Sch/U)_h$ qui satisfait
(1), (2), (3) et (4) du lemme \ref{lemma-refine-check-h},
et est donc un faisceau. Nous omettons quelques détails.

\medskip\noindent
Soit $\{U_i \to U\}_{i \in I}$ un recouvrement de $(\Sch/S)_h$.
Soit $(x_i, \varphi_{ij})$ une donnée de descente dans $\mathcal{S}$
relative à la famille $\{U_i \to U\}_{i \in I}$, voir
Champs, définition \ref{stacks-definition-descent-data}.
Considérons la règle $F$ qui, à $V/U$ dans $(\Sch/U)_h$, associe
l'ensemble des paires $(y, \psi_i)$, où $y$ est un objet de $\mathcal{S}_V$
et $\psi_i : y|_{U_i \times_U V} \to x_i|_{U_i \times_U V}$
est un morphisme de $\mathcal{S}$ au-dessus de $U_i \times_U V$
tel que
$$
\varphi_{ij}|_{U_i \times_U U_j \times_U V}
\circ
\psi_i|_{U_i \times_U U_j \times_U V}
=
\psi_j|_{U_i \times_U U_j \times_U V}
$$
à isomorphisme près. Puisque nous disposons déjà de la descente des morphismes, il est
clair que $F(V/U)$ est soit vide, soit un singleton. D'autre part,
$F(U_{i_0}/U)$ est non vide, car il contient
$(x_{i_0}, \varphi_{i_0i})$. Puisque notre but est de montrer que $F(U/U)$
est non vide, il suffit de montrer que $F$ est un faisceau sur $(\Sch/U)_h$.
Pour ce faire, nous pouvons utiliser le critère du lemme \ref{lemma-refine-check-h}.
Cependant, nos hypothèses (1), (2), (3) impliquent (en traçant certains
diagrammes commutatifs que nous omettons) que les propriétés (1), (2), (3) et (4)
du lemme \ref{lemma-refine-check-h} sont satisfaites pour $F$.

\medskip\noindent
Nous omettons de vérifier que, si $\mathcal{S}$ est un champ en groupoïdes,
alors (1), (2) et (3) sont satisfaites.
\end{proof}





\section{Normalisation faible absolue et recouvrements h}
\label{section-h-and-weak-normalization}

\noindent
Dans cette section, nous utilisons les critères obtenus dans la
section \ref{section-blow-up-h}
pour exhiber quelques faisceaux h et nous relions la faisceautisation h
du faisceau structural à la normalisation faible absolue.
Nous aurons besoin pour cela du lemme élémentaire suivant.

\begin{lemma}
\label{lemma-funny-blow-up}
Soient $Z, X, X', E$\newline les schémas d'un carré de quasi-éclatement comme dans
l'exemple \ref{example-two-generators}.
Alors $H^p(X', \mathcal{O}_{X'}) = 0$ pour $p > 0$ et\newline
$\Gamma(X, \mathcal{O}_X) \to \Gamma(X', \mathcal{O}_{X'})$
est un homomorphisme surjectif d'anneaux dont le noyau est un idéal de carré nul.
\end{lemma}

\begin{proof}
Supposons d'abord que $A = \mathbf{Z}[f_1, f_2]$ est l'anneau de polynômes.
Dans ce cas, notre carré de quasi-éclatement est l'éclatement de $X = \Spec(A)$
le long du sous-schéma fermé $Z$ et, en fait, $X' \subset \mathbf{P}^1_X$
est un diviseur de Cartier effectif découpé par la section globale
$f_2T_0 - f_1 T_1$ de $\mathcal{O}_{\mathbf{P}^1_X}(1)$.
Nous avons donc une résolution
$$
0 \to
\mathcal{O}_{\mathbf{P}^1_X}(-1) \to
\mathcal{O}_{\mathbf{P}^1_X} \to
\mathcal{O}_{X'} \to
0
$$
En utilisant la description de la cohomologie donnée dans
Cohomologie des schémas, section
\ref{coherent-section-cohomology-projective-space},
il s'ensuit que, dans ce cas,
$\Gamma(X, \mathcal{O}_X) \to \Gamma(X', \mathcal{O}_{X'})$
est un isomorphisme et $H^1(X', \mathcal{O}_{X'}) = 0$.

\medskip\noindent
Ensuite, nous observons que tout diagramme comme dans l'exemple \ref{example-two-generators}
est le changement de base du diagramme du paragraphe précédent par
l'homomorphisme d'anneaux $\mathbf{Z}[f_1, f_2] \to A$. Par conséquent, par Compléments sur les morphismes, lemmes
\ref{more-morphisms-lemma-check-h1-fibre-zero},
\ref{more-morphisms-lemma-h1-fibre-zero}, et
\ref{more-morphisms-lemma-h1-fibre-zero-check-h0-kappa},
nous concluons que $H^1(X', \mathcal{O}_{X'})$ est nul en général
et que l'homomorphisme
$H^0(X, \mathcal{O}_X) \to H^0(X', \mathcal{O}_{X'})$ est surjectif en général.

\medskip\noindent
Ensuite, dans le cas général, étudions le noyau. Si
$a \in A$ s'envoie sur zéro, alors, en examinant les cartes affines,
nous voyons que
$$
a = (f_1x - f_2)(a_0 + a_1x + \ldots + a_rx^r)\text{ dans }A[x]
$$
pour certains $r \geq 0$ et $a_0, \ldots, a_r \in A$, et de même
$$
a = (f_1 - f_2y)(b_0 + b_1y + \ldots + b_s y^s)\text{ dans }A[y]
$$
pour certains $s \geq 0$ et $b_0, \ldots, b_s \in A$. Cela signifie que nous
avons
$$
a = f_2 a_0,\ f_1 a_0 = f_2 a_1,\ \ldots,\ f_1 a_r = 0,
\ a = f_1 b_0,\ f_2 b_0 = f_1 b_1,\ \ldots,\ f_2 b_s = 0
$$
Si $(a', r', a'_i, s', b'_j)$ est un second système de ce type, alors nous avons
$$
aa' = f_1f_2a_0b'_0 = f_1f_2a_1b'_1 = f_1f_2a_2b'_2 = \ldots = 0
$$
comme souhaité.
\end{proof}

\noindent
Pour une $\mathbf{F}_p$-algèbre $A$,
nous posons $\colim_F A$ égal à la colimite du système
$$
A \xrightarrow{F} A \xrightarrow{F} A \xrightarrow{F} \ldots
$$
où $F : A \to A$, $a \mapsto a^p$ est l'endomorphisme de Frobenius.

\begin{lemma}
\label{lemma-h-sheaf-colim-F}
Soit $p$ un nombre premier. Soit $S$ un schéma sur $\mathbf{F}_p$.
Soit $(\Sch/S)_h$ un site comme dans la définition \ref{definition-big-small-h}.
Il existe un unique faisceau $\mathcal{F}$ sur $(\Sch/S)_h$ tel que
$$
\mathcal{F}(X) = \colim_F \Gamma(X, \mathcal{O}_X)
$$
pour tout objet quasi-compact et quasi-séparé $X$ de $(\Sch/S)_h$.
\end{lemma}

\begin{proof}
Notons $\mathcal{F}$ la faisceautisation de Zariski
du foncteur
$$
X \longrightarrow \colim_F \Gamma(X, \mathcal{O}_X)
$$
Pour les schémas quasi-compacts et quasi-séparés $X$,
nous avons\newline $\mathcal{F}(X) = \colim_F \Gamma(X, \mathcal{O}_X)$
par Faisceaux, lemme \ref{sheaves-lemma-directed-colimits-sections}
et le fait que $\mathcal{O}$ est un faisceau pour la topologie de Zariski.
Il suffit donc de montrer que $\mathcal{F}$ est un faisceau h.
Pour le démontrer, nous vérifions les conditions (1), (2), (3) et (4) du
lemme \ref{lemma-refine-check-h}.
La condition (1) est satisfaite parce que nous avons effectué une faisceautisation
de Zariski (presque superflue). La condition (2) est satisfaite parce que
$\mathcal{O}$ est un faisceau fppf (Descente, lemme
\ref{descent-lemma-sheaf-condition-holds}) et que, si
$A$ est l'égalisateur de deux homomorphismes $B \to C$ de $\mathbf{F}_p$-algèbres,
alors $\colim_F A$ est l'égalisateur des deux homomorphismes
$\colim_F B \to \colim_F C$.

\medskip\noindent
Nous vérifions la condition (3). Soient $A, f, J$ comme dans
l'exemple \ref{example-one-generator}.
Nous devons montrer que
$$
\colim_F A = \colim_F A/J \times_{\colim_F A/fA + J} \colim_F A/fA
$$
Cela se ramène à la question d'algèbre suivante : supposons que $a', a'' \in A$
soient tels que $F^n(a' - a'') \in fA + J$. Trouver $a \in A$ et $m \geq 0$
tels que $a - F^m(a') \in J$ et $a - F^m(a'') \in fA$, et montrer que
la paire $(a, m)$ est uniquement déterminée à un remplacement près de la
forme $(a, m) \mapsto (F(a), m + 1)$.
Pour cela, il suffit d'écrire $F^n(a' - a'') = f h + g$ avec $h \in A$ et $g \in J$,
de poser $a = F^n(a') - g = F^n(a'') + fh$ et de poser $m = n$.
Pour voir l'unicité, supposons que $(a_1, m_1)$ soit une seconde solution.
Par un remplacement de la forme donnée ci-dessus, nous pouvons supposer $m = m_1$.
Alors nous voyons que $a - a_1 \in J$ et $a - a_1 \in fA$.
Puisque $J$ est annulé par une puissance de $f$, nous voyons que
$a - a_1$ est un élément nilpotent. Par conséquent, $F^k(a - a_1)$ est nul
pour un certain $k$ assez grand. Ainsi, après d'autres remplacements, nous obtenons
$a = a_1$.

\medskip\noindent
Nous vérifions la condition (4). Soient $X, X', Z, E$ comme dans
l'exemple \ref{example-two-generators}. Par le
lemme \ref{lemma-funny-blow-up}, nous voyons que
$$
\mathcal{F}(X) = \colim_F \Gamma(X, \mathcal{O}_X)
\longrightarrow
\colim_F \Gamma(X', \mathcal{O}_{X'}) = \mathcal{F}(X')
$$
est bijectif. Puisque $E = \mathbf{P}^1_Z$ dans ce cas, nous
voyons aussi que $\mathcal{F}(Z) \to \mathcal{F}(E)$ est bijectif.
La conclusion vaut donc également dans ce cas.
\end{proof}

\noindent
Soit $p$ un nombre premier. Pour une $\mathbf{F}_p$-algèbre $A$,
nous posons $\lim_F A$ égal à la limite du système inverse
$$
\ldots \xrightarrow{F} A \xrightarrow{F} A \xrightarrow{F} A
$$
où $F : A \to A$, $a \mapsto a^p$ est l'endomorphisme de Frobenius.

\begin{lemma}
\label{lemma-h-sheaf-lim-F}
Soit $p$ un nombre premier. Soit $S$ un schéma sur $\mathbf{F}_p$.
Soit $(\Sch/S)_h$ un site comme dans la définition \ref{definition-big-small-h}.
La règle
$$
\mathcal{F}(X) = \lim_F \Gamma(X, \mathcal{O}_X)
$$
définit un faisceau sur $(\Sch/S)_h$.
\end{lemma}

\begin{proof}
Pour montrer que $\mathcal{F}$ est un faisceau, vérifions les conditions
(1), (2), (3) et (4) du lemme \ref{lemma-refine-check-h}.
La condition (1) est satisfaite parce que les limites de faisceaux sont des faisceaux
et que $\mathcal{O}$ est un faisceau de Zariski. La condition (2) est satisfaite parce que
$\mathcal{O}$ est un faisceau fppf (Descente, lemme
\ref{descent-lemma-sheaf-condition-holds}) et que, si $A$ est l'égalisateur
de deux homomorphismes $B \to C$ de $\mathbf{F}_p$-algèbres, alors $\lim_F A$
est l'égalisateur des deux homomorphismes $\lim_F B \to \lim_F C$.

\medskip\noindent
Nous vérifions la condition (3). Soient $A, f, J$ comme dans
l'exemple \ref{example-one-generator}.
Nous devons montrer que
\begin{align*}
\lim_F A
& \to
\lim_F A/J \times_{\lim_F A/fA + J} \lim_F A/fA \\
& =
\lim_F (A/J \times_{A/fA + J} A/fA) \\
& =
\lim_F A/(fA \cap J)
\end{align*}
est bijectif. Puisque $J$ est annulé par une puissance de $f$, nous voyons que
$\mathfrak a = fA \cap J$ est un idéal nilpotent, c'est-à-dire qu'il existe un
$n$ tel que $\mathfrak a^n = 0$. Il est immédiat
de vérifier que, dans ce cas, $\lim_F A \to \lim_F A/\mathfrak a$
est bijectif.

\medskip\noindent
Nous vérifions la condition (4). Soient $X, X', Z, E$ comme dans
l'exemple \ref{example-two-generators}. Par le
lemme \ref{lemma-funny-blow-up} et le même argument que ci-dessus,
nous voyons que
$$
\mathcal{F}(X) = \lim_F \Gamma(X, \mathcal{O}_X)
\longrightarrow
\lim_F \Gamma(X', \mathcal{O}_{X'}) = \mathcal{F}(X')
$$
est bijectif. Puisque $E = \mathbf{P}^1_Z$ dans ce cas, nous
voyons aussi que $\mathcal{F}(Z) \to \mathcal{F}(E)$ est bijectif.
La conclusion vaut donc également dans ce cas.
\end{proof}

\noindent
Dans le lemme suivant, nous utilisons la normalisation faible absolue $X^{awn}$
d'un schéma $X$, voir
Morphismes, section \ref{morphisms-section-seminormalization}.

\begin{lemma}
\label{lemma-weak-normalization-ph-sheaf}
Soit $(\Sch/S)_{ph}$ un site comme dans
Topologies, définition \ref{topologies-definition-big-small-ph}.
La règle
$$
X \longmapsto \Gamma(X^{awn}, \mathcal{O}_{X^{awn}})
$$
est un faisceau sur $(\Sch/S)_{ph}$.
\end{lemma}

\begin{proof}
Pour montrer que $\mathcal{F}$ est un faisceau, vérifions les conditions
(1) et (2) de Topologies, lemme \ref{topologies-lemma-characterize-sheaf}.
La condition (1) est satisfaite parce que la formation de $X^{awn}$ commute aux
recouvrements ouverts, voir Morphismes, lemme \ref{morphisms-lemma-seminormalization}
et sa démonstration.

\medskip\noindent
Soit $\pi : Y \to X$ un morphisme propre surjectif. Nous devons montrer
que l'égalisateur des deux homomorphismes
$$
\Gamma(Y^{awn}, \mathcal{O}_{Y^{awn}}) \to
\Gamma((Y \times_X Y)^{awn}, \mathcal{O}_{(Y \times_X Y)^{awn}})
$$
est égal à $\Gamma(X^{awn}, \mathcal{O}_{X^{awn}})$. Soit $f$ un
élément de cet égalisateur. Considérons alors le morphisme
$$
f : Y^{awn} \longrightarrow \mathbf{A}^1_X
$$
Puisque $Y^{awn} \to X$ est universellement fermé, l'image schématique
$Z$ de $f$ est un sous-schéma fermé de $\mathbf{A}^1_X$, propre sur $X$,
et $f : Y^{awn} \to Z$ est surjectif.
Voir Morphismes, lemme \ref{morphisms-lemma-scheme-theoretic-image-is-proper}.
Ainsi, $Z \to X$ est fini
(Morphismes, lemme \ref{morphisms-lemma-finite-proper})
et surjectif.

\medskip\noindent
Soit $k$ un corps et soient $z_1, z_2 : \Spec(k) \to Z$ deux
morphismes égalisés par $Z \to X$. Nous affirmons que $z_1 = z_2$.
Il suffit de montrer que les images $\lambda_i = z_i^*f \in k$ coïncident
(puisque le faisceau structural de $Z$ est engendré par $f$
sur le faisceau structural de $X$). Pour le voir, nous
choisissons une extension de corps
$K/k$ et des morphismes $y_1, y_2 : \Spec(K) \to Y^{awn}$ tels que
$z_i \circ (\Spec(K) \to \Spec(k)) = f \circ y_i$. Cela est possible
par la surjectivité du morphisme $Y^{awn} \to Z$. Choisissons une extension
algébriquement close $\Omega/k$ de très grand cardinal.
Pour tous les homomorphismes de $k$-algèbres $\sigma_i : K \to \Omega$,
nous obtenons
$$
\Spec(\Omega)
\xrightarrow{\sigma_1, \sigma_2}
\Spec(K \otimes_k K)
\xrightarrow{y_1, y_2}
Y^{awn} \times_X Y^{awn}
$$
Puisque le morphisme canonique
$(Y \times_X Y)^{awn} \to Y^{awn} \times_X Y^{awn}$
est un homéomorphisme universel et puisque $\Omega$ est algébriquement clos,
nous pouvons relever de manière unique la composition ci-dessus en un morphisme
$\Spec(\Omega) \to (Y \times_X Y)^{awn}$. Puisque $f$ appartient à l'égalisateur
ci-dessus, cela montre que $\sigma_1(\lambda_1) = \sigma_2(\lambda_2)$.
Un lemme facile sur les extensions de corps montre que cela implique
$\lambda_1 = \lambda_2$ ; les détails sont omis.

\medskip\noindent
Nous concluons que $Z \to X$ est universellement injectif, c'est-à-dire que
$Z \to X$ est injectif sur les points et induit des
extensions de corps résiduels purement inséparables
(Morphismes, lemme \ref{morphisms-lemma-universally-injective}).
En définitive, nous concluons que $Z \to X$ est un homéomorphisme universel, voir
Morphismes, lemme \ref{morphisms-lemma-universal-homeomorphism}.

\medskip\noindent
Soit $g : X^{awn} \to Z$ le morphisme obtenu par la propriété universelle
de $X^{awn}$. Alors $Y^{awn} \to X^{awn} \to Z$ et $f : Y^{awn} \to Z$
sont deux morphismes au-dessus de $X$. Par la propriété universelle de
$Y^{awn} \to Y$, les deux morphismes correspondants
$Y^{awn} \to Y \times_X Z$ au-dessus de $Y$ doivent être égaux. Cela implique
que $g \circ \pi^{wan} = f$ comme morphismes vers $\mathbf{A}^1_X$,
et nous concluons que $g \in \Gamma(X^{awn}, \mathcal{O}_{X^{awn}})$
est l'élément que nous cherchions.
\end{proof}

\begin{lemma}
\label{lemma-weak-normalization-h-sheaf}
Soit $S$ un schéma. Choisissons un site $(\Sch/S)_h$
comme dans la définition \ref{definition-big-small-h}.
La règle
$$
X \longmapsto \Gamma(X^{awn}, \mathcal{O}_{X^{awn}})
$$
est la faisceautisation du « faisceau structural » $\mathcal{O}$
sur $(\Sch/S)_h$. Il en va de même pour la topologie ph.
\end{lemma}

\begin{proof}
Dans le lemme \ref{lemma-weak-normalization-ph-sheaf},
nous avons vu que la règle $\mathcal{F}$ du lemme
définit un faisceau pour la topologie ph, et donc, a fortiori,
un faisceau pour la topologie h. Il existe clairement un morphisme canonique
de préfaisceaux d'anneaux $\mathcal{O} \to \mathcal{F}$.
Pour terminer la démonstration, il suffit de montrer
\begin{enumerate}
\item si $f \in \mathcal{O}(X)$ s'envoie sur zéro dans $\mathcal{F}(X)$,
alors il existe un recouvrement h $\{X_i \to X\}$ tel que $f|_{X_i} = 0$, et
\item étant donné $f \in \mathcal{F}(X)$, il existe un recouvrement h
$\{X_i \to X\}$ tel que $f|_{X_i}$ soit l'image de $f_i \in \mathcal{O}(X_i)$.
\end{enumerate}
Soit $f$ comme en (1). Alors $f|_{X^{awn}} = 0$. Cela signifie que $f$
est localement nilpotent. Ainsi, si $X' \subset X$ est le sous-schéma fermé
découpé par $f$, alors $X' \to X$ est une immersion fermée surjective de
présentation finie. Par conséquent, $\{X' \to X\}$ est le recouvrement h recherché.
Soit $f$ comme en (2). Après avoir remplacé $X$ par les membres d'un
recouvrement ouvert affine, nous pouvons supposer que $X = \Spec(A)$ est affine.
Alors $f \in A^{awn}$, voir
Morphismes, lemme \ref{morphisms-lemma-seminormalization-ring}.
Par Morphismes, lemme \ref{morphisms-lemma-universal-homeo-limit},
nous pouvons trouver un homomorphisme d'anneaux $A \to B$ de présentation finie
tel que $\Spec(B) \to \Spec(A)$ soit un homéomorphisme universel
et que $f$ soit l'image d'un élément $b \in B$ par
l'homomorphisme canonique $B \to A^{awn}$. Alors
$\{\Spec(B) \to \Spec(A)\}$ est un recouvrement h, et nous concluons.
L'assertion concernant la topologie ph se démontre de la même manière
(ou peut se déduire de l'assertion pour la topologie h).
\end{proof}

\noindent
Soit $p$ un nombre premier. Une $\mathbf{F}_p$-algèbre $A$ est
dite {\it parfaite} si l'homomorphisme $F : A \to A$, $x \mapsto x^p$,
est un automorphisme de $A$.

\begin{lemma}
\label{lemma-perfect-weankly-normal}
Soit $p$ un nombre premier. Une $\mathbf{F}_p$-algèbre $A$ est
absolument faiblement normale si et seulement si elle est parfaite.
\end{lemma}

\begin{proof}
Il résulte immédiatement de la condition (2)(b)\newline dans
Morphismes, définition \ref{morphisms-definition-seminormal-ring}
que, si $A$ est absolument faiblement normale, alors elle est parfaite.

\medskip\noindent
Supposons $A$ parfaite. Soient $x, y \in A$ tels que $x^3 = y^2$.
Si $p > 3$, alors nous pouvons écrire $p = 2n + 3m$ pour certains $n, m > 0$.
Choisissons $a, b \in A$ tels que $a^p = x$ et $b^p = y$.
En posant $c = a^n b^m$, nous avons
$$
c^{2p} = x^{2n} y^{2m} = x^{2n + 3m} = x^p
$$
et donc $c^2 = x$. De même, $c^3 = y$. Si $p = 2$, alors
écrivons $x = a^2$ pour obtenir $a^6 = y^2$, ce qui implique $a^3 = y$.
Si $p = 3$, alors écrivons $y = a^3$ pour obtenir $x^3 = a^6$, ce qui
implique $x = a^2$.

\medskip\noindent
Supposons $x, y \in A$ avec $\ell^\ell x = y^\ell$ pour un certain nombre
premier $\ell$. Si $\ell \not = p$, alors $a = y/\ell$ vérifie
$a^\ell = x$ et $\ell a = y$. Si $\ell = p$, alors
$y = 0$ et $x = a^p$ pour un certain $a$.
\end{proof}

\begin{lemma}
\label{lemma-char-p}
Soit $p$ un nombre premier.
\begin{enumerate}
\item Si $A$ est une $\mathbf{F}_p$-algèbre, alors $\colim_F A = A^{awn}$.
\item Si $S$ est un schéma sur $\mathbf{F}_p$, alors la
faisceautisation h de $\mathcal{O}$ associe à tout schéma quasi-compact
et quasi-séparé $X$ l'anneau $\colim_F \Gamma(X, \mathcal{O}_X)$.
\end{enumerate}
\end{lemma}

\begin{proof}
Démonstration de (1). Observons que $A \to \colim_F A$ induit un
homéomorphisme universel sur les spectres par
Algèbre, lemme \ref{algebra-lemma-p-ring-map}.
Il suffit donc de montrer que $B = \colim_F A$ est
absolument faiblement normal, voir
Morphismes, lemme \ref{morphisms-lemma-seminormalization-ring}.
Notons que l'homomorphisme d'anneaux $F : B \to B$ est un automorphisme ;
autrement dit, $B$ est un anneau parfait. Par conséquent, le
lemme \ref{lemma-perfect-weankly-normal} s'applique.

\medskip\noindent
Démonstration de (2). Cela résulte de (1) et des
lemmes \ref{lemma-h-sheaf-colim-F} et
\ref{lemma-weak-normalization-h-sheaf}
en raisonnant localement sur des ouverts affines.
\end{proof}






\section{Descente des fibrés vectoriels en caractéristique positive}
\label{section-descent-vectorbundles-h}

\noindent
Une référence pour cette section est \cite{Witt-Grass}.

\medskip\noindent
Pour un schéma $S$, notons $\textit{Vect}(S)$ la catégorie
des $\mathcal{O}_S$-modules localement libres de rang fini.
Soit $p$ un nombre premier. Soit $S$ un schéma quasi-compact et quasi-séparé
sur $\mathbf{F}_p$. Dans cette section, nous travaillerons avec la
catégorie
$$
\colim_F \textit{Vect}(S) =
\colim
\left(
\textit{Vect}(S) \xrightarrow{F^*}
\textit{Vect}(S) \xrightarrow{F^*}
\textit{Vect}(S) \xrightarrow{F^*} \ldots
\right)
$$
où $F : S \to S$ est le morphisme de Frobenius absolu. Concrètement,
un objet de cette catégorie est une paire $(\mathcal{E}, n)$,
où $\mathcal{E}$ est un $\mathcal{O}_S$-module localement libre de rang fini
et $n \geq 0$ est un entier. Pour les morphismes, nous prenons
$$
\Hom_{\colim_F \textit{Vect}(S)}((\mathcal{E}, n), (\mathcal{G}, m)) =
\colim_N \Hom_S(F^{N - n, *}\mathcal{E}, F^{N - m, *}\mathcal{G})
$$
où $F : S \to S$ est le morphisme de Frobenius absolu de $S$.
Ainsi, l'objet $(\mathcal{E}, n)$ est isomorphe à l'objet
$(F^*\mathcal{E}, n + 1)$.

\begin{lemma}
\label{lemma-colim-F-Vect}
Soit $p$ un nombre premier. Soit $S$ un schéma quasi-compact et quasi-séparé
sur $\mathbf{F}_p$. La catégorie $\colim_F \textit{Vect}(S)$
est équivalente à la catégorie des modules localement libres de rang fini
sur le faisceau d'anneaux $\colim_F \mathcal{O}_S$ sur $S$.
\end{lemma}

\begin{proof}
Omis.
\end{proof}

\begin{lemma}
\label{lemma-vector-bundle-I}
Soit $p$ un nombre premier. Considérons un carré de quasi-éclatement $X, X', Z, E$
en caractéristique $p$ comme dans l'exemple \ref{example-one-generator}.
Alors le foncteur
$$
\colim_F \textit{Vect}(X)
\longrightarrow
\colim_F \textit{Vect}(Z)
\times_{\colim_F \textit{Vect}(E)}
\colim_F \textit{Vect}(X')
$$
est une équivalence.
\end{lemma}

\begin{proof}
Soient $A, f, J$ comme dans l'exemple \ref{example-one-generator}.
Puisque tous nos schémas sont affines et que nous disposons du
Hom interne dans la catégorie des fibrés vectoriels, la pleine fidélité
du foncteur résulte du fait que
$$
\colim P \otimes_{A, F^N} A =
\colim P \otimes_{A, F^N} A/J
\times_{\colim P \otimes_{A, F^N} A/fA + J}
\colim P \otimes_{A, F^N} A/fA
$$
pour un $A$-module projectif de type fini $P$. En écrivant $P$ comme facteur direct
d'un module libre fini, cela résulte du cas où $P$ est libre
fini. Ce cas se ramène immédiatement au cas $P = A$. Le cas
$P = A$ résulte du lemme \ref{lemma-h-sheaf-colim-F}
(en fait, nous avons démontré ce cas directement dans la démonstration de ce lemme).

\medskip\noindent
Surjectivité essentielle. Nous obtenons ici le problème d'algèbre suivant.
Supposons que $P_1$ soit un $A/J$-module projectif de type fini,
que $P_2$ soit un $A/fA$-module projectif de type fini, et que
$$
\varphi :
P_1 \otimes_{A/J} A/fA + J
\longrightarrow
P_2 \otimes_{A/fA} A/fA + J
$$
soit un isomorphisme. Objectif : montrer qu'il existe un $N$, un
$A$-module projectif de type fini $P$, un isomorphisme
$\varphi_1 : P \otimes_A A/J \to P_1 \otimes_{A/J, F^N} A/J$,
et un isomorphisme
$\varphi_2 : P \otimes_A A/fA \to P_2 \otimes_{A/fA, F^N} A/fA$
compatible avec $\varphi$ de manière évidente.
On peut le voir comme suit. Observons d'abord que
$$
A/(J \cap fA) = A/J \times_{A/fA + J} A/fA
$$
Par conséquent, par Compléments sur l'algèbre, lemme
\ref{more-algebra-lemma-finitely-presented-module-over-fibre-product},
il existe un module projectif de type fini $P'$ sur
$A/(J \cap fA)$ muni d'isomorphismes
$\varphi'_1 : P' \otimes_A A/J \to P_1$ et
$\varphi_2 : P' \otimes_A A/fA \to P_2$
compatibles avec $\varphi$. Puisque $J$ est un idéal de type fini et
de torsion pour une puissance de $f$, nous voyons que $J \cap fA$ est un
idéal nilpotent. Ainsi, pour un certain $N$, il existe une factorisation
$$
A \xrightarrow{\alpha} A/(J \cap fA) \xrightarrow{\beta} A
$$
de $F^N$. En posant $P = P' \otimes_{A/(J \cap fA), \beta} A$,
nous concluons.
\end{proof}

\begin{lemma}
\label{lemma-vector-bundle-II}
Soit $p$ un nombre premier. Considérons un carré de quasi-éclatement $X, X', Z, E$
en caractéristique $p$ comme dans l'exemple \ref{example-two-generators}.
Alors le foncteur
$$
G :
\colim_F \textit{Vect}(X)
\longrightarrow
\colim_F \textit{Vect}(Z)
\times_{\colim_F \textit{Vect}(E)}
\colim_F \textit{Vect}(X')
$$
est une équivalence.
\end{lemma}

\begin{proof}
Pleine fidélité. Supposons que $(\mathcal{E}, n)$ et
$(\mathcal{F}, m)$ soient des objets de $\colim_F \textit{Vect}(X)$.
Soit $(a, b) : G(\mathcal{E}, n) \to G(\mathcal{F}, m)$
un morphisme dans le membre de droite. Nous pouvons choisir $N \gg 0$ et
considérer $a$ comme un homomorphisme
$a : F^{N - n, *}\mathcal{E}|_Z \to F^{N - m, *}\mathcal{F}|_Z$
et $b$ comme un homomorphisme
$b : F^{N - n, *}\mathcal{E}|_{X'} \to F^{N - m, *}\mathcal{F}|_{X'}$
qui coïncident sur $E$.
Choisissons un recouvrement fini par des ouverts affines
$X = X_1 \cup \ldots \cup X_n$ tel que $\mathcal{E}|_{X_i}$
et $\mathcal{F}|_{X_i}$ soient des $\mathcal{O}_{X_i}$-modules libres de rang fini.
Pour chaque $i$, le changement de base
$$
\xymatrix{
E_i \ar[r] \ar[d] & X'_i \ar[d] \\
Z_i \ar[r] & X_i
}
$$
est un autre carré de quasi-éclatement comme dans l'exemple \ref{example-two-generators}.
Pour ces carrés, nous savons que
$$
\colim_F H^0(X_i, \mathcal{O}_{X_i}) =
\colim_F H^0(Z_i, \mathcal{O}_{Z_i})
\times_{\colim_F H^0(E_i, \mathcal{O}_{E_i})}
\colim_F H^0(X'_i, \mathcal{O}_{X'_i})
$$
par le lemme \ref{lemma-h-sheaf-colim-F} (voir la démonstration du lemme).
Ainsi, après avoir augmenté $N$, nous pouvons supposer que
les homomorphismes $a|_{Z_i}$ et $b|_{X'_i}$ proviennent
d'homomorphismes $c_i : F^{N - n, *}\mathcal{E}|_{X_i} \to F^{N - m, *}\mathcal{F}|_{X_i}$.
Après avoir éventuellement augmenté $N$, nous pouvons supposer que $c_i$ et $c_j$
coïncident sur $X_i \cap X_j$. Ces homomorphismes se recollent donc pour donner le
morphisme recherché $(\mathcal{E}, n) \to (\mathcal{F}, m)$
dans le membre de gauche.

\medskip\noindent
Surjectivité essentielle. Soit $(\mathcal{F}, \mathcal{G}, \varphi)$ un
triplet constitué de
un $\mathcal{O}_Z$-module localement libre de rang fini $\mathcal{F}$,
un $\mathcal{O}_{X'}$-module localement libre de rang fini $\mathcal{G}$, et
un isomorphisme $\varphi : \mathcal{F}|_E \to \mathcal{G}|_E$.
Nous devons montrer qu'après avoir remplacé ce triplet par un tiré en arrière par une
puissance du Frobenius, il provient d'un $\mathcal{O}_X$-module localement libre de rang fini.

\medskip\noindent
Réduction au cas noethérien ; nous conseillons vivement au lecteur de sauter ce paragraphe.
Rappelons que\newline $X = \Spec(A)$ et $Z = \Spec(A/(f_1, f_2))$,
$X' = \text{Proj}(A[T_0, T_1]/(f_2T_0 - f_1T_1))$, et
$E = \mathbf{P}^1_Z$. Par Limites, lemme
\ref{limits-lemma-descend-invertible-modules},
nous pouvons trouver une $\mathbf{F}_p$-sous-algèbre de type fini
$A_0 \subset A$ contenant $f_1$ et $f_2$ telle que le triplet
$(\mathcal{F}, \mathcal{G}, \varphi)$ provienne de données définies sur
$X_0 = \Spec(A_0)$ et $Z_0 = \Spec(A_0/(f_1, f_2))$,
$X_0' = \text{Proj}(A_0[T_0, T_1]/(f_2T_0 - f_1T_1))$, et
$E_0 = \mathbf{P}^1_{Z_0}$. Nous pouvons donc supposer que nos schémas
sont noethériens.

\medskip\noindent
Supposons $X$ noethérien. Nous pouvons choisir un recouvrement fini par des ouverts affines
$X = X_1 \cup \ldots \cup X_n$ tel que $\mathcal{F}|_{Z \cap X_i}$ soit libre.
Puisque nous pouvons recoller les objets de $\colim_F \textit{Vect}(X)$
pour la topologie de Zariski (lemme \ref{lemma-colim-F-Vect}), et
puisque nous connaissons déjà la
pleine fidélité sur $X_i$ et $X_i \cap X_j$ (voir le premier paragraphe
de la démonstration), il suffit de démontrer l'existence sur chaque $X_i$.
Cela nous ramène au cas traité dans le paragraphe suivant.

\medskip\noindent
Supposons $X$ noethérien et $\mathcal{F} = \mathcal{O}_Z^{\oplus r}$.
En utilisant $\varphi$, nous obtenons un isomorphisme
$\mathcal{O}_E^{\oplus r} \to \mathcal{G}|_E$.
Soit $I = (f_1, f_2) \subset A$.
Soit $\mathcal{I} \subset \mathcal{O}_{X'}$
le faisceau d'idéaux de $E$ ; il est globalement engendré par $f_1$ et $f_2$.
Pour tout $n$, il existe une surjection
$$
(\mathcal{I}^n/\mathcal{I}^{n + 1})^{\oplus r} =
\mathcal{I}^n/\mathcal{I}^{n + 1} \otimes_{\mathcal{O}_E}
\mathcal{G}|_E \longrightarrow
\mathcal{I}^n\mathcal{G}/\mathcal{I}^{n + 1}\mathcal{G}
$$
Le premier groupe de cohomologie de ce module est donc nul.
Nous utilisons ici le fait que $E = \mathbf{P}^1_Z$ et que, par conséquent, son faisceau
structural, et même tout module quasi-cohérent globalement engendré,
a un $H^1$ nul. Comparer avec Compléments sur les morphismes, lemme
\ref{more-morphisms-lemma-globally-generated-vanishing}.
Alors, en utilisant les suites exactes courtes
$$
0 \to  \mathcal{I}^n\mathcal{G}/\mathcal{I}^{n + 1}\mathcal{G} \to
\mathcal{G}/\mathcal{I}^{n + 1}\mathcal{G} \to
\mathcal{G}/\mathcal{I}^n\mathcal{G} \to 0
$$
et en raisonnant par récurrence, nous voyons que
$$
\lim H^0(X', \mathcal{G}/\mathcal{I}^n\mathcal{G})
\to
H^0(E, \mathcal{G}|_E) = H^0(E, \mathcal{O}_E^{\oplus r}) =
A/I^{\oplus r}
$$
est surjectif. Par le théorème des fonctions formelles
(Cohomologie des schémas, théorème \ref{coherent-theorem-formal-functions}),
cela implique que
$$
H^0(X', \mathcal{G}) \to
H^0(E, \mathcal{G}|_E) = H^0(E, \mathcal{O}_E^{\oplus r}) =
A/I^{\oplus r}
$$
est surjectif. Nous pouvons donc choisir un homomorphisme
$\alpha : \mathcal{O}_{X'}^{\oplus r} \to \mathcal{G}$
compatible avec la trivialisation donnée
de $\mathcal{G}|_E$. Ainsi, $\alpha$ est un isomorphisme sur
un voisinage ouvert de $E$ dans $X'$. Tout point
de $Z$ possède donc un voisinage ouvert affine sur lequel
nous pouvons résoudre le problème. Puisque $X' \setminus E \to X \setminus Z$
est un isomorphisme, il en va de même pour les points de $X$ qui ne sont pas dans $Z$.
Un autre argument de recollement de Zariski achève donc la démonstration.
\end{proof}

\begin{proposition}
\label{proposition-h-descent-vector-bundles-p}
Soit $p$ un nombre premier. Soit $S$ un schéma de caractéristique $p$.
Alors la catégorie fibrée en groupoïdes
$$
p : \mathcal{S} \longrightarrow (\Sch/S)_h
$$
dont la catégorie fibre au-dessus de $U$ est la catégorie
des $\colim_F \mathcal{O}_U$-modules localement libres de rang fini sur $U$
est un champ en groupoïdes. De plus, si $U$ est quasi-compact
et quasi-séparé, alors $\mathcal{S}_U$ est $\colim_F \textit{Vect}(U)$.
\end{proposition}

\begin{proof}
La dernière assertion est le contenu du lemme \ref{lemma-colim-F-Vect}.
Pour démontrer la proposition, nous vérifierons les conditions (1), (2) et (3) du
lemme \ref{lemma-refine-check-h-stack}.

\medskip\noindent
La condition (1) est satisfaite parce que, par définition, nous avons le recollement
pour la topologie de Zariski.

\medskip\noindent
Pour voir la condition (2), supposons que $f : X \to Y$ soit un morphisme
plat, propre, surjectif et de présentation finie sur $S$, avec $Y$ affine.
Puisque $Y, X, X \times_Y X$ sont quasi-compacts et quasi-séparés,
nous pouvons utiliser la description des catégories fibres donnée dans
l'énoncé de la proposition. Il suffit alors clairement de
montrer que
$$
\textit{Vect}(Y)
\longrightarrow
\textit{Vect}(X)
\times_{\textit{Vect}(X \times_Y X)}
\textit{Vect}(X)
$$
est une équivalence (car cela impliquera la même chose pour les colimites).
Cela résulte immédiatement de la descente fppf des modules localement libres de rang fini, voir
Descente, proposition \ref{descent-proposition-fpqc-descent-quasi-coherent} et
lemme \ref{descent-lemma-finite-locally-free-descends}.

\medskip\noindent
La condition (3) est le contenu des
lemmes \ref{lemma-vector-bundle-I} et \ref{lemma-vector-bundle-II}.
\end{proof}

\begin{lemma}
\label{lemma-trivial-fibres-dvr}
Soit $f : X \to S$ un morphisme propre à fibres géométriquement connexes,
où $S$ est le spectre d'un anneau de valuation discrète. Notons $\eta \in S$
le point générique et notons $X_n \subset X$ le sous-schéma fermé
découpé par la $n$-ième puissance d'une uniformisante sur $S$.
Il existe alors
un entier $n$ tel que l'assertion suivante soit vraie : tout
$\mathcal{O}_X$-module localement libre de rang fini $\mathcal{E}$
tel que $\mathcal{E}|_{X_\eta}$ et $\mathcal{E}|_{X_n}$
soient libres, est libre.
\end{lemma}

\begin{proof}
Nous nous ramenons d'abord au cas où $X \to S$ possède une section. Écrivons $S = \Spec(A)$.
Choisissons un point fermé $\xi$ de $X_\eta$. Choisissons une extension
d'anneaux de valuation discrète $A \subset B$ telle que le corps des fractions
de $B$ soit $\kappa(\xi)$. Cela est possible par Krull--Akizuki
(Algèbre, lemme \ref{algebra-lemma-integral-closure-Dedekind})
et par le fait que $\kappa(\xi)$ est une extension finie du
corps des fractions de $A$.
Par le critère valuatif de propreté
(Morphismes, lemme \ref{morphisms-lemma-characterize-proper}),
nous obtenons un point à valeurs dans $B$, $\tau : \Spec(B) \to X$,
qui induit une section $\sigma : \Spec(B) \to X_B$.
Pour un $\mathcal{O}_X$-module localement libre de rang fini $\mathcal{E}$,
notons $\mathcal{E}_B$ son tiré en arrière sur $X_B$.
Par changement de base plat
(Cohomologie des schémas, lemme \ref{coherent-lemma-flat-base-change-cohomology}),
nous voyons que $H^0(X_B, \mathcal{E}_B) = H^0(X, \mathcal{E}) \otimes_A B$.
Ainsi, si $\mathcal{E}_B$ est libre de rang $r$, alors les sections de
$H^0(X, \mathcal{E})$ engendrent le $B$-module libre
$\tau^*\mathcal{E} = \sigma^*\mathcal{E}_B$.
En particulier, nous pouvons trouver $r$ sections globales $s_1, \ldots, s_r$
de $\mathcal{E}$ qui engendrent $\tau^*\mathcal{E}$. Alors
$$
s_1, \ldots, s_r :
\mathcal{O}_X^{\oplus r}
\longrightarrow
\mathcal{E}
$$
est un homomorphisme de $\mathcal{O}_X$-modules localement libres de rang $r$,
et son tiré en arrière sur $X_B$ est un homomorphisme de $\mathcal{O}_{X_B}$-modules libres
dont la restriction est un isomorphisme en un point de chaque fibre.
En prenant le déterminant, nous obtenons une fonction
$g \in \Gamma(X_\eta, \mathcal{O}_{X_B})$
qui est inversible en un point de chaque fibre.
Comme les fibres sont propres et connexes, nous voyons que $g$
doit être inversible (détails omis ; indication : utiliser Variétés, lemme
\ref{varieties-lemma-proper-geometrically-reduced-global-sections}).
Il suffit donc de démontrer le lemme pour le changement de base $X_B \to \Spec(B)$.

\medskip\noindent
Supposons que nous ayons une section $\sigma : S \to X$. Soit $\mathcal{E}$
un $\mathcal{O}_X$-module localement libre de rang fini que nous supposons
libre sur la fibre générique et sur $X_n$ (nous choisirons $n$ plus tard).
Choisissons un isomorphisme $\sigma^*\mathcal{E} = \mathcal{O}_S^{\oplus r}$.
Considérons l'homomorphisme
$$
K = R\Gamma(X, \mathcal{E}) \longrightarrow
R\Gamma(S, \sigma^*\mathcal{E}) = A^{\oplus r}
$$
dans $D(A)$. En raisonnant comme ci-dessus, nous voyons que $\mathcal{E}$ est libre si (et seulement si)
l'homomorphisme induit $H^0(K) = H^0(X, \mathcal{E}) \to A^{\oplus r}$ est surjectif.

\medskip\noindent
Posons $L = R\Gamma(X, \mathcal{O}_X^{\oplus r})$ et observons que
l'homomorphisme correspondant $L \to A^{\oplus r}$ possède la propriété voulue.
Observons que $K \otimes_A Q(A) \cong L \otimes_A Q(A)$
par changement de base plat et par l'hypothèse que $\mathcal{E}$
est libre sur la fibre générique. Soit $\pi \in A$ une uniformisante. Observons que
$$
K \otimes_A^\mathbf{L} A/\pi^m A =
R\Gamma(X, \mathcal{E} \xrightarrow{\pi^m} \mathcal{E})
$$
et de même pour $L$.
Notons $\mathcal{E}_{tors} \subset \mathcal{E}$ le sous-faisceau cohérent des
sections supportées par la fibre spéciale, et procédons de même pour les autres
$\mathcal{O}_X$-modules. Choisissons $k > 0$ tel que
$(\mathcal{O}_X)_{tors} \to \mathcal{O}_X/\pi^k \mathcal{O}_X$
soit injectif (Cohomologie des schémas, lemme \ref{coherent-lemma-Artin-Rees}).
Puisque $\mathcal{E}$ est localement libre, nous voyons
que $\mathcal{E}_{tors} \subset \mathcal{E}/\pi^k\mathcal{E}$.
Alors, pour $n \geq m + k$, nous avons des isomorphismes
\begin{align*}
(\mathcal{E} \xrightarrow{\pi^m} \mathcal{E})
& \cong
(\mathcal{E}/\pi^k\mathcal{E} \xrightarrow{\pi^m}
\mathcal{E}/\pi^{k + m}\mathcal{E}) \\
& \cong
(\mathcal{O}_X^{\oplus r}/\pi^k\mathcal{O}_X^{\oplus r} \xrightarrow{\pi^m}
\mathcal{O}_X^{\oplus r}/\pi^{k + m}\mathcal{O}_X^{\oplus r}) \\
& \cong
(\mathcal{O}_X^{\oplus r} \xrightarrow{\pi^m} \mathcal{O}_X^{\oplus r})
\end{align*}
dans $D(\mathcal{O}_X)$. Cela détermine un isomorphisme
$$
K \otimes_A^\mathbf{L} A/\pi^m A \cong L \otimes_A^\mathbf{L} A/\pi^m A
$$
dans $D(A)$ (valable lorsque $n \geq m + k$). Observons que ces isomorphismes
sont compatibles au passage à l'image inverse par $\sigma$ ; en particulier,
nous concluons que
$K \otimes_A^\mathbf{L} A/\pi^m A \to (A/\pi^m A)^{\oplus r}$
induit une surjection sur les modules de cohomologie en degré $0$ (puisque
c'est vrai pour $L$).
Puisque $A$ est un anneau de valuation discrète, nous avons
$$
K \cong \bigoplus H^i(K)[-i]
\quad\text{et}\quad
L \cong
\bigoplus H^i(L)[-i]
$$
dans $D(A)$. Voir Compléments sur l'algèbre, exemple
\ref{more-algebra-example-finite-injective-finite-global-dimension}.
Les groupes de cohomologie $H^i(K) = H^i(X, \mathcal{E})$ et
$H^i(L) = H^i(X, \mathcal{O}_X)^{\oplus r}$
sont des $A$-modules de type fini d'après Cohomologie des schémas, lemme
\ref{coherent-lemma-proper-over-affine-cohomology-finite}.
Par Compléments sur l'algèbre, lemme
\ref{more-algebra-lemma-generalized-valuation-ring-modules},
ces modules sont des sommes directes de modules cycliques.
Nous avons vu ci-dessus que le rang $\beta_i$ de la
partie libre de $H^i(K)$ et de $H^i(L)$ est le même.
Observons ensuite que
$$
H^i(L \otimes_A^\mathbf{L} A/\pi^m A) =
H^i(L)/\pi^m H^i(L) \oplus H^{i + 1}(L)[\pi^m]
$$
et de même pour $K$. Soit $e$ le plus grand entier
tel que $A/\pi^eA$ apparaisse comme facteur direct de $H^i(X, \mathcal{O}_X)$,
ou, de manière équivalente, de $H^i(L)$, pour un certain $i$. En prenant alors $m = e + 1$,
nous voyons que $H^i(L \otimes_A^\mathbf{L} A/\pi^m A)$ est une somme directe de
$\beta_i$ copies de $A/\pi^m A$ et d'autres modules cycliques,
chacun annulé par $\pi^e$. Par le même raisonnement pour $K$
et l'isomorphisme
$K \otimes_A^\mathbf{L} A/\pi^m A \cong L \otimes_A^\mathbf{L} A/\pi^m A$,
il s'ensuit que $H^i(K)$
ne peut avoir aucun facteur direct cyclique de la forme $A/\pi^l A$
avec $l > e$. (Il s'ensuit également que $K$ est isomorphe à $L$
comme objet de $D(A)$, mais nous n'en aurons pas besoin.)
La seule manière dont l'homomorphisme
$$
H^0(K \otimes^\mathbf{L}_A A/\pi^{e + 1} A) =
H^0(K)/\pi^{e + 1}H^0(K) \oplus H^1(K)[\pi^{e + 1}]
\longrightarrow
(A/\pi^{e + 1} A)^{\oplus r}
$$
puisse être surjectif est donc qu'il soit surjectif sur le
premier facteur direct. C'est ce que nous voulions montrer.
(Pour être précis, l'entier $n$ de l'énoncé du
lemme, s'il existe une section $\sigma$,
doit être égal à $k + e + 1$, où $k$ et $e$ sont comme ci-dessus
et ne dépendent que de $X$.)
\end{proof}

\begin{lemma}
\label{lemma-trivial-over-dvrs}
Soit $f : X \to S$ un morphisme de schémas.
Soit $\mathcal{E}$ un $\mathcal{O}_X$-module localement libre de rang fini.
Supposons que
\begin{enumerate}
\item $f$ est plat et propre et $\mathcal{O}_S = f_*\mathcal{O}_X$,
\item $S$ est un schéma noethérien normal,
\item le tiré en arrière de $\mathcal{E}$ sur $X \times_S \Spec(\mathcal{O}_{S, s})$
est libre pour tout point de codimension $1$, noté $s \in S$.
\end{enumerate}
Alors $\mathcal{E}$ est isomorphe au tiré en arrière d'un
$\mathcal{O}_S$-module localement libre de rang fini.
\end{lemma}

\begin{proof}
Nous allons démontrer que le morphisme canonique
$$
\Phi : f^*f_*\mathcal{E} \longrightarrow \mathcal{E}
$$
est un isomorphisme. Par changement de base plat (Cohomologie des schémas, lemme
\ref{coherent-lemma-flat-base-change-cohomology})
et, d'après les hypothèses (1) et (3), nous voyons que
son tiré en arrière sur $X \times_S \Spec(\mathcal{O}_{S, s})$
est un isomorphisme pour tout point de codimension $1$, noté $s \in S$.
Par Diviseurs, lemme \ref{divisors-lemma-check-isomorphism-via-depth-and-ass},
il suffit de démontrer que $\text{profondeur}((f^*f_*\mathcal{E})_x) \geq 2$
pour tout point $x \in X$ s'envoyant sur un point $s \in S$ de codimension $\geq 2$.
Puisque $f$ est plat et que
$(f^*f_*\mathcal{E})_x = (f_*\mathcal{E})_s \otimes_{\mathcal{O}_{S, s}}
\mathcal{O}_{X, x}$, il suffit de démontrer que
$\text{profondeur}((f_*\mathcal{E})_s) \geq 2$, voir
Algèbre, lemme \ref{algebra-lemma-apply-grothendieck}.
Puisque $S$ est un schéma noethérien normal
et que $\dim(\mathcal{O}_{S, s}) \geq 2$,
nous avons $\text{profondeur}(\mathcal{O}_{S, s}) \geq 2$, voir
Propriétés, lemme \ref{properties-lemma-criterion-normal}.
Nous obtenons donc le résultat voulu d'après
Diviseurs, lemme \ref{divisors-lemma-depth-pushforward}.
\end{proof}

\noindent
Nous pouvons utiliser les résultats ci-dessus pour démontrer l'énoncé
miraculeux suivant.

\begin{theorem}
\label{theorem-pullback-trivial-fibres}
Soit $p$ un nombre premier. Soit $Y$ un schéma quasi-compact et quasi-séparé
sur $\mathbf{F}_p$.
Soit $f : X \to Y$ un morphisme propre, surjectif et de présentation finie
à fibres géométriquement connexes.
Alors le foncteur
$$
\colim_F \textit{Vect}(Y) \longrightarrow \colim_F \textit{Vect}(X)
$$
est pleinement fidèle et son image essentielle se décrit comme suit.
Soit $\mathcal{E}$ un $\mathcal{O}_X$-module localement libre de rang fini.
Supposons que, pour tout $y \in Y$, il existe des entiers $n_y, r_y \geq 0$
tels que
$$
F^{n_y, *}\mathcal{E}|_{X_{y, red}}
\cong
\mathcal{O}_{X_{y, red}}^{\oplus r_y}
$$
Alors, pour un certain $n \geq 0$, le tiré en arrière par la $n$-ième puissance de Frobenius
$F^{n, *}\mathcal{E}$ est le tiré en arrière d'un
$\mathcal{O}_Y$-module localement libre de rang fini.
\end{theorem}

\begin{proof}
Démonstration de la pleine fidélité. Puisque les fibrés vectoriels sur $Y$ sont localement
triviaux, cela se ramène à l'assertion que
$$
\colim_F \Gamma(Y, \mathcal{O}_Y)
\longrightarrow
\colim_F \Gamma(X, \mathcal{O}_X)
$$
est bijectif. Puisque $\{X \to Y\}$ est un recouvrement h, cela
résultera du lemme \ref{lemma-h-sheaf-colim-F} si nous pouvons montrer que les deux homomorphismes
$$
\colim_F \Gamma(X, \mathcal{O}_X)
\longrightarrow
\colim_F \Gamma(X \times_Y X, \mathcal{O}_{X \times_Y X})
$$
sont égaux. Soit $g \in \Gamma(X, \mathcal{O}_X)$,
et notons $g_1$ et $g_2$ les deux tirés en arrière de $g$ à $X \times_Y X$.
Puisque $X_{y, red}$ est géométriquement connexe, nous
voyons que\newline $H^0(X_{y, red}, \mathcal{O}_{X_{y, red}})$ est
une extension purement inséparable de $\kappa(y)$, voir
Variétés, lemme
\ref{varieties-lemma-proper-geometrically-reduced-global-sections}.
Ainsi, $g^q|_{X_{y, red}}$ provient d'un élément de
$\kappa(y)$ pour une certaine puissance de $p$, $q$ (qui peut dépendre de $y$).
Il s'ensuit que $g_1^q$ et $g_2^q$ s'envoient sur le même
élément du corps résiduel en tout point de
$(X \times_Y X)_y = X_y \times_y X_y$.
Par conséquent, la restriction de $g_1 - g_2$ à $(X \times_Y X)_{red}$ est nulle.
Ainsi, $(g_1 - g_2)^n = 0$ pour un certain $n$ que nous pouvons choisir
égal à une puissance de $p$, comme souhaité.

\medskip\noindent
Description de l'image essentielle. Soit $\mathcal{E}$ comme dans l'énoncé
de la proposition. Nous nous ramenons d'abord au cas noethérien.

\medskip\noindent
Soit $y \in Y$ un point, considéré comme un morphisme
$y \to Y$ du spectre du corps résiduel vers $Y$.
Nous pouvons écrire $y \to Y$ comme une limite filtrante de morphismes $Y_i \to Y$ de
présentation finie avec $Y_i$ affine. (Il est préférable de le démontrer
soi-même, mais cela résulte aussi formellement de
Limites, lemme \ref{limits-lemma-relative-approximation} et
\ref{limits-lemma-limit-affine}.)
Pour chaque $i$, posons $Z_i = Y_i \times_Y X$. Alors $X_y = \lim Z_i$
et $X_{y, red} = \lim Z_{i, red}$.
Par Limites, lemme \ref{limits-lemma-descend-modules-finite-presentation},
nous pouvons trouver un $i$ tel que
$F^{n_y, *}\mathcal{E}|_{Z_{i, red}} \cong
\mathcal{O}_{Z_{i, red}}^{\oplus r_y}$.
Fixons $i$.
Nous avons $Z_{i, red} = \lim Z_{i, j}$, où $Z_{i, j} \to Z_i$
est un épaississement de présentation finie (Limites, lemme
\ref{limits-lemma-closed-is-limit-closed-and-finite-presentation}).
En utilisant le même lemme que précédemment, nous pouvons trouver un $j$ tel que
$F^{n_y, *}\mathcal{E}|_{Z_{i, j}} \cong \mathcal{O}_{Z_{i, j}}^{\oplus r_y}$.
Nous concluons que, pour tout $y \in Y$, il existe un morphisme
$Y_y \to Y$ de présentation finie dont l'image contient $y$
et un épaississement $Z_y \to Y_y \times_Y X$ tel que
$F^{n_y, *}\mathcal{E}|_{Z_y} \cong \mathcal{O}_{Z_y}^{\oplus r_y}$.
Observons que l'image de $Y_y \to Y$ est constructible
(Morphismes, lemme \ref{morphisms-lemma-chevalley}).
Puisque $Y$ est quasi-compact pour la topologie constructible
(Topologie, lemme \ref{topology-lemma-constructible-hausdorff-quasi-compact} et
Propriétés, lemme \ref{properties-lemma-quasi-compact-quasi-separated-spectral}),
nous concluons qu'il existe un nombre fini de morphismes
$$
Y_1 \to Y,\ Y_2 \to Y,\ \ldots,\ Y_N \to Y
$$
de présentation finie tels que $Y = \bigcup \Im(Y_a \to Y)$
ensemblistement et tels que, pour chaque $a \in \{1, \ldots, N\}$,
il existe des entiers $n_a, r_a \geq 0$ et un épaississement
$Z_a \subset Y_a \times_Y X$ de présentation finie tel que
$F^{n_a, *}\mathcal{E}|_{Z_a} \cong \mathcal{O}_{Z_a}^{\oplus r_a}$.

\medskip\noindent
Sous cette forme, la condition se descend à une approximation
noethérienne absolue. Nous conseillons vivement au lecteur de sauter
ce paragraphe. Écrivons d'abord $Y = \lim_{i \in I} Y_i$ comme une limite cofiltrante
de schémas de type fini sur $\mathbf{F}_p$ à morphismes de transition
affines (Limites, lemme \ref{limits-lemma-relative-approximation}).
Ensuite, nous pouvons supposer que nous avons des morphismes propres $f_i : X_i \to Y_i$
dont le changement de base à $Y$ redonne $f : X \to Y$, voir
Limites, lemme \ref{limits-lemma-descend-finite-presentation}.
Après avoir augmenté $i$, nous pouvons supposer qu'il existe un
$\mathcal{O}_{X_i}$-module localement libre de rang fini $\mathcal{E}_i$ dont le tiré en arrière
sur $X$ est isomorphe à $\mathcal{E}$, voir
Limites, lemme \ref{limits-lemma-descend-invertible-modules}.
Choisissons $0 \in I$ et notons $E \subset Y_0$ le sous-ensemble constructible
où les fibres géométriques de $f_0$ sont connexes, voir
Compléments sur les morphismes, lemme
\ref{more-morphisms-lemma-nr-geom-connected-components-constructible}.
Alors $Y \to Y_0$ s'envoie dans $E$, voir
Compléments sur les morphismes, lemme
\ref{more-morphisms-lemma-base-change-fibres-geometrically-connected}.
Ainsi, $Y_i \to Y_0$ s'envoie dans $E$ pour $i \gg 0$, voir
Limites, lemme \ref{limits-lemma-limit-contained-in-constructible}.
Nous voyons donc que les fibres de $f_i$ sont géométriquement connexes
pour $i \gg 0$. Par Limites, lemme \ref{limits-lemma-descend-finite-presentation},
pour $i$ assez grand, nous pouvons trouver des morphismes
$Y_{i, a} \to Y_i$ de type fini dont le changement de base à $Y$
redonne $Y_a \to Y$, $a  \in \{1, \ldots, N\}$.
Après avoir éventuellement augmenté $i$, nous pouvons trouver des épaississements
$Z_{i, a} \subset Y_{i, a} \times_{Y_i} X_i$ dont le changement de base
à $Y_a \times_Y X$ redonne $Z_a$ (même référence que précédemment,
combinée avec
Limites, lemmes
\ref{limits-lemma-descend-closed-immersion-finite-presentation} et
\ref{limits-lemma-descend-surjective}).
Puisque $Z_a = \lim Z_{i, a}$, nous constatons qu'après avoir augmenté $i$, nous pouvons supposer
$F^{n_a, *}\mathcal{E}_i|_{Z_{i, a}} \cong
\mathcal{O}_{Z_{i, a}}^{\oplus r_a}$, voir
Limites, lemme \ref{limits-lemma-descend-modules-finite-presentation}.
Enfin, après avoir augmenté $i$ une dernière fois, nous pouvons supposer que
$\coprod Y_{i, a} \to Y_i$ est surjectif d'après
Limites, lemme \ref{limits-lemma-descend-surjective}.
À ce stade, toutes les hypothèses sont satisfaites pour $X_i \to Y_i$
et $\mathcal{E}_i$, et nous voyons qu'il suffit de démontrer le résultat
pour $X_i \to Y_i$ et $\mathcal{E}_i$.

\medskip\noindent
Supposons $Y$ de type fini sur $\mathbf{F}_p$.
Pour démontrer le résultat, nous allons raisonner par récurrence sur $\dim(Y)$.
Nous cherchons un objet de
$\colim_F \textit{Vect}(Y)$ dont le tiré en arrière soit
l'objet de $\colim_F \textit{Vect}(X)$ déterminé par $\mathcal{E}$.
Par la pleine fidélité déjà démontrée et en vertu de la
proposition \ref{proposition-h-descent-vector-bundles-p},
il suffit de construire une descente de $\mathcal{E}$
après avoir remplacé $Y$ par les membres d'un recouvrement h
et $X$ par le changement de base correspondant. Cela signifie
que nous pouvons remplacer $Y$ par un schéma muni d'un morphisme propre et surjectif
vers $Y$, pourvu que cela n'augmente pas la dimension de $Y$.
Si $T \subset T'$ est un épaississement de schémas de type fini
sur $\mathbf{F}_p$, alors
$\colim_F \textit{Vect}(T) = \colim_F \textit{Vect}(T')$
car $\{T \to T'\}$ est un recouvrement h tel que $T \times_{T'} T = T$.
Si $T' \to T$ est un homéomorphisme universel de schémas
de type fini sur $\mathbf{F}_p$, alors
$\colim_F \textit{Vect}(T) = \colim_F \textit{Vect}(T')$
car $\{T \to T'\}$ est un recouvrement h tel que la diagonale
$T \subset T \times_{T'} T$ est un épaississement.

\medskip\noindent
D'après les remarques générales faites ci-dessus, nous remplaçons
$X$ par sa réduction et pouvons supposer que $X$ est réduit.
Considérons la factorisation de Stein $X \to Y' \to Y$, voir
Compléments sur les morphismes, théorème
\ref{more-morphisms-theorem-stein-factorization-Noetherian}.
Alors $Y' \to Y$ est un homéomorphisme universel
de schémas de type fini sur $\mathbf{F}_p$.
D'après ce qui précède, nous pouvons remplacer $Y$ par $Y'$.
Nous pouvons donc supposer $f_*\mathcal{O}_X = \mathcal{O}_Y$
et que $Y$ est réduit. Cela nous ramène au cas traité
dans le paragraphe suivant.

\medskip\noindent
Supposons $Y$ réduit et $f_*\mathcal{O}_X = \mathcal{O}_Y$
sur un sous-schéma ouvert dense de $Y$.
Alors $X \to Y$ est plat sur un sous-schéma
ouvert dense $V \subset Y$, voir
Morphismes, proposition \ref{morphisms-proposition-generic-flatness-reduced}.
Par le lemme \ref{lemma-flat-after-blowing-up},
il existe un éclatement $V$-admissible $Y' \to Y$ tel que
le transformé strict $X'$ de $X$ soit plat sur $Y'$.
Observons que $\dim(Y') = \dim(Y)$, puisque $Y$ et $Y'$ possèdent
un sous-schéma ouvert dense commun. Par Compléments sur les morphismes, lemme
\ref{more-morphisms-lemma-proper-flat-nr-geom-conn-comps-lower-semicontinuous}
et le fait que $V \subset Y'$ est dense,
toutes les fibres de $f' : X' \to Y'$ sont géométriquement connexes.
Nous avons encore $(f'_*\mathcal{O}_{X'})|_V = \mathcal{O}_V$.
Écrivons
$$
Y' \times_Y X = X' \cup E \times_Y X
$$
où $E \subset Y'$ est le diviseur exceptionnel de l'éclatement.
Par les remarques générales ci-dessus, il suffit de démontrer l'existence
pour $Y' \times_Y X \to Y'$ et la restriction de $\mathcal{E}$
à $Y' \times_Y X$.
Supposons que nous trouvions un objet $\xi'$ dans $\colim_F \textit{Vect}(Y')$
dont le tiré en arrière soit la restriction de $\mathcal{E}$ à $X'$ (considérée
comme un objet de la catégorie colimite).
Par récurrence sur $\dim(Y)$, nous pouvons trouver un objet $\xi''$ dans
$\colim_F \textit{Vect}(E)$ dont le tiré en arrière soit la restriction de
$\mathcal{E}$ à $E \times_Y X$. Alors la pleine fidélité
détermine un unique isomorphisme $\xi'|_E \to \xi''$
compatible avec les identifications données avec la restriction
de $\mathcal{E}$ à $E \times_{Y'} X'$. Puisque
$$
\{E \times_Y X \to Y' \times_Y X, X' \to Y' \times_Y X\}
$$
est un recouvrement h donné par une paire d'immersions fermées avec
$$
(E \times_Y X) \times_{(Y' \times_Y X)} X' = E \times_{Y'} X'
$$
nous concluons que le tiré en arrière de $\xi'$ est la restriction de
$\mathcal{E}$ à $Y' \times_Y X$. Il suffit donc de trouver
$\xi'$, et nous nous ramenons au cas traité dans le paragraphe suivant.

\medskip\noindent
Supposons $Y$ réduit, $f$ plat et $f_*\mathcal{O}_X = \mathcal{O}_Y$
sur un sous-schéma ouvert dense de $Y$. Dans ce cas, nous considérons la
normalisation $Y^\nu \to Y$ (Morphismes, section
\ref{morphisms-section-normalization}). C'est un
morphisme fini surjectif
(Morphismes, lemme \ref{morphisms-lemma-nagata-normalization} et
\ref{morphisms-lemma-ubiquity-nagata}) qui est un isomorphisme
sur un ouvert dense. Par conséquent, d'après nos remarques générales, nous pouvons
remplacer $Y$ par $Y^\nu$ et $X$ par $Y^\nu \times_Y X$.
Après ce remplacement, nous voyons que $\mathcal{O}_Y = f_*\mathcal{O}_X$
(car la factorisation de Stein doit être un isomorphisme
dans ce cas ; petit détail omis).

\medskip\noindent
Supposons que $Y$ soit noethérien normal, que $f$ soit plat et que
$f_*\mathcal{O}_X = \mathcal{O}_Y$. Après avoir remplacé $\mathcal{E}$
par un tiré en arrière par une puissance convenable de Frobenius, nous pouvons supposer que $\mathcal{E}$
est trivial sur les fibres schématiques de $f$ aux points génériques
des composantes irréductibles de $Y$ (car
$\colim_F \textit{Vect}(-)$ est invariant par les homéomorphismes
universels, voir ci-dessus). Comme dans les arguments précédents
(lors de la réduction au cas noethérien), nous concluons qu'il existe un
sous-schéma ouvert dense $V \subset Y$ tel que $\mathcal{E}|_{f^{-1}(V)}$ soit libre.
Soit $Z \subset Y$ un sous-schéma fermé tel que
$Y = V \amalg Z$ ensemblistement. Soient $z_1, \ldots, z_t \in Z$
les points génériques des composantes irréductibles de $Z$
de codimension $1$. Alors $A_i = \mathcal{O}_{Y, z_i}$ est
un anneau de valuation discrète. Soit $n_i$ l'entier obtenu dans le
lemme \ref{lemma-trivial-fibres-dvr} pour le schéma $X_{A_i}$ sur $A_i$.
Après avoir remplacé $\mathcal{E}$ par un tiré en arrière par une puissance
convenable de Frobenius, nous pouvons supposer que $\mathcal{E}$ est libre sur
$X_{A_i/\mathfrak m_i^{n_i}}$ (là encore parce que la catégorie colimite
est invariante par homéomorphismes universels, voir ci-dessus).
Alors le lemme \ref{lemma-trivial-fibres-dvr} nous dit que
$\mathcal{E}$ est libre sur $X_{A_i}$.
Nous concluons enfin en appliquant le lemme \ref{lemma-trivial-over-dvrs}.
\end{proof}










\section{Éclatement de complexes}
\label{section-blowup-complexes}

\noindent
Cette section établit des formes normales pour les objets parfaits de la catégorie
dérivée après éclatement.

\begin{lemma}
\label{lemma-fitting-ideals-complex}
Soit $X$ un schéma. Soit $E \in D(\mathcal{O}_X)$ pseudo-cohérent.
Pour tous $p, k \in \mathbf{Z}$, il existe un faisceau quasi-cohérent
d'idéaux de type fini $\text{Fit}_{p, k}(E) \subset \mathcal{O}_X$
ayant la propriété suivante : pour tout ouvert $U \subset X$
tel que $E|_U$ soit isomorphe à
$$
\ldots \to
\mathcal{O}_U^{\oplus n_{b - 2}}
\xrightarrow{d_{b - 2}}
\mathcal{O}_U^{\oplus n_{b - 1}}
\xrightarrow{d_{b - 1}}
\mathcal{O}_U^{\oplus n_b} \to 0 \to \ldots
$$
la restriction $\text{Fit}_{p, k}(E)|_U$ est engendrée par les
mineurs de la matrice de $d_p$ de taille
$$
- k + n_{p + 1} - n_{p + 2} + \ldots + (-1)^{b - p + 1} n_b
$$
Convention : l'idéal engendré par les mineurs $r \times r$
est $\mathcal{O}_U$ si $r \leq 0$, et l'idéal engendré par les
mineurs $r \times r$, lorsque $r > \min(n_p, n_{p + 1})$, est nul.
\end{lemma}

\begin{proof}
Observons que $E$ a, localement sur $X$, la forme indiquée dans le lemme ; voir
la section \ref{more-algebra-section-pseudo-coherent} de Compléments d'algèbre,
la section \ref{cohomology-section-pseudo-coherent} de Cohomologie et
la section \ref{perfect-section-spell-out} de Catégories dérivées de schémas.
Il suffit donc de démontrer que l'idéal des mineurs est indépendant
du représentant choisi. Pour cela, il suffit de raisonner
sur les anneaux locaux. Sur un anneau local $(R, \mathfrak m, \kappa)$,
considérons un complexe borné supérieurement
$$
F^\bullet :
\ldots \to
R^{\oplus n_{b - 2}}
\xrightarrow{d_{b - 2}}
R^{\oplus n_{b - 1}}
\xrightarrow{d_{b - 1}}
R^{\oplus n_b} \to 0 \to \ldots
$$
Notons $\text{Fit}_{k, p}(F^\bullet) \subset R$ l'idéal engendré
par les mineurs d'ordre $k - n_{p + 1} + n_{p + 2} - \ldots + (-1)^{b - p} n_b$
de la matrice de $d_p$. Supposons qu'un coefficient matriciel d'une
différentielle de $F^\bullet$ soit inversible. Choisissons alors le plus grand
entier $i$ tel que $d_i$ ait un coefficient matriciel inversible.
D'après le lemme \ref{algebra-lemma-add-trivial-complex} d'Algèbre,
le complexe $F^\bullet$ est isomorphe à la somme directe d'un complexe trivial
$\ldots \to 0 \to R \to R \to 0 \to \ldots$, dont les termes non nuls
sont en degrés $i$ et $i + 1$, et d'un complexe $(F')^\bullet$.
Nous laissons au lecteur le soin de vérifier que
$\text{Fit}_{p, k}(F^\bullet) = \text{Fit}_{p, k}((F')^\bullet)$ ;
c'est ici qu'intervient la formule donnant la taille des mineurs.
Si $(F')^\bullet$ possède une autre différentielle ayant un coefficient
matriciel inversible, nous recommençons, et ainsi de suite.
En poursuivant de cette manière, nous obtenons finalement un complexe $(F^\infty)^\bullet$
dont toutes les différentielles ont des matrices à coefficients
dans $\mathfrak m$. Il peut être nécessaire d'effectuer une infinité
d'étapes, mais, pour toute borne inférieure fixée, seul un nombre fini de ces
étapes modifie le complexe dans les degrés $\geq $ cette borne.
Ainsi, la « limite » $(F^\infty)^\bullet$ est un complexe borné supérieurement
bien défini de modules libres de rang fini ; elle est munie
d'un quasi-isomorphisme $(F^\infty)^\bullet \to F^\bullet$
vers le complexe de départ, et
$\text{Fit}_{p, k}(F^\bullet) = \text{Fit}_{p, k}((F^\infty)^\bullet)$.
Puisque le complexe $(F^\infty)^\bullet$ est unique à isomorphisme près d'après
le lemme
\ref{more-algebra-lemma-lift-pseudo-coherent-from-residue-field} de Compléments d'algèbre,
la démonstration est achevée.
\end{proof}

\begin{lemma}
\label{lemma-blowup-complex}
Soit $X$ un schéma. Soit $E \in D(\mathcal{O}_X)$ parfait.
Soit $U \subset X$ un sous-schéma ouvert schématiquement dense
tel que $H^i(E|_U)$ soit localement libre de rang constant $r_i$
pour tout $i \in \mathbf{Z}$.
Alors il existe un éclatement $U$-admissible $b : X' \to X$ tel que
$H^i(Lb^*E)$ soit un $\mathcal{O}_{X'}$-module parfait
de dimension de Tor $\leq 1$ pour tout $i \in \mathbf{Z}$.
\end{lemma}

\begin{proof}
Nous allons construire et étudier l'éclatement localement sur des ouverts affines. Supposons donc que
$V \subset X$ soit un sous-schéma ouvert affine tel que
$E|_V$ puisse être représenté par le complexe
$$
\mathcal{O}_V^{\oplus n_a} \xrightarrow{d_a}
\ldots \xrightarrow{d_{b - 1}} \mathcal{O}_V^{\oplus n_b}
$$
Posons $k_ i = r_{i + 1} - r_{i + 2} + \ldots + (-1)^{b - i + 1}r_b$.
Un calcul, que nous omettons, montre que, sur $U \cap V$, le rang de
$d_i$ est
$$
\rho_i = - k_i + n_{i + 1} - n_{i + 2} + \ldots + (-1)^{b - i + 1}n_b
$$
au sens où le conoyau de $d_i$ est localement
libre de rang $n_{i + 1} - \rho_i$. Soit
$\mathcal{I}_i \subset \mathcal{O}_V$ l'idéal engendré par les mineurs
$\rho_i \times \rho_i$ de la matrice de $d_i$.

\medskip\noindent
D'une part, en comparant avec le lemme \ref{lemma-fitting-ideals-complex},
nous voyons que l'idéal $\mathcal{I}_i$ correspond à l'idéal global
$\text{Fit}_{i, k_i}(E)$, dont on a démontré qu'il était
indépendant du choix du complexe représentant $E|_V$.
D'autre part, $\mathcal{I}_i$ est le $(n_{i + 1} - \rho_i)$-ième
idéal de Fitting de $\Coker(d_i)$. Gardons ce fait à l'esprit.

\medskip\noindent
Soit $b : X' \to X$ l'éclatement le long du produit des idéaux
$\text{Fit}_{i, k_i}(E)$ ; cette définition a un sens puisque, localement sur $X$,
presque tous ces idéaux sont égaux à l'idéal unité (voir ci-dessus).
Cet éclatement domine les éclatements $b_i : X'_i \to X$ le long des
idéaux $\text{Fit}_{i, k_i}(E)$ ; voir
le lemme \ref{divisors-lemma-blowing-up-two-ideals} de Diviseurs.
D'après le lemme \ref{divisors-lemma-blowup-fitting-ideal} de Diviseurs,
chaque $b_i$ est un éclatement $U$-admissible. Il s'ensuit que
$b$ est un éclatement $U$-admissible (nous omettons un petit détail ; comparer avec
la démonstration du lemme \ref{divisors-lemma-dominate-admissible-blowups} de Diviseurs).
Enfin, $U$ demeure un sous-schéma ouvert schématiquement dense de $X'$.
Après avoir remplacé $X$ par $X'$, nous nous trouvons donc dans la situation
étudiée au paragraphe suivant.

\medskip\noindent
Supposons que $\text{Fit}_{i, k_i}(E)$ soit un idéal inversible pour tout $i$.
Choisissons un ouvert affine $V$ et un complexe de modules libres de rang fini
représentant $E|_V$ comme ci-dessus. Il résulte du
lemme \ref{divisors-lemma-blowup-fitting-ideal} de Diviseurs,
que $\Coker(d_i)$ est de dimension de Tor $\leq 1$. Ainsi,
$\Im(d_i)$ est localement libre de rang fini, en tant que noyau d'une application
d'un module localement libre de rang fini vers un module de présentation finie
et de dimension de Tor $\leq 1$. Par conséquent, $\Ker(d_i)$
est lui aussi localement libre de rang fini (par le même argument). La suite
exacte courte
$$
0 \to \Im(d_{i - 1}) \to \Ker(d_i) \to H^i(E)|_V \to 0
$$
donne alors le résultat voulu, ce qui achève la démonstration.
\end{proof}

\begin{lemma}
\label{lemma-blowup-complex-integral}
Soit $X$ un schéma intègre. Soit $E \in D(\mathcal{O}_X)$ parfait.
Alors il existe un ouvert non vide $U \subset X$
tel que $H^i(E|_U)$ soit localement libre de rang constant $r_i$
pour tout $i \in \mathbf{Z}$, ainsi qu'un éclatement $U$-admissible
$b : X' \to X$ tel que $H^i(Lb^*E)$ soit un
$\mathcal{O}_{X'}$-module parfait de dimension de Tor $\leq 1$ pour tout $i \in \mathbf{Z}$.
\end{lemma}

\begin{proof}
Nous invitons vivement le lecteur à trouver sa propre démonstration de
l'existence de $U$. Soit $\eta \in X$ le point générique.
La restriction de $E$ à $\eta$ est isomorphe dans $D(\kappa(\eta))$ à un
complexe fini $V^\bullet$ d'espaces vectoriels de dimension finie dont les
différentielles sont nulles. Posons $r_i = \dim_{\kappa(\eta)} V^i$.
Alors l'objet parfait $E'$ de $D(\mathcal{O}_X)$ représenté par le complexe
dont les termes sont $\mathcal{O}_X^{\oplus r_i}$ et les différentielles nulles
devient isomorphe à $E$ après restriction à $\eta$.
Par conséquent, d'après le lemme
\ref{perfect-lemma-descend-relatively-perfect} de Catégories dérivées de schémas,
il existe un voisinage ouvert $U$ de $\eta$ tel que $E|_U$ et $E'|_U$
soient isomorphes. Cela démontre la première assertion. La seconde résulte de la
première et du lemme \ref{lemma-blowup-complex}, puisque tout ouvert non vide est
schématiquement dense dans le schéma intègre $X$.
\end{proof}

\begin{remark}
\label{remark-when-you-have-a-complex}
Soit $X$ un schéma. Soit $E \in D(\mathcal{O}_X)$ un objet parfait tel que
$H^i(E)$ soit un $\mathcal{O}_X$-module parfait de dimension de Tor $\leq 1$
pour tout $i \in \mathbf{Z}$. Cette propriété permet parfois de ramener
des questions sur $E$ à des questions sur $H^i(E)$. Par exemple, supposons que
$$
\mathcal{E}^a \xrightarrow{d^a} \ldots
\xrightarrow{d^{b - 2}} \mathcal{E}^{b - 1}
\xrightarrow{d^{b - 1}} \mathcal{E}^b
$$
soit un complexe borné de $\mathcal{O}_X$-modules localement libres
de rang fini représentant $E$. Alors $\Im(d^i)$ et $\Ker(d^i)$ sont des
$\mathcal{O}_X$-modules localement libres de rang fini pour tout $i$. En effet, supposons par
récurrence que cela soit vrai pour tous les indices strictement supérieurs à $i$.
Nous pouvons alors commencer par utiliser la suite exacte courte
$$
0 \to \Im(d^i) \to \Ker(d^{i + 1}) \to H^{i + 1}(E) \to 0
$$
et l'hypothèse selon laquelle $H^{i + 1}(E)$ est parfait de dimension de Tor $\leq 1$
pour conclure que $\Im(d^i)$ est localement libre de rang fini.
Le même argument, appliqué ensuite à la suite exacte courte
$$
0 \to \Ker(d^i) \to \mathcal{E}^i \to \Im(d^i) \to 0
$$
montre que $\Ker(d^i)$ est localement libre de rang fini.
Il en résulte que les triangles distingués
$$
\tau_{\leq k - 1}E \to \tau_{\leq k}E \to H^k(E)[-k] \to
(\tau_{\leq k - 1}E)[1]
$$
sont représentés par les suites exactes courtes suivantes de
complexes bornés de modules localement libres de rang fini
$$
\begin{matrix}
& &
& &
& &
0 \\
& &
& &
& &
\downarrow \\
\mathcal{E}^a & \to &
\ldots & \to &
\mathcal{E}^{k - 2} & \to &
\Ker(d^{k - 1}) \\
\downarrow & &
& &
\downarrow & &
\downarrow \\
\mathcal{E}^a & \to &
\ldots & \to &
\mathcal{E}^{k - 2} & \to &
\mathcal{E}^{k - 1} & \to &
\Ker(d^k) \\
& &
& &
& &
\downarrow & &
\downarrow \\
& &
& &
& &
\Im(d^{k - 1}) & \to &
\Ker(d^k) \\
& &
& &
& &
\downarrow \\
& &
& &
& &
0
\end{matrix}
$$
Ici, les complexes sont les lignes, et les zéros « évidents »
sont omis dans le diagramme.
\end{remark}





\section{Éclatement de modules parfaits}
\label{section-blowup-perfect-modules}

\noindent
Cette section tente d'établir des formes normales pour les modules parfaits de dimension de Tor
$\leq 1$ après éclatement. Nous n'y parvenons que partiellement.

\begin{lemma}
\label{lemma-blowup-pd1-derived}
Soit $X$ un schéma. Soit $\mathcal{F}$ un
$\mathcal{O}_X$-module parfait de dimension de Tor $\leq 1$. Pour tout
éclatement $b : X' \to X$, on a $Lb^*\mathcal{F} = b^*\mathcal{F}$,
et $b^*\mathcal{F}$ est un $\mathcal{O}_X$-module parfait
de dimension de Tor $\leq 1$.
\end{lemma}

\begin{proof}
Nous pouvons supposer que $X = \Spec(A)$ est affine et que le $A$-module $M$
correspondant à $\mathcal{F}$ admet une présentation
$$
0 \to A^{\oplus m} \to A^{\oplus n} \to M \to 0
$$
Supposons que $I \subset A$ soit un idéal et que $a \in I$. Rappelons que
l'algèbre affine d'éclatement $A[\frac{I}{a}]$ est un sous-anneau de $A_a$.
Comme la localisation est exacte, $A_a^{\oplus m} \to A_a^{\oplus n}$
est injective. Par suite, $A[\frac{I}{a}]^{\oplus m} \to A[\frac{I}{a}]^{\oplus n}$
est elle aussi injective. Cela démontre le lemme.
\end{proof}

\begin{lemma}
\label{lemma-blowup-pd1}
Soit $X$ un schéma. Soit $\mathcal{F}$ un $\mathcal{O}_X$-module parfait
de dimension de Tor $\leq 1$. Soit $U \subset X$ un ouvert schématiquement
dense tel que $\mathcal{F}|_U$ soit localement libre de rang constant
$r$. Alors il existe un éclatement $U$-admissible $b : X' \to X$ tel qu'il
existe une suite exacte courte canonique
$$
0 \to \mathcal{K} \to b^*\mathcal{F} \to \mathcal{Q} \to 0
$$
où $\mathcal{Q}$ est localement libre de rang $r$ et
$\mathcal{K}$ est un $\mathcal{O}_X$-module parfait
de dimension de Tor $\leq 1$ dont la restriction à $U$ est nulle.
\end{lemma}

\begin{proof}
Combiner le lemme \ref{divisors-lemma-blowup-fitting-ideal} de Diviseurs et
le lemme \ref{lemma-blowup-pd1-derived}.
\end{proof}

\begin{lemma}
\label{lemma-canonical-blowup-torsion-pd1}
Soit $X$ un schéma. Soit $\mathcal{F}$ un $\mathcal{O}_X$-module parfait
de dimension de Tor $\leq 1$. Soit $U \subset X$ un ouvert tel que
$\mathcal{F}|_U = 0$. Alors il existe un éclatement $U$-admissible
$$
b : X' \to X
$$
tel que $\mathcal{F}' = b^*\mathcal{F}$ soit muni de deux filtrations
canoniques localement finies
$$
0 = F^0 \subset F^1 \subset F^2 \subset \ldots \subset \mathcal{F}'
\quad\text{et}\quad
\mathcal{F}' = F_1 \supset F_2 \supset F_3 \supset \ldots \supset 0
$$
telles que, pour tout $n \geq 1$, il existe un diviseur de Cartier effectif
$D_n \subset X'$ tel que
$$
F^i/F^{i - 1}
\quad\text{et}\quad
F_i/F_{i + 1}
$$
sont localement libres de rang $i$ sur $D_i$.
\end{lemma}

\begin{proof}
Choisissons un ouvert affine $V \subset X$ tel qu'il existe une présentation
$$
0 \to \mathcal{O}_V^{\oplus n} \xrightarrow{A} \mathcal{O}_V^{\oplus n} \to
\mathcal{F} \to 0
$$
pour un certain $n$ et une certaine matrice $A$. L'idéal le long duquel nous allons éclater
est le produit des idéaux de Fitting $\text{Fit}_k(\mathcal{F})$
pour $k \geq 0$. Cette définition a un sens puisque, dans la situation affine
ci-dessus, $\text{Fit}_k(\mathcal{F})|_V = \mathcal{O}_V$
pour $k > n$. Il est clair qu'il s'agit d'un éclatement $U$-admissible. D'après
le lemme \ref{divisors-lemma-blowing-up-two-ideals} de Diviseurs,
nous voyons que, sur $X'$, les idéaux $\text{Fit}_k(\mathcal{F})$ sont inversibles.
Nous sommes donc ramenés au cas étudié au
paragraphe suivant.

\medskip\noindent
Supposons que $\text{Fit}_k(\mathcal{F})$ soit un idéal inversible pour
$k \geq 0$. Si $E_k \subset X$ est le diviseur de Cartier
effectif défini par $\text{Fit}_k(\mathcal{F})$
pour $k \geq 0$, alors les diviseurs de Cartier effectifs
$D_k$ de l'énoncé du lemme vérifieront
$$
E_k = D_{k + 1} + 2 D_{k + 2} + 3 D_{k + 3} + \ldots
$$
Cette somme a un sens puisque la famille $D_k$ sera localement finie.
De plus, elle détermine de manière unique les diviseurs de Cartier effectifs $D_k$ ;
il suffit donc de construire $D_k$ localement.

\medskip\noindent
Choisissons un ouvert affine $V \subset X$ et une présentation de $\mathcal{F}|_V$
comme ci-dessus. Nous allons construire les diviseurs et les filtrations par récurrence
sur l'entier $n$ de la présentation. Posons $D_k|_V = \emptyset$
pour $k > n$ et $D_n|V = E_{n - 1}|_V$.
Après avoir restreint $V$, nous pouvons supposer que
$\text{Fit}_{n - 1}(\mathcal{F})|_V$ est engendré par un
seul non-diviseur de zéro $f \in \Gamma(V, \mathcal{O}_V)$.
Comme $\text{Fit}_{n - 1}(\mathcal{F})|_V$ est l'idéal engendré
par les coefficients de $A$, il existe une matrice
$A'$ à coefficients dans $\Gamma(V, \mathcal{O}_V)$ telle que $A = fA'$.
Définissons $\mathcal{F}'$ sur $V$ par la suite exacte courte
$$
0 \to \mathcal{O}_V^{\oplus n} \xrightarrow{A'} \mathcal{O}_V^{\oplus n} \to
\mathcal{F}' \to 0
$$
Puisque les coefficients de $A'$ engendrent l'idéal unité de
$\Gamma(V, \mathcal{O}_V)$, le module $\mathcal{F}'$ admet localement sur $V$
une présentation où $n$ est diminué de
$1$ ; voir le lemme \ref{algebra-lemma-add-trivial-complex} d'Algèbre.
Notons en outre que
$f^{n - k}\text{Fit}_k(\mathcal{F}') = \text{Fit}_k(\mathcal{F})|_V$
pour $k = 0, \ldots, n$. Ainsi, $\text{Fit}_k(\mathcal{F}')$
est un idéal inversible pour tout $k$. Par récurrence,
il existe donc des diviseurs de Cartier effectifs $D'_k \subset V$
tels que $\mathcal{F}'$ admette deux filtrations canoniques comme dans l'énoncé
du lemme. Posons alors $D_k|_V = D'_k$ pour $k = 1, \ldots, n - 1$.
Les égalités entre diviseurs de Cartier effectifs
affichées ci-dessus sont vérifiées avec ces choix. Il reste
à construire les filtrations. En effet,
nous disposons des suites exactes courtes
$$
0 \to
\mathcal{O}_{D_n \cap V}^{\oplus n} \to
\mathcal{F} \to \mathcal{F}' \to 0
\quad\text{et}\quad
0 \to
\mathcal{F}' \to \mathcal{F} \to
\mathcal{O}_{D_n \cap V}^{\oplus n} \to 0
$$
qui proviennent des deux factorisations $A = A'f = f A'$ de $A$.
Ces suites sont canoniques parce que, dans la première,
le sous-module est $\Ker(f : \mathcal{F} \to \mathcal{F})$
et, dans la seconde, le module quotient est
$\Coker(f : \mathcal{F} \to \mathcal{F})$.
\end{proof}

\begin{lemma}
\label{lemma-blowup-map-pd1}
Soit $X$ un schéma. Soit $\varphi : \mathcal{F} \to \mathcal{G}$
un homomorphisme de $\mathcal{O}_X$-modules parfaits de dimension de Tor $\leq 1$.
Soit $U \subset X$ un ouvert schématiquement dense
tel que $\mathcal{F}|_U = 0$ et $\mathcal{G}|_U = 0$.
Alors il existe un éclatement $U$-admissible $b : X' \to X$ tel que
le noyau, l'image et le conoyau de $b^*\varphi$ soient des
$\mathcal{O}_{X'}$-modules parfaits de dimension de Tor $\leq 1$.
\end{lemma}

\begin{proof}
Les hypothèses montrent que l'objet $(\mathcal{F} \to \mathcal{G})$
de $D(\mathcal{O}_X)$ est parfait. Nous obtenons donc un éclatement $U$-admissible
qui convient au conoyau et au noyau grâce aux lemmes \ref{lemma-blowup-complex}
et \ref{lemma-blowup-pd1-derived} (le second permettant de déterminer la forme du complexe
après image inverse). L'image est le noyau du conoyau ;
elle est donc elle aussi parfaite de dimension de Tor $\leq 1$.
\end{proof}











\section{Un opérateur introduit par Berthelot et Ogus}
\label{section-eta}

\noindent
Lire d'abord la section \ref{cohomology-section-eta} de Cohomologie.

\medskip\noindent
Soit $X$ un schéma. Soit $D \subset X$ un diviseur de Cartier effectif.
Soit $\mathcal{I} = \mathcal{I}_D \subset \mathcal{O}_X$
le faisceau d'idéaux de $D$ ; voir
la section
\ref{divisors-section-effective-Cartier-invertible} de Diviseurs.
On peut manifestement appliquer les considérations de la section
\ref{cohomology-section-eta} de Cohomologie à $X$ et $\mathcal{I}$.

\begin{lemma}
\label{lemma-eta-stalks}
Soit $X$ un schéma. Soit $D \subset X$ un diviseur de Cartier effectif
de faisceau d'idéaux $\mathcal{I} \subset \mathcal{O}_X$.
Soit $\mathcal{F}^\bullet$ un complexe de
$\mathcal{O}_X$-modules quasi-cohérents tel que $\mathcal{F}^i$ soit
sans $\mathcal{I}$-torsion pour tout $i$. Alors
$\eta_\mathcal{I}\mathcal{F}^\bullet$ est un complexe de
$\mathcal{O}_X$-modules quasi-cohérents. De plus,
si $U = \Spec(A) \subset X$ est un ouvert affine et $D \cap U = V(f)$,
alors $\eta_f(\mathcal{F}^\bullet(U))$ est canoniquement isomorphe
à $(\eta_\mathcal{I}\mathcal{F}^\bullet)(U)$.
\end{lemma}

\begin{proof}
Omis.
\end{proof}

\begin{lemma}
\label{lemma-Leta}
Soit $X$ un schéma. Soit $D \subset X$ un diviseur de Cartier effectif
de faisceau d'idéaux $\mathcal{I} \subset \mathcal{O}_X$.
Le foncteur $L\eta_\mathcal{I} : D(\mathcal{O}_X) \to D(\mathcal{O}_X)$ du
lemme \ref{cohomology-lemma-Leta} de Cohomologie envoie
$D_\QCoh(\mathcal{O}_X)$ dans lui-même.
De plus, si $X = \Spec(A)$ est affine et $D = V(f)$,
alors le foncteur $L\eta_f$ sur $D(A)$ défini dans le
lemme \ref{more-algebra-lemma-Leta} de Compléments d'algèbre
et le foncteur $L\eta_\mathcal{I}$ sur $D_\QCoh(\mathcal{O}_X)$
se correspondent par l'équivalence du lemme
\ref{perfect-lemma-affine-compare-bounded} de Catégories dérivées de schémas.
\end{lemma}

\begin{proof}
Omis.
\end{proof}





\section{Éclatement de complexes, II}
\label{section-blowup-complexes-II}

\noindent
Le contenu de cette section sera utilisé pour construire une version de la
construction du graphe de Macpherson dans la section \ref{section-blowup-complexes-III}.
 
\begin{situation}
\label{situation-complex-and-divisor}
Ici $X$ est un schéma, $D \subset X$ est un diviseur de Cartier effectif
de faisceau d'idéaux $\mathcal{I} \subset \mathcal{O}_X$, et
$M$ est un objet parfait de $D(\mathcal{O}_X)$.
\end{situation}

\noindent
Soit $(X, D, M)$ un triplet comme dans la
situation \ref{situation-complex-and-divisor}.
Considérons un ouvert affine $U = \Spec(A) \subset X$
tel que
\begin{enumerate}
\item $D \cap U = V(f)$ pour un non-diviseur de zéro $f \in A$, et
\item il existe un complexe borné $M^\bullet$ de
$A$-modules libres de rang fini représentant $M|_U$ (via l'équivalence de
Catégories dérivées des schémas, lemme
\ref{perfect-lemma-affine-compare-bounded}).
\end{enumerate}
Nous dirons que $(U, A, f, M^\bullet)$ est une
{\it carte affine} pour $(X, D, M)$.
Considérons les idéaux $I_i(M^\bullet, f) \subset A$ définis dans
Plus sur l'algèbre, section \ref{more-algebra-section-perfect-eta}.
Disons que $(X, S, M)$ est un {\it bon triplet} si, pour tout $x \in D$,
il existe une carte affine $(U, A, f, M^\bullet)$
avec $x \in U$ et telle que les idéaux $I_i(M^\bullet, f)$ soient principaux pour tout
$i \in \mathbf{Z}$.

\begin{lemma}
\label{lemma-pullback-triple-ideals-good}
Dans la situation \ref{situation-complex-and-divisor}, soit $h : Y \to X$ un
morphisme de schémas tel que l'image inverse $E = h^{-1}D$ de $D$
soit définie (Diviseurs, définition
\ref{divisors-definition-pullback-effective-Cartier-divisor}).
Soit $(U, A, f, M^\bullet)$ une carte affine pour $(X, D, M)$.
Soit $V = \Spec(B) \subset Y$ un ouvert affine avec $h(V) \subset U$.
Notons $g \in B$ l'image de $f \in A$.
Alors
\begin{enumerate}
\item $(V, B, g, M^\bullet \otimes_A B)$ est une carte affine pour $(Y, E, Lh^*M)$,
\item $I_i(M^\bullet, f)B = I_i(M^\bullet \otimes_A B, g)$ dans $B$, et
\item si $(X, D, M)$ est un bon triplet, alors
$(Y, E, Lh^*M)$ est un bon triplet.
\end{enumerate}
\end{lemma}

\begin{proof}
La première assertion résulte des observations suivantes :
$g$ est un non-diviseur de zéro dans $B$ qui définit $E \cap V \subset V$
et $M^\bullet \otimes_A B$ représente $M^\bullet \otimes_A^\mathbf{L} B$
et représente donc l'image inverse de $M$ sur $V$ d'après
Catégories dérivées des schémas, lemme
\ref{perfect-lemma-quasi-coherence-pullback}.
La partie (2) résulte de la partie (1) et de
Plus sur l'algèbre, lemme \ref{more-algebra-lemma-eta-base-change-pre}.
En combinant ceci avec
Plus sur l'algèbre, lemme \ref{more-algebra-lemma-eta-base-change-pre},
nous concluons que la deuxième assertion du lemme est vraie.
\end{proof}

\begin{lemma}
\label{lemma-good-triple}
Soient $X, D, \mathcal{I}, M$ comme dans la
situation \ref{situation-complex-and-divisor}.
Si $(X, D, M)$ est un bon triplet, alors
$L\eta_\mathcal{I}M$ est un objet parfait
de $D(\mathcal{O}_X)$.
\end{lemma}

\begin{proof}
Traduction de
Plus sur l'algèbre, lemme \ref{more-algebra-lemma-eta-locally-free}.
Pour effectuer la traduction, utiliser le lemme \ref{lemma-Leta}.
\end{proof}

\begin{lemma}
\label{lemma-good-triple-bdd-loc-free}
Soient $X, D, \mathcal{I}, M$ comme dans la
situation \ref{situation-complex-and-divisor}.
Supposons que $(X, D, M)$ soit un bon triplet.
S'il existe un complexe localement borné $\mathcal{M}^\bullet$
de $\mathcal{O}_X$-modules localement libres de rang fini représentant $M$,
alors il existe un complexe localement borné $\mathcal{Q}^\bullet$
de $\mathcal{O}_{X'}$-modules localement libres de rang fini représentant
$L\eta_\mathcal{I}M$.
\end{lemma}

\begin{proof}
D'après Cohomologie, lemme \ref{cohomology-lemma-Leta},
le complexe $\mathcal{Q}^\bullet = \eta_\mathcal{I}\mathcal{M}^\bullet$
représente $L\eta_\mathcal{I}M$. Pour vérifier que ce complexe
est localement borné et formé de modules localement libres de rang fini, nous pouvons travailler
localement sur des ouverts affines. La bornitude est alors claire.
Choisissons une carte affine $(U, A, f, M^\bullet)$ pour $(X, D, M)$
telle que les idéaux $I_i(M^\bullet, f)$ soient principaux et telle que
$\mathcal{M}^i|_U$ soit libre de rang fini pour chaque $i$. Notre hypothèse
que $(X, D, M)$ est un bon triplet nous permet de le faire. En posant
$N^i = \Gamma(U, \mathcal{M}^i|_U)$, nous obtenons un complexe borné
$N^\bullet$ de $A$-modules libres de rang fini représentant le
même objet de $D(A)$ que le complexe $M^\bullet$ (d'après
Catégories dérivées des schémas, lemme
\ref{perfect-lemma-affine-compare-bounded}).
Alors $I_i(N^\bullet, f)$ est un idéal principal pour tout $i$
d'après Plus sur l'algèbre, lemme \ref{more-algebra-lemma-ideal-well-defined}.
Par conséquent, le complexe $\eta_fN^\bullet$ est un complexe borné
de $A$-modules localement libres de rang fini. Comme
$\mathcal{Q}^i|_U$ est le $\mathcal{O}_U$-module quasi-cohérent
correspondant à $\eta_fN^i$ d'après le lemme \ref{lemma-eta-stalks}, nous concluons.
\end{proof}

\begin{lemma}
\label{lemma-complex-and-divisor-eta-pull}
Dans la situation \ref{situation-complex-and-divisor}, soit $h : Y \to X$
un morphisme de schémas tel que l'image inverse $E = h^{-1}D$
soit définie. Si $(X, D, M)$ est un bon triplet, alors
$$
Lh^*(L\eta_\mathcal{I}M) = L\eta_\mathcal{J}(Lh^*M)
$$
dans $D(\mathcal{O}_Y)$, où $\mathcal{J}$ est le faisceau d'idéaux de $E$.
\end{lemma}

\begin{proof}
Traduction de
Plus sur l'algèbre, lemme \ref{more-algebra-lemma-eta-base-change}.\newline
Utiliser les lemmes \ref{lemma-eta-stalks} et \ref{lemma-Leta}
pour effectuer la traduction.
\end{proof}

\begin{lemma}
\label{lemma-complex-and-divisor-blowup-pre}
Dans la situation \ref{situation-complex-and-divisor},
il existe un unique morphisme $b : X' \to X$ tel que
\begin{enumerate}
\item l'image inverse $D' = b^{-1}D$ soit définie et que
$(X', D', M')$ soit un bon triplet, où $M' = Lb^*M$, et
\item pour tout morphisme de schémas $h : Y \to X$ tel que
l'image inverse $E = h^{-1}D$ soit définie et que $(Y, E, Lh^*M)$
soit un bon triplet, il existe une unique factorisation de $h$ à travers $b$.
\end{enumerate}
De plus, pour toute carte affine $(U, A, f, M^\bullet)$, la restriction
$b^{-1}(U) \to U$ est l'éclatement dans le produit des idéaux
$I_i(M^\bullet, f)$ et, pour tout ouvert quasi-compact $W \subset X$, la
restriction $b|_{b^{-1}(W)} : b^{-1}(W) \to W$ est, relativement à $W \setminus D$, un
éclatement admissible.
\end{lemma}

\begin{proof}
La preuve consiste simplement à éclater localement $X$ dans les idéaux produits
$I_i(M^\bullet, f)$ pour toute carte affine $(U, A, f, M^\bullet)$.
Les premiers lemmes de Plus sur l'algèbre, section
\ref{more-algebra-section-perfect-eta}, montrent que ceci est bien défini.
La propriété universelle (2) résulte alors de la
propriété universelle de l'éclatement.
On trouvera les détails ci-dessous.

\medskip\noindent
Soit $U, A, f, M^\bullet$ une carte affine pour $(X, D, M)$.
Tous les idéaux $I_i(M^\bullet, f)$ sauf un nombre fini sont
égaux à $A$, il est donc légitime de considérer
$$
I = \prod\nolimits_i I_i(M^\bullet, f)
$$
et c'est un idéal de type fini de $A$. Notons
$$
b_U : U' \to U
$$
l'éclatement de $U$ dans $I$. Alors $b_U^{-1}(U \cap D)$ est défini d'après
Diviseurs, lemme
\ref{divisors-lemma-blow-up-pullback-effective-Cartier}.
Rappelons que $f^{r_i} \in I_i(M^\bullet, f)$ et que, par conséquent,
$b_U$ est un éclatement $(U \setminus D)$-admissible.
D'après Diviseurs, lemme \ref{divisors-lemma-blowing-up-two-ideals},
pour chaque $i$, le morphisme $b_U$ se factorise en $U' \to U'_i \to U$,
où $U'_i \to U$ est l'éclatement dans $I_i(M^\bullet, f)$
et $U' \to U'_i$ est un autre éclatement. Il s'ensuit que
l'image inverse $I_i(M^\bullet, f)\mathcal{O}_{U'}$
de $I_i(M^\bullet, f)$ sur $U'$ est un faisceau d'idéaux inversible, voir
Diviseurs, lemmes \ref{divisors-lemma-blow-up-pullback-effective-Cartier} et
\ref{divisors-lemma-blowing-up-gives-effective-Cartier-divisor}.
Il s'ensuit que $(U', b^{-1}D, Lb^*M|_U)$ est un bon triplet, voir
le lemme \ref{lemma-pullback-triple-ideals-good}
pour le comportement des idéaux $I_i(-,-)$ par image inverse.
Enfin, nous affirmons que $b_U : U' \to U$ possède la propriété universelle
mentionnée dans la partie (2) de l'énoncé du lemme. En effet, supposons que
$h : Y \to U$ soit un morphisme de schémas tel que
l'image inverse $E = h^{-1}(D \cap U)$ soit définie et que $(Y, E, Lh^*M)$
soit un bon triplet. Alors $Y$ est recouvert par des cartes affines
$(V, B, g, N^\bullet)$ telles que $I_i(N^\bullet, g)$ soit
un idéal inversible pour chaque $i$. Alors $g$ et l'image de
$f$ dans $B$ diffèrent par une unité, puisqu'ils découpent tous deux le diviseur de Cartier
effectif $E \cap V$. Nous pouvons donc supposer que $g$ est l'image de $f$ d'après
Plus sur l'algèbre, lemme \ref{more-algebra-lemma-eta-change-unit}.
Alors $I_i(N^\bullet, g)$ est isomorphe à $I_i(M^\bullet \otimes_A B, g)$
comme $B$-module d'après Plus sur l'algèbre, lemme
\ref{more-algebra-lemma-ideal-well-defined}.
Ainsi $I_i(M^\bullet \otimes_A B, g) = I_i(M^\bullet, f)B$
(lemme \ref{lemma-pullback-triple-ideals-good})
est un $B$-module inversible.
Par conséquent, l'idéal $IB$ est inversible. Il s'ensuit que $I\mathcal{O}_Y$
est inversible. Nous obtenons donc une unique factorisation de $h$ à travers
$b_U$ d'après Diviseurs, lemme \ref{divisors-lemma-universal-property-blowing-up}.

\medskip\noindent
Soit $\mathcal{B}$ l'ensemble des ouverts affines $U \subset X$
pour lesquels il existe une carte affine $(U, A, f, M^\bullet)$
pour $(X, D, M)$. Alors $\mathcal{B}$ est une base de la topologie
de $X$ ; détails omis. Pour $U \in \mathcal{B}$, nous avons
le morphisme $b_U : U' \to U$ construit ci-dessus, qui satisfait
la propriété universelle sur $U$. Si $U_1 \subset U_2 \subset X$
appartiennent tous deux à $\mathcal{B}$, alors $b_{U_1} : U'_1 \to U_1$
est canoniquement isomorphe à
$$
b_{U_2}|_{b_{U_2}^{-1}(U_1)} : b_{U_2}^{-1}(U_1) \longrightarrow U_1
$$
par la propriété universelle. Autrement dit, nous obtenons un isomorphisme
$U'_1 \to b_{U_2}^{-1}(U_1)$ sur $U_1$. Ces isomorphismes satisfont
la condition de cocycle (de nouveau par la propriété universelle) et donc, d'après
Constructions, lemme \ref{constructions-lemma-relative-glueing},
nous obtenons un morphisme $b : X' \to X$ dont la restriction à chaque
$U$ de $\mathcal{B}$ est isomorphe à $U' \to U$.
Alors le morphisme $b : X' \to X$ satisfait les propriétés (1) et (2)
de l'énoncé du lemme, car ces propriétés peuvent être vérifiées
localement (détails omis).

\medskip\noindent
Il nous reste à démontrer la dernière assertion du lemme.
Soit $W \subset X$ un ouvert quasi-compact.
Choisissons un recouvrement fini $W = U_1 \cup \ldots \cup U_T$
tel que, pour chaque $1 \leq t \leq T$, il existe une carte affine
$(U_t, A_t, f_t, M_t^\bullet)$. Nous utiliserons ci-dessous le fait que, pour tout
ouvert affine $V = \Spec(B) \subset U_t \cap U_{t'}$, nous avons
(a) les images de $f_t$ et $f_{t'}$ dans $B$ diffèrent par une unité, et
(b) les complexes $M_t^\bullet \otimes_A B$ et $M_{t'} \otimes_A B$
définissent des objets isomorphes de $D(B)$. Pour $i \in \mathbf{Z}$, posons
$$
N_i = \max\nolimits_{t = 1, \ldots, T}
\left(\sum\nolimits_{j \geq i} (-1)^{j - i} rk(M_t^j)\right)
$$
Alors $N_t - \sum\nolimits_{j \geq i} (-1)^{j - i} rk(M_t^j) \geq 0$
et nous pouvons considérer les idéaux
$$
I_{t, i} =
f_t^{N_i - \sum\nolimits_{j \geq i} (-1)^{j - i} rk(M_t^j)}
I_i(M_t^\bullet, f_t)
\subset A_t
$$
Il résulte de Plus sur l'algèbre, lemmes
\ref{more-algebra-lemma-eta-change-unit} et
\ref{more-algebra-lemma-ideal-well-defined},
que les idéaux $I_{t, i}$ se recollent en un idéal quasi-cohérent de type fini
$\mathcal{I}_i \subset \mathcal{O}_W$. De plus, tous ces idéaux sauf
un nombre fini sont égaux à $\mathcal{O}_W$.
Clairement, le morphisme $X' \to X$ construit ci-dessus
se restreint à l'éclatement de $W$ dans le produit
des idéaux $\mathcal{I}_i$. Ceci achève la preuve.
\end{proof}

\begin{lemma}
\label{lemma-complex-and-divisor-blowup}
Dans la situation \ref{situation-complex-and-divisor}, soit $b : X' \to X$
le morphisme du lemme \ref{lemma-complex-and-divisor-blowup-pre}.
Considérons le diviseur de Cartier effectif $D' = b^{-1}D$ de faisceau d'idéaux
$\mathcal{I}' \subset \mathcal{O}_{X'}$. Alors $Q = L\eta_{\mathcal{I}'}Lb^*M$
est un objet parfait de $D(\mathcal{O}_{X'})$.
\end{lemma}

\begin{proof}
Cela résulte des lemmes \ref{lemma-complex-and-divisor-blowup-pre}
et \ref{lemma-good-triple}.
\end{proof}

\begin{lemma}
\label{lemma-complex-and-divisor-blowup-base-change}
Dans la situation \ref{situation-complex-and-divisor}, soit $h : Y \to X$
un morphisme de schémas tel que l'image inverse $E = h^{-1}D$
soit définie. Soient $b : X' \to X$, resp.\ $c : Y' \to Y$, construits dans le
lemme \ref{lemma-complex-and-divisor-blowup-pre} pour
$D \subset X$ et $M$, resp.\ $E \subset Y$ et $Lh^*M$.
Alors $Y'$ est le transformé strict de $Y$ relativement à $b : X' \to X$
(voir la preuve pour une formulation précise) et
$$
L\eta_{\mathcal{J}'}L(h \circ c)^*M = L(Y' \to X')^*Q
$$
où $Q = L\eta_{\mathcal{I}'}Lb^*M$ est comme dans le
lemme \ref{lemma-complex-and-divisor-blowup}.
En particulier, si $(Y, E, Lh^*M)$ est un bon triplet et si
$k : Y \to X'$ est l'unique morphisme tel que
$h = b \circ k$, alors $L\eta_\mathcal{J}Lh^*M = Lk^*Q$.
\end{lemma}

\begin{proof}
Notons $E' = c^{-1}E$. Alors $(Y', E', L(h \circ c)^*M)$ est un bon
triplet. Par conséquent, d'après la propriété universelle du
lemme \ref{lemma-complex-and-divisor-blowup-pre},
il existe un unique morphisme
$$
h' : Y' \longrightarrow X'
$$
tel que $b \circ h' = h \circ c$. En particulier, il existe un morphisme
$(h', c) : Y' \to X' \times_X Y$. Nous affirmons que, si $W \subset X$ est un ouvert
quasi-compact tel que $b^{-1}(W) \to W$ soit un éclatement,
ce morphisme identifie $Y'|_W$ au transformé strict de $Y_W$
relativement à $b^{-1}(W) \to W$. Pour vérifier ce fait,
il suffit de raisonner localement sur $W$ ; nous pouvons donc démontrer
l'assertion sur une carte affine. C'est ce que nous faisons au paragraphe suivant.

\medskip\noindent
Soit $(U, A, f, M^\bullet)$ une carte affine pour $(X, D, M)$.
Rappelons que, d'après la preuve du lemme \ref{lemma-complex-and-divisor-blowup},
la restriction de $b : X' \to X$ à $U$ est l'éclatement
de $U = \Spec(A)$ dans le produit des idéaux $I_i(M^\bullet, f)$.
Maintenant, si $V = \Spec(B) \subset Y$ est un ouvert affine avec $h(V) \subset U$,
alors $(V, B, g, M^\bullet \otimes_A B)$ est une carte affine pour
$(Y, E, Lh^*M)$, où $g \in B$ est l'image de $f$, voir le
lemme \ref{lemma-pullback-triple-ideals-good}.
Ainsi, la restriction de $c : Y' \to Y$ à $V$
est l'éclatement dans le produit des idéaux
$I_i(M^\bullet, f)B$, c'est-à-dire que le morphisme $c : Y' \to Y$ au-dessus de
$h^{-1}(U)$ est l'éclatement de $h^{-1}(U)$ dans
l'idéal $\prod I_i(M^\bullet, f) \mathcal{O}_{h^{-1}(U)}$.
Puisque cela vaut aussi pour le transformé strict, nous voyons que
notre assertion sur les transformés stricts est vraie.

\medskip\noindent
Cela étant dit, l'égalité
$L\eta_{\mathcal{J}'}L(h \circ c)^*M = L(Y' \to X')^*Q$
résulte du
lemme \ref{lemma-complex-and-divisor-eta-pull}.
La dernière assertion en est un cas particulier
(à savoir, le cas où $c = \text{id}_Y$ et $k = h'$).
\end{proof}

\begin{lemma}
\label{lemma-complex-and-divisor-blowup-good}
Dans la situation \ref{situation-complex-and-divisor}, soit $W \subset X$
le sous-schéma ouvert maximal sur lequel les faisceaux de cohomologie
de $M$ sont localement libres. Alors le morphisme $b : X' \to X$
du lemme \ref{lemma-complex-and-divisor-blowup-pre} est un isomorphisme
au-dessus de $W$.
\end{lemma}

\begin{proof}
Cela est vrai car, pour toute carte affine $(U, A, f, M^\bullet)$
avec $U \subset W$, les idéaux $I_i(M^\bullet, f)$ sont localement
engendrés par une puissance de $f$ d'après
Plus sur l'algèbre, lemme \ref{more-algebra-lemma-eta-cohomology-free}.
Puisque $f$ est un non-diviseur de zéro,
l'éclatement $b^{-1}(U) \to U$ est un isomorphisme.
\end{proof}

\begin{lemma}
\label{lemma-complex-and-divisor-blowup-good-T}
Soient $X, D, \mathcal{I}, M$ comme dans la
situation \ref{situation-complex-and-divisor}.
Si $(X, D, M)$ est un bon triplet, alors il existe une
immersion fermée
$$
i : T \longrightarrow D
$$
de présentation finie ayant les propriétés suivantes :
\begin{enumerate}
\item $T$ contient schématiquement $D \cap W$,
où $W \subset X$ est l'ouvert maximal sur lequel les
faisceaux de cohomologie de $M$ sont localement libres,
\item les faisceaux de cohomologie de $Li^*L\eta_\mathcal{I}M$
sont localement libres, et
\item pour tout point $t \in T$ d'image $x = i(t) \in W$, le rang
de $H^i(M)_x$ sur $\mathcal{O}_{X, x}$ et le rang
de $H^i(Li^*L\eta_\mathcal{I}M)_t$ sur $\mathcal{O}_{T, t}$ sont égaux.
\end{enumerate}
\end{lemma}

\begin{proof}
Soit $(U, A, f, M^\bullet)$ une carte affine pour $(X, D, M)$\newline
telle que $I_i(M^\bullet, f)$ soit un idéal principal pour tout $i \in \mathbf{Z}$.
Nous définissons alors $T \cap U \subset D \cap U$ comme le sous-schéma fermé
défini par l'idéal
$$
J(M^\bullet, f) = \sum J_i(M^\bullet, f) \subset A/fA
$$
étudié dans Plus sur l'algèbre, lemmes
\ref{more-algebra-lemma-eta-vanishing-beta-plus-pre} et
\ref{more-algebra-lemma-eta-vanishing-beta-plus} ;
d'après le second lemme, nous voyons que $T \cap U \to D \cap U$
est donné par l'homomorphisme d'anneaux $A/fA \to C$ qui y est étudié.
Puisque $(X, D, M)$ est un bon triplet, nous pouvons recouvrir $X$ par
des cartes affines de cette forme et, d'après le premier des deux
lemmes, cette construction se recolle. Nous obtenons donc un sous-schéma
fermé $T \subset D$ qui, sur les bonnes cartes affines ci-dessus,
est donné par l'idéal $J(M^\bullet, f)$. Les propriétés
(1) et (2) résultent alors du second lemme. Détails omis.
Petite observation pour aider le lecteur :
puisque $\eta_fM^\bullet$ est un complexe de modules localement libres
d'après Plus sur l'algèbre, lemme \ref{more-algebra-lemma-eta-locally-free},
nous voyons que $Li^*L\eta_\mathcal{I}M|_{T \cap U}$
est représenté par le complexe $\eta_fM^\bullet \otimes_A C$
de $C$-modules. L'assertion (3) sur les rangs résulte de
Cohomologie, lemme \ref{cohomology-lemma-eta-cohomology-locally-free}.
\end{proof}

\begin{lemma}
\label{lemma-complex-and-divisor-blowup-T}
Dans la situation \ref{situation-complex-and-divisor}. Soient $b : X' \to X$
et $D'$ comme dans le lemme \ref{lemma-complex-and-divisor-blowup-pre}. Soit
$Q = L\eta_{\mathcal{I}'}Lb^*M$ comme dans le
lemme \ref{lemma-complex-and-divisor-blowup}.
Soit $W \subset X$ l'ouvert maximal sur lequel $M$ a des
modules de cohomologie localement libres.
Alors il existe une immersion fermée $i : T \to D'$ de présentation finie
telle que
\begin{enumerate}
\item $D' \cap b^{-1}(W) \subset T$ schématiquement,
\item $Li^*Q$ ait des faisceaux de cohomologie localement libres, et
\item pour $t \in T$ s'envoyant sur $w \in W$, le rang
de $H^i(Li^*Q)_t$ sur $\mathcal{O}_{T, t}$ soit égal au
rang de $H^i(M)_x$ sur $\mathcal{O}_{X, x}$.
\end{enumerate}
\end{lemma}

\begin{proof}
Le lemme \ref{lemma-complex-and-divisor-blowup-good} nous dit que
$b$ est un isomorphisme au-dessus de $W$. Ainsi $b^{-1}(W) \subset X'$
est contenu dans l'ouvert maximal $W' \subset X'$ sur lequel $Lb^*M$
a des faisceaux de cohomologie localement libres. Les assertions du lemme elles-mêmes
sont alors une conséquence immédiate du
lemme \ref{lemma-complex-and-divisor-blowup-good-T} appliqué
à $(X', D', Lb^*M)$ et des autres lemmes mentionnés dans l'énoncé.
\end{proof}

\begin{lemma}
\label{lemma-complex-and-divisor-blowup-T-ranks}
Dans la situation \ref{situation-complex-and-divisor}. Soient $b : X' \to X$,
$D' \subset X'$ et $Q$ comme dans le
lemme \ref{lemma-complex-and-divisor-blowup}.
Soit $\rho = (\rho_i)_{i \in \mathbf{Z}}$ une famille d'entiers.
Soit $W(\rho) \subset X$ le sous-schéma ouvert maximal sur lequel
$H^i(M)$ est localement libre de rang $\rho_i$ pour tout $i$.
Soit $i : T \to D'$ comme dans le lemme \ref{lemma-complex-and-divisor-blowup-T}.
Alors il existe un sous-schéma ouvert et fermé $T(\rho) \subset T$
contenant schématiquement $D' \cap b^{-1}(W(\rho))$
tel que $H^i(Li^*Q|_{T(\rho)})$ soit localement libre de rang $\rho_i$
pour tout $i$.
\end{lemma}

\begin{proof}
Soit $T(\rho) \subset T$ le sous-schéma ouvert et fermé
sur lequel\newline $H^i(Li^*Q)$ est de rang $\rho_i$ pour tout $i$. L'assertion
résulte immédiatement de l'assertion du
lemme \ref{lemma-complex-and-divisor-blowup-T}
sur les rangs des modules de cohomologie.
\end{proof}

\begin{lemma}
\label{lemma-complex-and-divisor-blowup-T-represented-bdd-loc-free}
Dans la situation \ref{situation-complex-and-divisor}. Soient $b : X' \to X$,
$D' \subset X'$ et $Q$ comme dans le
lemme \ref{lemma-complex-and-divisor-blowup}.
S'il existe un complexe localement borné $\mathcal{M}^\bullet$
de $\mathcal{O}_X$-modules localement libres de rang fini représentant $M$,
alors il existe un complexe localement borné $\mathcal{Q}^\bullet$
de $\mathcal{O}_{X'}$-modules localement libres de rang fini représentant $Q$.
\end{lemma}

\begin{proof}
Rappelons que $Q = L\eta_{\mathcal{I}'}Lb^*M$, où $\mathcal{I}'$ est le
faisceau d'idéaux du diviseur de Cartier effectif $D'$.
Le complexe localement borné
$(\mathcal{M}')^\bullet = b^*\mathcal{M}^\bullet$ de
$\mathcal{O}_{X'}$-modules localement libres de rang fini
représente $Lb^*M$. Ainsi, le lemme résulte du
lemme \ref{lemma-good-triple-bdd-loc-free}.
\end{proof}

\begin{lemma}
\label{lemma-complex-and-divisor-derived}
Soit $X$ un schéma et soit $D \subset X$ un diviseur de Cartier effectif. Soit
$M \in D(\mathcal{O}_X)$ un objet parfait. Soit $W \subset X$ l'ouvert maximal
sur lequel les faisceaux de cohomologie $H^i(M)$ sont localement libres.
Il existe un morphisme propre $b : X' \longrightarrow X$
et un objet $Q$ de $D(\mathcal{O}_{X'})$ ayant les propriétés suivantes :
\begin{enumerate}
\item $b : X' \to X$ est un isomorphisme au-dessus de $X \setminus D$,
\item $b : X' \to X$ est un isomorphisme au-dessus de $W$,
\item $D' = b^{-1}D$ est un diviseur de Cartier effectif,
\item $Q = L\eta_{\mathcal{I}'}Lb^*M$, où $\mathcal{I}'$
est le faisceau d'idéaux de $D'$,
\item $Q$ est un objet parfait de $D(\mathcal{O}_{X'})$,
\item il existe une immersion fermée $i : T \to D'$ de présentation finie
telle que
\begin{enumerate}
\item $D' \cap b^{-1}(W) \subset T$ schématiquement,
\item $Li^*Q$ ait des faisceaux de cohomologie localement libres de rang fini,
\item pour $t \in T$ d'image $w \in W$, le rang
de $H^i(Li^*Q)_t$ sur $\mathcal{O}_{T, t}$ soit égal au
rang de $H^i(M)_x$ sur $\mathcal{O}_{X, x}$,
\end{enumerate}
\item pour toute carte affine $(U, A, f, M^\bullet)$ pour $(X, D, M)$,
la restriction de $b$ à $U$ est l'éclatement de $U = \Spec(A)$
dans l'idéal $I = \prod I_i(M^\bullet, f)$, et
\item pour toute carte affine $(V, B, g, N^\bullet)$ pour $(X', D', Lb^*N)$
telle que $I_i(N^\bullet, g)$ soit principal, nous avons
\begin{enumerate}
\item $Q|_V$ correspond à $\eta_gN^\bullet$,
\item $T \subset V \cap D'$ correspond à l'idéal
$J(N^\bullet, g) = \sum J_i(N^\bullet, g) \subset B/gB$
étudié dans
Plus sur l'algèbre, lemme \ref{more-algebra-lemma-eta-vanishing-beta-plus}.
\end{enumerate}
\item Si $M$ peut être représenté par un complexe localement borné
de $\mathcal{O}_X$-modules localement libres de rang fini, alors $Q$ peut
être représenté par un complexe borné de
$\mathcal{O}_{X'}$-modules localement libres de rang fini.
\end{enumerate}
\end{lemma}

\begin{proof}
Cet énoncé rassemble les informations obtenues dans les
lemmes \ref{lemma-pullback-triple-ideals-good},
\ref{lemma-good-triple},
\ref{lemma-complex-and-divisor-eta-pull},
\ref{lemma-complex-and-divisor-blowup-pre},
\ref{lemma-complex-and-divisor-blowup},
\ref{lemma-complex-and-divisor-blowup-base-change},
\ref{lemma-complex-and-divisor-blowup-good},
\ref{lemma-complex-and-divisor-blowup-good-T},
\ref{lemma-complex-and-divisor-blowup-T}, et
\ref{lemma-complex-and-divisor-blowup-T-represented-bdd-loc-free}.
\end{proof}










\section{Éclatement de complexes, III}
\label{section-blowup-complexes-III}

\noindent
Dans cette section, nous donnons une ``version algébrique'' de la version de la
construction du graphe de Macpherson donnée dans \cite[Section 18.1]{F}.

\medskip\noindent
Soit $X$ un schéma. Soit $E$ un objet parfait de $D(\mathcal{O}_X)$.
Soit $U \subset X$ le sous-schéma ouvert maximal tel que
$E|_U$ ait des faisceaux de cohomologie localement libres.

\medskip\noindent
Considérons le diagramme commutatif
$$
\xymatrix{
\mathbf{A}^1_X \ar[r] \ar[rd] &
\mathbf{P}^1_X \ar[d]^p &
(\mathbf{P}^1_X)_\infty \ar[l] \ar[ld] \\
& X \ar@/_1em/[ur]_\infty
}
$$
Rappelons ici que $\mathbf{A}^1 = D_+(T_0)$ est le premier ouvert
affine standard de $\mathbf{P}^1$ et que $\infty = V_+(T_0)$ est le
diviseur de Cartier effectif complémentaire, et que le diagramme ci-dessus
est l'image inverse de ces schémas sur $X$. Observons que
$\infty : X \to (\mathbf{P}^1_X)_\infty$ est un isomorphisme.
Alors
$$
(\mathbf{P}^1_X, (\mathbf{P}^1_X)_\infty, Lp^*E)
$$
est un triplet comme dans la situation \ref{situation-complex-and-divisor}
de la section \ref{section-blowup-complexes-II}.
Soient
$$
b : W \longrightarrow \mathbf{P}^1_X\quad\text{et}\quad
W_\infty = b^{-1}((\mathbf{P}^1_X)_\infty)
$$
l'éclatement et le diviseur de Cartier effectif construits
à partir de ce triplet dans le
lemme \ref{lemma-complex-and-divisor-blowup-pre}.
Nous notons aussi
$$
Q = L\eta_{\mathcal{I}} Lb^*M = L\eta_\mathcal{I} L(p \circ b)^*E
$$
l'objet parfait de $D(\mathcal{O}_W)$
considéré dans le lemme \ref{lemma-complex-and-divisor-blowup}.
Ici $\mathcal{I} \subset \mathcal{O}_W$ est le faisceau d'idéaux
de $W_\infty$.

\begin{lemma}
\label{lemma-graph-construction}
La construction ci-dessus a les propriétés suivantes :
\begin{enumerate}
\item $b$ est un isomorphisme au-dessus de $\mathbf{P}^1_U \cup \mathbf{A}^1_X$,
\item la restriction de $Q$ à $\mathbf{A}^1_X$
est égale à l'image inverse de $E$,
\item il existe une immersion fermée $i : T \to W_\infty$
de présentation finie telle que $(W_\infty \to X)^{-1}U \subset T$
schématiquement et telle que $Li^*Q$ ait des faisceaux de cohomologie
localement libres,
\item pour $t \in T$ d'image $u \in U$, le
rang de $H^i(Li^*Q)_t$ sur $\mathcal{O}_{T, t}$ est égal au rang
de $H^i(M)_u$ sur $\mathcal{O}_{U, u}$,
\item si $E$ peut être représenté par un complexe localement borné de
$\mathcal{O}_X$-modules localement libres de rang fini, alors $Q$ peut être représenté
par un complexe localement borné de $\mathcal{O}_W$-modules localement libres de rang fini.
\end{enumerate}
\end{lemma}

\begin{proof}
Cela résulte immédiatement des résultats de la
section \ref{section-blowup-complexes-II} ; pour un énoncé
qui rassemble tout ce qui est nécessaire, voir le
lemme \ref{lemma-complex-and-divisor-derived}.
\end{proof}











\begin{multicols}{2}[\section{Autres chapitres}]
\noindent
Préliminaires
\begin{enumerate}
\item \hyperref[introduction-section-phantom]{Introduction}
\item \hyperref[conventions-section-phantom]{Conventions}
\item \hyperref[sets-section-phantom]{Théorie des ensembles}
\item \hyperref[categories-section-phantom]{Catégories}
\item \hyperref[topology-section-phantom]{Topologie}
\item \hyperref[sheaves-section-phantom]{Faisceaux sur les espaces}
\item \hyperref[sites-section-phantom]{Sites et faisceaux}
\item \hyperref[stacks-section-phantom]{Champs}
\item \hyperref[fields-section-phantom]{Corps}
\item \hyperref[algebra-section-phantom]{Algèbre commutative}
\item \hyperref[brauer-section-phantom]{Groupes de Brauer}
\item \hyperref[homology-section-phantom]{Algèbre homologique}
\item \hyperref[derived-section-phantom]{Catégories dérivées}
\item \hyperref[simplicial-section-phantom]{Méthodes simpliciales}
\item \hyperref[more-algebra-section-phantom]{Compléments d'algèbre}
\item \hyperref[smoothing-section-phantom]{Lissage des morphismes d'anneaux}
\item \hyperref[modules-section-phantom]{Faisceaux de modules}
\item \hyperref[sites-modules-section-phantom]{Modules sur les sites}
\item \hyperref[injectives-section-phantom]{Objets injectifs}
\item \hyperref[cohomology-section-phantom]{Cohomologie des faisceaux}
\item \hyperref[sites-cohomology-section-phantom]{Cohomologie sur les sites}
\item \hyperref[dga-section-phantom]{Algèbre différentielle graduée}
\item \hyperref[dpa-section-phantom]{Algèbre à puissances divisées}
\item \hyperref[sdga-section-phantom]{Faisceaux différentiels gradués}
\item \hyperref[hypercovering-section-phantom]{Hyperrecouvrements}
\end{enumerate}
Schémas
\begin{enumerate}
\setcounter{enumi}{25}
\item \hyperref[schemes-section-phantom]{Schémas}
\item \hyperref[constructions-section-phantom]{Constructions de schémas}
\item \hyperref[properties-section-phantom]{Propriétés des schémas}
\item \hyperref[morphisms-section-phantom]{Morphismes de schémas}
\item \hyperref[coherent-section-phantom]{Cohomologie des schémas}
\item \hyperref[divisors-section-phantom]{Diviseurs}
\item \hyperref[limits-section-phantom]{Limites de schémas}
\item \hyperref[varieties-section-phantom]{Variétés}
\item \hyperref[topologies-section-phantom]{Topologies sur les schémas}
\item \hyperref[descent-section-phantom]{Descente}
\item \hyperref[perfect-section-phantom]{Catégories dérivées de schémas}
\item \hyperref[more-morphisms-section-phantom]{Compléments sur les morphismes}
\item \hyperref[flat-section-phantom]{Compléments sur la platitude}
\item \hyperref[groupoids-section-phantom]{Schémas en groupoïdes}
\item \hyperref[more-groupoids-section-phantom]{Compléments sur les schémas en groupoïdes}
\item \hyperref[etale-section-phantom]{Morphismes étales de schémas}
\end{enumerate}
Thèmes de théorie des schémas
\begin{enumerate}
\setcounter{enumi}{41}
\item \hyperref[chow-section-phantom]{Homologie de Chow}
\item \hyperref[intersection-section-phantom]{Théorie de l'intersection}
\item \hyperref[pic-section-phantom]{Schémas de Picard des courbes}
\item \hyperref[weil-section-phantom]{Théories cohomologiques de Weil}
\item \hyperref[adequate-section-phantom]{Modules adéquats}
\item \hyperref[dualizing-section-phantom]{Complexes dualisants}
\item \hyperref[duality-section-phantom]{Dualité pour les schémas}
\item \hyperref[discriminant-section-phantom]{Discriminants et différentes}
\item \hyperref[derham-section-phantom]{Cohomologie de de Rham}
\item \hyperref[local-cohomology-section-phantom]{Cohomologie locale}
\item \hyperref[algebraization-section-phantom]{Géométrie algébrique et formelle}
\item \hyperref[curves-section-phantom]{Courbes algébriques}
\item \hyperref[resolve-section-phantom]{Résolution des surfaces}
\item \hyperref[models-section-phantom]{Réduction semi-stable}
\item \hyperref[functors-section-phantom]{Foncteurs et morphismes}
\item \hyperref[equiv-section-phantom]{Catégories dérivées de variétés}
\item \hyperref[pione-section-phantom]{Groupes fondamentaux de schémas}
\item \hyperref[etale-cohomology-section-phantom]{Cohomologie étale}
\item \hyperref[crystalline-section-phantom]{Cohomologie cristalline}
\item \hyperref[proetale-section-phantom]{Cohomologie pro-étale}
\item \hyperref[relative-cycles-section-phantom]{Cycles relatifs}
\item \hyperref[more-etale-section-phantom]{Compléments de cohomologie étale}
\item \hyperref[trace-section-phantom]{La formule des traces}
\end{enumerate}
Espaces algébriques
\begin{enumerate}
\setcounter{enumi}{64}
\item \hyperref[spaces-section-phantom]{Espaces algébriques}
\item \hyperref[spaces-properties-section-phantom]{Propriétés des espaces algébriques}
\item \hyperref[spaces-morphisms-section-phantom]{Morphismes d'espaces algébriques}
\item \hyperref[decent-spaces-section-phantom]{Espaces algébriques décents}
\item \hyperref[spaces-cohomology-section-phantom]{Cohomologie des espaces algébriques}
\item \hyperref[spaces-limits-section-phantom]{Limites d'espaces algébriques}
\item \hyperref[spaces-divisors-section-phantom]{Diviseurs sur les espaces algébriques}
\item \hyperref[spaces-over-fields-section-phantom]{Espaces algébriques sur des corps}
\item \hyperref[spaces-topologies-section-phantom]{Topologies sur les espaces algébriques}
\item \hyperref[spaces-descent-section-phantom]{Descente et espaces algébriques}
\item \hyperref[spaces-perfect-section-phantom]{Catégories dérivées d'espaces}
\item \hyperref[spaces-more-morphisms-section-phantom]{Compléments sur les morphismes d'espaces}
\item \hyperref[spaces-flat-section-phantom]{Platitude sur les espaces algébriques}
\item \hyperref[spaces-groupoids-section-phantom]{Groupoïdes dans les espaces algébriques}
\item \hyperref[spaces-more-groupoids-section-phantom]{Compléments sur les groupoïdes dans les espaces}
\item \hyperref[bootstrap-section-phantom]{Amorçage}
\item \hyperref[spaces-pushouts-section-phantom]{Sommes amalgamées d'espaces algébriques}
\end{enumerate}
Thèmes de géométrie
\begin{enumerate}
\setcounter{enumi}{81}
\item \hyperref[spaces-chow-section-phantom]{Groupes de Chow des espaces}
\item \hyperref[groupoids-quotients-section-phantom]{Quotients de groupoïdes}
\item \hyperref[spaces-more-cohomology-section-phantom]{Compléments sur la cohomologie des espaces}
\item \hyperref[spaces-simplicial-section-phantom]{Espaces simpliciaux}
\item \hyperref[spaces-duality-section-phantom]{Dualité pour les espaces}
\item \hyperref[formal-spaces-section-phantom]{Espaces algébriques formels}
\item \hyperref[restricted-section-phantom]{Algébrisation des espaces formels}
\item \hyperref[spaces-resolve-section-phantom]{Résolution des surfaces revisitée}
\end{enumerate}
Théorie des déformations
\begin{enumerate}
\setcounter{enumi}{89}
\item \hyperref[formal-defos-section-phantom]{Théorie formelle des déformations}
\item \hyperref[defos-section-phantom]{Théorie des déformations}
\item \hyperref[cotangent-section-phantom]{Le complexe cotangent}
\item \hyperref[examples-defos-section-phantom]{Problèmes de déformation}
\end{enumerate}
Champs algébriques
\begin{enumerate}
\setcounter{enumi}{93}
\item \hyperref[algebraic-section-phantom]{Champs algébriques}
\item \hyperref[examples-stacks-section-phantom]{Exemples de champs}
\item \hyperref[stacks-sheaves-section-phantom]{Faisceaux sur les champs algébriques}
\item \hyperref[criteria-section-phantom]{Critères de représentabilité}
\item \hyperref[artin-section-phantom]{Axiomes d'Artin}
\item \hyperref[quot-section-phantom]{Espaces de Quot et de Hilbert}
\item \hyperref[stacks-properties-section-phantom]{Propriétés des champs algébriques}
\item \hyperref[stacks-morphisms-section-phantom]{Morphismes de champs algébriques}
\item \hyperref[stacks-limits-section-phantom]{Limites de champs algébriques}
\item \hyperref[stacks-cohomology-section-phantom]{Cohomologie des champs algébriques}
\item \hyperref[stacks-perfect-section-phantom]{Catégories dérivées de champs}
\item \hyperref[stacks-introduction-section-phantom]{Introduction aux champs algébriques}
\item \hyperref[stacks-more-morphisms-section-phantom]{Compléments sur les morphismes de champs}
\item \hyperref[stacks-geometry-section-phantom]{La géométrie des champs}
\end{enumerate}
Thèmes de théorie des espaces de modules
\begin{enumerate}
\setcounter{enumi}{107}
\item \hyperref[moduli-section-phantom]{Champs de modules}
\item \hyperref[moduli-curves-section-phantom]{Espaces de modules de courbes}
\end{enumerate}
Divers
\begin{enumerate}
\setcounter{enumi}{109}
\item \hyperref[examples-section-phantom]{Exemples}
\item \hyperref[exercises-section-phantom]{Exercices}
\item \hyperref[guide-section-phantom]{Guide bibliographique}
\item \hyperref[desirables-section-phantom]{Desiderata}
\item \hyperref[coding-section-phantom]{Style de programmation}
\item \hyperref[obsolete-section-phantom]{Obsolète}
\item \hyperref[fdl-section-phantom]{Licence de documentation libre GNU}
\item \hyperref[index-section-phantom]{Index généré automatiquement}
\end{enumerate}
\end{multicols}



\bibliography{my}
\bibliographystyle{amsalpha}

\end{document}
